% Generated by build_chapter.py from the GY1982 files. Do not edit by hand; edit the files and rebuild.
\chapter{The academic year 1982--83}\label{ch:1982}
\noindent This chapter covers all 37 supplied catalog records, from
6~September~1982 to 18~April~1983: 36 mathematical records and the year heading
579. The occurrence index contains 335 items. Catalog page ranges are
retained as supplied; item locators follow the actual handwriting. On the two
April meeting dates, the source puts 594 before 593 and 596 before 595.
\GSpage{P001115}, although included in the input packet, begins Deligne's
5~September~1983 talk and belongs to the following academic year. The present
mathematical coverage therefore ends at \GSpage{P001114}.

\paragraph{The year's connected lines.}
Affine Lie algebras, their Heisenberg and vertex realizations, the Virasoro
algebra, and semi-infinite constructions connect Kirillov, Fuchs, Feigin,
Neretin, and the string lectures. Zamolodchikov and Polyakov move between
classical constraints, random surfaces, conformal symmetry, and Liouville
theory; Schwarz and Kobzarev supply the gauge-theoretic and defect background.
The companion keeps formal path integrals distinct from later mathematical
constructions. In particular, algebraic quantization of cyclic string charges
has an established existence result
\citep{enriquezhalbout2010quantization,bahnsmeinecke2011string}; it is not
presented as a newly unresolved existence question.

The geometric line runs through orbital integrals, canonical relations,
moment maps, and localization. Beilinson's disk and gluing discussion has an
unusually strong documentary match: his later paper explicitly follows
S.~I.~Gelfand's 1982 notes \citep[introduction]{beilinson1987glue}.
Manin's superflags and Gelfand's continuous Grassmannians have more uneven
levels of formulation. In the spring, the finite-gap and Hamiltonian line
joins Krichever's perturbation theory, Novikov's many-valued actions,
Semenov-Tian-Shansky's factorization, and Wilson's Grassmannian construction
of integrable hierarchies \citep{krichever1983hessians,semenov1983classical,segalwilson1985loop}.
Wodzicki's spectral zeta functions and boundary index theory create a bridge
to the final arithmetic discussions without identifying spectral and
arithmetic zeta functions indiscriminately.

\paragraph{What the classifications mean.}
The counts describe occurrences, including recurring exposition, rather than
independent discoveries. A publication class and literal correctness are
separate: an old theorem can be copied with a wrong sign, and a contemporary
paper can repair a false notebook formula.
\begin{center}
\begin{tabular}{@{}lrrrrrrr@{}}
\toprule
Class & A & B & C & D & E & F & Total\\
\midrule
Items & 47 & 29 & 3 & 0 & 214 & 42 & 335\\
\bottomrule
\end{tabular}
\end{center}
Here A denotes near-contemporary publication by the speaker; B, a later
established development by others; C, a specified partial resolution; D,
documented lack of development; E, prior knowledge; and F, insufficient
identification. No class D is assigned merely because a search failed.
The separate assessment axes and source audit distinguish direct reading
of a result from bibliographic metadata checked.

\paragraph{Questions and bounded development.}
The problem index has 19 entries, including exact recurrences.
The genuine remaining tasks include the global four-unitary free energy
(GP-1982-02), general continuation for curve zeta functions (GP-1982-42), and
actual automorphy of arbitrary curve cohomology (GP-1982-43).
For the matrix integral, the development proves finite-rank formulas and
uses the small-potential theorem of
\citet[Theorem~0.1]{collinsguionnetmaurel2009asymptotics} to justify
\(F_4(\beta)=-4\beta+6\beta^2+O(\beta^3)\) locally, in the notebook's
negative-exponent convention. An independent exact index contraction checks
the coefficients. This does not evaluate the integral at every coupling.
In arithmetic, the genus-two continuation theorem over totally real fields
\citep[Theorems~1.1.2--1.1.3]{boxercalegarigeepilloni2021surfaces}
must be distinguished from actual modularity with additional hypotheses.
The broad remarks on continuous Grassmannian topology, arbitrary abelian
theta functions, and group-free localization remain explicitly underspecified.

\paragraph{Continuities across years.}
The earlier occurrence index supplies exact returns to compact-group weight
multiplicities (GP-1981-41), eta regularity (GP-1981-43), and spectral
constructions without self-adjointness (GP-1981-44). The finite-gap and KP
lectures continue the 1973--77 sequence; the superspace discussion continues
1978-455 and 1981-537. These subject links are kept distinct from an exact
recurrence of a question. All present-status judgments are dated
November~2026; uncertain readings and unresolved identifications remain
visible rather than being completed by historical guesswork.

\GStalk{561}{1982-09-06}{A. A. Kirillov}{P001026-P001031}{Affine sl(2), the principal Heisenberg algebra, and vertex operators}
Kirillov explains the Lepowsky--Wilson realization of the basic affine \(\mathfrak{sl}_2\)-module. Gelfand interrupts with generators and relations, root monomials, Heisenberg subalgebras and possible analogies with paths, \(p\)-adic algebra and quadratic forms. The actual mathematical entry occupies P001026--P001029; the catalogue range overlaps Zamolodchikov, whose heading begins P001030. The loop cocycle, central generator and oscillator indices require explicit normalization corrections. The 1978 published construction makes the principal line expository in September 1982.

\begin{GSitem}{1982-561.1}{E}{Cartan generators and the Serre quotient}
\GSfield{Kind; source}{construction; \GSpage{P001026-P001027}; I. M. Gelfand}
\GSfield{The notes}{Generators \(e_i,h_i,f_i\) satisfy \([e_i,f_j]=\delta_{ij}h_i\), \([h_i,e_j]=a_{ij}e_j\), \([h_i,f_j]=-a_{ij}f_j\), and commuting \(h_i\). Gelfand says \ru{Корней можно не учить}. The extra relations are written as ordinary powers, and the off-diagonal inequality is strict.}
\GSfield{Statement}{For a generalized Cartan matrix \(A\), require \(a_{ii}=2\), \(a_{ij}\in\mathbb Z_{\le0}\) for \(i\ne j\), and \(a_{ij}=0\iff a_{ji}=0\). The derived Kac--Moody algebra is generated by these elements subject to the displayed brackets and \((\operatorname{ad}e_i)^{1-a_{ij}}e_j=(\operatorname{ad}f_i)^{1-a_{ij}}f_j=0\), \(i\ne j\). Root vectors are obtained from iterated Lie brackets, not associative products.}
\GSfield{Reconstruction}{The generalized Cartan conditions, adjoint notation, and distinction between the derived algebra and its Cartan extension are editorial qualifications.}
\GSfield{Fate}{The finite-growth construction predates the talk: \citet[Chapter II]{kac1968graded}; the standard presentation is organized in \citet[Chapter 1]{kac1983infinite}. The notebook compresses several statements into an informal uniqueness assertion.}
\GSfield{Present status}{As of November 2026: the qualified statement is established; the notebook is an exposition, not a newly open question.}
\GSfield{Assessment}{literal statement: false as written; publication match: partial; verification: publisher metadata checked; development: annotated; confidence: high}
\end{GSitem}
\begin{GSitem}{1982-561.2}{E}{Finite and affine types and the rank-two examples}
\GSfield{Kind; source}{theorem; \GSpage{P001027-P001028}}
\GSfield{The notes}{Finite-dimensionality is identified with \(A>0\); affine type with \(A\ge0\), rank \(n-1\). Rank-two examples are \(A_2,B_2,G_2\). The first affine matrix is written \(\left(\begin{smallmatrix}2&-2\\-2&-2\end{smallmatrix}\right)\), then corrected implicitly in Kirillov\textquotesingle s repetition.}
\GSfield{Statement}{For an indecomposable symmetrizable generalized Cartan matrix \(A\), choose positive diagonal \(D\) with \(DA\) symmetric. Finite type is equivalent to positive definiteness of \(DA\); affine type to positive semidefiniteness of corank one. The affine matrices here are \(\left(\begin{smallmatrix}2&-2\\-2&2\end{smallmatrix}\right)\) and \(\left(\begin{smallmatrix}2&-4\\-1&2\end{smallmatrix}\right)\).}
\GSfield{Reconstruction}{Indecomposability, symmetrizability and the positive diagonal symmetrizer are added; the lower-right affine entry is emended to +2.}
\GSfield{Fate}{Finite/affine classification is old material in \citet[Chapter II]{kac1968graded}; \citet[Chapter 4]{kac1983infinite} gives the generalized Cartan formulation. Positivity of a nonsymmetric matrix requires interpretation.}
\GSfield{Present status}{As of November 2026: the qualified statement is established; the notebook is an exposition, not a newly open question.}
\GSfield{Assessment}{literal statement: false as written; publication match: partial; verification: publisher metadata checked; development: annotated; confidence: high}
\end{GSitem}
\begin{GSitem}{1982-561.3}{E}{The Heisenberg central extension of a symplectic space}
\GSfield{Kind; source}{definition; \GSpage{P001028}}
\GSfield{The notes}{\(E\oplus\mathbb R\mathbf1\), with \([x,y]=\langle x,y\rangle\mathbf1\), is introduced as the Heisenberg algebra.}
\GSfield{Statement}{If \((E,\omega)\) is a real symplectic vector space, define \([(x,a),(y,b)]=(0,\omega(x,y))\). Bilinearity, antisymmetry and centrality give a Lie algebra; for finite-dimensional nondegenerate \(E\) its center is exactly \(\mathbb R\mathbf1\).}
\GSfield{Reconstruction}{Nondegeneracy and the finite-dimensional qualification for the center are explicit here.}
\GSfield{Fate}{This is the standard central-extension model underlying the oscillator realization in \citet[pp.~43--53]{lepowskywilson1978construction}; the definition itself needs no conjectural step.}
\GSfield{Present status}{As of November 2026: the qualified statement is established; the notebook is an exposition, not a newly open question.}
\GSfield{Assessment}{literal statement: correct; publication match: partial; verification: publisher metadata checked; development: annotated; confidence: high}
\end{GSitem}
\begin{GSitem}{1982-561.4}{E}{Loop realization and the sign of the central cocycle}
\GSfield{Kind; source}{construction; \GSpage{P001028}}
\GSfield{The notes}{The notes use \([A,B]=AB-BA+\operatorname{Res}\operatorname{tr}(A\,dB)\mathbf1\), with \(e_1=E_{12}\), \(f_1=E_{21}\), \(e_2=tE_{21}\), \(f_2=t^{-1}E_{12}\), and \(h_2=-h_1+\mathbf1\).}
\GSfield{Statement}{On \(\mathfrak{sl}_2(\mathbb C[t,t^{-1}])\oplus\mathbb C K\), the cocycle \(\operatorname{Res}\operatorname{tr}((dA)B)\) yields \([e_2,f_2]=-h_1+K\). Using \(\operatorname{Res}\operatorname{tr}(A\,dB)\) instead yields \(-h_1-K\). Either convention defines the central extension, but the displayed bracket and \(h_2\) must agree.}
\GSfield{Reconstruction}{The residue convention is fixed explicitly; replacing the cocycle sign is an emendation, not a different historical claim.}
\GSfield{Fate}{This is the loop realization of the affine algebra, already used in the construction summarized by \citet{lepowskywilson1978construction}; see also \citet[Chapter 7]{kac1983infinite}. The sign discrepancy is checked directly by differentiating \(t^{-1}\).}
\GSfield{Present status}{As of November 2026: the qualified statement is established; the notebook is an exposition, not a newly open question.}
\GSfield{Assessment}{literal statement: false as written; publication match: partial; verification: publisher metadata checked; development: annotated; confidence: high}
\end{GSitem}
\begin{GSitem}{1982-561.5}{F}{Root monomials and the proposal to make paths continuous}
\GSfield{Kind; source}{question; \GSpage{P001028}; I. M. Gelfand}
\GSfield{The notes}{Weights count occurrences of the \(e_j\) in a monomial. A staircase path is drawn, followed by \ru{как сделать, чтобы пути были непрерывны?}}
\GSfield{Statement}{A Lie word has multidegree given by its generator multiplicities. The question about continuous paths specifies neither a topology, a limiting parameter, nor which algebraic operations the path model must preserve; no unique mathematical problem can be recovered.}
\GSfield{Reconstruction}{Only the multigrading is definite. No Littelmann-path or stochastic-path theorem is imported.}
\GSfield{Fate}{Comparison with the published affine construction and Kac\textquotesingle s representation framework yields no exact target. Later path models cannot be identified with this sentence merely from the word ``paths\textquotesingle\textquotesingle.}
\GSfield{Present status}{As of November 2026: the multigrading is standard, while the continuous-path request remains underdetermined as a historical item.}
\GSfield{Assessment}{literal statement: not determinable; publication match: none found; verification: unresolved lead; development: annotated; confidence: low}
\end{GSitem}
\begin{GSitem}{1982-561.6}{E}{The basic highest-weight module}
\GSfield{Kind; source}{construction; \GSpage{P001029}}
\GSfield{The notes}{A cyclic vector satisfies \(e_1v_0=e_2v_0=0\), \(h_1v_0=0\), \(h_2v_0=v_0\). Its space is \(\mathbb C[x_1,x_3,x_5,\ldots]\).}
\GSfield{Statement}{For the affine algebra with \(K=h_1+h_2\), the irreducible integrable highest-weight module of weight \(\Lambda\) with \(\Lambda(h_1)=0\), \(\Lambda(h_2)=1\) has level one and the principal bosonic realization on the polynomial algebra in positive odd variables. The Verma module with these highest weights must first be quotiented by its maximal proper submodule.}
\GSfield{Reconstruction}{Irreducibility, integrability, level and the quotient qualification make the notebook’s “most economical” precise.}
\GSfield{Fate}{Exactly the odd-variable realization is identified by the original abstract of \citet{lepowskywilson1978construction}, published August 1978. The distinction between Verma and irreducible integrable modules is essential.}
\GSfield{Present status}{As of November 2026: the qualified statement is established; the notebook is an exposition, not a newly open question.}
\GSfield{Assessment}{literal statement: correct with missing hypotheses; publication match: exact; verification: publisher metadata checked; development: annotated; confidence: high}
\end{GSitem}
\begin{GSitem}{1982-561.7}{E}{Nonconjugate loop Cartans and the square-root obstruction}
\GSfield{Kind; source}{claim; \GSpage{P001029}}
\GSfield{The notes}{After quotienting by the center, diagonal loops and matrices \(\left(\begin{smallmatrix}0&f(t)\\tf(t)&0\end{smallmatrix}\right)\) form nonconjugate commutative subalgebras. Pointwise diagonalization is possible, but continuous diagonalization is not.}
\GSfield{Statement}{For \(t\in S^1\), \(B(t)=\left(\begin{smallmatrix}0&1\\t&0\end{smallmatrix}\right)\) has distinct eigenvalues \(\pm\sqrt t\). A continuous periodic diagonalization with labeled diagonal entries would give a continuous square root of \(t\), impossible since its winding number would be one half. Over \(\mathbb C[t,t^{-1}]\) this likewise obstructs conjugation to the split diagonal Cartan.}
\GSfield{Reconstruction}{The domain S1, complex coefficient field, and nonvanishing choice f=1 are specified; no obstruction is claimed on a simply connected arc.}
\GSfield{Fate}{The principal Heisenberg in \citet{lepowskywilson1978construction} implements this distinction. The winding-number explanation supplied here is an elementary reconstruction of the specific notebook example.}
\GSfield{Present status}{As of November 2026: the qualified statement is established; the notebook is an exposition, not a newly open question.}
\GSfield{Assessment}{literal statement: correct; publication match: partial; verification: publisher metadata checked; development: annotated; confidence: high}
\end{GSitem}
\begin{GSitem}{1982-561.8}{F}{A p-adic analogy without a specified statement}
\GSfield{Kind; source}{remark; \GSpage{P001029}; I. M. Gelfand}
\GSfield{The notes}{Gelfand says that ``this algebra\textquotesingle\textquotesingle arises naturally in the \(p\)-adic case.}
\GSfield{Statement}{The sentence does not specify a local field, a group, a building, or a representation. Affine root systems occur in several inequivalent \(p\)-adic constructions; the intended assertion is not determined.}
\GSfield{Fate}{The nearby loop-algebra context and earlier affine/building entries were compared, but they do not select one statement. The remark is not evidence of an unsolved classification.}
\GSfield{Present status}{As of November 2026: the historical referent remains unidentified.}
\GSfield{Assessment}{literal statement: not determinable; publication match: none found; verification: unresolved lead; development: annotated; confidence: low}
\end{GSitem}
\begin{GSitem}{1982-561.9}{F}{Quadratic forms and an unspecified count of 2 to the n}
\GSfield{Kind; source}{remark; \GSpage{P001029}; I. M. Gelfand}
\GSfield{The notes}{Gelfand mentions Mishchenko and quadratic forms: \ru{количество \(2^n\), но кан. виды не описаны}.}
\GSfield{Statement}{No ground field, coefficient ring, equivalence relation or meaning of \(n\) is supplied. Consequently the count \(2^n\) is not a well-defined classification claim.}
\GSfield{Fate}{Neither the cited Lepowsky--Wilson setting nor the supplied earlier-year index supplies these missing data. No later quadratic-form classification is equated with this allusion.}
\GSfield{Present status}{As of November 2026: unidentified; a scan or other record specifying the forms would be needed before research development.}
\GSfield{Assessment}{literal statement: not determinable; publication match: none found; verification: unresolved lead; development: annotated; confidence: low}
\end{GSitem}
\begin{GSitem}{1982-561.10}{E}{Principal Heisenberg modes and the Fock representation}
\GSfield{Kind; source}{construction; \GSpage{P001029}}
\GSfield{The notes}{\(B_{2k+1}=\left(\begin{smallmatrix}0&t^k\\t^{k+1}&0\end{smallmatrix}\right)\), with \([B_{2k+1},B_{2l+1}]=\delta(k+l)(2l+1)\mathbf1\), is followed by multiplication/differentiation operators and an irreducibility assertion.}
\GSfield{Statement}{For the printed cocycle \(c(A,B)=\operatorname{Res}\operatorname{tr}(A\,dB)\), direct multiplication gives \([B_{2k+1},B_{2l+1}]=(2l+1)\delta_{k+l+1,0}K\). Abstractly, modes with \([a_m,a_n]=m\delta_{m+n,0}\) act for positive odd \(m\) by \(a_m=\partial_{x_m}\), \(a_{-m}=m x_m\) on polynomials; the vacuum is 1 and the representation is irreducible.}
\GSfield{Reconstruction}{The Fock-mode notation is normalized separately from the notebook\textquotesingle s inconsistent signs; positive modes annihilate the highest-weight vacuum.}
\GSfield{Fate}{The Heisenberg restriction is the standard oscillator component of \citet{lepowskywilson1978construction}. The missing shift in the Kronecker delta and mixed variable indices must be repaired independently of its normalization.}
\GSfield{Present status}{As of November 2026: the qualified statement is established; the notebook is an exposition, not a newly open question.}
\GSfield{Assessment}{literal statement: false as written; publication match: partial; verification: publisher metadata checked; development: annotated; confidence: high}
\end{GSitem}
\begin{GSitem}{1982-561.11}{E}{Exponential modes realize the remaining affine generators}
\GSfield{Kind; source}{method; \GSpage{P001029}}
\GSfield{The notes}{The notes form \(L=e^{\sum x_k}e^{4\sum\partial_{x_k}}=\sum L_m\), with weighted homogeneous components, and express \(C_{2k+1}\) and \(H_{2k}\) through them.}
\GSfield{Statement}{The principal realization extends the Heisenberg action on \(\mathbb C[x_1,x_3,\ldots]\) to the basic affine module using coefficients of a normally ordered formal exponential field. Each coefficient acts on a polynomial through finitely many derivative terms; the multiplication part is interpreted in the graded completion. The coefficients and grading must be chosen together to satisfy the affine brackets.}
\GSfield{Reconstruction}{The formal-series domain and completion are added; no unverified numerical normalization is promoted to an exact operator formula.}
\GSfield{Fate}{This is the specific construction described in \citet[pp.~43--53]{lepowskywilson1978construction}. Publisher abstract verification establishes the match; it does not verify the notebook\textquotesingle s coefficient 4 or recover its omitted spectral parameter.}
\GSfield{Present status}{As of November 2026: the qualified statement is established; the notebook is an exposition, not a newly open question.}
\GSfield{Assessment}{literal statement: correct with missing hypotheses; publication match: exact; verification: publisher metadata checked; development: annotated; confidence: medium}
\end{GSitem}

\GStalk{562}{1982-09-06}{A. B. Zamolodchikov}{P001029-P001034}{Classical string constraints and cyclic invariant charges}
The entry proper begins on P001030 and ends with the first lines of P001033. It moves from the Nambu--Goto action and Dirac constraints through two Witt algebras to oscillator quantization and the Virasoro anomaly. It then asks for the quantum algebra of cyclic iterated integrals. The classical construction, the asserted invariance of naively normally ordered operators, abstract deformation quantization, and completeness are distinct items with different later fates.

\begin{GSitem}{1982-562.1}{E}{Lorentzian area action and its sign}
\GSfield{Kind; source}{definition; \GSpage{P001030}}
\GSfield{The notes}{The target metric is \(\operatorname{diag}(1,-1,\ldots,-1)\), while the written integrand is \[\sqrt{\dot X^2(X^{\prime})^2-(\dot X\cdot X^{\prime})^2}.\]}
\GSfield{Statement}{For a timelike immersed worldsheet with induced metric \(h_{ab}=\partial_aX\cdot\partial_bX\), a real action, with \(T>0\), is \[\begin{aligned} S&=-T\int\sqrt{-\det h}\,d\tau\,d\sigma\\ &=-T\int\sqrt{(\dot X\cdot X^{\prime})^2-\dot X^2(X^{\prime})^2}\,d\tau\,d\sigma.\end{aligned}\] Degenerate worldsheets need separate treatment.}
\GSfield{Reconstruction}{The tension, timelike immersion and real square-root convention are supplied.}
\GSfield{Fate}{This is established string mechanics; \citet[Section~1.2]{tong2009string} supplies an explicitly normalized reference. The square-root sign follows directly from Lorentzian signature.}
\GSfield{Present status}{As of November 2026: the qualified statement is established; the notebook is an exposition, not a newly open question.}
\GSfield{Assessment}{literal statement: false as written; publication match: partial; verification: read the result; development: annotated; confidence: high}
\end{GSitem}
\begin{GSitem}{1982-562.2}{E}{Singular Legendre transform and first-class constraints}
\GSfield{Kind; source}{construction; \GSpage{P001030-P001031}}
\GSfield{The notes}{The canonical Hamiltonian vanishes. Constraints are first \(p^2-(X^{\prime})^2=0,\ p\cdot X^{\prime}=0\), then \((p\pm X^{\prime})^2=0\); a Hamiltonian is formed from arbitrary multipliers.}
\GSfield{Statement}{For the Lorentzian area action, \(p^2+T^2(X^{\prime})^2=0\) and \(p\cdot X^{\prime}=0\). The canonical Hamiltonian vanishes by degree-one velocity homogeneity. The total Hamiltonian is a linear combination of these first-class constraints; lapse and shift encode worldsheet reparametrization.}
\GSfield{Reconstruction}{Tension T is restored and the conflicting sign is repaired.}
\GSfield{Fate}{Dirac developed the framework in August 1949 lectures and \citet[Sections~2--4]{dirac1950generalized}. The string specialization is checked by differentiating the action.}
\GSfield{Present status}{As of November 2026: the qualified statement is established; the notebook is an exposition, not a newly open question.}
\GSfield{Assessment}{literal statement: false as written; publication match: partial; verification: read the result; development: annotated; confidence: high}
\end{GSitem}
\begin{GSitem}{1982-562.3}{E}{Two classical Witt algebras and Fourier conventions}
\GSfield{Kind; source}{theorem; \GSpage{P001031}}
\GSfield{The notes}{Modes of \((p+X^{\prime})^2\) and \((p-X^{\prime})^2\) close as commuting Witt algebras. \(X\) and \(p\) use opposite Fourier signs, but their mode bracket uses \(\delta_{n+m,0}\).}
\GSfield{Statement}{The smeared chiral constraints have the brackets of vector fields on \(S^1\), with opposite orientation conventions for the chiralities and zero mixed bracket. If \(X=\sum x_ne^{in\sigma}\), \(p=\sum p_me^{-im\sigma}\), normalization aside \(\{x_n,p_m\}=\delta_{nm}\); same-sign expansions give \(\delta_{n+m,0}\).}
\GSfield{Reconstruction}{The algebraic claim and the inconsistent Fourier convention are separated; factors of \(2\pi\) depend on normalization.}
\GSfield{Fate}{This is the classical constraint calculation underlying \citet[Sections~1.4 and 2.1]{tong2009string}. Fourier orthogonality gives the delta correction. First-class closure also proves preservation of the constraints by multiplier evolution.}
\GSfield{Present status}{As of November 2026: the qualified statement is established; the notebook is an exposition, not a newly open question.}
\GSfield{Assessment}{literal statement: false as written; publication match: partial; verification: read the result; development: annotated; confidence: high}
\end{GSitem}
\begin{GSitem}{1982-562.4}{E}{Oscillators, Fock space and the Virasoro central term}
\GSfield{Kind; source}{construction; \GSpage{P001031-P001032}}
\GSfield{The notes}{Oscillators combine \(p_n,x_n\), a vacuum is annihilated for \(n\ge0\), and quadratic modes acquire \(D(n^3-n)/12\). Individual Fourier modes are called Hermitian.}
\GSfield{Statement}{Normalize \([\alpha_m^\mu,\alpha_n^\nu]=m\eta^{\mu\nu}\delta_{m+n,0}\), \((\alpha_n^\mu)^\dagger=\alpha_{-n}^\mu\), \(\alpha_n|0;p\rangle=0\) for \(n>0\), and \(L_n=\tfrac12\sum_m:\alpha_{n-m}\cdot\alpha_m:\). On finite excitations, \([L_m,L_n]=(m-n)L_{m+n}+\tfrac D{12}(m^3-m)\delta_{m+n,0}\). The zero mode is momentum, not an oscillator annihilator.}
\GSfield{Reconstruction}{A common algebraic domain, Minkowski contraction, one-half factor and zero-mode sector are explicit.}
\GSfield{Fate}{The oscillator construction predates 1982; \citet[Section~2.1]{tong2009string} makes conventions explicit. Reality, mode factors and the zero-mode qualification are distinct defects.}
\GSfield{Present status}{As of November 2026: the qualified statement is established; the notebook is an exposition, not a newly open question.}
\GSfield{Assessment}{literal statement: false as written; publication match: partial; verification: read the result; development: annotated; confidence: high}
\end{GSitem}
\begin{GSitem}{1982-562.5}{E}{Incompatibility of all quantum Virasoro constraints}
\GSfield{Kind; source}{claim; \GSpage{P001032}}
\GSfield{The notes}{The notes say that all \(L_n\Psi=0\) are incompatible with the central term.}
\GSfield{Statement}{If \(c\ne0\) and every \(L_n v=0\) on a common invariant domain, then \(0=[L_2,L_{-2}]v=(4L_0+c/2)v=(c/2)v\), hence \(v=0\). This is a statement before canceling matter and ghost central charges.}
\GSfield{Reconstruction}{The elementary proof and the distinction from the total anomaly-free BRST complex are editorial.}
\GSfield{Fate}{The claim follows immediately from the known Virasoro relation; \citet[Section~2.1.2]{tong2009string} discusses the physical constraint prescription.}
\GSfield{Present status}{As of November 2026: the qualified statement is established; the notebook is an exposition, not a newly open question.}
\GSfield{Assessment}{literal statement: correct with missing hypotheses; publication match: partial; verification: read the result; development: annotated; confidence: high}
\end{GSitem}
\begin{GSitem}{1982-562.6}{E}{Half-constraints and physical matrix elements}
\GSfield{Kind; source}{method; \GSpage{P001032}}
\GSfield{The notes}{Only \(L_n\Psi=0\) for \(n\ge0\) are imposed; every constraint is then said to have vanishing physical matrix elements.}
\GSfield{Statement}{With \(L_n^\dagger=L_{-n}\), \(L_n v=0\) for \(n>0\) and \((L_0-a)v=0\) imply \(\langle u,L_nv\rangle=0\) for \(n\ne0\) and \(\langle u,(L_0-a)v\rangle=0\) between physical vectors. These conditions alone prove neither positivity nor consistency in arbitrary dimension.}
\GSfield{Reconstruction}{The intercept and limits of the matrix-element conclusion are restored.}
\GSfield{Fate}{This is the old covariant prescription with the intercept suppressed; see \citet[Section~2.1.2]{tong2009string}. Conjugation proves the negative-mode assertion.}
\GSfield{Present status}{As of November 2026: the qualified statement is established; the notebook is an exposition, not a newly open question.}
\GSfield{Assessment}{literal statement: correct with missing hypotheses; publication match: partial; verification: read the result; development: annotated; confidence: high}
\end{GSitem}
\begin{GSitem}{1982-562.7}{B}{Cyclic iterated integrals as classical string observables}
\GSfield{Kind; source}{construction; \GSpage{P001032-P001033}}
\GSfield{The notes}{Ordered products of \(y^\mu=p^\mu+iX^{\prime\mu}\) are integrated over cyclically ordered points. Their \(Q\)s are said to commute with constraints and momentum; the classical bracket has length \(n+m-2\).}
\GSfield{Statement}{For a closed classical string put \(u_\pm=p\pm TX^{\prime}\). Sum ordered iterated integrals over cyclic permutations of tensor indices. These charges are independent of the starting point, invariant under orientation-preserving reparametrizations and conserved on shell. Their span, shuffle multiplication and induced Poisson bracket form a Poisson algebra; the bracket contracts one index from each word.}
\GSfield{Reconstruction}{The real Lorentzian currents replace the unexplained printed \(i\). The length rule is schematic: coefficients, contractions and shuffle relations remain necessary.}
\GSfield{Fate}{This matches the Pohlmeyer--Rehren construction, explicitly given in \citet[Equations~1.1--1.2 and Section~2]{bahnsmeinecke2011string}. Its first cited publication is Pohlmeyer 1982; accessible evidence did not date circulation before this September talk. Later algebraic development is by others, hence B.}
\GSfield{Present status}{As of November 2026: the classical cyclic-charge Poisson algebra is established.}
\GSfield{Assessment}{literal statement: correct with missing hypotheses; publication match: exact; verification: read the result; development: annotated; confidence: high}
\end{GSitem}
\begin{GSitem}{1982-562.8}{B}{Normal ordering does not preserve all charge relations}
\GSfield{Kind; source}{claim; \GSpage{P001032}}
\GSfield{The notes}{After adding normal-ordering colons, the notes assert \([Q,L_n]=0\).}
\GSfield{Statement}{Expand classical charges in oscillators and normally order formally. Their constraint commutators and defining relations still require verification. Naive canonical normal ordering does not in general yield the invariant algebra with only observable quantum corrections.}
\GSfield{Reconstruction}{The universal quantum assertion is separated from classical invariance and alternative quantizations.}
\GSfield{Fate}{\citet[Sections~4--5]{bahns2004string} finds anomalous terms which cannot be removed by changing target dimension. This negative result concerns that prescription, not every deformation quantization.}
\GSfield{Present status}{As of November 2026: the unqualified canonical normal-ordering assertion fails.}
\GSfield{Assessment}{literal statement: false as written; publication match: exact; verification: read the result; development: annotated; confidence: high}
\end{GSitem}
\begin{GSitem}{1982-562.9}{B}{The quantum algebra of cyclic charges}
\GSfield{Kind; source}{question; \GSpage{P001032-P001033}}
\GSfield{The notes}{Zamolodchikov asks for the mutual commutation relations of the \(Q\)s, after stating the classical length rule.}
\GSfield{Statement}{For the cyclic-charge Poisson algebra \(\mathfrak h\), construct an associative formal deformation over \(\mathbb C[[\hbar]]\) with zeroth product the shuffle product and first commutator \(i\hbar\{\, ,\,\}\), and describe its quantum relations. The notes impose no Hilbert-space representation requirement.}
\GSfield{Reconstruction}{Formal deformation and the semiclassical condition specify the algebraic question. Preferred explicit presentations and positive representations are stronger requirements.}
\GSfield{Fate}{\citet[Theorem~1.1]{enriquezhalbout2010quantization}, circulated in 2008 and published in 2010, constructs quasi-Lie-bialgebra quantization. \citet[Proposition~15 and p.~20]{bahnsmeinecke2011string} explicitly deduces existence for \(\mathfrak h\) without quadratic generation. This substantial algebraic answer does not validate the notebook operators.}
\GSfield{Present status}{As of November 2026: abstract formal quantization exists; the question is not catalogued as an open existence problem.}
\GSfield{Assessment}{literal statement: correct with missing hypotheses; publication match: partial; verification: read the result; development: annotated; confidence: high; problem \hyperref[prob:GP-1982-01]{GP-1982-01}}
\end{GSitem}
\begin{GSitem}{1982-562.10}{B}{Completeness of invariant charges}
\GSfield{Kind; source}{claim; \GSpage{P001033}}
\GSfield{The notes}{The continuation calls \(\{Q\}\) a complete set of invariants commuting with \(P^\mu\).}
\GSfield{Statement}{Completeness requires a specified class of worldsheets and an equivalence by translations and reparametrization. It does not assert free algebraic independence or generation of every momentum-commuting function without regularity restrictions.}
\GSfield{Reconstruction}{Completeness is interpreted geometrically, without promoting the abbreviated assertion to an unrestricted theorem.}
\GSfield{Fate}{\citet[pp.~177--186]{pohlmeyerrehren1988completeness} treats the exact geometric question (received April 1987, published June 1988). Its opened abstract establishes the match; \citet[Section~1]{bahns2004string} explains auxiliary boosts and translations along total momentum. The original full hypotheses were inaccessible.}
\GSfield{Present status}{As of November 2026: a geometric completeness theorem exists for its stated class; the notebook omits necessary scope and auxiliary data.}
\GSfield{Assessment}{literal statement: correct with missing hypotheses; publication match: exact; verification: indirect; development: annotated; confidence: medium}
\end{GSitem}

\GStalk{563}{1982-09-13}{A. S. Schwarz}{P001032-P001038}{Gauge symmetry, topological strings and mirror particles}
The actual entry occupies P001033--P001036. A discussion of subjects to study leads to gauge covariantization, symmetry breaking and vacuum topology. Schwarz then describes strings whose holonomy reverses electric charge or exchanges ordinary and mirror matter. The printed group-theoretic criterion lacks a crucial stabilizer hypothesis. The final quantum-stability question and astrophysical anecdotes do not specify enough data for an independently identifiable theorem.

\begin{GSitem}{1982-563.1}{E}{Gauge covariantization and the connection transformation}
\GSfield{Kind; source}{construction; \GSpage{P001033}}
\GSfield{The notes}{For \(G\subset SO(n)\), replace \(\partial_\mu\) by \(\partial_\mu+A_\mu\), use \(F_{\mu\nu}=\partial_\mu A_\nu-\partial_\nu A_\mu+[A_\mu,A_\nu]\), and add a curvature-square term. The notes give \(\phi\mapsto g\phi\), \(A\mapsto gAg^{-1}+dg\,g^{-1}\).}
\GSfield{Statement}{For a connection \(D=d+A\) acting on a left representation and \(\phi^g=g\phi\), covariance \(D^g(g\phi)=gD\phi\) requires \(A^g=gAg^{-1}-dg\,g^{-1}\). Then \(F^g=gFg^{-1}\), and invariant inner products give gauge-invariant matter and Yang--Mills kinetic terms.}
\GSfield{Reconstruction}{The left action and sign of the covariant derivative are kept exactly as in the notebook; only the inconsistent connection sign is emended.}
\GSfield{Fate}{The local gauge principle is the content of \citet[pp.~191--195]{yangmills1954conservation}. The sign repair follows by expanding \(d(g\phi)\), independently of convention choices in that paper.}
\GSfield{Present status}{As of November 2026: the qualified statement is established; the notebook is an exposition, not a newly open question.}
\GSfield{Assessment}{literal statement: false as written; publication match: partial; verification: read the result; development: annotated; confidence: high}
\end{GSitem}
\begin{GSitem}{1982-563.2}{E}{Perturbative renormalizability after gauging}
\GSfield{Kind; source}{claim; \GSpage{P001034}}
\GSfield{The notes}{The statement attributed to 't Hooft says that if \(L\) is renormalizable, its localization is also renormalizable.}
\GSfield{Statement}{In four dimensions, Yang--Mills theories with compact gauge group and power-counting renormalizable matter interactions admit perturbative renormalization compatible with gauge identities when gauge anomalies cancel. Spontaneous breaking by scalar vacuum expectation values does not by itself destroy this property. Gauging an arbitrary global symmetry without anomaly checks is not a valid universal implication.}
\GSfield{Reconstruction}{Dimension, interaction degree, anomaly cancellation and perturbative scope are supplied.}
\GSfield{Fate}{\citet[Introduction and Sections~2--5]{thooft1971massive} treats the gauge identities and explicit scalar models in 1971. This is perturbative renormalizability, not construction of a nonperturbative continuum probability measure.}
\GSfield{Present status}{As of November 2026: the qualified statement is established; the notebook is an exposition, not a newly open question.}
\GSfield{Assessment}{literal statement: correct with missing hypotheses; publication match: partial; verification: read the result; development: annotated; confidence: high}
\end{GSitem}
\begin{GSitem}{1982-563.3}{E}{Vacuum manifold and symmetry breaking}
\GSfield{Kind; source}{definition; \GSpage{P001034}}
\GSfield{The notes}{Constant minimizers of \(U\) form \(\mathcal R\). For \(U(\phi)=\lambda(|\phi|^2-a^2)^2\), the vacuum set is a sphere; \(H_\phi\) is the subgroup fixing a vacuum.}
\GSfield{Statement}{For \(\lambda>0,a>0\) and \(\phi\in\mathbb R^n\), the minima are \(S^{n-1}_a\). When a compact group \(G\) acts transitively on a vacuum component, it is \(G/H_\phi\). This vacuum orbit is a space of Higgs field values; gauge-equivalent constant configurations need not be distinct physical vacua.}
\GSfield{Reconstruction}{Positivity, transitivity, vacuum component and gauge-equivalence distinctions are made explicit.}
\GSfield{Fate}{This is standard Higgs-model geometry; \citet[Section~3]{thooft1971massive} explicitly analyzes the scalar minimum and residual gauge symmetry.}
\GSfield{Present status}{As of November 2026: the qualified statement is established; the notebook is an exposition, not a newly open question.}
\GSfield{Assessment}{literal statement: correct with missing hypotheses; publication match: partial; verification: read the result; development: annotated; confidence: high}
\end{GSitem}
\begin{GSitem}{1982-563.4}{F}{Grand unification and the affine-group proposal}
\GSfield{Kind; source}{remark; \GSpage{P001034}; I. M. Gelfand}
\GSfield{The notes}{The low-energy group is \(SU(3)\times SU(2)\times U(1)\); candidates for a larger group include \(SU(5),SO(10),E_8\). Gelfand proposes an affine group, emphasizing finite rank and an invariant form.}
\GSfield{Statement}{An embedding of the indicated Lie algebra into a larger gauge algebra is only group-theoretic input. A unification model additionally requires matter representations, symmetry-breaking fields, interactions and anomaly conditions. The affine suggestion specifies none of these and therefore does not define a unique field theory or existence problem.}
\GSfield{Fate}{The finite-group mirror construction subsequently published by \citet{fateevfrolovschwarz1982mirror} is an identifiable nearby model, not a realization of the unspecified affine proposal. No exact affine-model target was found in the author bibliography or prior occurrence index.}
\GSfield{Present status}{As of November 2026: the proposal remains underdetermined; it is not evidence of an unsolved general unification theorem.}
\GSfield{Assessment}{literal statement: not determinable; publication match: thematic only; verification: unresolved lead; development: annotated; confidence: medium}
\end{GSitem}
\begin{GSitem}{1982-563.5}{E}{Monopole sectors from the sphere at spatial infinity}
\GSfield{Kind; source}{construction; \GSpage{P001034}}
\GSfield{The notes}{Finite-energy fields on \(\mathbb R^3\) approach \(\mathcal R\) at infinity, defining \(S^2\to\mathcal R\), hence an element of \(\pi_2(\mathcal R)\) and monopoles.}
\GSfield{Statement}{For a Higgs model with positive potential vanishing exactly on a vacuum manifold, impose asymptotic boundary conditions ensuring that the normalized Higgs field extends continuously to the sphere at infinity. Its homotopy class lies in \(\pi_2(\mathcal R)\). Nontrivial classes distinguish sectors but do not alone prove that a finite-energy minimizing solution exists in every class.}
\GSfield{Reconstruction}{Explicit boundary regularity and the distinction between a topological sector and a solution are added.}
\GSfield{Fate}{The topology is established background to the gauge-string analysis summarized in \citet{schwarz1982charge}; the notebook compresses the analytic existence issue into the topological classification.}
\GSfield{Present status}{As of November 2026: the qualified statement is established; the notebook is an exposition, not a newly open question.}
\GSfield{Assessment}{literal statement: correct with missing hypotheses; publication match: partial; verification: publisher metadata checked; development: annotated; confidence: high}
\end{GSitem}
\begin{GSitem}{1982-563.6}{E}{Strings and fundamental-group sectors}
\GSfield{Kind; source}{construction; \GSpage{P001034-P001035}}
\GSfield{The notes}{A string is independent of the longitudinal coordinate and has finite energy per unit length. A linking circle maps to \(\mathcal R\); stability is followed by a question mark.}
\GSfield{Statement}{For a translationally invariant configuration on \(\mathbb R^2\times\mathbb R\), require \(\int_{\mathbb R^2}\mathcal E\,dx\,dy<\infty\) and a limiting Higgs map \(S^1_\infty\to\mathcal R\). Its class in \(\pi_1(\mathcal R)\) obstructs deformation to vacuum through admissible finite-energy configurations. This is topological stability, not a proof of spectral, nonlinear or quantum stability.}
\GSfield{Reconstruction}{Boundary conditions, energy per unit length and the meaning of stability are separated.}
\GSfield{Fate}{This is the vacuum-topology mechanism used in \citet{schwarz1982charge}. The scanned linking-circle diagram represents a meridian of the string, not a worldsheet boundary.}
\GSfield{Present status}{As of November 2026: the qualified statement is established; the notebook is an exposition, not a newly open question.}
\GSfield{Assessment}{literal statement: correct with missing hypotheses; publication match: partial; verification: publisher metadata checked; development: annotated; confidence: high}
\end{GSitem}
\begin{GSitem}{1982-563.7}{E}{Charge-reversing holonomy and the missing disconnected stabilizer}
\GSfield{Kind; source}{claim; \GSpage{P001035}}
\GSfield{The notes}{If an element \(\alpha\in\widetilde G\) conjugates the electromagnetic \(U(1)\) by inversion, the notes assert existence of a topologically nontrivial charge-reversing string. The unbroken group is written \(K\times U(1)\).}
\GSfield{Statement}{Let \(\widetilde G\) be simply connected and \(H\) its unbroken subgroup. A string class corresponds to a component of \(H\), since \(\pi_1(\widetilde G/H)\cong\pi_0(H)\). A representative holonomy \(\alpha\in H\) can reverse charge if \(\alpha Q\alpha^{-1}=-Q\). Merely finding \(\alpha\) in \(\widetilde G\) is insufficient: connected \(H\) gives no such fundamental-group string. Field-equation existence needs model-specific energy analysis.}
\GSfield{Reconstruction}{The homotopy-sequence calculation and the condition alpha in H outside its identity component repair the stated criterion.}
\GSfield{Fate}{The mechanism was announced in \citet{schwarz1981nonlocal}; \citet{schwarz1982charge} explicitly assumes a disconnected unbroken subgroup. The 1981 announcement makes this E despite the later detailed paper.}
\GSfield{Present status}{As of November 2026: the Alice-string mechanism is established; the printed sufficient condition is incomplete.}
\GSfield{Assessment}{literal statement: false as written; publication match: exact; verification: publisher metadata checked; development: annotated; confidence: high}
\end{GSitem}
\begin{GSitem}{1982-563.8}{A}{Ordinary-to-mirror particle transport around a string}
\GSfield{Kind; source}{claim; \GSpage{P001035}}
\GSfield{The notes}{There are unified models containing mirror particles and strings around which a particle becomes its mirror partner. The two sectors interact weakly through heavy bosons and share gravity; mirror particles are not antiparticles.}
\GSfield{Statement}{Choose a gauge model whose disconnected residual symmetry exchanges two matter sectors. Holonomy in that component transports a state in one sector to its mirror partner. Decoupling of local nongravitational interactions at low energy is a separate statement about heavy mediator masses and couplings, not a consequence of topology alone.}
\GSfield{Reconstruction}{Holonomy exchange and low-energy coupling suppression are distinct assumptions.}
\GSfield{Fate}{\citet[pp.~427--432]{fateevfrolovschwarz1982mirror} publishes the specific model in December 1982, after this talk, and its abstract states precisely the sector exchange. An earlier preprint date was not verified. Classification A records the speaker's subsequent publication, with abstract-level verification.}
\GSfield{Present status}{As of November 2026: the proposed gauge-theory mechanism was published; this does not establish the existence of mirror matter in nature.}
\GSfield{Assessment}{literal statement: correct with missing hypotheses; publication match: exact; verification: publisher metadata checked; development: annotated; confidence: high}
\end{GSitem}
\begin{GSitem}{1982-563.9}{F}{Quantum stability of the proposed strings}
\GSfield{Kind; source}{question; \GSpage{P001036}}
\GSfield{The notes}{\ru{Неясно, сущ. ли квантово стабильные нити}; classical stability is asserted.}
\GSfield{Statement}{A quantum-stability question must specify the Lagrangian, coupling regime, boundary sector and whether stability means an exact conserved topological charge, absence of tunneling, a positive fluctuation spectrum or a lifetime. The notes do not select one meaning or one model.}
\GSfield{Fate}{The Schwarz and Fateev--Frolov--Schwarz publications identify nearby classical constructions; their accessible abstracts do not supply the missing quantum target. No exact quantum-stability theorem can responsibly be assigned to this sentence.}
\GSfield{Present status}{As of November 2026: underdetermined, rather than a newly formulated universal open problem.}
\GSfield{Assessment}{literal statement: not determinable; publication match: none found; verification: unresolved lead; development: annotated; confidence: medium}
\end{GSitem}
\begin{GSitem}{1982-563.10}{F}{Mirror-photon visibility and the quasar anecdote}
\GSfield{Kind; source}{remark; \GSpage{P001035-P001036}}
\GSfield{The notes}{The diagram suggests that a string crossing a line of sight can turn photons into mirror photons. Khlopov mentions a preprint about a quasar whose luminosity changed sharply over a century.}
\GSfield{Statement}{A sector-exchanging holonomy can change which asymptotic gauge field a transported state couples to. Predicting an observed light curve requires an explicit spacetime string configuration, propagation and source model. The notebook gives no quasar name, preprint author list or measurable parameters.}
\GSfield{Fate}{The mirror-particle paper identifies the particle-transport mechanism only. The observational anecdote could not be uniquely identified from the named author and sparse description; it is not treated as experimental confirmation.}
\GSfield{Present status}{As of November 2026: the particular astrophysical claim is unidentifiable from this record.}
\GSfield{Assessment}{literal statement: not determinable; publication match: thematic only; verification: unresolved lead; development: annotated; confidence: low}
\end{GSitem}

\GStalk{564}{1982-09-20}{Ya. G. Sinai}{P001036-P001039}{Period doubling, renormalization and period tripling}
On P001037 Sinai proposes physics topics and then discusses Feigenbaum universality; Migdal adds further agenda items. Sinai's mathematics ends midway down P001038. The entry gives period-doubling scaling, robustness of suitable families, a suggested flow extension, complex period tripling and the renormalization fixed-point equation. Numerical approximation is distinguished from a validated proof, and the broad post-accumulation question is not treated as a uniquely defined open problem.

\begin{GSitem}{1982-564.1}{E}{Normalized even families do not by themselves force a cascade}
\GSfield{Kind; source}{definition; \GSpage{P001037-P001038}}
\GSfield{The notes}{Families satisfy \(f_\mu:[-1,1]\to[-1,1]\), \(f_\mu(-x)=f_\mu(x)\), \(f_\mu(0)=1\); \(1-\mu x^2\) is the example. The period-doubling sequence is then asserted.}
\GSfield{Statement}{A period-doubling universality theorem needs a unimodal family with nondegenerate quadratic critical point, appropriate regularity and a parameter curve crossing the renormalization stable manifold transversely. Evenness and normalization alone admit constant families and do not imply any bifurcation.}
\GSfield{Reconstruction}{Unimodality, quadratic criticality and parameter transversality are added; the constant-family counterexample is elementary.}
\GSfield{Fate}{\citet[Definition of M and Theorem~5]{lanford1982feigenbaum} specifies the analytic framework and verifies transverse crossing for the quadratic family. His April 1981 preprint predates the talk.}
\GSfield{Present status}{As of November 2026: the qualified statement is established; the notebook is an exposition, not a newly open question.}
\GSfield{Assessment}{literal statement: correct with missing hypotheses; publication match: partial; verification: read the result; development: annotated; confidence: high}
\end{GSitem}
\begin{GSitem}{1982-564.2}{E}{Universal period-doubling parameter scaling}
\GSfield{Kind; source}{claim; \GSpage{P001038}}
\GSfield{The notes}{At successive parameters \(\mu_n\), stable periods double; \((\mu_n-\mu_{n+1})/(\mu_{n-1}-\mu_n)\to\delta^{-1}\), with \(\delta\approx4.6992\).}
\GSfield{Statement}{For an admissible transverse period-doubling family in the quadratic universality class, the scaling constant is the expanding eigenvalue of the period-doubling renormalization derivative, numerically \(\delta\approx4.669201609\). The asymptotic \(\mu_n-\mu_\infty\sim C\delta^{-n}\) requires the relevant bifurcation-locus transversality, not just existence of a fixed point.}
\GSfield{Reconstruction}{The admissible-family qualifier and distinction between the eigenvalue theorem and full parameter asymptotic are explicit.}
\GSfield{Fate}{\citet[Theorems~3--5]{lanford1982feigenbaum} proves the core fixed-point and hyperbolicity assertions but explicitly notes an unproved transversality refinement after Theorem~4. The entry is expository; the incorrect decimal is visible in the scan.}
\GSfield{Present status}{As of November 2026: the qualified statement is established; the notebook is an exposition, not a newly open question.}
\GSfield{Assessment}{literal statement: false as written; publication match: partial; verification: read the result; development: annotated; confidence: high}
\end{GSitem}
\begin{GSitem}{1982-564.3}{E}{Why universality is robust for transverse families}
\GSfield{Kind; source}{claim; \GSpage{P001038}}
\GSfield{The notes}{The set of families with the stated bifurcation sequence is said to be open.}
\GSfield{Statement}{Near a hyperbolic renormalization fixed point with one unstable direction, its stable manifold has codimension one. A one-parameter family crossing that manifold transversely retains this local crossing under sufficiently small perturbations in the chosen differentiable topology. Infinite renormalizability itself is not an open property of individual maps.}
\GSfield{Reconstruction}{Local family transversality replaces an incorrect reading as an open set of infinitely renormalizable maps.}
\GSfield{Fate}{\citet[Theorem~3 and the discussion following it]{lanford1982feigenbaum} supplies the codimension-one stable manifold. The openness statement concerns families and their crossing, as the handwritten insertion \ru{сем-в} indicates.}
\GSfield{Present status}{As of November 2026: the qualified statement is established; the notebook is an exposition, not a newly open question.}
\GSfield{Assessment}{literal statement: correct with missing hypotheses; publication match: partial; verification: read the result; development: annotated; confidence: high}
\end{GSitem}
\begin{GSitem}{1982-564.4}{F}{Suggested extension from maps to differential equations}
\GSfield{Kind; source}{claim; \GSpage{P001038}}
\GSfield{The notes}{For \(\dot x_i=f(x_1,\ldots,x_r,\mu)\), the notes state the same scaling and an attracting periodic orbit whose basin excludes only finitely many closed trajectories. A sequence of nested loops is drawn.}
\GSfield{Statement}{Period doubling of a smooth flow is studied through a Poincare return map near a periodic orbit. To import a one-dimensional renormalization theorem one needs a specified reduction and contraction in transverse directions. The notes do not give these assumptions; a global basin assertion for arbitrary vector fields is false.}
\GSfield{Reconstruction}{A Poincare-section formulation is a proposed interpretation, not a theorem contained in the notes.}
\GSfield{Fate}{The map theorem of \citet{lanford1982feigenbaum} does not establish this general flow claim. The named ODE, phase-space region and reduction are absent, so no exact subsequent theorem is assigned.}
\GSfield{Present status}{As of November 2026: the unrestricted flow statement is false; its intended restricted family cannot be identified.}
\GSfield{Assessment}{literal statement: false as written; publication match: thematic only; verification: unresolved lead; development: annotated; confidence: medium}
\end{GSitem}
\begin{GSitem}{1982-564.5}{A}{Two-parameter period tripling and complex scaling constants}
\GSfield{Kind; source}{claim; \GSpage{P001038}}
\GSfield{The notes}{With two parameters, period tripling can occur and a universal constant has been found.}
\GSfield{Statement}{For holomorphic quadratic-like families depending on a complex parameter, period tripling motivates a renormalization fixed point with a complex expanding eigenvalue; the real and imaginary parameters form a two-parameter family. Numerical evidence for the fixed point, spectrum and extension to nearby real two-dimensional maps is distinct from a rigorous theorem for all such families.}
\GSfield{Reconstruction}{The complex-analytic setting and numerical character are recovered from the exact later publication.}
\GSfield{Fate}{\citet[pp.~187--188]{golbergsinaikhanin1983tripling}, received 28 June 1982 and published in 1983, is the exact speaker-coauthored match. It describes polynomial truncations and numerical eigenvalues. Its numerical evidence is not relabeled as a complete proof; the receipt date is not asserted as public circulation.}
\GSfield{Present status}{As of November 2026: the recorded numerical discovery has an exact publication. The broader real-family universality claim is not certified by that paper; no unrestricted resolution is claimed.}
\GSfield{Assessment}{literal statement: correct with missing hypotheses; publication match: exact; verification: read the result; development: annotated; confidence: high}
\end{GSitem}
\begin{GSitem}{1982-564.6}{E}{The renormalization fixed-point equation}
\GSfield{Kind; source}{construction; \GSpage{P001038}}
\GSfield{The notes}{\(-\alpha^{-1}f(f(\alpha x))=f(x)\) is given; the notes say renormalization arises from linearizing it.}
\GSfield{Statement}{For an even normalized unimodal map with \(a=-f(1)>0\), define \(Rf(x)=-a^{-1}f(f(-ax))\). A fixed point satisfies the notebook equation after matching the evenness and scale convention. Renormalization is the nonlinear operator \(R\); its derivative at a fixed point governs local stability and universal scaling.}
\GSfield{Reconstruction}{The scale is determined by normalization and the roles of R and its derivative are separated.}
\GSfield{Fate}{\citet[Sections~1--2, Theorems~1 and 3]{lanford1982feigenbaum} proves existence of an analytic fixed point and one-dimensional expanding spectrum. Linearization analyzes the renormalization operator; it does not define the original nonlinear map.}
\GSfield{Present status}{As of November 2026: the qualified statement is established; the notebook is an exposition, not a newly open question.}
\GSfield{Assessment}{literal statement: correct with missing hypotheses; publication match: partial; verification: read the result; development: annotated; confidence: high}
\end{GSitem}
\begin{GSitem}{1982-564.7}{E}{Degree-ten approximation and a computer-assisted proof}
\GSfield{Kind; source}{method; \GSpage{P001038}}
\GSfield{The notes}{The equation is said to be solved using a polynomial of degree ten and numerical computation.}
\GSfield{Statement}{A degree-ten polynomial is an approximate fixed point. A proof can combine a Newton-type map, explicit analytic tail estimates and interval arithmetic bounding rounding errors, so that a contraction theorem establishes a nearby exact analytic solution. A finite polynomial fit alone proves neither existence nor hyperbolicity.}
\GSfield{Reconstruction}{The contraction and error-control steps are supplied from the matched proof.}
\GSfield{Fate}{This is precisely the method outlined in \citet[Section~3]{lanford1982feigenbaum}, already circulated in April 1981. The distinction between approximation and validated bounds explains the notebook compression.}
\GSfield{Present status}{As of November 2026: the qualified statement is established; the notebook is an exposition, not a newly open question.}
\GSfield{Assessment}{literal statement: correct with missing hypotheses; publication match: partial; verification: read the result; development: annotated; confidence: high}
\end{GSitem}
\begin{GSitem}{1982-564.8}{F}{Stochastic behavior beyond period-doubling accumulation}
\GSfield{Kind; source}{question; \GSpage{P001038}}
\GSfield{The notes}{The question asks what happens to the vector field for \(\mu>\mu_\infty\). For maps, an ergodic invariant measure and mixing are asserted as the onset of stochasticity.}
\GSfield{Statement}{For the quadratic family there is a positive-measure set of parameters with an absolutely continuous invariant probability measure; this is not an assertion for every parameter above the accumulation point. Ergodicity alone is inadequate as a chaos criterion, since a periodic-orbit measure is ergodic. The unspecified vector-field family cannot be assigned a unique post-accumulation theorem.}
\GSfield{Reconstruction}{The definite interval-map comparison is separated from the underdetermined question about flows.}
\GSfield{Fate}{\citet[Theorem~2]{jakobson1980invariant} announces the positive-parameter-measure result before this talk, including robustness under \(C^3\) perturbations of families. It does not answer the unformulated flow question or eliminate periodic windows.}
\GSfield{Present status}{As of November 2026: the map comparison has a pre-talk theorem; no definite unresolved flow problem can be extracted.}
\GSfield{Assessment}{literal statement: correct with missing hypotheses; publication match: partial; verification: read the result; development: annotated; confidence: medium}
\end{GSitem}

\GStalk{565}{1982-09-20}{A. A. Migdal}{P001037-P001042}{Stochastic quantization and reduction to large random matrices}
The mathematical entry starts on the lower half of P001038 after Sinai and ends on P001040. It develops Langevin quantization, factorization of normalized traces, quenched momenta and reduction of a matrix field to finitely many matrices. The definite calculation requested is the all-coupling large-\(N\) free energy of four unitary matrices. The manuscript retains a negative exponent and all ordered index pairs; those conventions are essential to the development. The one-matrix external-field integral at the end is a separate already studied model.

\begin{GSitem}{1982-565.1}{E}{Planar counting and the matrix-field action}
\GSfield{Kind; source}{method; \GSpage{P001037-P001038}}
\GSfield{The notes}{The agenda states \(u_\infty=sp_\infty=O_\infty\). Migdal considers a \(U(N)\)-valued or \(\mathfrak u(N)\)-valued field, action \(S=(N/\lambda)\int\operatorname{tr}[\tfrac12(\partial\varphi)^2+U(\varphi)]\), and large-\(N\) non-Gaussian fluctuations.}
\GSfield{Statement}{For Hermitian matrix fields with a stable regularized Euclidean action and fixed coupling, ribbon-graph perturbation theory organizes connected vacuum diagrams by Euler characteristic: an orientable genus-\(g\) diagram contributes \(N^{2-2g}\). Orthogonal and symplectic matrix theories also have nonorientable corrections. For anti-Hermitian fields a positive kinetic form is \(-\operatorname{tr}(\partial\varphi)^2\), so the printed sign requires a convention change.}
\GSfield{Reconstruction}{A regulator, stable action, Hermitian convention and perturbative meaning are specified; the three classical series are not literally identified.}
\GSfield{Fate}{\citet[Sections 2--3]{thooft1974planar}, received in December 1973 and published in 1974, establishes the diagrammatic expansion. The agenda equality is at most a leading-planar analogy, not an isomorphism of infinite classical Lie algebras.}
\GSfield{Present status}{As of November 2026: the qualified statement is established; the notebook is an exposition, not a newly open question.}
\GSfield{Assessment}{literal statement: false as written; publication match: partial; verification: read the result; development: annotated; confidence: high}
\end{GSitem}
\begin{GSitem}{1982-565.2}{E}{Langevin equilibrium and stochastic quantization}
\GSfield{Kind; source}{method; \GSpage{P001039}}
\GSfield{The notes}{The equation \(\partial_\tau\varphi+\delta S/\delta\varphi=J\) is assigned unit white-noise covariance and is claimed to reproduce \(Z^{-1}\int Qe^{-S}D\varphi\) by a long-time average.}
\GSfield{Statement}{For a finite regulator, \(dX=-\nabla S(X)\,d\tau+\sqrt2\,dB_\tau\) has invariant density \(Z^{-1}e^{-S}\) when \(Z<\infty\); ergodicity is additionally needed for almost-sure time averages. Unit noise covariance instead gives density proportional to \(e^{-2S}\). A continuum limit requires separate construction and renormalization.}
\GSfield{Reconstruction}{The finite-dimensional statement and direct equilibrium check are editorial; the historical identification relies on the later paper abstract.}
\GSfield{Fate}{The Parisi--Wu method was published in 1981 and is the explicit starting point of \citet[abstract]{alfarosakita1983quenched}. Here the normalization is checked directly from the finite-dimensional stationary Fokker--Planck equation. No construction of four-dimensional continuum Yang--Mills measure follows.}
\GSfield{Present status}{As of November 2026: the qualified statement is established; the notebook is an exposition, not a newly open question.}
\GSfield{Assessment}{literal statement: false as written; publication match: partial; verification: indirect; development: annotated; confidence: high}
\end{GSitem}
\begin{GSitem}{1982-565.3}{E}{Large-N factorization and the meaning of self-averaging}
\GSfield{Kind; source}{claim; \GSpage{P001039}}
\GSfield{The notes}{For \(G=N^{-1}\operatorname{tr}(\varphi(x_1)\cdots\varphi(x_n))\), the notes assert self-averaging up to \(O(1/N)\), after taking stochastic time large.}
\GSfield{Statement}{In the planar expansion of a regularized invariant matrix theory, connected correlators of two normalized single traces are \(O(N^{-2})\). When this estimate holds uniformly for the actual probability measure, standard deviations are \(O(N^{-1})\). Planar graph counting alone does not establish uniform stochastic-time convergence or an exchange of the limits \(N\to\infty\), \(\tau\to\infty\), and cutoff removal.}
\GSfield{Reconstruction}{The distinction between covariance and standard deviation and between formal and convergent limits is added.}
\GSfield{Fate}{The topological mechanism is in \citet[Sections 2--3]{thooft1974planar}; factorization is used in \citet[abstract]{eguchikawai1982reduction}, published April 1982. The qualification between formal expansion and a measure theorem is necessary.}
\GSfield{Present status}{As of November 2026: the qualified statement is established; the notebook is an exposition, not a newly open question.}
\GSfield{Assessment}{literal statement: correct with missing hypotheses; publication match: partial; verification: read the result; development: annotated; confidence: high}
\end{GSitem}
\begin{GSitem}{1982-565.4}{E}{Quenched momenta and conjugated noise}
\GSfield{Kind; source}{construction; \GSpage{P001039-P001040}}
\GSfield{The notes}{Diagonal random momenta \(p_\mu\) are introduced, with \(\varphi(x)=e^{ip\cdot x}\bar\varphi e^{-ip\cdot x}\) and similarly conjugated matrix noise. The density is \(e^{-p^2/\Lambda^2}\), said to have variance \(1/\Lambda^2\).}
\GSfield{Statement}{In a regulated planar expansion, color-index momentum differences \(p_i-p_j\) represent line momenta. Averaging sampled diagonal momenta with a specified cutoff density approximates the momentum integrals. For one real coordinate with normalized density proportional to \(e^{-p^2/\Lambda^2}\), the variance is \(\Lambda^2/2\). Conjugated finite-matrix noise is not entrywise spacetime white noise as an exact finite-\(N\) identity.}
\GSfield{Reconstruction}{The planar and cutoff meaning is made explicit; the elementary Gaussian variance is corrected.}
\GSfield{Fate}{\citet[abstract]{alfarosakita1983quenched} gives the same stochastic route; its July 1982 preprint preceded this September talk, although publication followed in February 1983. This identifies an expository method, without certifying a general nonperturbative equivalence.}
\GSfield{Present status}{As of November 2026: the qualified statement is established; the notebook is an exposition, not a newly open question.}
\GSfield{Assessment}{literal statement: false as written; publication match: exact; verification: indirect; development: annotated; confidence: high}
\end{GSitem}
\begin{GSitem}{1982-565.5}{A}{The reduced Langevin equation and formal matrix measure}
\GSfield{Kind; source}{construction; \GSpage{P001040}}
\GSfield{The notes}{Equations (15)--(18) replace spatial derivatives by commutators with \(p_\mu\), write \(U(\varphi)\) in the drift, and then integrate \(e^{+\operatorname{tr}([p,\varphi]^2+U_\Lambda(\varphi))}\) over a set labeled \(U(N)\), without a partition-function denominator.}
\GSfield{Statement}{If \(\varphi(x)=e^{ipx}\Phi e^{-ipx}\), then \(-\partial^2\varphi=e^{ipx}\sum_\mu[p_\mu,[p_\mu,\Phi]]e^{-ipx}\). The drift from \(\operatorname{tr}U(\Phi)\) contains \(U^{\prime}(\Phi)\). A reduced expectation must specify whether \(\Phi\) is Hermitian, anti-Hermitian or unitary, a stable action, a measure on that domain, normalization, and the separate averaging of quenched momenta.}
\GSfield{Reconstruction}{The displayed commutator identity is checked directly. The action derivative, domain and probability normalization are repairs, not silent replacements of the manuscript.}
\GSfield{Fate}{The stochastic prescription is documented by \citet{alfarosakita1983quenched}; the speaker published the related QCD reduction in \citet[pp.~425--428]{migdal1982reduction} later in 1982. Only institutional metadata and abstracts were accessible, so the exact displayed equations are not certified as the published result.}
\GSfield{Present status}{As of November 2026: a historical reduction program is identified; the notebook equations themselves do not define the claimed probability measure.}
\GSfield{Assessment}{literal statement: false as written; publication match: partial; verification: publisher metadata checked; development: annotated; confidence: high}
\end{GSitem}
\begin{GSitem}{1982-565.6}{C}{All-coupling free energy of four unitary matrices}
\GSfield{Kind; source}{problem; \GSpage{P001040}}
\GSfield{The notes}{Equation (19) asks for \[\lim_{N\to\infty}N^{-2}\log\int\exp\!\left[-\beta N\sum_{\mu,\nu=1}^4\operatorname{tr}(U_\mu U_\nu U_\mu^\dagger U_\nu^\dagger)\right]\prod_{\mu=1}^4dU_\mu.\] The minus sign and the complete ordered double sum are visible in the scan.}
\GSfield{Statement}{With normalized Haar measures on \(U(N)\), set \[Z_N(\beta)=\int_{U(N)^4}\exp\!\left[-\beta N\sum_{\mu,\nu=1}^4\operatorname{tr}(U_\mu U_\nu U_\mu^\dagger U_\nu^\dagger)\right]dU.\] Determine \(F(\beta)=\lim_{N\to\infty}N^{-2}\log Z_N(\beta)\) for real \(\beta\), including existence, analytic branches and transition values. The sum is real because ordered off-diagonal terms are conjugate pairs; diagonal terms contribute \(4N\).}
\GSfield{Reconstruction}{Normalized Haar measure, real coupling and the precise all-coupling objective are supplied. Diagonal terms and exponent sign are retained.}
\GSfield{Fate}{\citet[Theorem~0.1]{collinsguionnetmaurel2009asymptotics}, first circulated in 2006 and published in 2009, gives convergence and analytic expansion for sufficiently small potential. The chapter computes the first two derivatives and \(F(\beta)=-4\beta+6\beta^2+O(\beta^3)\). The original \citet{eguchikawai1982reduction} uses different sign and summation conventions. Searches through November 2026 found no exact global solution for the displayed model; \citet[Section~5]{maeta2026regularized} remains a numerical small-model study and proposes larger reduced models for later work. The bounded calculation in \hyperref[dev:1982-migdal]{the development section} proves the first two local coefficients in the exact notebook convention.}
\GSfield{Present status}{As of November 2026: small-coupling existence and expansion are established; a global all-coupling formula and phase description for the exact recorded convention were not located in the primary literature searched.}
\GSfield{Assessment}{literal statement: correct with missing hypotheses; publication match: partial; verification: read the result; development: developed in this chapter; confidence: medium; problem \hyperref[prob:GP-1982-02]{GP-1982-02}}
\GSfield{Worth developing}{Worth developing: a global characterization would connect a rigorously controlled local free energy with possible large-matrix phase transitions in a definite finite-dimensional integral.}
\end{GSitem}
\begin{GSitem}{1982-565.7}{E}{The one-unitary external-field integral}
\GSfield{Kind; source}{claim; \GSpage{P001040}}
\GSfield{The notes}{Equation (20) is \(\int_{U(N)}e^{\beta N\operatorname{tr}(Uh+h^\dagger U^\dagger)}dU\), with arbitrary \(h\). It is said to be solved by a name resembling \ru{Пиризен}, Gross, and then Fateev.}
\GSfield{Statement}{For fixed complex \(N\times N\) source \(h\), unitary invariance reduces the normalized Haar integral to the singular values of \(h\). The large-\(N\) external-field problem requires a specified limiting distribution and scaling of these singular values; it is distinct from the interacting four-matrix integral.}
\GSfield{Reconstruction}{The likely name identification is editorial. The finite-\(N\) reduction to singular values follows from Haar invariance.}
\GSfield{Fate}{\citet[pp.~120--124]{brezingross1980external}, published November 1980, is the precise Brézin--Gross identification. The opened author bibliography and publisher abstract support the large-\(N\) match; they do not verify an arbitrary finite-\(N\) closed formula or the unidentified Fateev extension.}
\GSfield{Present status}{As of November 2026: the qualified statement is established; the notebook is an exposition, not a newly open question.}
\GSfield{Assessment}{literal statement: correct with missing hypotheses; publication match: exact; verification: publisher metadata checked; development: annotated; confidence: medium}
\end{GSitem}
\begin{GSitem}{1982-565.8}{F}{The unspecified Fateev extension of the external-field solution}
\GSfield{Kind; source}{remark; \GSpage{P001040}}
\GSfield{The notes}{The final sentence adds \ru{затем Фатеев} without title, formula, parameter restriction or reference.}
\GSfield{Statement}{The note does not distinguish a finite-\(N\) formula, a new asymptotic regime, or a different external-field model. It therefore does not determine a separate mathematical claim attributable to Fateev.}
\GSfield{Fate}{The Brézin--Gross primary bibliographic match was checked, together with searches for Fateev and the unitary external-field integral. No exact Fateev publication matching this abbreviated sentence was verified.}
\GSfield{Present status}{As of November 2026: the attribution remains unidentified; it is not entered as an open problem.}
\GSfield{Assessment}{literal statement: not determinable; publication match: none found; verification: unresolved lead; development: annotated; confidence: low}
\end{GSitem}

\GStalk{566}{1982-09-27}{D. B. Fuchs}{P001040-P001043}{Virasoro Verma modules and semi-infinite exterior powers}
The actual entry occupies P001041--P001042. Fuchs presents the central extension of circle vector fields, Verma modules, tensor-density modules and their semi-infinite exterior powers, then the Kac determinant criterion and a second oscillator construction. The explicitly cited Feigin--Fuchs paper had appeared in the second issue of 1982. The notebook reverses two signs, omits commutator factors, and writes operators which are not the required Heisenberg modes; these defects are separated from the valid representation-theoretic content.

\begin{GSitem}{1982-566.1}{E}{The Witt algebra and its one-dimensional central extension}
\GSfield{Kind; source}{definition; \GSpage{P001041}}
\GSfield{The notes}{\([e_i,e_j]=(j-i)e_{i+j}\), and the central correction at \(i+j=0\) is \((j^3-j)z/12\). The extension is called unique.}
\GSfield{Statement}{The Laurent-polynomial vector fields \(e_i=z^{i+1}\partial_z\) have the displayed bracket. Up to rescaling the central generator and equivalence, there is one nontrivial one-dimensional central extension, the Virasoro algebra. For smooth circle vector fields one uses continuous Lie-algebra cochains; uniqueness excludes the split extension.}
\GSfield{Reconstruction}{Nontriviality, equivalence and the continuous-cochain qualification are supplied.}
\GSfield{Fate}{The degree-two class is in \citet[Theorem~1]{gelfandfuchs1968circle} (1968); the exact mode convention appears in \citet[Section~3.1]{feiginfuchs1982invariant}.}
\GSfield{Present status}{As of November 2026: the qualified statement is established; the notebook is an exposition, not a newly open question.}
\GSfield{Assessment}{literal statement: correct with missing hypotheses; publication match: exact; verification: read the result; development: annotated; confidence: high}
\end{GSitem}
\begin{GSitem}{1982-566.2}{E}{Verma modules, PBW basis and partition-number dimensions}
\GSfield{Kind; source}{construction; \GSpage{P001041}}
\GSfield{The notes}{The vacuum satisfies \(e_{-1}v=e_{-2}v=0,\ e_0v=hv,\ zv=cv\); the ordered positive-mode monomials are required to be a basis.}
\GSfield{Statement}{The induced Verma module \(V_{h,c}\) has basis \(e_{i_1}\cdots e_{i_k}v\), \(i_1\ge\cdots\ge i_k>0\). Its level-\(n\) dimension is the partition number \(p(n)\), hence graded character \(\prod_{m\ge1}(1-q^m)^{-1}\). Since \(e_{-1},e_{-2}\) generate the negative subalgebra, their annihilation conditions imply every \(e_{-j}v=0\).}
\GSfield{Reconstruction}{The PBW basis condition is essential: a highest-weight quotient alone need not be a Verma module.}
\GSfield{Fate}{\citet[Section~3.1]{feiginfuchs1982invariant} gives this exact convention and universal property before the talk.}
\GSfield{Present status}{As of November 2026: the qualified statement is established; the notebook is an exposition, not a newly open question.}
\GSfield{Assessment}{literal statement: correct; publication match: exact; verification: read the result; development: annotated; confidence: high}
\end{GSitem}
\begin{GSitem}{1982-566.3}{E}{Tensor densities with two parameters}
\GSfield{Kind; source}{construction; \GSpage{P001041}}
\GSfield{The notes}{\(e_if_j=[j+\mu-(i+1)\lambda]f_{i+j}\), interpreted as \(f_j=z^{j+\mu}(dz)^{-\lambda}\).}
\GSfield{Statement}{The Lie derivative of a density of weight \(-\lambda\) gives \(e_i(z^{j+\mu}(dz)^{-\lambda})=[j+\mu-(i+1)\lambda]z^{i+j+\mu}(dz)^{-\lambda}\). For nonintegral \(\mu\), this is an algebraic Laurent module or a density with monodromy, not necessarily an ordinary single-valued function on the circle.}
\GSfield{Reconstruction}{The meaning of a nonintegral exponent is made explicit.}
\GSfield{Fate}{The exact algebraic module is defined in \citet[Section~2.1]{feiginfuchs1982invariant}; the displayed action can also be verified directly by the Lie-derivative rule.}
\GSfield{Present status}{As of November 2026: the qualified statement is established; the notebook is an exposition, not a newly open question.}
\GSfield{Assessment}{literal statement: correct with missing hypotheses; publication match: exact; verification: read the result; development: annotated; confidence: high}
\end{GSitem}
\begin{GSitem}{1982-566.4}{E}{Semi-infinite wedges and normal-ordering of the zero mode}
\GSfield{Kind; source}{construction; \GSpage{P001041-P001042}}
\GSfield{The notes}{Take wedges differing finitely from \(\cdots\wedge f_{-3}\wedge f_{-2}\wedge f_{-1}\). Nonzero modes act by a finite surviving sum, but the naive zero-mode sum diverges; the notes set \(e_0=[e_{-1},e_1]\) and compare \([e_{-2},e_2]\).}
\GSfield{Statement}{Fix the charge-zero tail. For \(i\ne0\), repeated factors make all but finitely many Leibniz terms vanish. The Virasoro relations require \(e_0=\tfrac12[e_{-1},e_1]\) and \(z=2([e_{-2},e_2]-4e_0)\). The central scalar records the finite discrepancy after regularizing the divergent diagonal sum.}
\GSfield{Reconstruction}{Charge, regularization and the missing factors are explicit.}
\GSfield{Fate}{\citet[Section~3.3, Proposition~3.1]{feiginfuchs1982invariant} supplies the exact wedge construction. The commutator coefficients follow directly from the notebook algebra.}
\GSfield{Present status}{As of November 2026: the qualified statement is established; the notebook is an exposition, not a newly open question.}
\GSfield{Assessment}{literal statement: false as written; publication match: exact; verification: read the result; development: annotated; confidence: high}
\end{GSitem}
\begin{GSitem}{1982-566.5}{E}{Highest weight and central charge of the wedge representation}
\GSfield{Kind; source}{theorem; \GSpage{P001042}}
\GSfield{The notes}{The notebook has \(e_0=h+n\), \(h=\tfrac12\mu(\mu-2\lambda+1)\), followed by a question mark.}
\GSfield{Statement}{For the stated tail ending at \(f_{-1}\), the normal-ordered wedge module has \(h=\tfrac12\mu(\mu-2\lambda-1)\) and \(c=-2(6\lambda^2+6\lambda+1)\); a vector of level \(n\) has eigenvalue \(h+n\). Other charge sectors replace \(\mu\) by \(\mu+\text{charge}\).}
\GSfield{Reconstruction}{The omitted central charge is restored from the explicitly cited paper and the tail convention fixes the sign.}
\GSfield{Fate}{These constants are printed in \citet[Section~3.3, Proposition~3.1, Russian p.~59]{feiginfuchs1982invariant}. The notebook sign is a faithful mathematical error, already marked doubtful by the recorder.}
\GSfield{Present status}{As of November 2026: the qualified statement is established; the notebook is an exposition, not a newly open question.}
\GSfield{Assessment}{literal statement: false as written; publication match: exact; verification: read the result; development: annotated; confidence: high}
\end{GSitem}
\begin{GSitem}{1982-566.6}{E}{When the wedge realization is a Verma module}
\GSfield{Kind; source}{claim; \GSpage{P001042}}
\GSfield{The notes}{\ru{Если неприводим, то явл. модулем Верма}.}
\GSfield{Statement}{If the charge-zero wedge module is irreducible, its nonzero vacuum generates it. The universal highest-weight map from \(V_{h,c}\) is then surjective, and both level-\(n\) spaces have dimension \(p(n)\); it is an isomorphism. Without irreducibility or cyclicity, equal characters alone do not make the canonical map an isomorphism.}
\GSfield{Reconstruction}{The universal-map and graded-dimension argument is written explicitly.}
\GSfield{Fate}{The equal-character comparison and canonical map are in \citet[Section~3.3, pp.~59--60]{feiginfuchs1982invariant}. The short proof supplied here spells out the notebook inference.}
\GSfield{Present status}{As of November 2026: the qualified statement is established; the notebook is an exposition, not a newly open question.}
\GSfield{Assessment}{literal statement: correct; publication match: exact; verification: read the result; development: annotated; confidence: high}
\end{GSitem}
\begin{GSitem}{1982-566.7}{E}{The reducibility locus of Virasoro Verma modules}
\GSfield{Kind; source}{question; \GSpage{P001042}}
\GSfield{The notes}{Kac is credited with the criterion \(\Phi_{r,s}(h,c)=0\) for positive integers \(r,s\), but both factors use a minus sign before \((c-13)\).}
\GSfield{Statement}{Choose \(t\in\mathbb C^\times\) with \(c=13-6(t+t^{-1})\). Then \(V_{h,c}\) is reducible exactly when \(h=h_{r,s}(t)=((rt-s)^2-(t-1)^2)/(4t)\) for some \(r,s\ge1\). The first singular-vector level is the least \(rs\) giving the weight. The determinant at level \(n\) has factors \((h-h_{r,s})^{p(n-rs)}\), up to a nonzero normalization.}
\GSfield{Reconstruction}{The parameterized form avoids the notebook polynomial sign error and makes the first-level qualification precise.}
\GSfield{Fate}{Kac announced the result in 1979; \citet[Theorem~3.1 and Section~4]{feiginfuchs1982invariant} states the determinant and supplies a proof before this talk. The polynomial printed there has plus signs before \((c-13)\).}
\GSfield{Present status}{As of November 2026: the qualified statement is established; the notebook is an exposition, not a newly open question.}
\GSfield{Assessment}{literal statement: false as written; publication match: exact; verification: read the result; development: annotated; confidence: high; problem \hyperref[prob:GP-1982-03]{GP-1982-03}}
\end{GSitem}
\begin{GSitem}{1982-566.8}{E}{The central-charge-zero specialization}
\GSfield{Kind; source}{claim; \GSpage{P001042}}
\GSfield{The notes}{\(c=0\Rightarrow h=(m^2-1)/24\).}
\GSfield{Statement}{At \(c=0\), take \(t=3/2\). Reducible highest weights are exactly \(((3r-2s)^2-1)/24\), \(r,s>0\), equivalently \((m^2-1)/24\) with \(m\in\mathbb Z\). The assertion describes reducibility, not every allowed highest weight at \(c=0\).}
\GSfield{Reconstruction}{The missing condition and integer range are supplied; the arithmetic equivalence is checked here.}
\GSfield{Fate}{This is direct substitution in the criterion of \citet[Theorem~3.1]{feiginfuchs1982invariant}. Every integer is \(3r-2s\) with positive \(r,s\) after adding a sufficiently large multiple of \((2,3)\).}
\GSfield{Present status}{As of November 2026: the qualified statement is established; the notebook is an exposition, not a newly open question.}
\GSfield{Assessment}{literal statement: correct with missing hypotheses; publication match: exact; verification: read the result; development: annotated; confidence: high}
\end{GSitem}
\begin{GSitem}{1982-566.9}{E}{Bosonic polynomial realization and the incorrect printed modes}
\GSfield{Kind; source}{construction; \GSpage{P001042}}
\GSfield{The notes}{The second construction sets \(F_k=\sum_i x_i\partial_{x_{i+k}}\) and builds quadratic \(e_i\), with coefficient \(\alpha=\mu-i(\lambda+\tfrac12)\).}
\GSfield{Statement}{For a corrected oscillator realization, let \([a_m,a_n]=m\delta_{m+n,0}\), \(a_n=n\partial_{x_n}\), \(a_{-n}=x_n\) for \(n>0\), and \(a_0=p\). Then \(L_n=\tfrac12\sum_k:a_{n-k}a_k:-Q(n+1)a_n\) has \(c=1-12Q^2\), \(h=\tfrac12p(p-2Q)\). Relabel \(e_i=L_{-i}\) to match the notes. The printed \(F_k\) annihilate the constant polynomial for every \(k\), so cannot satisfy a nonzero Heisenberg central bracket.}
\GSfield{Reconstruction}{The displayed normalized Heisenberg realization is a replacement, not a transcription correction. It is not attributed directly to the semi-infinite-wedge formula in Feigin--Fuchs.}
\GSfield{Fate}{The background-charge bosonic construction predates the talk through the older work cited by \citet[Theorem~2]{neretin1983highest}; that 1983 paper gives explicit corrected operators. Priority is verified indirectly, while the printed failure and the normalized replacement can be checked algebraically.}
\GSfield{Present status}{As of November 2026: the qualified statement is established; the notebook is an exposition, not a newly open question.}
\GSfield{Assessment}{literal statement: false as written; publication match: partial; verification: indirect; development: annotated; confidence: medium}
\end{GSitem}

\GStalk{567}{1982-09-27}{A. B. Zamolodchikov}{P001041-P001047}{From random paths to Polyakov sums over surfaces}
The entry begins on P001043 and ends above the 4 October heading on P001046. Classical area and conformal coordinates lead into three path regularizations: lattice walks, fixed-length random rods, and Gaussian increments. Their properly normalized diffusive limits motivate the intrinsic-metric integral for surfaces. Several formulas omit scaling factors, and the passage from fixed-metric Gaussian fields to a sum over geometries remains a formal construction.

\begin{GSitem}{1982-567.1}{E}{Area, conformal coordinates and the minimal-surface equation}
\GSfield{Kind; source}{construction; \GSpage{P001043}}
\GSfield{The notes}{\(G_{ab}=\partial_aX\cdot\partial_bX\), \(A=\int\sqrt{\det G}\); in a gauge \(G_{ab}=\rho\delta_{ab}\), the equations are \(\partial_a\partial_aX=0\), \(\partial_1X\cdot\partial_2X=0\), and \(|\partial_1X|^2=|\partial_2X|^2\).}
\GSfield{Statement}{For a smooth immersed Euclidean surface, local isothermal coordinates make the induced metric conformally flat. In these coordinates the mean-curvature-zero equation is harmonicity of \(X\), together with conformality. Branch points and global conformal moduli require additional treatment. A stationary minimal immersion need not be a global area minimizer.}
\GSfield{Reconstruction}{Locality, immersion and the distinction between stationary and globally minimizing surfaces are explicit.}
\GSfield{Fate}{The same classical formulation is recalled in \citet[Equations~(3)--(11)]{polyakov1981quantum}, published July 1981.}
\GSfield{Present status}{As of November 2026: the qualified statement is established; the notebook is an exposition, not a newly open question.}
\GSfield{Assessment}{literal statement: correct with missing hypotheses; publication match: exact; verification: read the result; development: annotated; confidence: high}
\end{GSitem}
\begin{GSitem}{1982-567.2}{E}{Propagation of conformality from the boundary}
\GSfield{Kind; source}{claim; \GSpage{P001043}}
\GSfield{The notes}{The two conformality conditions need only hold on the boundary.}
\GSfield{Statement}{Let \(X:\overline\Omega\to\mathbb R^d\) be \(C^1\) on a bounded smooth planar closure and harmonic in \(\Omega\). Its Hopf function \(H=X_z\cdot X_z\) is holomorphic. If both conformality equations hold on every boundary component, then \(H=0\) on the boundary, hence throughout \(\Omega\) by the maximum principle. Where \(dX\ne0\), \(X\) is a minimal immersion.}
\GSfield{Reconstruction}{We add the bounded smooth planar domain, regularity and harmonicity assumptions, and exclude branch points when claiming immersion.}
\GSfield{Fate}{This elementary implication belongs to the classical conformal formulation discussed in \citet[Equations~(5)--(11)]{polyakov1981quantum}. The proof here \ours supplies the boundary regularity and maximum-principle argument absent from the notes.}
\GSfield{Present status}{As of November 2026: the qualified statement is established; the notebook is an exposition, not a newly open question.}
\GSfield{Assessment}{literal statement: correct with missing hypotheses; publication match: partial; verification: read the result; development: annotated; confidence: high}
\end{GSitem}
\begin{GSitem}{1982-567.3}{F}{What measure defines the sum over surfaces?}
\GSfield{Kind; source}{question; \GSpage{P001043}}
\GSfield{The notes}{\(\sum_{\text{surfaces}}e^{-\mu_0^2A}\), followed by \ru{Мера --- ?}}
\GSfield{Statement}{The expression needs an ensemble, topology, embedding or mapping convention, parametrization quotient, boundary data, cutoff and normalization. The bare area weight alone does not determine a measure on a space of surfaces.}
\GSfield{Reconstruction}{No unspecified lattice-surface or Liouville existence problem is silently substituted.}
\GSfield{Fate}{\citet[Equation~(14)]{polyakov1981quantum} proposes an intrinsic-metric formalism. It does not identify a unique measure for every interpretation of the notebook sum.}
\GSfield{Present status}{As of November 2026: the general measure question is underdetermined; precise subclasses must be assessed separately.}
\GSfield{Assessment}{literal statement: not determinable; publication match: partial; verification: read the result; development: annotated; confidence: medium}
\end{GSitem}
\begin{GSitem}{1982-567.4}{E}{Reparametrization invariance of curve length}
\GSfield{Kind; source}{construction; \GSpage{P001043}}
\GSfield{The notes}{Length \(L=\int|\dot X|\,dt\) is invariant under changes of parameter; minimizing gives \(\ddot X=0\), with \(|\dot X|=1\).}
\GSfield{Statement}{For a regular curve with fixed endpoints, the Euler--Lagrange equation is \(d(\dot X/|\dot X|)/dt=0\). It becomes \(\ddot X=0\) after choosing arclength, with constant unit tangent. The interval length must equal the endpoint distance for a unit-speed minimizing segment.}
\GSfield{Reconstruction}{The order of variation and gauge choice, regularity and endpoint compatibility are supplied.}
\GSfield{Fate}{This is the elementary point-particle analogue used in \citet[Equation~(1)]{polyakov1981quantum}; the precise gauge-fixed equation follows by direct variation.}
\GSfield{Present status}{As of November 2026: the qualified statement is established; the notebook is an exposition, not a newly open question.}
\GSfield{Assessment}{literal statement: correct with missing hypotheses; publication match: partial; verification: read the result; development: annotated; confidence: high}
\end{GSitem}
\begin{GSitem}{1982-567.5}{E}{The lattice-walk geometric series and its convergence range}
\GSfield{Kind; source}{method; \GSpage{P001044}}
\GSfield{The notes}{Replace \(\mathbb R^d\) by a lattice, give every step length one, and sum \(e^{-mN}\); the notes say the sum converges and is explicitly computable.}
\GSfield{Statement}{For unrestricted nearest-neighbor walks, the Fourier multiplier for the unnormalized step count is \(A(p)=2\sum_{j=1}^d\cos p_j\). The Green series is \((1-e^{-m}A(p))^{-1}\) when the Neumann series converges, in particular for \(2de^{-m}<1\). At the critical weight convergence of a fixed-endpoint Green function depends on dimension; below critical mass the unqualified convergence assertion fails.}
\GSfield{Reconstruction}{The walk model, normalization and sufficient uniform convergence condition are specified; no necessity claim is made at a single endpoint.}
\GSfield{Fate}{The path-sum framework predates the seminar; compare \citet{polyakov1981quantum}. The finite-step multiplier, norm bound and geometric series are calculated directly here.}
\GSfield{Present status}{As of November 2026: the qualified statement is established; the notebook is an exposition, not a newly open question.}
\GSfield{Assessment}{literal statement: false as written; publication match: partial; verification: read the result; development: annotated; confidence: high}
\end{GSitem}
\begin{GSitem}{1982-567.6}{E}{Fixed-length random rods and their endpoint distribution}
\GSfield{Kind; source}{construction; \GSpage{P001044}}
\GSfield{The notes}{\(G_N(R)=\int\prod d n_i\,\delta(\sum n_i-R/\Delta)\); Fourier transformation gives \(C(p,\Delta)^N\), and a recurrence is requested.}
\GSfield{Statement}{With normalized rotation-invariant measure \(\sigma\) on \(S^{d-1}\), the endpoint density is \(p_N(R)=\int\delta(R-\Delta\sum n_i)\prod d\sigma(n_i)=\Delta^{-d}G_N(R)\). It satisfies \(p_{N+1}(R)=\int p_N(R-\Delta n)d\sigma(n)\) and \(\widehat p_N=C^N\), where \(C=\int e^{i\Delta p\cdot n}d\sigma(n)\). For \(d\ge2\), \(C=\Gamma(d/2)(2/(\Delta|p|))^{d/2-1}J_{d/2-1}(\Delta|p|)\), continuously extended at zero.}
\GSfield{Reconstruction}{Probability normalization, the missing delta-function Jacobian and the recurrence are supplied.}
\GSfield{Fate}{This is a direct finite-dimensional derivation of the notebook regularization. Its historical role is the path analogue of \citet[Equations~(1)--(2)]{polyakov1981quantum}.}
\GSfield{Present status}{As of November 2026: the qualified statement is established; the notebook is an exposition, not a newly open question.}
\GSfield{Assessment}{literal statement: false as written; publication match: partial; verification: read the result; development: annotated; confidence: high}
\end{GSitem}
\begin{GSitem}{1982-567.7}{E}{Diffusive scaling, mass tuning and the quadratic momentum term}
\GSfield{Kind; source}{method; \GSpage{P001044-P001045}}
\GSfield{The notes}{Equations (19)--(20) sum the step characteristic function and then use an exponent \(a+b\vec p-m\), before obtaining a propagator \(1/(p^2+m_R^2)\).}
\GSfield{Statement}{For isotropic unit steps, \(C(p,\Delta)=1-\Delta^2|p|^2/(2d)+O(\Delta^4|p|^4)\). Put \(a_\Delta=\Delta^2/(2d)\) and choose per-step weight \(q_\Delta=1-m_R^2a_\Delta+o(a_\Delta)\). Then \(a_\Delta\sum_{N\ge0}(q_\Delta C)^N\to(m_R^2+|p|^2)^{-1}\) for \(m_R>0\). Unnormalized step measures require additional entropy compensation in \(q_\Delta\).}
\GSfield{Reconstruction}{A definite scaling limit replaces the undefined constants a,b and clarifies mass renormalization.}
\GSfield{Fate}{The limit follows by expansion and a geometric series, within the path framework of \citet{polyakov1981quantum}. A linear isotropic momentum term cannot produce it.}
\GSfield{Present status}{As of November 2026: the qualified statement is established; the notebook is an exposition, not a newly open question.}
\GSfield{Assessment}{literal statement: false as written; publication match: partial; verification: read the result; development: annotated; confidence: high}
\end{GSitem}
\begin{GSitem}{1982-567.8}{E}{Gaussian increments, Wiener measure and proper time}
\GSfield{Kind; source}{method; \GSpage{P001045}}
\GSfield{The notes}{The rods become springs, their lengths are called Gaussian, and a formal integral \(e^{-\frac12\int\dot X^2}\) over paths is called Wiener measure.}
\GSfield{Statement}{Gaussian displacement vectors with covariance \(\Delta I\), each normalized by \((2\pi\Delta)^{-d/2}\), converge to Brownian motion with generator \(\tfrac12\Delta_x\). The endpoint heat kernel has Fourier multiplier \(e^{-\tau|p|^2/2}\). Integrating against \(e^{-a\tau}\,d\tau\), \(a>0\), yields \((a+|p|^2/2)^{-1}\). Brownian paths have infinite ordinary arclength almost surely; the kinetic action is symbolic rather than a finite pathwise function.}
\GSfield{Reconstruction}{Displacements replace Gaussian lengths; normalization, generator and the symbolic action convention are explicit.}
\GSfield{Fate}{The Wiener-functional tradition precedes the seminar; \citet{kac1949distributions} is verified bibliographically only. The Gaussian convolution and Laplace transform above are checked directly. Matching the previous convention requires rescaling time and mass.}
\GSfield{Present status}{As of November 2026: the qualified statement is established; the notebook is an exposition, not a newly open question.}
\GSfield{Assessment}{literal statement: false as written; publication match: partial; verification: publisher metadata checked; development: annotated; confidence: high}
\end{GSitem}
\begin{GSitem}{1982-567.9}{E}{The Euclidean massive Green function}
\GSfield{Kind; source}{claim; \GSpage{P001045}}
\GSfield{The notes}{\(\widehat G(p)=z/(p^2+m^2)\), whose inverse is called a Hankel function.}
\GSfield{Statement}{With Fourier inverse \((2\pi)^{-d}\int e^{ip\cdot R}\,dp\), for \(m>0\), \(R\ne0\), the kernel is \(G(R)=(2\pi)^{-d/2}(m/|R|)^{d/2-1}K_{d/2-1}(m|R|)\). A Hankel-function expression requires analytic continuation to an imaginary argument or a Lorentzian prescription; the decaying Euclidean expression uses modified Bessel \(K\).}
\GSfield{Reconstruction}{The transform convention, massive regime and relation to the Hankel terminology are stated.}
\GSfield{Fate}{This follows by integrating the Gaussian heat kernel in proper time, the framework associated with \citet{kac1949distributions}. The normalization and function type are an editorial calculation.}
\GSfield{Present status}{As of November 2026: the qualified statement is established; the notebook is an exposition, not a newly open question.}
\GSfield{Assessment}{literal statement: correct with missing hypotheses; publication match: partial; verification: publisher metadata checked; development: annotated; confidence: high}
\end{GSitem}
\begin{GSitem}{1982-567.10}{E}{A spring mesh gives a Gaussian field at fixed worldsheet geometry}
\GSfield{Kind; source}{construction; \GSpage{P001045}}
\GSfield{The notes}{A square spring network gives \(\int DX\exp[-\tfrac12\int(\partial X)^2]\); a regular grid is said to be insufficient for general surfaces.}
\GSfield{Statement}{For a fixed two-dimensional mesh and boundary conditions, a quadratic spring energy defines a finite-dimensional Gaussian measure after removing or fixing zero modes. Its continuum limit is a Gaussian free field at the chosen conformal structure. Varying embedded vertices on a fixed mesh does not by itself specify a sum over intrinsic metrics or topologies.}
\GSfield{Reconstruction}{Boundary/zero-mode conditions and fixed intrinsic geometry are added.}
\GSfield{Fate}{The distinction motivates \citet[Equations~(6), (14)]{polyakov1981quantum}. The notebook drawings do not define a unique discretization-independent surface measure.}
\GSfield{Present status}{As of November 2026: the qualified statement is established; the notebook is an exposition, not a newly open question.}
\GSfield{Assessment}{literal statement: correct with missing hypotheses; publication match: partial; verification: read the result; development: annotated; confidence: high}
\end{GSitem}
\begin{GSitem}{1982-567.11}{E}{Why the bosonic critical dimension is 26}
\GSfield{Kind; source}{claim; \GSpage{P001045}}
\GSfield{The notes}{The fixed-metric quadratic action gives the \ru{прав. ответ} in target dimension 26.}
\GSfield{Statement}{In bosonic Polyakov gauge fixing, \(D\) free scalar matter fields contribute central charge \(D\), and the reparametrization ghosts contribute \(-26\). The local Weyl anomaly cancels at \(D=26\). This does not prove convergence over topologies, absence of tachyons, or existence of an unrestricted surface probability measure.}
\GSfield{Reconstruction}{The anomaly statement replaces the unqualified phrase correct answer.}
\GSfield{Fate}{The coefficient is calculated in \citet[Equations~(23)--(25)]{polyakov1981quantum}, before the talk.}
\GSfield{Present status}{As of November 2026: the qualified statement is established; the notebook is an exposition, not a newly open question.}
\GSfield{Assessment}{literal statement: correct with missing hypotheses; publication match: exact; verification: read the result; development: annotated; confidence: high}
\end{GSitem}
\begin{GSitem}{1982-567.12}{E}{The independent-metric action and cosmological term}
\GSfield{Kind; source}{construction; \GSpage{P001045-P001046}}
\GSfield{The notes}{The final integral introduces \(g_{ab}\) and \(+\lambda\int\sqrt g\) in the exponent, then says minimizing recovers the earlier equations.}
\GSfield{Statement}{The classical action \(S_P=\tfrac12\int\sqrt g\,g^{ab}\partial_aX\cdot\partial_bX\), varied in \(X\) and \(g\), imposes harmonicity and vanishing stress tensor. For a nondegenerate induced metric this makes \(g\) conformal to it and recovers area. A nonzero classical cosmological term changes the trace equation: in two dimensions the kinetic stress tensor is traceless, so that term must vanish. In the formal quantum measure it is instead a cutoff-dependent counterterm.}
\GSfield{Reconstruction}{The trace check explicitly separates the classical and quantum roles of lambda.}
\GSfield{Fate}{\citet[Equations~(6)--(7), (14), p.~209]{polyakov1981quantum} separates these roles. Gauge volume, moduli and boundary conditions also need specification.}
\GSfield{Present status}{As of November 2026: the qualified statement is established; the notebook is an exposition, not a newly open question.}
\GSfield{Assessment}{literal statement: false as written; publication match: exact; verification: read the result; development: annotated; confidence: high}
\end{GSitem}

\GStalk{568}{1982-10-04}{A. A. Kirillov}{P001045-P001047}{Springer resolution and limiting orbital integrals}
The entry occupies the middle and lower part of P001046 and the first lines of P001047. Kirillov reports Ginzburg\textquotesingle s work after explaining adjoint orbits and Springer resolution. The 1983 paper identifies the nilpotent-orbit limit formula precisely for complex semisimple algebras regarded as real. Gelfand then asks how a singular hyperbola integral recovers the value of a test function at the crossing.

\begin{GSitem}{1982-568.1}{E}{Parameters of regular adjoint orbits}
\GSfield{Kind; source}{definition; \GSpage{P001046}}
\GSfield{The notes}{\(\Omega_c\) is parametrized by \(\ell=\operatorname{rank}G\) invariants, illustrated by \(SL(2,\mathbb R)\) hyperboloids and traces of powers.}
\GSfield{Statement}{For a complex semisimple algebra, invariant polynomials parametrize the adjoint quotient, and regular semisimple fibers are single complex adjoint orbits. For real groups, invariant values alone can combine several real orbits. In \(\mathfrak{sl}_2(\mathbb R)\), the negative and positive values of the invariant quadratic form give different hyperboloid types; the two sheets need not be one orbit.}
\GSfield{Reconstruction}{The field, regular semisimple restriction and real-orbit distinction are supplied.}
\GSfield{Fate}{The complex-as-real scope used by \citet[Sections~1--2]{ginzburg1983integrals} is narrower than the notebook\textquotesingle s initial arbitrary Lie group. The rank-one real qualification follows directly from the adjoint quadratic form.}
\GSfield{Present status}{As of November 2026: the qualified statement is established; the notebook is an exposition, not a newly open question.}
\GSfield{Assessment}{literal statement: correct with missing hypotheses; publication match: partial; verification: read the result; development: annotated; confidence: high}
\end{GSitem}
\begin{GSitem}{1982-568.2}{E}{Springer incidence resolution}
\GSfield{Kind; source}{construction; \GSpage{P001046}}
\GSfield{The notes}{For the unipotent variety \(N\), let \(\mathcal B\) be the Borel variety and \(\mathcal M=\{(n,B):n\in N,\ n\in B\}\); projection to \(N\) is called a desingularization.}
\GSfield{Statement}{For a connected complex semisimple algebraic group in characteristic zero, \(\widetilde N=G\times^B U_B\to N\) is proper and birational, with smooth source; its fiber over \(n\) is the Borel subgroups containing \(n\). In Lie-algebra form \(G\times^B\mathfrak n_B\simeq T^*(G/B)\) resolves the nilpotent cone. Properness and birationality, not just equal dimension and generic bijectivity, belong to the resolution statement.}
\GSfield{Reconstruction}{The statement specifies a complex algebraic group and properness. It also distinguishes the group construction from the Lie algebra construction.}
\GSfield{Fate}{This is the preexisting Springer framework cited by \citet{springer1976trigonometric} and used explicitly in \citet[Section~9]{ginzburg1983integrals}. The exact earlier incidence-resolution occurrence is 1971-231.2; 1971-231.5 concerns the related Springer-fiber dimension.}
\GSfield{Present status}{As of November 2026: the qualified statement is established; the notebook is an exposition, not a newly open question.}
\GSfield{Assessment}{literal statement: correct with missing hypotheses; publication match: exact; verification: indirect; development: annotated; confidence: high}
\end{GSitem}
\begin{GSitem}{1982-568.3}{E}{The real cylinder-to-cone illustration}
\GSfield{Kind; source}{construction; \GSpage{P001046}}
\GSfield{The notes}{The drawing maps a cylinder to a double cone by \((x,y,z)\mapsto(xz,yz,z)\), collapsing its middle circle; it is identified with the \(SL(2,\mathbb R)\) example.}
\GSfield{Statement}{On \(x^2+y^2=1\), this map lands in \(X^2+Y^2=Z^2\), is one-to-one away from the cone vertex, and collapses exactly \(S^1\times\{0\}\). It models the topology of the real rank-one Springer map, whose exceptional fiber is \(\mathbb{RP}^1\simeq S^1\). The complex exceptional fiber is instead \(\mathbb{CP}^1\); the real drawing is not a description of all complex algebraic data.}
\GSfield{Reconstruction}{The source cylinder equation and the real versus complex interpretation are supplied.}
\GSfield{Fate}{The topology is checked directly from the displayed map and the flag fiber, within the Springer framework used by \citet[Sections~3,9]{ginzburg1983integrals}.}
\GSfield{Present status}{As of November 2026: the qualified statement is established; the notebook is an exposition, not a newly open question.}
\GSfield{Assessment}{literal statement: correct with missing hypotheses; publication match: partial; verification: read the result; development: annotated; confidence: high}
\end{GSitem}
\begin{GSitem}{1982-568.4}{B}{Nilpotent orbital integrals from regular semisimple ones}
\GSfield{Kind; source}{theorem; \GSpage{P001046}}
\GSfield{The notes}{\(\psi(c)=\int_{\Omega_c}f\), with coherently chosen measures; \(\int_{\Omega_0}f=(P\psi)(0)\), \(P\) a differential operator.}
\GSfield{Statement}{For a complex semisimple Lie algebra regarded as a real Lie algebra and a Schwartz test function \(f\), each nilpotent orbital integral is a constant multiple of a constant-coefficient differential operator applied at zero to the normalized Harish--Chandra orbital transform on a Cartan subalgebra. The operator polynomial is obtained from the corresponding Springer homology representation. Raw integrals in arbitrary quotient coordinates do not automatically satisfy the abbreviated identity.}
\GSfield{Reconstruction}{The group class, Schwartz decay, Cartan normalization, measure constants and orbit-dependent operator are restored.}
\GSfield{Fate}{\citet[Formula~(1), Sections~1--7]{ginzburg1983integrals} proves the Barbasch--Vogan conjectural formula in this scope. Presented 19 October 1982, received 23 November 1982 and published in 1983, it matches Kirillov\textquotesingle s report of another mathematician\textquotesingle s work; class B records that authorship.}
\GSfield{Present status}{As of November 2026: the identified complex-group limit formula is established; no claim about all real reductive groups is inferred.}
\GSfield{Assessment}{literal statement: correct with missing hypotheses; publication match: exact; verification: read the result; development: annotated; confidence: high}
\end{GSitem}
\begin{GSitem}{1982-568.5}{E}{Recovering the crossing value from hyperbola integrals}
\GSfield{Kind; source}{question; \GSpage{P001047}; I. M. Gelfand}
\GSfield{The notes}{Gelfand asks to recover \(f(0)\) from integrals over \(xy=c\), and refers to an expansion of \(\delta(xy-c)\) in the first volume of Generalized Functions.}
\GSfield{Statement}{For \(f\in C_c^\infty(\mathbb R^2)\), set \(I_f(c)=\int f(x,y)\delta(xy-c)\,dx\,dy\), \(c\ne0\). Then \(I_f(c)=-2f(0,0)\log|c|+O(1)\), so \(f(0,0)=\lim_{c\to0}I_f(c)/(-2\log|c|)\). This uses the delta/coarea measure, not ordinary arclength. Restricting to a single positive branch halves the leading coefficient.}
\GSfield{Reconstruction}{Compact support, delta measure and the exact leading coefficient are supplied; this is an elementary reconstruction of the recorded question.}
\GSfield{Fate}{The notes themselves identify the preexisting distribution-theoretic setting of \citet{gelfandshilov1964generalized}; only its bibliographic record was accessible. The stated leading term is verified here \ours by writing \(I_f(c)=\int f(x,c/x)\,dx/|x|\), splitting near zero, and estimating the smooth remainder.}
\GSfield{Present status}{As of November 2026: the qualified statement is established; the notebook is an exposition, not a newly open question.}
\GSfield{Assessment}{literal statement: correct with missing hypotheses; publication match: partial; verification: publisher metadata checked; development: annotated; confidence: high; problem \hyperref[prob:GP-1982-04]{GP-1982-04}}
\end{GSitem}
\begin{GSitem}{1982-568.6}{E}{Why the singular limit contains logarithms}
\GSfield{Kind; source}{method; \GSpage{P001047}}
\GSfield{The notes}{The text speaks of expanding \(\delta(xy-c)\) in powers of \(c\), then taking the hyperbola limit to obtain integrals over straight lines.}
\GSfield{Statement}{At the critical level \(c=0\), the pullback distribution cannot be handled by an ordinary Taylor expansion: \(d(xy)=0\) at the crossing. Asymptotics involve logarithmic terms and finite-part distributions on the axes. Writing a finite coefficient through axis integrals requires a finite-part prescription; changing that prescription is compensated by a multiple of the point delta distribution. The coefficient of the fixed distribution family itself is not arbitrary.}
\GSfield{Reconstruction}{The need for logarithmic/finite-part terms is an explicit mathematical correction, without claiming an unread full asymptotic formula.}
\GSfield{Fate}{The source allusion is to \citet{gelfandshilov1964generalized}; the leading logarithm above directly rules out an unqualified power-series or finite unrenormalized limit. The exact book passage was not located.}
\GSfield{Present status}{As of November 2026: the qualified statement is established; the notebook is an exposition, not a newly open question.}
\GSfield{Assessment}{literal statement: false as written; publication match: partial; verification: publisher metadata checked; development: annotated; confidence: high}
\end{GSitem}

\GStalk{569}{1982-10-04}{A. B. Zamolodchikov}{P001045-P001051}{From path entropy to the Polyakov metric integral and Liouville theory}
The actual entry starts halfway through P001047 and ends on P001049. Zamolodchikov reviews the proper-time representation of a particle path sum, then the fixed-topology Polyakov metric integral, quotienting by reparametrizations, functional metrics, conformal gauge and the Liouville action. The uncertainty about summing over handles is separate from the fixed-topology calculation. Several expressions are formal prescriptions: neither an infinite-dimensional Riemannian metric nor division by a gauge-group volume by itself constructs a measure.

\begin{GSitem}{1982-569.1}{E}{Proper time and the particle Green function}
\GSfield{Kind; source}{method; \GSpage{P001047}}
\GSfield{The notes}{Equations (2)--(2\textquotesingle) replace the length-weighted sum by an integral over \(\tau>0\) and Brownian paths with kinetic term \(\frac12\int\dot x^2dt\), giving \(\tau^{-d/2}\exp[-m^2\tau-R^2/(2\tau)]\).}
\GSfield{Statement}{For Brownian motion with generator \(\frac12\Delta\) on \(\mathbb R^d\), the heat kernel is \(p_\tau(R)=(2\pi\tau)^{-d/2}e^{-|R|^2/(2\tau)}\). For \(m^2>0\), its Laplace transform is the resolvent kernel of \(m^2-\frac12\Delta\), with Fourier multiplier \((m^2+|p|^2/2)^{-1}\). The length-weighted lattice sum reaches this expression only after a specified diffusive limit and mass renormalization.}
\GSfield{Reconstruction}{The normalized heat kernel, generator and positive-mass hypothesis are supplied. The formula does not assign ordinary arc length to Brownian paths.}
\GSfield{Fate}{This is the particle comparison underlying \citet[Eqs.~(1)--(5)]{polyakov1981quantum}; it repeats the preceding lecture in a consistent proper-time normalization.}
\GSfield{Present status}{As of November 2026: the qualified statement is established; the notebook is an exposition, not a newly open question.}
\GSfield{Assessment}{literal statement: correct with missing hypotheses; publication match: partial; verification: read the result; development: annotated; confidence: high}
\end{GSitem}
\begin{GSitem}{1982-569.2}{E}{Entropy subtraction and the renormalized mass}
\GSfield{Kind; source}{method; \GSpage{P001047}}
\GSfield{The notes}{The number of length-\(L\) paths is written \(f(L)\sim L^\beta e^{\varkappa L}\), with \(\varkappa\) cutoff-dependent and \(\beta\) cutoff-independent; the notes conclude \(m^2=m_0-\varkappa\).}
\GSfield{Statement}{For a specified lattice-walk ensemble the exponential growth constant shifts the bare length penalty to \(m_0-\varkappa\). A Brownian continuum limit additionally rescales the number of steps as inverse step length squared and fixes the diffusion coefficient. Consequently the numerical identity \(m^2=m_0-\varkappa\) is schematic until the units of \(L\) and \(\tau\) and the cutoff scaling are specified. Polynomial exponents depend on dimension, endpoints and the walk ensemble.}
\GSfield{Reconstruction}{The identification is kept as a scaling prescription. No self-avoiding-walk or arbitrary-regulator universality claim is supplied.}
\GSfield{Fate}{\citet[Eqs.~(1)--(5)]{polyakov1981quantum} uses the same random-walk analogy. The entropy subtraction is standard; the notes do not state a universality theorem encompassing different path ensembles.}
\GSfield{Present status}{As of November 2026: the qualified statement is established; the notebook is an exposition, not a newly open question.}
\GSfield{Assessment}{literal statement: correct with missing hypotheses; publication match: partial; verification: read the result; development: annotated; confidence: high}
\end{GSitem}
\begin{GSitem}{1982-569.3}{F}{The remark about singular infinite-dimensional measures}
\GSfield{Kind; source}{remark; \GSpage{P001047}}
\GSfield{The notes}{A margin says:\par\ru{В \(\infty\)-мерном пр-ве почти \(\forall\) две меры взаимно сингулярны}.}
\GSfield{Statement}{There is no specified class of measures or measure on that class, so the phrase almost every pair has no definite quantifier. Infinite-dimensional Gaussian measures can be equivalent or mutually singular depending on their means and covariance operators; dimension alone does not decide the comparison.}
\GSfield{Reconstruction}{No particular Gaussian-equivalence criterion is attributed to the speaker.}
\GSfield{Fate}{The statement cannot be assigned a precise probability theorem from the notebook. It serves as a warning that changing a continuum regularization can change the measure class.}
\GSfield{Present status}{As of November 2026: an underdetermined heuristic, not an identified open problem.}
\GSfield{Assessment}{literal statement: not determinable; publication match: none found; verification: unresolved lead; development: annotated; confidence: low}
\end{GSitem}
\begin{GSitem}{1982-569.4}{E}{The fixed-topology Polyakov metric integral}
\GSfield{Kind; source}{construction; \GSpage{P001048}}
\GSfield{The notes}{Equation (5) uses maps \(x(\xi)\), an independent positive metric \(g\), a Gaussian kinetic action and the area term \(\mu_0^2\int\sqrt g\); the topology is explicitly fixed.}
\GSfield{Statement}{For a fixed oriented surface \(\Sigma\), the formal bosonic-string action is \(S[x,g]=\frac12\int_\Sigma |dx|_g^2\,dA_g+\mu_0^2\int_\Sigma dA_g\). Maps and metrics are integrated modulo the chosen diffeomorphism gauge; conformal anomaly and metric moduli remain part of the definition. Equality with a combinatorial sum over embedded surfaces requires a regulator and a scaling-limit theorem, neither specified here.}
\GSfield{Reconstruction}{Fixed topology, positivity and formal status are made explicit.}
\GSfield{Fate}{\citet[Eqs.~(3)--(10)]{polyakov1981quantum}, published 23 July 1981, gives this metric-integral formulation. The notebook reports an existing formal construction, not an established equality of all possible random-surface measures.}
\GSfield{Present status}{As of November 2026: the qualified statement is established; the notebook is an exposition, not a newly open question.}
\GSfield{Assessment}{literal statement: correct with missing hypotheses; publication match: exact; verification: read the result; development: annotated; confidence: high}
\end{GSitem}
\begin{GSitem}{1982-569.5}{F}{Could surfaces with many handles dominate?}
\GSfield{Kind; source}{question; \GSpage{P001048}}
\GSfield{The notes}{The margin asks about the relative contributions of handles and holes, suggesting that \ru{Пенистые пов.} with many handles may dominate.}
\GSfield{Statement}{An all-topology sum requires a class of surfaces, boundary data, a regulator, a weight for each genus, and a summation or convergence criterion. These are not fixed, so the statement does not select a definite theorem about large-genus asymptotics or a nonperturbative string partition function.}
\GSfield{Reconstruction}{No genus fugacity or limiting ensemble is invented.}
\GSfield{Fate}{Zamolodchikov\textquotesingle s \citet{zamolodchikov1982entropy} is a genuine 1982 publication on random-surface entropy; its existence does not identify a solution of this unspecified all-genus question. The modern fixed-sphere construction \citet{davidkupiainenrhodesvargas2016liouville} likewise has a different scope.}
\GSfield{Present status}{As of November 2026: the historical question remains underdetermined; no open-problem identifier is assigned.}
\GSfield{Assessment}{literal statement: not determinable; publication match: thematic only; verification: publisher metadata checked; development: annotated; confidence: medium}
\end{GSitem}
\begin{GSitem}{1982-569.6}{E}{Functional metrics and the trace direction}
\GSfield{Kind; source}{definition; \GSpage{P001048}}
\GSfield{The notes}{Equation (6) gives the two formal quadratic expressions \[\|\delta x\|^2=\int\sqrt g\,|\delta x|^2,\] \[\|\delta g\|^2=\int\sqrt g\,(g^{aa\prime}g^{bb\prime}+c\,g^{ab}g^{a\prime b\prime})\delta g_{ab}\delta g_{a\prime b\prime}.\]}
\GSfield{Statement}{For a symmetric variation \(h\) of a two-dimensional positive metric, write \(h=h_0+\frac12(\operatorname{tr}_g h)g\). The second integrand equals \(|h_0|_g^2+(c+\frac12)(\operatorname{tr}_g h)^2\). It defines a positive inner product on all metric variations exactly when \(c>-\frac12\); equality is degenerate and smaller values give an indefinite trace direction.}
\GSfield{Reconstruction}{The omitted condition on c is supplied. A positive field-space norm still does not create a countably additive flat volume measure.}
\GSfield{Fate}{The invariant functional-metric prescription is in \citet[Eqs.~(6)--(10)]{polyakov1981quantum}. The positivity condition follows by the displayed trace decomposition, an elementary check \ours.}
\GSfield{Present status}{As of November 2026: the qualified statement is established; the notebook is an exposition, not a newly open question.}
\GSfield{Assessment}{literal statement: correct with missing hypotheses; publication match: partial; verification: read the result; development: annotated; confidence: high}
\end{GSitem}
\begin{GSitem}{1982-569.7}{E}{Diffeomorphisms and the covector connection sign}
\GSfield{Kind; source}{construction; \GSpage{P001048}}
\GSfield{The notes}{Equation (7) gives \(\delta g_{ab}=\nabla_a\varepsilon_b+\nabla_b\varepsilon_a\), but defines \(\nabla_a\varepsilon_b=\partial_a\varepsilon_b+\Gamma^c_{ab}\varepsilon_c\).}
\GSfield{Statement}{For the Levi--Civita Christoffel symbols, a covector satisfies \(\nabla_a\varepsilon_b=\partial_a\varepsilon_b-\Gamma^c_{ab}\varepsilon_c\). With \(\varepsilon_b=g_{bc}\varepsilon^c\), the metric Lie derivative is \((\mathcal L_\varepsilon g)_{ab}=\nabla_a\varepsilon_b+\nabla_b\varepsilon_a\). The plus sign belongs to the connection on contravariant components.}
\GSfield{Reconstruction}{The usual Levi--Civita convention is fixed; the inverse coordinate-change convention could reverse the whole Lie derivative, but not only its Christoffel term.}
\GSfield{Fate}{This is the gauge variation used in \citet[Eqs.~(7)--(10)]{polyakov1981quantum}. The discrepancy is a faithful notebook sign error, confirmed on the scan.}
\GSfield{Present status}{As of November 2026: the qualified statement is established; the notebook is an exposition, not a newly open question.}
\GSfield{Assessment}{literal statement: false as written; publication match: partial; verification: read the result; development: annotated; confidence: high}
\end{GSitem}
\begin{GSitem}{1982-569.8}{B}{Does dividing by gauge volume define the measure?}
\GSfield{Kind; source}{claim; \GSpage{P001048}}
\GSfield{The notes}{After \(d\mu(x,g)=\mathcal Dx\,\mathcal Dg/\mathcal D\omega\), the notes say \ru{Теперь интеграл (5) определён хорошо}.}
\GSfield{Statement}{Functional metrics and a formal quotient are not by themselves a measure construction. One must control a regulator, zero modes, gauge stabilizers, renormalization and metric moduli. A precise later result is the construction of spacelike Liouville correlation functions on the sphere for \(0<\gamma<2\), \(\mu>0\), \(Q=2/\gamma+\gamma/2\), and insertions satisfying \(\alpha_i<Q\), \(\sum_i\alpha_i>2Q\). This is a specified model, not the unrestricted quotient integral.}
\GSfield{Reconstruction}{The later theorem is stated separately from the 1982 formal claim.}
\GSfield{Fate}{\citet[Theorem~3.2 and Definition~3.6]{davidkupiainenrhodesvargas2016liouville}, first circulated 27 October 2014 and published in 2016, gives the indicated rigorous construction. The bare uninserted sphere integral has a divergent zero-mode direction; the notebook\textquotesingle s assertion is not justified by (6)--(9).}
\GSfield{Present status}{As of November 2026: the stated subcritical sphere theory is constructed. The notes do not define a universal measure problem beyond that scope.}
\GSfield{Assessment}{literal statement: false as written; publication match: partial; verification: read the result; development: annotated; confidence: high}
\end{GSitem}
\begin{GSitem}{1982-569.9}{E}{Conformal gauge, residual automorphisms and moduli}
\GSfield{Kind; source}{method; \GSpage{P001048}}
\GSfield{The notes}{The representative is \(g_{ab}=\rho\,\delta_{ab}\), followed by \ru{Остаётся конформная группа (конечномерная)}.}
\GSfield{Statement}{Every smooth two-dimensional metric is locally conformally flat. Globally write \(g=e^\varphi\widehat g(m)\) using a reference metric in its conformal class and its moduli \(m\). For a fixed compact Riemann surface the global conformal automorphism group is finite-dimensional; germs of local conformal transformations are not. A single globally Euclidean chart does not cover an arbitrary closed surface.}
\GSfield{Reconstruction}{The distinction between local coordinates, global automorphisms and moduli is added.}
\GSfield{Fate}{\citet[Eqs.~(11)--(12) and discussion of moduli]{polyakov1981quantum} uses conformal gauge with global qualifications. The later sphere construction \citet[Section~3.2]{davidkupiainenrhodesvargas2016liouville} treats the residual Möbius action explicitly.}
\GSfield{Present status}{As of November 2026: the qualified statement is established; the notebook is an exposition, not a newly open question.}
\GSfield{Assessment}{literal statement: correct with missing hypotheses; publication match: partial; verification: read the result; development: annotated; confidence: high}
\end{GSitem}
\begin{GSitem}{1982-569.10}{E}{Matter and ghost determinants are different operators}
\GSfield{Kind; source}{claim; \GSpage{P001048}}
\GSfield{The notes}{After Gaussian integration the notes say the determinants are identical and write \([\det\Delta_g]^{(26-d)/2}\) in equation (11).}
\GSfield{Statement}{Writing \(\Delta_g\) for the nonnegative scalar Laplacian, integrating \(d\) scalar matter fields gives a factor \((\det{}^{\prime}\Delta_g)^{-d/2}\), with the translation zero modes treated separately. Gauge fixing contributes the determinant of the vector-to-traceless-tensor operator \(P\varepsilon=\nabla\varepsilon+(\nabla\varepsilon)^t-g\,\operatorname{div}\varepsilon\), together with zero-mode and moduli factors. The ghost conformal anomaly has coefficient \(-26\); it is not an exact identity of its determinant with the thirteenth power of the scalar Laplacian determinant.}
\GSfield{Reconstruction}{Primed determinants and the actual ghost operator replace the literal scalar-determinant identification.}
\GSfield{Fate}{\citet[Eqs.~(13)--(24)]{polyakov1981quantum} distinguishes the operators before calculating their conformal variations. Equation (11) captures an anomaly coefficient while suppressing determinant-line, zero-mode and global factors.}
\GSfield{Present status}{As of November 2026: the qualified statement is established; the notebook is an exposition, not a newly open question.}
\GSfield{Assessment}{literal statement: false as written; publication match: partial; verification: read the result; development: annotated; confidence: high}
\end{GSitem}
\begin{GSitem}{1982-569.11}{E}{The induced Liouville action and its scope}
\GSfield{Kind; source}{construction; \GSpage{P001049}}
\GSfield{The notes}{Equation (12) sets \(\rho=e^\varphi\) and display the exponent \((d-26)S_L/(48\pi)\), with \(S_L=\int[\frac12(\partial\varphi)^2+m^2e^\varphi]\,d^2\xi\).}
\GSfield{Statement}{The Weyl-anomaly calculation induces a local kinetic term for the conformal factor and a renormalized area term. On a curved reference surface it also includes a background-curvature coupling. The rigorous subcritical spacelike normalization uses a renormalized exponential, \(Q=2/\gamma+\gamma/2\), and central charge \(c_L=1+6Q^2>25\). Cancellation with matter requires \(c_{\mathrm m}+c_L-26=0\); this regime corresponds to \(c_{\mathrm m}<1\), not every integer target dimension \(d<26\).}
\GSfield{Reconstruction}{The curved-background term, quantum renormalization and restricted matter-charge range are modern qualifications.}
\GSfield{Fate}{The formal induced action was already in \citet[Eqs.~(25)--(31)]{polyakov1981quantum}. The quantum parameter shift and rigorous Weyl anomaly are verified in \citet[Introduction and Theorem~3.11]{davidkupiainenrhodesvargas2016liouville}. These later results clarify a specified branch without turning the notebook\textquotesingle s formal expression into a theorem for all dimensions.}
\GSfield{Present status}{As of November 2026: the formal anomaly mechanism is established, and the stated spacelike sphere realization is rigorous; no blanket all-dimension claim is made.}
\GSfield{Assessment}{literal statement: correct with missing hypotheses; publication match: exact; verification: read the result; development: annotated; confidence: high}
\end{GSitem}
\begin{GSitem}{1982-569.12}{E}{The field-dependent norm of the conformal factor}
\GSfield{Kind; source}{definition; \GSpage{P001049}}
\GSfield{The notes}{Equation (13) gives \(\|\delta\varphi\|^2=\int e^\varphi(\delta\varphi)^2\,d^2\xi\) as the metric behind \(d\mu(\varphi)\).}
\GSfield{Statement}{For \(g=e^\varphi\widehat g\), the induced trace-variation norm is proportional to \(\int_\Sigma e^\varphi(\delta\varphi)^2dA_{\widehat g}\), with a positive proportionality factor depending on the trace coefficient in (6). It is a field-dependent formal \(L^2\) metric. Its associated symbolic volume must not be replaced by flat \(\mathcal D\varphi\) without calculating a regulator-dependent Jacobian.}
\GSfield{Reconstruction}{The proportionality constant and background volume form are restored.}
\GSfield{Fate}{\citet[Eqs.~(29)--(31)]{polyakov1981quantum} explicitly retains this nontrivial metric in the Liouville integral. Modern probabilistic constructions begin from a Gaussian free field and renormalized multiplicative chaos instead of a literal Riemannian volume in infinitely many dimensions.}
\GSfield{Present status}{As of November 2026: the qualified statement is established; the notebook is an exposition, not a newly open question.}
\GSfield{Assessment}{literal statement: correct with missing hypotheses; publication match: exact; verification: read the result; development: annotated; confidence: high}
\end{GSitem}

\GStalk{570}{1982-10-11}{V. L. Ginzburg}{P001049-P001052}{Coisotropic reduction, canonical projectors and double fibrations}
The actual entry occupies P001050 and the upper part of P001051, before Fröhlich. The catalogue attribution to V. L. Ginzburg is retained; the notebook gives only V. Ginzburg and explicitly presents the material after Guillemin--Sternberg. The full speaker identity has not been independently resolved from another dated seminar source. V. A. Ginzburg’s 1983 paper is used for mathematical comparison, without identifying its author as this speaker. The talk relates a regular coisotropic quotient to a symmetric idempotent canonical relation and an operator algebra, then Shmelyov gives the conormal embedding condition for a double fibration.

\begin{GSitem}{1982-570.1}{E}{The characteristic foliation and regular symplectic reduction}
\GSfield{Kind; source}{construction; \GSpage{P001050}}
\GSfield{The notes}{\(X\) is complex, \(\Sigma\subset T^*X\) is conic and coisotropic, and \(\Sigma\to X\) is a fibration. The kernel of the restricted symplectic form is said to be integrable; its quotient \(M\) is assumed to be a fibration.}
\GSfield{Statement}{For a smooth coisotropic \(\Sigma\) in a symplectic manifold \((P,\omega)\), \(\ker(\omega|_\Sigma)=(T\Sigma)^\omega\) is integrable. If its leaf space \(M\) is smooth and the quotient \(q:\Sigma\to M\) is a surjective submersion with connected fibers, there is a unique symplectic form \(\omega_M\) satisfying \(q^*\omega_M=\omega|_\Sigma\). Closed leaves alone do not guarantee a smooth Hausdorff quotient.}
\GSfield{Reconstruction}{Smoothness and the quotient-submersion hypothesis are made explicit; the closed-form/Frobenius argument is the standard qualification.}
\GSfield{Fate}{The notes expressly attribute the framework to earlier Guillemin--Sternberg work. \citet[Theorem~1]{ginzburg1983symplectic} uses the same characteristic reduction; its reference~3 records the author\textquotesingle s 1981 preprint. The 1983 publication date therefore cannot establish novelty in October 1982.}
\GSfield{Present status}{As of November 2026: the qualified statement is established; the notebook is an exposition, not a newly open question.}
\GSfield{Assessment}{literal statement: correct with missing hypotheses; publication match: partial; verification: read the result; development: annotated; confidence: high}
\end{GSitem}
\begin{GSitem}{1982-570.2}{E}{The Lagrangian equivalence relation of a coisotropic quotient}
\GSfield{Kind; source}{construction; \GSpage{P001050}}
\GSfield{The notes}{\(\Lambda=\{(\xi_1,\xi_2)\in\Sigma^2:\xi_1\sim\xi_2\}\) is Lagrangian and satisfies \(\Lambda^t=\Lambda\), \(\Lambda\circ\Lambda=\Lambda\).}
\GSfield{Statement}{With the preceding regular quotient, \(\Lambda=\Sigma\times_M\Sigma\) is Lagrangian in \(P\times\overline P\), where the second symplectic form is \(-\omega\). If \(\dim P=2n\) and \(\operatorname{codim}\Sigma=k\), then \(\dim\Lambda=2n\); equality of the two pullbacks of \(q^*\omega_M\) gives isotropy. Symmetry and idempotence are the ordinary equivalence-relation identities.}
\GSfield{Reconstruction}{The reversed symplectic sign and the regular fiber product are essential additions.}
\GSfield{Fate}{This is the canonical-relation form of the reduction used in \citet[Section~2]{ginzburg1983symplectic}. The dimension and pullback check above is supplied explicitly \ours; the source does not present it as a new 1983 theorem.}
\GSfield{Present status}{As of November 2026: the qualified statement is established; the notebook is an exposition, not a newly open question.}
\GSfield{Assessment}{literal statement: correct with missing hypotheses; publication match: partial; verification: read the result; development: annotated; confidence: high}
\end{GSitem}
\begin{GSitem}{1982-570.3}{F}{The proposed converse for symmetric idempotent relations}
\GSfield{Kind; source}{conjecture; \GSpage{P001051}}
\GSfield{The notes}{\ru{По-видимому, обратно}: a Lagrangian \(\Lambda\) with \(\Lambda^t=\Lambda\) and \(\Lambda\circ\Lambda=\Lambda\) should arise from \(\Sigma\).}
\GSfield{Statement}{Under regularity assumptions on both projections, a smooth support \(\Sigma\), a smooth quotient and connected equivalence classes, the relation is the characteristic-leaf relation of a coisotropic support. The notebook does not say whether singular supports, immersed relations, disconnected classes or non-Hausdorff quotients are allowed, so an unrestricted global converse is not a uniquely specified assertion.}
\GSfield{Reconstruction}{A regular local reading is described separately; no singular global theorem is imported.}
\GSfield{Fate}{\citet{ginzburg1983symplectic} gives a related coisotropic realization theorem for Lagrangian families, not the unrestricted converse written here. Searches for canonical projectors and symmetric idempotent relations did not identify an exact source matching all omitted global choices.}
\GSfield{Present status}{As of November 2026: the regular interpretation belongs to standard symplectic reduction, while the notebook does not determine a distinct global open problem.}
\GSfield{Assessment}{literal statement: not determinable; publication match: partial; verification: read the result; development: annotated; confidence: medium}
\end{GSitem}
\begin{GSitem}{1982-570.4}{E}{When Fourier integral kernels form an algebra}
\GSfield{Kind; source}{claim; \GSpage{P001051}}
\GSfield{The notes}{Fourier integral operators associated with \(\Lambda\) are denoted \(\mathcal E_\Lambda\); they form a ring because \(\Lambda\circ\Lambda=\Lambda\).}
\GSfield{Statement}{Closure under composition requires an admissible operator class: appropriate proper support, a clean canonical fiber product and the corresponding excess/order convention. In the transverse case, properly supported Fourier integral operators compose with summed orders when the canonical projection is proper. An idempotent set-theoretic relation alone does not ensure that arbitrary distribution kernels can be multiplied and integrated.}
\GSfield{Reconstruction}{Support, cleanliness and filtration qualifications are supplied, without identifying the notation with a unique sheaf of microdifferential operators.}
\GSfield{Fate}{The preexisting transverse theorem is \citet[Theorems~4.2.2--4.2.3]{hormander1971fourier}, published in 1971. The notebook\textquotesingle s characteristic quotient generally calls for its clean-composition extension; it does not specify a complete analytic or complex-microlocal calculus.}
\GSfield{Present status}{As of November 2026: the qualified statement is established; the notebook is an exposition, not a newly open question.}
\GSfield{Assessment}{literal statement: correct with missing hypotheses; publication match: partial; verification: read the result; development: annotated; confidence: high}
\end{GSitem}
\begin{GSitem}{1982-570.5}{E}{The product model: pseudodifferential in one factor, smoothing in the other}
\GSfield{Kind; source}{construction; \GSpage{P001051}}
\GSfield{The notes}{For \(X=X_1\times X_2\), the formula is \(\Lambda=\operatorname{diag}(T^*X_1)\times0\); operators are pseudodifferential in \(X_1\), smoothing in \(X_2\). The preceding prose describes covectors normal to \(X_1\), along \(X_2\).}
\GSfield{Statement}{The formula and operator description agree with \(\Sigma=T^*X_1\times0_{X_2}\): its covectors annihilate tangent vectors to the \(X_2\) fibers. Its characteristic leaves vary in \(X_2\), and \(\Lambda=\operatorname{diag}(T^*X_1)\times(0_{X_2}\times0_{X_2})\). Here each zero denotes an entire zero section, not a single zero covector at one base point. Product kernels are pseudodifferential in the first variables and smooth in both second-factor variables.}
\GSfield{Reconstruction}{The zero sections are expanded and the normal-direction wording is emended to match the two independent displayed descriptions.}
\GSfield{Fate}{The distinction follows directly from the conormal and kernel definitions in \citet[Chapters~II--III]{hormander1971fourier}. The two factor labels in the notebook\textquotesingle s verbal normal-direction description are reversed relative to its formula and operator conclusion.}
\GSfield{Present status}{As of November 2026: the qualified statement is established; the notebook is an exposition, not a newly open question.}
\GSfield{Assessment}{literal statement: false as written; publication match: partial; verification: read the result; development: annotated; confidence: high}
\end{GSitem}
\begin{GSitem}{1982-570.6}{E}{Shmelyov’s conormal embedding condition}
\GSfield{Kind; source}{method; \GSpage{P001051}; Shmelyov}
\GSfield{The notes}{Shmelyov draws a double fibration \(Z\to X,Y\) and requires the induced map \(N^*Z\to T^*Y\) to be an embedding.}
\GSfield{Statement}{For a smooth incidence submanifold \(Z\subset X\times Y\), remove zero covectors and use the signed conormal as the canonical relation of the integral transform \(R:C_c^\infty(X)\to C^\infty(Y)\). Injective immersion of its projection to the data cotangent bundle is the Bolker condition. With appropriate properness and a nonvanishing weight, \(R^*R\) is pseudodifferential and elliptic on visible covectors; complete visibility gives global microlocal ellipticity. This alone does not imply exact injectivity for every smooth weight.}
\GSfield{Reconstruction}{The incidence embedding, zero-section removal, signs, properness and visibility are specified.}
\GSfield{Fate}{\citet[Introduction and Section~2]{homanzhou2017radon} explicitly attributes the FIO and Bolker normal-operator theorem to earlier Guillemin--Sternberg work. Its modern injectivity and stability theorems impose additional analytic or generic hypotheses. Thus the condition in the notes is expository; later stronger results are kept separate.}
\GSfield{Present status}{As of November 2026: the qualified statement is established; the notebook is an exposition, not a newly open question.}
\GSfield{Assessment}{literal statement: correct with missing hypotheses; publication match: exact; verification: indirect; development: annotated; confidence: high}
\end{GSitem}

\GStalk{571}{1982-10-11}{J. Fröhlich}{P001050-P001056}{Nonminimal monopoles, four-manifold topology, and Witten’s Morse deformation}
The entry begins midway down P001051 after Ginzburg and ends on P001054. Fröhlich reports Taubes, Donaldson and Witten, with a final announcement about Enss. The notes distinguish small from zero eigenvalues, but their deformed differential changes sign between pages, their oscillator loses diagonal terms, and its ground state loses the Gaussian. The reconstructions keep the monopole sector, regularity of instanton moduli and long-range scattering modifiers explicit.

\begin{GSitem}{1982-571.1}{E}{The massless Yang–Mills–Higgs energy and magnetic charge}
\GSfield{Kind; source}{definition; \GSpage{P001051}}
\GSfield{The notes}{For \(SU(2)\) with a three-component adjoint scalar, \(E=\int(|D_A\varphi|^2+F^2)\,dx\) on \(\mathbb R^3\). A monopole has \(\deg(\varphi|_{S^2_\infty})=1\); minimum-energy solutions are known.}
\GSfield{Statement}{Choose a positive invariant inner product on \(\mathfrak{su}_2\), coupling one, finite energy and \(|\Phi|\to v>0\). The normalized asymptotic Higgs field defines \(k=\deg(\Phi/|\Phi|:S^2_\infty\to S^2)\). For \(E=\frac12\int(|B|^2+|D_A\Phi|^2)\), square completion gives \(E\ge4\pi v|k|\), attained by \(B=\pm D_A\Phi\). The absent Higgs potential identifies the Prasad--Sommerfield limit.}
\GSfield{Reconstruction}{Asymptotic norm, positive inner product, factor one-half and magnetic field notation are added.}
\GSfield{Fate}{This is the older Bogomolny framework underlying \citet[abstract and setup]{taubes1982nonminimal1}. The square completion fixes the normalization; no claim that every minimizer is explicitly produced by twistors is needed.}
\GSfield{Present status}{As of November 2026: the qualified statement is established; the notebook is an exposition, not a newly open question.}
\GSfield{Assessment}{literal statement: correct with missing hypotheses; publication match: partial; verification: publisher metadata checked; development: annotated; confidence: high}
\end{GSitem}
\begin{GSitem}{1982-571.2}{E}{Nonminimal stationary points and the monopole–antimonopole minimax}
\GSfield{Kind; source}{question; \GSpage{P001051}}
\GSfield{The notes}{After the charge-one definition the notes ask for other stationary points, answer yes, cite Taubes, draw a noncontractible loop and compare the solution with a monopole--antimonopole pair. The displayed maximum and minimum are reversed.}
\GSfield{Statement}{Does the finite-energy \(SU(2)\) Yang--Mills--Higgs functional in the massless-Higgs limit have critical points outside both Bogomolny equations? The loop method seeks \(c=\inf_{\gamma\in\Gamma}\sup_{t\in S^1}E(\gamma(t))\), with \(\Gamma\) a fixed nontrivial loop class in the gauge quotient and compactness proved separately. A monopole--antimonopole pair has total charge zero; it does not itself answer a fixed charge-one question.}
\GSfield{Reconstruction}{The definite existence question is separated from the ambiguous fixed-sector reading; the variational order is emended.}
\GSfield{Fate}{\citet[abstract]{taubes1982nonminimal2} published existence in September 1982, before this talk. The later each-sector multiplicity theorem \citep{taubes1985minmax} requires the corrected argument in \citet{taubes1989erratum}. These papers were checked at publisher-abstract and metadata level.}
\GSfield{Present status}{As of November 2026: nonminimal finite-energy solutions are established in this massless-Higgs model; sector details are not inferred from the pair picture.}
\GSfield{Assessment}{literal statement: false as written; publication match: partial; verification: publisher metadata checked; development: annotated; confidence: high; problem \hyperref[prob:GP-1982-05]{GP-1982-05}}
\end{GSitem}
\begin{GSitem}{1982-571.3}{E}{A unit instanton on a simply connected homology four-sphere}
\GSfield{Kind; source}{question; \GSpage{P001052}}
\GSfield{The notes}{For compact \(M^4\) with \(\pi_1(M)=1\), \(H^2(M)=0\), and a conformal structure, is there an \(SU(2)\) connection with \(F=\pm *F\) and characteristic class one? The answer is yes, citing Taubes, CMP 1981.}
\GSfield{Statement}{For a smooth closed oriented simply connected four-manifold with \(H^2(M;\mathbb R)=0\) and a smooth metric, the charge-one \(SU(2)\) bundle admits an anti-self-dual connection. Equivalently reverse orientation and use self-duality with \(c_2[M]=-1\). The Hodge star on two-forms is conformally invariant.}
\GSfield{Reconstruction}{We specify closedness, smoothness, orientation and cohomology coefficients, as well as the sign relating charge to self-duality.}
\GSfield{Fate}{Taubes existence is invoked explicitly in \citet[Section~3, reference~6]{donaldson1983selfdual}. The identifiable original is JDG 17 (1982), 139--170, rather than CMP 1981. Verification is indirect through Donaldson; the original full proof was inaccessible.}
\GSfield{Present status}{As of November 2026: existence is established; an instanton does not identify the manifold diffeomorphically with the standard sphere.}
\GSfield{Assessment}{literal statement: correct with missing hypotheses; publication match: exact; verification: indirect; development: annotated; confidence: high; problem \hyperref[prob:GP-1982-06]{GP-1982-06}}
\end{GSitem}
\begin{GSitem}{1982-571.4}{B}{The five-dimensional instanton moduli calculation}
\GSfield{Kind; source}{claim; \GSpage{P001052}}
\GSfield{The notes}{The parameter space is a five-dimensional manifold, as for the sphere.}
\GSfield{Statement}{The charge-one anti-self-dual complex has virtual real dimension \(8-3(1-b_1+b_2^+)=5\). The moduli space is smooth of this dimension near irreducible unobstructed connections. A suitable generic perturbation gives regularity; the index alone does not give a smooth moduli space for every unperturbed metric.}
\GSfield{Reconstruction}{Virtual dimension, actual regularity and stabilizers are distinguished.}
\GSfield{Fate}{\citet[Lemma~1]{donaldson1983selfdual} supplies generic perturbation, using the earlier Atiyah--Hitchin--Singer dimension formula.}
\GSfield{Present status}{As of November 2026: the identification and limits are as stated above.}
\GSfield{Assessment}{literal statement: correct with missing hypotheses; publication match: partial; verification: read the result; development: annotated; confidence: high}
\end{GSitem}
\begin{GSitem}{1982-571.5}{B}{The instanton end and a five-dimensional cobordism}
\GSfield{Kind; source}{construction; \GSpage{P001052}}
\GSfield{The notes}{The parameter space \(B\) has boundary diffeomorphic to \(M^4\).}
\GSfield{Statement}{The charge-one moduli space has a bubbling end parametrized by concentration point and scale, \(M\times(0,\varepsilon)\). After regular perturbation and compactification, absence of reducibles in the \(H^2=0\) case yields a compact oriented five-manifold with boundary \(M\). The boundary is not already contained in the smooth space of connections.}
\GSfield{Reconstruction}{Bubbling compactification and perturbation are restored.}
\GSfield{Fate}{\citet[Lemma~2]{donaldson1983selfdual} gives the product end and compact complement. Received 9 August 1982, published January 1983; an exact earlier public circulation date was not found.}
\GSfield{Present status}{As of November 2026: the identification and limits are as stated above.}
\GSfield{Assessment}{literal statement: correct with missing hypotheses; publication match: exact; verification: read the result; development: annotated; confidence: high}
\end{GSitem}
\begin{GSitem}{1982-571.6}{E}{Oriented four-dimensional cobordism and the signature}
\GSfield{Kind; source}{theorem; \GSpage{P001052}}
\GSfield{The notes}{\(M^4\) is a boundary iff its Pontryagin class vanishes, equivalently its intersection-form signature vanishes.}
\GSfield{Statement}{For smooth closed oriented \(M^4\), oriented null-cobordism is equivalent to \(\sigma(M)=0\), or \(\langle p_1(TM),[M]\rangle=0\). Indeed \(\Omega_4^{SO}\simeq\mathbb Z\), generated by \(\mathbb{CP}^2\), and \(\langle p_1,[M]\rangle=3\sigma(M)\). For connected oriented \(M\), the number also detects the top-degree integral class.}
\GSfield{Reconstruction}{Oriented cobordism and the characteristic number specify the shorthand.}
\GSfield{Fate}{This is old material: \citet[Theorem~IV.13 and Corollary~IV.14]{thom1954globales}.}
\GSfield{Present status}{As of November 2026: the qualified statement is established; the notebook is an exposition, not a newly open question.}
\GSfield{Assessment}{literal statement: correct with missing hypotheses; publication match: exact; verification: read the result; development: annotated; confidence: high}
\end{GSitem}
\begin{GSitem}{1982-571.7}{B}{Donaldson’s diagonalization theorem versus a manifold classification}
\GSfield{Kind; source}{theorem; \GSpage{P001052}}
\GSfield{The notes}{A positive-definite form implies that the manifold is an \ru{обобщ. проект. плоск. \(CP^2\)}; hence its integral form is a sum of squares.}
\GSfield{Statement}{If \(M\) is smooth, closed, oriented and simply connected with positive-definite intersection form, then \(Q_M\simeq_{\mathbb Z}\operatorname{diag}(1,\ldots,1)\). The theorem does not assert \(M\cong\#_{b_2}\mathbb{CP}^2\) by diffeomorphism.}
\GSfield{Reconstruction}{The conclusion is restricted to integral forms.}
\GSfield{Fate}{The exact conclusion is \citet[Theorem, p.~81]{donaldson1983selfdual}. Fröhlich reports Donaldson; the manifold-equality wording exceeds that theorem.}
\GSfield{Present status}{As of November 2026: the identification and limits are as stated above.}
\GSfield{Assessment}{literal statement: false as written; publication match: exact; verification: read the result; development: annotated; confidence: high}
\end{GSitem}
\begin{GSitem}{1982-571.8}{B}{Reducible instantons and cones on the complex projective plane}
\GSfield{Kind; source}{construction; \GSpage{P001052}}
\GSfield{The notes}{For the positive-definite case, \(B\) has singularities that are cones over \(\mathbb{CP}^2\).}
\GSfield{Statement}{In this charge-one self-dual problem, suitable perturbations leave isolated reducibles modeled on \(\mathbb C^3/S^1=C(\mathbb{CP}^2)\). Their number is \(n(Q)=\frac12\#\{a:Q(a)=1\}\). Removing cone neighborhoods gives a cobordism between \(M\) and \(n(Q)\) projective planes; signature and the norm-one-vector bound force diagonalization.}
\GSfield{Reconstruction}{Charge, regularity, reducibility and the lattice count are supplied.}
\GSfield{Fate}{The local model and count appear in \citet[Sections~3--4]{donaldson1983selfdual}.}
\GSfield{Present status}{As of November 2026: the identification and limits are as stated above.}
\GSfield{Assessment}{literal statement: correct with missing hypotheses; publication match: exact; verification: read the result; development: annotated; confidence: high}
\end{GSitem}
\begin{GSitem}{1982-571.9}{E}{The weak Morse inequalities and their hypotheses}
\GSfield{Kind; source}{theorem; \GSpage{P001052}}
\GSfield{The notes}{For isolated critical points, \(M_p\) counts points with \(p\) negative Hessian eigenvalues; \(M_p\ge B_p\).}
\GSfield{Statement}{For a smooth Morse function on a smooth closed manifold, let \(m_p\) count its nondegenerate critical points of index \(p\). Then \(m_p\ge b_p(M;\mathbb R)\). Isolated degenerate critical points cannot in general be counted solely by their negative Hessian eigenvalues.}
\GSfield{Reconstruction}{Closedness and nondegeneracy are restored.}
\GSfield{Fate}{These classical inequalities are the starting point of \citet[Section~2]{witten1982supersymmetry}.}
\GSfield{Present status}{As of November 2026: the qualified statement is established; the notebook is an exposition, not a newly open question.}
\GSfield{Assessment}{literal statement: correct with missing hypotheses; publication match: partial; verification: read the result; development: annotated; confidence: high}
\end{GSitem}
\begin{GSitem}{1982-571.10}{B}{Conjugating the de Rham differential}
\GSfield{Kind; source}{construction; \GSpage{P001052-P001053}}
\GSfield{The notes}{\(d_t=e^{-tf}de^{tf}\), \(d_t^2=0\), and unchanged cohomology are followed by \(d_t=d-t\,df\wedge\).}
\GSfield{Statement}{The conjugation gives \(d_t=d+t\,df\wedge\), \(d_t^*=d^*+t\,\iota_{\nabla f}\). Multiplication by \(e^{tf}\) is a chain isomorphism \((\Omega^*,d_t)\to(\Omega^*,d)\), proving nilpotence and unchanged cohomology. The two later minuses instead correspond to \(t\mapsto-t\).}
\GSfield{Reconstruction}{The consistent sign and chain isomorphism replace the loose appeal to homotopy.}
\GSfield{Fate}{See \citet[Section~2, equations~(9)--(10)]{witten1982supersymmetry}. Received 26 March 1982, published December 1982; receipt is not a verified public-preprint date.}
\GSfield{Present status}{As of November 2026: the identification and limits are as stated above.}
\GSfield{Assessment}{literal statement: false as written; publication match: exact; verification: read the result; development: annotated; confidence: high}
\end{GSitem}
\begin{GSitem}{1982-571.11}{B}{Deformed Hodge representatives}
\GSfield{Kind; source}{construction; \GSpage{P001052-P001053}}
\GSfield{The notes}{Representatives satisfy \(d_t\omega=0\) and \((\omega,d_t\beta)=0\), hence \(\Box_t\omega=0\).}
\GSfield{Statement}{On a closed Riemannian manifold, \(\Delta_t=d_t^*d_t+d_td_t^*\) is nonnegative self-adjoint elliptic, with \(\langle\Delta_t\omega,\omega\rangle=\|d_t\omega\|^2+\|d_t^*\omega\|^2\). Deformed harmonic forms represent its cohomology, so \(\dim\ker(\Delta_t|_{\Omega^p})=b_p\) for every finite \(t\).}
\GSfield{Reconstruction}{Compactness and the analytic operator setting are explicit.}
\GSfield{Fate}{This is the Hodge step in \citet[Section~2, equation~(11)]{witten1982supersymmetry}.}
\GSfield{Present status}{As of November 2026: the identification and limits are as stated above.}
\GSfield{Assessment}{literal statement: correct with missing hypotheses; publication match: exact; verification: read the result; development: annotated; confidence: high}
\end{GSitem}
\begin{GSitem}{1982-571.12}{E}{Exterior algebra as a fermionic Fock space}
\GSfield{Kind; source}{construction; \GSpage{P001053-P001054}; A. A. Kirillov and Jürgen Fröhlich}
\GSfield{The notes}{Kirillov identifies \(a_i\) with contraction and \(a_i^*\) with exterior multiplication, correcting the initial description of \(a_i\). The mixed anticommutator is \(\delta_{ij}\), and \(a_i\Omega_0=0\).}
\GSfield{Statement}{For an orthonormal one-particle basis, \(\mathcal F_-=\Lambda^*\mathbb C^n\), \(a_i^*=e_i\wedge\), \(a_i=\iota_{e_i}\), \(\Omega_0=1\). Besides \(\{a_i,a_j^*\}=\delta_{ij}\), one has \(\{a_i,a_j\}=\{a_i^*,a_j^*\}=0\). Ordered creation products form the exterior basis; an unrestricted antisymmetric coefficient sum requires a matching factorial convention.}
\GSfield{Reconstruction}{The omitted anticommutators and coefficient convention are supplied.}
\GSfield{Fate}{The standard fermionic algebra is recalled in \citet[Section~2, before equation~(13)]{witten1982supersymmetry}.}
\GSfield{Present status}{As of November 2026: the qualified statement is established; the notebook is an exposition, not a newly open question.}
\GSfield{Assessment}{literal statement: false as written; publication match: partial; verification: read the result; development: annotated; confidence: high}
\end{GSitem}
\begin{GSitem}{1982-571.13}{B}{The Witten Laplacian and its Hessian term}
\GSfield{Kind; source}{construction; \GSpage{P001054}}
\GSfield{The notes}{\(\Box_t=\Box+t^2(df)^2+t\sum[a_i^*,a_j]\partial_i\partial_jf\).}
\GSfield{Statement}{For the plus-sign differential, in an orthonormal frame, \[\Delta_t=\Delta+t^2|df|^2+t\sum_{i,j}(\nabla^2f)_{ij}[a_i^*,a_j].\] The Hessian is covariant. On functions its last term is \(-t\operatorname{tr}(\nabla^2f)\), independently checked by applying \(d_t^*d_t\).}
\GSfield{Reconstruction}{Covariant derivatives and the nonnegative Laplacian convention are restored.}
\GSfield{Fate}{Compare \citet[Section~2, equation~(13)]{witten1982supersymmetry}. The zero-form check is supplied here \ours.}
\GSfield{Present status}{As of November 2026: the identification and limits are as stated above.}
\GSfield{Assessment}{literal statement: correct with missing hypotheses; publication match: exact; verification: read the result; development: annotated; confidence: high}
\end{GSitem}
\begin{GSitem}{1982-571.14}{B}{The normal harmonic oscillator at a Morse point}
\GSfield{Kind; source}{construction; \GSpage{P001054}}
\GSfield{The notes}{For \(f=f_0+\sum s_i\lambda_i y_i^2+O(|y|^3)\), the potential is \(4t^2\sum\lambda_i^2y_i^2\), but the Hessian sum uses \(i<j\) and marginal factors four.}
\GSfield{Statement}{For \(\lambda_i>0\), \(s_i=\pm1\), the quadratic model is \[H_t^{\rm quad}=\sum_i(-\partial_{y_i}^2+4t^2\lambda_i^2y_i^2)+2t\sum_i s_i\lambda_i[a_i^*,a_i].\] The Hessian has diagonal entries \(2s_i\lambda_i\); restricting to \(i<j\) loses them all. Since \([a_i^*,a_i]=2N_i-1\), the bosonic ground energy cancels for occupation \(N_i=0\) in positive directions and \(N_i=1\) in negative directions.}
\GSfield{Reconstruction}{Diagonal terms and their coefficient are emended.}
\GSfield{Fate}{See \citet[Section~2, equations~(15)--(18)]{witten1982supersymmetry}. The notebook normalization gives the factor two, independently checked \ours.}
\GSfield{Present status}{As of November 2026: the identification and limits are as stated above.}
\GSfield{Assessment}{literal statement: false as written; publication match: exact; verification: read the result; development: annotated; confidence: high}
\end{GSitem}
\begin{GSitem}{1982-571.15}{B}{The missing Gaussian in the local ground state}
\GSfield{Kind; source}{claim; \GSpage{P001054}}
\GSfield{The notes}{The operator allegedly annihilates \(dy_{n-p+1}\wedge\cdots\wedge dy_n\).}
\GSfield{Statement}{The unique square-integrable zero mode of the quadratic model, up to scalar, is \[\omega_{0,t}=e^{-t\sum_i\lambda_i y_i^2}dy_{n-p+1}\wedge\cdots\wedge dy_n.\] The bare wedge is neither annihilated nor square-integrable on \(\mathbb R^n\). Direct differentiation verifies \(H_t^{\rm quad}\omega_{0,t}=0\); all other oscillator/Fock occupations have positive energy.}
\GSfield{Reconstruction}{The scalar Gaussian is an essential emendation.}
\GSfield{Fate}{The degree selection matches \citet[Section~2, equation~(18)]{witten1982supersymmetry}. The Gaussian substitution is checked here \ours.}
\GSfield{Present status}{As of November 2026: the identification and limits are as stated above.}
\GSfield{Assessment}{literal statement: false as written; publication match: exact; verification: read the result; development: annotated; confidence: high}
\end{GSitem}
\begin{GSitem}{1982-571.16}{B}{Small eigenvalues, quasimodes and the Morse inequalities}
\GSfield{Kind; source}{method; \GSpage{P001054}}
\GSfield{The notes}{The notes explicitly seek small rather than necessarily zero eigenvalues; each critical point gives an approximate eigenfunction.}
\GSfield{Statement}{On a closed manifold with Morse \(f\), for large \(t\) the low-energy spectral subspace of \(\Delta_t\) in degree \(p\) has dimension \(m_p\), containing the \(b_p\)-dimensional kernel. Cut-off local Gaussians are quasimodes; their small eigenvalues need not vanish. The low-energy subcomplex yields strong and weak Morse inequalities.}
\GSfield{Reconstruction}{A spectral cutoff, the large-parameter regime and the subcomplex make the shorthand precise.}
\GSfield{Fate}{\citet[Section~2]{witten1982supersymmetry} presents the semiclassical mechanism. Its original asymptotic argument is distinguished from a complete modern analytic proof.}
\GSfield{Present status}{As of November 2026: the identification and limits are as stated above.}
\GSfield{Assessment}{literal statement: correct with missing hypotheses; publication match: exact; verification: read the result; development: annotated; confidence: high}
\end{GSitem}
\begin{GSitem}{1982-571.17}{B}{Gelfand’s replacement of cycle modifications by form modifications}
\GSfield{Kind; source}{remark; \GSpage{P001054}; I. M. Gelfand}
\GSfield{The notes}{Gelfand says \ru{вместо перестроек в циклах --- перестройки в формах}, connecting differential operators with Morse modifications.}
\GSfield{Statement}{The remark describes the passage from handle attachments and cycles to spectral localization of the de Rham complex near critical points. It proposes no theorem beyond the surrounding construction.}
\GSfield{Fate}{This is a conceptual reading of \citet[Section~2]{witten1982supersymmetry}.}
\GSfield{Present status}{As of November 2026: the identification and limits are as stated above.}
\GSfield{Assessment}{literal statement: correct; publication match: partial; verification: read the result; development: annotated; confidence: high}
\end{GSitem}
\begin{GSitem}{1982-571.18}{B}{Supersymmetry as geometry of field configurations}
\GSfield{Kind; source}{method; \GSpage{P001054}}
\GSfield{The notes}{\(d_t\) is a supercharge; an expression with \(\delta/\delta\varphi(x)\) and \(ta^*(x)\varphi^n(x)\) is compared with \(d+t\,df\wedge\). In two dimensions fermions are forms on bosonic configuration space.}
\GSfield{Statement}{Finite-dimensional supersymmetric quantum mechanics uses \(d_t,d_t^*\). For a theory on a spatial circle, the analogous configuration space is a space of maps. The form picture is formal until domains, a measure and regularized operators are specified. The functional derivative displayed in the notes alone does not define a quantum field operator.}
\GSfield{Reconstruction}{The finite-dimensional result is separated from the formal infinite-dimensional analogy.}
\GSfield{Fate}{\citet[Sections~4--5]{witten1982supersymmetry} develops the loop-space analogy and explicitly leaves analytic foundations incomplete.}
\GSfield{Present status}{As of November 2026: the identification and limits are as stated above.}
\GSfield{Assessment}{literal statement: correct with missing hypotheses; publication match: partial; verification: read the result; development: annotated; confidence: high}
\end{GSitem}
\begin{GSitem}{1982-571.19}{F}{The proposed higher-dimensional gauge-field interpretation}
\GSfield{Kind; source}{question; \GSpage{P001054}}
\GSfield{The notes}{For \(d>2\), gauge fields may be needed for an analogous interpretation, which is unavailable. A parenthesis calls the Poincaré group \ru{простая}.}
\GSfield{Statement}{No supersymmetry algebra, target, signature, gauge group or proposed complex is specified, so no single unresolved problem is determined. If \ru{простая} means algebraically simple, it is false because translations form a proper normal subgroup; in context it may mean uncomplicated.}
\GSfield{Fate}{\citet[Section~5]{witten1982supersymmetry} discusses Kähler scalar theories and an unclear gauge-theory structure. This thematic match does not identify the missing target.}
\GSfield{Present status}{As of November 2026: the notebook does not specify a definite higher-dimensional problem whose present openness can be assessed.}
\GSfield{Assessment}{literal statement: not determinable; publication match: thematic only; verification: read the result; development: annotated; confidence: medium}
\end{GSitem}
\begin{GSitem}{1982-571.20}{B}{Enss’s announcement of three-body Coulomb asymptotic completeness}
\GSfield{Kind; source}{theorem; \GSpage{P001054}}
\GSfield{The notes}{Enss is credited with asymptotic completeness for three particles with Coulomb interaction.}
\GSfield{Statement}{For three nonrelativistic particles in \(\mathbb R^3\) with positive finite masses and pair Coulomb potentials, after removing center-of-mass motion, the modified channel wave operators exist and their orthogonal ranges exhaust the continuous spectral subspace. Channels include three separated particles and each bound pair plus a third particle. Long-range Dollard modifiers account for Coulomb phases.}
\GSfield{Reconstruction}{Quantum mechanics, channel projections and long-range comparison dynamics are supplied.}
\GSfield{Fate}{The 1983 chapter \citet[abstract]{enss1983threebody}, from a 1981 conference, is explicitly short-range. The Coulomb scope is verified in \citet[Section~1, equations~(1.12)--(1.13), Section~4]{enss1989longrange}, which removes earlier implicit subsystem conditions. The precise 1982 announced version was not located.}
\GSfield{Present status}{As of November 2026: three-body Coulomb completeness is established; the search does not certify that every detail of the 1989 version was already proved in October 1982.}
\GSfield{Assessment}{literal statement: correct with missing hypotheses; publication match: partial; verification: read the result; development: annotated; confidence: high; problem \hyperref[prob:GP-1982-07]{GP-1982-07}}
\end{GSitem}

\GStalk{572}{1982-10-18}{J. Fröhlich}{P001054-P001059}{Random lattice surfaces, roughening, and the proposed Regge continuum limit}
The mathematical entry occupies P001055--P001057. The first model is an integer-valued solid-on-solid height graph with a rigorous roughening result; the later questions concern more general plaquette surfaces and their topology. A closely matching 1982 primary lecture text makes this distinction explicit. Boundary-independent exponential growth was developed in a 1983 paper by the speaker and collaborators, but neither a power-law correction nor a general continuum-limit theorem follows from it. The final Regge curvature theorem is kept separate from the proposed convergence of random-metric integrals.

\begin{GSitem}{1982-572.1}{E}{Four motivations for statistical mechanics of surfaces}
\GSfield{Kind; source}{remark; \GSpage{P001055}}
\GSfield{The notes}{Applications listed are crystal growth, phase interfaces, lattice gauge theories and dual-resonance strings.}
\GSfield{Statement}{These motivate different ensembles: height graphs approximate interfaces, gauge expansions produce plaquette configurations, and string functional integrals sum over maps and metrics. Sharing a surface interpretation does not identify their measures or critical exponents.}
\GSfield{Fate}{The same four motivations appear in \citet[Introduction]{frohlichpfisterspencer1982surfaces}.}
\GSfield{Present status}{As of November 2026: the qualified statement is established; the notebook is an exposition, not a newly open question.}
\GSfield{Assessment}{literal statement: correct; publication match: exact; verification: read the result; development: annotated; confidence: high}
\end{GSitem}
\begin{GSitem}{1982-572.2}{E}{The solid-on-solid measure with zero boundary heights}
\GSfield{Kind; source}{definition; \GSpage{P001055-P001056}}
\GSfield{The notes}{For finite \(\Lambda\subset\mathbb Z^2\), \(h_p\in\mathbb Z\), \(h_p=0\) outside, add vertical walls and use \(e^{-\beta\operatorname{area}(S)}=e^{-\beta|\Lambda|}\prod_{\langle p,q\rangle}e^{-\beta|h_p-h_q|}\).}
\GSfield{Statement}{Let \(\mathcal H_\Lambda=\{h:\mathbb Z^2\to\mathbb Z:h=0\text{ outside }\Lambda\}\). Sum each unordered nearest-neighbor edge meeting \(\Lambda\) once and set \(H_\Lambda(h)=\sum_{\langle p,q\rangle}|h_p-h_q|\). For \(\beta>0\), \(Z_\Lambda=\sum_h e^{-\beta H_\Lambda(h)}<\infty\), and \(\mu_{\Lambda,\beta}=Z_\Lambda^{-1}e^{-\beta H_\Lambda}\). The constant horizontal area cancels in normalization; this model excludes overhangs.}
\GSfield{Reconstruction}{The bond-counting convention and finite-volume state space are explicit.}
\GSfield{Fate}{Compare \citet[Section~7, equation~(7.16)]{frohlichspencer1981transition}.}
\GSfield{Present status}{As of November 2026: the qualified statement is established; the notebook is an exposition, not a newly open question.}
\GSfield{Assessment}{literal statement: correct with missing hypotheses; publication match: exact; verification: read the result; development: annotated; confidence: high}
\end{GSitem}
\begin{GSitem}{1982-572.3}{E}{Localization and exponential height tails at large inverse temperature}
\GSfield{Kind; source}{theorem; \GSpage{P001056}}
\GSfield{The notes}{For large \(\beta\), \(\mathbb E h_p^2\) is bounded and \(\Pr(|h_p|=n)\sim e^{-{\rm const}\,n}\), attributed to Dobrushin.}
\GSfield{Statement}{For sufficiently large \(\beta\), zero-boundary SOS measures satisfy a volume-uniform bound \(\Pr(|h_p|\ge n)\le C_\beta e^{-c_\beta n}\), hence \(\sup_{\Lambda,p}\mathbb E h_p^2<\infty\). A Peierls contour argument proves such exponential bounds; the displayed \(\sim\) is not by itself a verified sharp tail asymptotic.}
\GSfield{Reconstruction}{An exponential bound replaces an unsupported asymptotic equivalence; the finite-volume zero-boundary convention is retained.}
\GSfield{Fate}{\citet[Section~7, Remark~2 after Theorem~7.3]{frohlichspencer1981transition} recalls the low-temperature argument. The notebook attribution to Dobrushin is broader than the exact SOS proof identified here.}
\GSfield{Present status}{As of November 2026: the qualified statement is established; the notebook is an exposition, not a newly open question.}
\GSfield{Assessment}{literal statement: correct with missing hypotheses; publication match: partial; verification: read the result; development: annotated; confidence: high}
\end{GSitem}
\begin{GSitem}{1982-572.4}{E}{Logarithmic height-difference fluctuations at small inverse temperature}
\GSfield{Kind; source}{theorem; \GSpage{P001056}}
\GSfield{The notes}{For small \(\beta\), \(c_2\log(1+|p-q|)\le\mathbb E(h_p-h_q)^2\le c_1\log(1+|p-q|)\).}
\GSfield{Statement}{In a thermodynamic gradient-state limit of the two-dimensional SOS model, for sufficiently small \(\beta>0\) there are positive finite \(c_-(\beta),c_+(\beta)\) bounding \(\mathbb E(h_p-h_q)^2\) above and below by the corresponding multiples of \(\log(1+|p-q|)\). This is a statement about differences; unpinned absolute heights need not have a finite-variance infinite-volume law.}
\GSfield{Reconstruction}{Gradient limits, coupling-dependent constants and the distinction between absolute and relative heights are stated.}
\GSfield{Fate}{The lower bound is \citet[Theorem~7.3 and equation~(7.18)]{frohlichspencer1981transition}; the matching upper bound is recorded in \citet[Theorem~4(2)]{frohlichpfisterspencer1982surfaces}. The 1981 publication predates this talk.}
\GSfield{Present status}{As of November 2026: the qualified statement is established; the notebook is an exposition, not a newly open question.}
\GSfield{Assessment}{literal statement: correct with missing hypotheses; publication match: exact; verification: read the result; development: annotated; confidence: high}
\end{GSitem}
\begin{GSitem}{1982-572.5}{E}{Microscopic height shifts versus an enhanced continuous symmetry}
\GSfield{Kind; source}{claim; \GSpage{P001056}}
\GSfield{The notes}{The product weight is said to have continuous symmetry \(h_p\mapsto h_p+\mathrm{const}\), broken by the boundary condition, and to gain symmetry on shrinking the squares.}
\GSfield{Statement}{For integer-valued heights the exact global shifts are \(h\mapsto h+k\), \(k\in\mathbb Z\); zero boundary values break even this action on a fixed finite-volume state space. Real shifts occur in a continuous-field representation or a scaling description, not as transformations of the original integer configurations. A scaling-limit claim additionally needs an observable topology and rescaling prescription.}
\GSfield{Reconstruction}{The discrete symmetry and the formal enhanced symmetry are separated.}
\GSfield{Fate}{\citet[Section~2, equations~(28)--(29)]{frohlichpfisterspencer1982surfaces} puts the real shift symmetry in a renormalized continuous-field representation.}
\GSfield{Present status}{As of November 2026: the qualified statement is established; the notebook is an exposition, not a newly open question.}
\GSfield{Assessment}{literal statement: false as written; publication match: partial; verification: read the result; development: annotated; confidence: high}
\end{GSitem}
\begin{GSitem}{1982-572.6}{E}{The passage from height graphs to general plaquette surfaces}
\GSfield{Kind; source}{construction; \GSpage{P001056}}
\GSfield{The notes}{\(Z_{\beta,\gamma}=\sum_{\partial S=\gamma}e^{-\beta|S|}\), and the analogous small-\(\beta\) questions are said not to be done.}
\GSfield{Statement}{To define a finite count, specify connected surfaces attached to a fixed finite boundary \(\gamma\), the ambient lattice, permitted intersections, plaquette multiplicities and topology. If disconnected components can be placed arbitrarily far away, even a fixed-area count can be infinite. Connected sets with each plaquette used at most once have at most exponential counting growth, by bounded-degree lattice-animal enumeration.}
\GSfield{Reconstruction}{Connectedness and multiplicity restrictions are necessary editorial qualifications; no unique general ensemble is silently selected.}
\GSfield{Fate}{\citet[Section~3, ensembles (a)--(c)]{frohlichpfisterspencer1982surfaces} explicitly distinguishes the alternatives. Its SOS theorem does not apply to all of them.}
\GSfield{Present status}{As of November 2026: the qualified statement is established; the notebook is an exposition, not a newly open question.}
\GSfield{Assessment}{literal statement: correct with missing hypotheses; publication match: partial; verification: read the result; development: annotated; confidence: high}
\end{GSitem}
\begin{GSitem}{1982-572.7}{F}{The asserted exponential times power counting asymptotic}
\GSfield{Kind; source}{claim; \GSpage{P001056}}
\GSfield{The notes}{\(a_\gamma(n)=\#\{S:\partial S=\gamma,|S|=n\}\sim C_\gamma^n n^{\delta(\gamma)}\), allegedly from general combinatorial theorems.}
\GSfield{Statement}{An exponential growth rate, where established, means \(\log a_\gamma(n)=n\log C+o(n)\) along admissible area values. It does not imply a polynomial prefactor: for example \(a_n=\lfloor C^ne^{\sqrt n}\rfloor\) has the same exponential rate and no form \(A C^n n^\delta\). The notebook does not fix the surface class or area periodicity needed for the stronger claim.}
\GSfield{Reconstruction}{The verified leading growth is separated from the unverified power-law refinement.}
\GSfield{Fate}{\citet[abstract]{durhuusfrohlichjonsson1983surfaces} states only exponential growth with an \(o(n)\) remainder for its two specified models. The counterexample to the inference is an elementary check \ours.}
\GSfield{Present status}{As of November 2026: the quoted general-combinatorics justification does not establish the displayed asymptotic; its exact model-dependent version is not determined by the notes.}
\GSfield{Assessment}{literal statement: not determinable; publication match: partial; verification: publisher metadata checked; development: annotated; confidence: high}
\end{GSitem}
\begin{GSitem}{1982-572.8}{A}{Boundary independence of the exponential growth rate}
\GSfield{Kind; source}{theorem; \GSpage{P001056}}
\GSfield{The notes}{The limit \(n^{-1}\log a_\gamma(n)\) is introduced and \(C_\gamma\) is asserted independent of \(\gamma\).}
\GSfield{Statement}{For the self-avoiding and planar random-surface models considered by Durhuus--Fröhlich--Jónsson, the admissible-area growth satisfies \(a_\gamma(n)=\exp(\beta_0n+o(n))\), with \(\beta_0\) independent of the fixed admissible boundary. A finite boundary modification explains the mechanism, but it must preserve the chosen surface class; it is not a theorem about every object called a surface.}
\GSfield{Reconstruction}{The theorem is restricted to the identified models and admissible areas.}
\GSfield{Fate}{\citet[abstract]{durhuusfrohlichjonsson1983surfaces} is a direct later match, published 31 October 1983 with the speaker as coauthor. Exact preprint date and full proof were not verified.}
\GSfield{Present status}{As of November 2026: the identification and limits are as stated above.}
\GSfield{Assessment}{literal statement: correct with missing hypotheses; publication match: exact; verification: publisher metadata checked; development: annotated; confidence: high}
\end{GSitem}
\begin{GSitem}{1982-572.9}{F}{Does the subexponential exponent depend only on boundary topology?}
\GSfield{Kind; source}{conjecture; \GSpage{P001056}}
\GSfield{The notes}{\(\delta(\gamma)\) should depend on topological properties of \(\gamma\); sketches contrast one and several boundary components.}
\GSfield{Statement}{Once a particular asymptotic \(a_\gamma(n)\sim A_\gamma C^n n^{\delta_\gamma}\) and a surface ensemble are fixed, one can ask whether \(\delta_\gamma\) depends only on the number and knot/link type of boundary components, or also on genus restrictions and embedding dimension. The notes do not choose these parameters and do not establish existence of \(\delta_\gamma\).}
\GSfield{Fate}{The exponential theorem in \citet{durhuusfrohlichjonsson1983surfaces} does not determine a power exponent. Searches of surface-counting and topology literature did not yield a theorem matching the unqualified conjecture.}
\GSfield{Present status}{As of November 2026: a universal topological exponent cannot be assigned without specifying the ensemble and proving that such an exponent exists.}
\GSfield{Assessment}{literal statement: not determinable; publication match: partial; verification: publisher metadata checked; development: annotated; confidence: medium}
\end{GSitem}
\begin{GSitem}{1982-572.10}{F}{The number of handles of a typical surface}
\GSfield{Kind; source}{question; \GSpage{P001057}}
\GSfield{The notes}{What happens when all topologies are allowed, and how many handles does a typical surface have?}
\GSfield{Statement}{A typical genus requires a normalized probability measure: for example a fixed-area ensemble, or a grand-canonical ensemble with specified area and genus weights. It also requires a class of actual surfaces, since a branched plaquette complex need not have a manifold genus. None of the limiting regimes or weights is fixed at this question.}
\GSfield{Fate}{The genus-weighted alternatives already appear in \citet[Section~3, equation~(35)]{frohlichpfisterspencer1982surfaces}. They illustrate the missing specification rather than resolve a unique notebook problem.}
\GSfield{Present status}{As of November 2026: the question is underdetermined; results for fixed-topology, tethered, branched or triangulated models cannot be substituted indiscriminately.}
\GSfield{Assessment}{literal statement: not determinable; publication match: thematic only; verification: read the result; development: annotated; confidence: high}
\end{GSitem}
\begin{GSitem}{1982-572.11}{E}{The partition-function abscissa from counting growth}
\GSfield{Kind; source}{theorem; \GSpage{P001057}}
\GSfield{The notes}{\(Z_{\beta,\gamma}<\infty\) for \(\beta>\beta_0\) and \(Z_{\beta,\gamma}=\infty\) for \(\beta<\beta_0\); the equality cases are overwritten.}
\GSfield{Statement}{For finite nonnegative counts \(a_\gamma(n)\), set \(\beta_0=\limsup_n n^{-1}\log a_\gamma(n)\). The root test gives convergence of \(\sum_n a_\gamma(n)e^{-\beta n}\) for \(\beta>\beta_0\) and divergence for \(\beta<\beta_0\). At equality both behaviors are possible. This statement remains valid without a polynomial correction or a true growth-rate limit.}
\GSfield{Reconstruction}{Finite counts and strict inequalities avoid unsupported critical-point claims.}
\GSfield{Fate}{This elementary power-series fact underlies the surface partition functions in \citet[Section~3, equations~(36)--(37)]{frohlichpfisterspencer1982surfaces}. The limsup formulation is supplied directly \ours.}
\GSfield{Present status}{As of November 2026: the qualified statement is established; the notebook is an exposition, not a newly open question.}
\GSfield{Assessment}{literal statement: correct with missing hypotheses; publication match: exact; verification: read the result; development: annotated; confidence: high}
\end{GSitem}
\begin{GSitem}{1982-572.12}{F}{Divergent normalization does not define a positive-density surface phase}
\GSfield{Kind; source}{claim; \GSpage{P001057}}
\GSfield{The notes}{When \(Z_{\beta,\gamma}=\infty\), the density of occupied plaquettes is positive, as for a Peano curve.}
\GSfield{Statement}{An infinite unnormalized sum does not define probabilities by \(Z^{-1}w\), and therefore by itself implies no density statement. A finite-volume ensemble and thermodynamic limit could define a dense phase, but these are additional choices. Divergence of a one-surface generating function and positive-density plaquette percolation are different statements.}
\GSfield{Fate}{The distinct ensembles in \citet[Section~3]{frohlichpfisterspencer1982surfaces} do not justify this implication. Its failure is logical: the probability measure used to state the conclusion is absent.}
\GSfield{Present status}{As of November 2026: the asserted implication remains invalid without a specified normalized finite-volume construction.}
\GSfield{Assessment}{literal statement: false as written; publication match: none found; verification: read the result; development: annotated; confidence: high}
\end{GSitem}
\begin{GSitem}{1982-572.13}{F}{The conjectured roughening interval for general surfaces}
\GSfield{Kind; source}{conjecture; \GSpage{P001057}}
\GSfield{The notes}{For \(\beta_0<\beta<\beta_R\), fluctuations should be logarithmic. For \(\beta>\beta_R\), the notes write roughening but describe a nearly flat shape.}
\GSfield{Statement}{The description at large inverse temperature is the smooth or rigid regime. Roughening means loss of this rigidity on decreasing \(\beta\). For general embedded plaquette surfaces, the existence of a distinct transition requires a chosen ensemble and a fluctuation observable; the height difference of the SOS graph is no longer automatically defined.}
\GSfield{Fate}{\citet[Section~3, equations~(42)--(45)]{frohlichpfisterspencer1982surfaces} formulates model-specific roughening questions. \citet{durhuusfrohlichjonsson1983surfaces} discusses, rather than proves, such a transition. The 1981 SOS theorem cannot settle this separate conjecture.}
\GSfield{Present status}{As of November 2026: the general-surface question lacks a unique ensemble/observable; the reversed phase terminology is a definite correction.}
\GSfield{Assessment}{literal statement: false as written; publication match: partial; verification: read the result; development: annotated; confidence: high}
\end{GSitem}
\begin{GSitem}{1982-572.14}{E}{The zero-color limit as an analytic interpolation}
\GSfield{Kind; source}{construction; \GSpage{P001057}}
\GSfield{The notes}{The surface theory is equivalent to \(SU(N)\) gauge theory as \(N\to0\), with Maritan and Omero in Physics Letters cited; the weight \(w_N(S)\) tends to \(w(S)\).}
\GSfield{Statement}{This means an analytic interpolation of suitably rescaled lattice-gauge expansions and Wilson observables in the color parameter. In the identified small-coupling construction the zero-color limit yields a self-avoiding surface generating function. It is not a limit through positive integer group ranks approaching zero, and it is not the unrestricted continuum Yang--Mills theory.}
\GSfield{Reconstruction}{Analytic continuation, field rescaling and the small-coupling restriction are restored.}
\GSfield{Fate}{\citet[Theorem~6(2), reference~28]{frohlichpfisterspencer1982surfaces} records the construction and cites the 1981 CERN preprint. The original Maritan--Omero paper was not obtained; verification of its result is indirect through this primary account.}
\GSfield{Present status}{As of November 2026: the qualified statement is established; the notebook is an exposition, not a newly open question.}
\GSfield{Assessment}{literal statement: correct with missing hypotheses; publication match: exact; verification: indirect; development: annotated; confidence: high}
\end{GSitem}
\begin{GSitem}{1982-572.15}{E}{The proposed edge-length measure and its scale divergence}
\GSfield{Kind; source}{construction; \GSpage{P001057}}
\GSfield{The notes}{The weight contains area, vertex angle defects, a denominator \(e^{-c\,\#\mathrm{handles}}\), and \(\prod_{\rm edges}d\ell/\ell\).}
\GSfield{Statement}{For a fixed closed triangulation, lengths must be positive and satisfy triangle inequalities. The displayed denominator equals a genus factor \(e^{cg}\), not suppression when \(c>0\). The bare measure is not normalizable over all lengths: write \(\ell_e=r u_e\), fix one shape coordinate, and restrict the other ratios to a compact nondegenerate open set. The product measure includes \(dr/r\), angle defects are scale invariant, and area is \(r^2A(u)\); hence the integral diverges at \(r=0\) for finite couplings. An area constraint, scale quotient or cutoff is essential.}
\GSfield{Reconstruction}{Admissible lengths and a fixed closed triangulation specify the counterexample; the notebook itself does not supply a normalization or regulator.}
\GSfield{Fate}{The Regge functional-integral setting is discussed in \citet{cheegermullerschrader1982lattice}. The explicit small-scale divergence is an editorial calculation \ours, independent of any continuum conjecture.}
\GSfield{Present status}{As of November 2026: the qualified statement is established; the notebook is an exposition, not a newly open question.}
\GSfield{Assessment}{literal statement: correct with missing hypotheses; publication match: partial; verification: publisher metadata checked; development: annotated; confidence: high}
\end{GSitem}
\begin{GSitem}{1982-572.16}{E}{Regge curvature concentrated on codimension-two hinges}
\GSfield{Kind; source}{definition; \GSpage{P001057}}
\GSfield{The notes}{\(R=\sum_f\operatorname{area}(f)\,\operatorname{defect}(f)\), with \(f\) a \((d-2)\)-simplex; the reference says Regge 64.}
\GSfield{Statement}{In a piecewise Euclidean \(d\)-manifold, an interior hinge \(f\) has deficit \(\epsilon_f=2\pi-\sum_{\sigma\supset f}\theta_{\sigma,f}\). In the convention \(\mathrm{Scal}=2K\) in dimension two, the integrated scalar-curvature measure is \(2\sum_f\operatorname{Vol}_{d-2}(f)\epsilon_f\). A convention absorbing the factor two is possible, but must be stated. Boundary terms require exterior-angle contributions.}
\GSfield{Reconstruction}{Deficits, normalization and the boundary qualification are supplied.}
\GSfield{Fate}{The foundational reference is \citet{regge1961relativity}, published in February 1961. This recurs from 1981-546.6; the formula is defined carefully in \citet[Section~2]{cheegermullerschrader1984curvature}.}
\GSfield{Present status}{As of November 2026: the qualified statement is established; the notebook is an exposition, not a newly open question.}
\GSfield{Assessment}{literal statement: correct with missing hypotheses; publication match: exact; verification: read the result; development: annotated; confidence: high}
\end{GSitem}
\begin{GSitem}{1982-572.17}{E}{Weak convergence of piecewise-flat curvature measures}
\GSfield{Kind; source}{theorem; \GSpage{P001057}}
\GSfield{The notes}{The hinge sum tends to \(\int\sqrt g\,R(x)\,d^dx\) as the triangulation approaches a smooth surface or manifold.}
\GSfield{Statement}{For geodesic triangulations approximating a smooth compact Riemannian manifold, with mesh tending to zero and uniform positive fatness of simplices, the corresponding piecewise-flat Lipschitz--Killing curvature measures converge weakly to their smooth counterparts. In particular the consistently normalized hinge sum converges to total scalar curvature. Arbitrarily thin simplices or pointwise hinge-curvature convergence are not covered.}
\GSfield{Reconstruction}{The approximation, fatness and measure topology are essential hypotheses.}
\GSfield{Fate}{The precise result is \citet[Theorem~5.1]{cheegermullerschrader1984curvature}. Its 1981 announcement and 1982 chapter precede this talk; the earlier occurrence is 1981-546.7. This differs from convergence of Regge field equations in GP-1981-21.}
\GSfield{Present status}{As of November 2026: the qualified statement is established; the notebook is an exposition, not a newly open question.}
\GSfield{Assessment}{literal statement: correct with missing hypotheses; publication match: exact; verification: read the result; development: annotated; confidence: high}
\end{GSitem}
\begin{GSitem}{1982-572.18}{F}{Does the displayed random-metric integral converge to Polyakov theory?}
\GSfield{Kind; source}{question; \GSpage{P001057}}
\GSfield{The notes}{After the curvature theorem, the notes claim that equation (14) tends to Polyakov string theory, followed by Exercise: Prove this!}
\GSfield{Statement}{Convergence of deterministic curvature actions does not prove convergence of probability measures or functional integrals. The proposed limit lacks a triangulation ensemble, scale normalization, treatment of degenerating lengths, gauge quotient, matter fields, coupling renormalization and a topology of convergence. Even the unregulated fixed-triangulation measure above diverges. These omissions prevent a unique theorem from being reconstructed.}
\GSfield{Fate}{Compare the conjectural discussion in \citet[Section~3]{frohlichpfisterspencer1982surfaces} with the deterministic theorem of \citet{cheegermullerschrader1984curvature} and Polyakov\textquotesingle s formal integral \citep{polyakov1981quantum}. None proves the claimed measure limit.}
\GSfield{Present status}{As of November 2026: the notebook formulation does not specify a continuum-limit problem; modern results for other random-surface measures do not settle this displayed unregulated integral.}
\GSfield{Assessment}{literal statement: not determinable; publication match: thematic only; verification: read the result; development: annotated; confidence: high}
\end{GSitem}

\GStalk{573}{1982-11-01}{A. B. Zamolodchikov}{P001057-P001061}{Primary fields, Virasoro degeneracy, and the Liouville spectrum}
The entry begins on P001058 and ends on P001059; the broader catalogue range includes the preceding surface lecture and following Beilinson record. Zamolodchikov asks how to compute primary-field matrix elements and which highest-weight modules have null states; Neretin supplies the reducibility criterion. The second half returns to the independent-metric string integral, the Liouville action, a cylinder evolution kernel and a conjectural Virasoro description of the spectrum. These are distinct questions: BPZ develops the correlator program in 1984, whereas a much later spectral theorem gives a rigorous real-coupling Liouville realization. Several notebook formulas require consistent mode conventions or substantive emendations.

\begin{GSitem}{1982-573.1}{E}{The Virasoro bracket and its missing cubic central factor}
\GSfield{Kind; source}{definition; \GSpage{P001058}}
\GSfield{The notes}{\([L_n,L_m]=(m-n)L_{n+m}+(c/12)\delta_{n+m,0}\mathbf1\). The scan confirms that no cubic mode factor is written.}
\GSfield{Statement}{Reserve \(L_n\) for the notebook and write \(\ell_n=-L_n\), so the modern highest weight is \(h=-\Delta\). Use \([\ell_n,\ell_m]=(n-m)\ell_{n+m}+\frac{c}{12}(n^3-n)\delta_{n+m,0}\mathbf1\), with \(\ell_n\) for \(n>0\) annihilating the highest-weight vector. In notebook generators this has the reverse Witt term \((m-n)L_{n+m}\) and the same cubic central term. The written constant delta term is not antisymmetric, already giving \([L_0,L_0]=c\mathbf1/12\). The neighboring primary law needs the corresponding overall sign change.}
\GSfield{Reconstruction}{Throughout the modern statements use \(\ell_n=-L_n\) and \(h=-\Delta\); retain the notebook symbols in quotations. Restoring the cubic factor is a mathematical emendation.}
\GSfield{Fate}{The central extension predates the talk: \citet[Theorem~1]{gelfandfuchs1968circle}. The convention matching the recorded primary commutator is printed in \citet[equation~(2.21)]{belavinpolyakovzamolodchikov1984infinite}.}
\GSfield{Present status}{As of November 2026: the qualified statement is established; the notebook is an exposition, not a newly open question.}
\GSfield{Assessment}{literal statement: false as written; publication match: exact; verification: read the result; development: annotated; confidence: high}
\end{GSitem}
\begin{GSitem}{1982-573.2}{E}{A cyclic highest-weight representation and its Verma cover}
\GSfield{Kind; source}{construction; \GSpage{P001058}}
\GSfield{The notes}{\(L_n|0\rangle=0\) for \(n>0\), \(L_0|0\rangle=\Delta|0\rangle\), and all vectors are generated by products of negative modes.}
\GSfield{Statement}{For fixed \(c,h\), the induced Verma module \(M(c,h)\) has basis \(\ell_{-n_1}\cdots\ell_{-n_k}|h\rangle\), \(n_1\ge\cdots\ge n_k\ge1\), and level-\(N\) dimension \(p(N)\). Every cyclic highest-weight representation with these parameters is its quotient. Here \(h=-\Delta\) under \(\ell_n=-L_n\). The recorded generation assertion does not say whether the quotient is already irreducible. A general highest-weight vector is not the conformally invariant vacuum: the latter also has \(h=0\) and \(\ell_{-1}|0\rangle=0\).}
\GSfield{Reconstruction}{Ordered PBW monomials, the universal quotient and distinct notation h for the state weight are supplied.}
\GSfield{Fate}{The Verma construction and determinant were established before November 1982; see \citet[Sections~3--4]{feiginfuchs1982invariant}. The vacuum distinction is explicit in \citet[equations~(2.23)--(2.27)]{belavinpolyakovzamolodchikov1984infinite}.}
\GSfield{Present status}{As of November 2026: the qualified statement is established; the notebook is an exposition, not a newly open question.}
\GSfield{Assessment}{literal statement: correct with missing hypotheses; publication match: exact; verification: read the result; development: annotated; confidence: high}
\end{GSitem}
\begin{GSitem}{1982-573.3}{E}{Primary covariance does not uniquely define the field}
\GSfield{Kind; source}{claim; \GSpage{P001058}}
\GSfield{The notes}{\([L_n,\varphi(z)]=z^{n+1}\varphi^{\prime}(z)+(n+1)\Delta z^n\varphi(z)\), followed by \ru{Он опрд. этой ф-лой однозначно}.}
\GSfield{Statement}{In the modern convention, a primary field of weight \(\delta\) obeys \([\ell_n,\varphi(z)]=[z^{n+1}\partial_z+(n+1)\delta z^n]\varphi(z)\). Because \(L_n=-\ell_n\), the notebook reverse-bracket generators require the negative of that right-hand side. The field weight need not equal the source-module weight \(h=-\Delta\). Covariance does not determine an operator uniquely: scalar multiplication preserves it, and module channels, multiplicities and normalization must be specified. For fixed generic source and target Verma modules, formal descendant matrix elements can be recursively fixed from a compatible leading matrix element.}
\GSfield{Reconstruction}{The covariance sign, field weight, source and target data and normalization are made consistent. No unrestricted existence or uniqueness of intertwiners is asserted.}
\GSfield{Fate}{\citet[equation~(3.9) and its following sentence]{belavinpolyakovzamolodchikov1984infinite} identifies the law with earlier dual-model vertex operators; its historical antecedent is verified indirectly. The failure of literal uniqueness follows directly by rescaling.}
\GSfield{Present status}{As of November 2026: primary covariance is standard, but the printed uniqueness claim lacks essential data.}
\GSfield{Assessment}{literal statement: false as written; publication match: exact; verification: indirect; development: annotated; confidence: high}
\end{GSitem}
\begin{GSitem}{1982-573.4}{A}{Computing primary correlators and the role of null states}
\GSfield{Kind; source}{question; \GSpage{P001058}}
\GSfield{The notes}{The first question asks how to calculate \(\langle0|\varphi(z_1)\cdots\varphi(z_n)|0\rangle\).}
\GSfield{Statement}{For specified Virasoro modules, primary weights, intertwining channels and leading three-point matrix elements, commute modes to determine descendant matrix elements and conformal-block coefficients. In a vacuum correlator, global Ward identities constrain the position dependence; three-point coefficients and four-point cross-ratio dependence require dynamical data. If a singular vector vanishes in a chosen irreducible quotient, its decoupling gives a differential equation for correlators. Without those extra choices the displayed expectation has no unique numerical answer.}
\GSfield{Reconstruction}{The additional channel and structure-constant data make the computational question definite. A classification of every conformal field theory is not inserted into the notes.}
\GSfield{Fate}{This is a close precursor of \citet[Sections~3--5]{belavinpolyakovzamolodchikov1984infinite}: the paper separates conformal blocks from structure constants and derives differential equations for degenerate fields. Received on 22 November 1983 and published in 1984, it is a later publication by the speaker, not a theorem already published in November 1982.}
\GSfield{Present status}{As of November 2026: the Ward-identity, block and null-state program is established; the data in the notebook do not select a unique full conformal field theory.}
\GSfield{Assessment}{literal statement: correct with missing hypotheses; publication match: partial; verification: read the result; development: annotated; confidence: high; problem \hyperref[prob:GP-1982-08]{GP-1982-08}}
\end{GSitem}
\begin{GSitem}{1982-573.5}{E}{The contravariant Gram form and positivity}
\GSfield{Kind; source}{construction; \GSpage{P001058}}
\GSfield{The notes}{\(L_{-n}^{+}=L_n\) and \(\langle0|0\rangle=1\) are said to recover the scalar product.}
\GSfield{Statement}{For real \(c,h\), these conditions determine the contravariant Hermitian form on the Verma module by moving positive modes to the right; for complex parameters use the algebraic bilinear form with the standard antiautomorphism. The form can be degenerate or indefinite. For example \(\|\ell_{-1}|h\rangle\|^2=2h\), so positivity does not follow for arbitrary \(h\). The radical is removed when forming the irreducible highest-weight quotient.}
\GSfield{Reconstruction}{Hermitian versus bilinear conventions and the absence of automatic positivity are made explicit.}
\GSfield{Fate}{\citet[Section~3 and Theorem~3.1]{feiginfuchs1982invariant} gives the determinant of this form before the talk. The displayed level-one norm is a direct check \ours.}
\GSfield{Present status}{As of November 2026: the qualified statement is established; the notebook is an exposition, not a newly open question.}
\GSfield{Assessment}{literal statement: correct with missing hypotheses; publication match: exact; verification: read the result; development: annotated; confidence: high}
\end{GSitem}
\begin{GSitem}{1982-573.6}{E}{Neretin on reducibility and the first singular level}
\GSfield{Kind; source}{question; \GSpage{P001058}; Yu. A. Neretin}
\GSfield{The notes}{Neretin gives a two-factor polynomial with \(-\Delta\) in both factors, then says degeneration occurs at level \(\alpha\beta\), for positive integers \(\alpha,\beta\).}
\GSfield{Statement}{For the reverse Witt bracket of the notebook, set \(\ell_n=-L_n\) and \(h=-\Delta\). In the corresponding central normalization, write \(c=13-6(t+t^{-1})\), \(t\ne0\). Reducibility means \(h=((rt-s)^2-(t-1)^2)/(4t)\) for some \(r,s\ge1\), and the first singular level is the least such \(rs\). Thus the recorded \(-\Delta\) is exactly \(h\), and its two-factor polynomial has the correct Kac form. At level two the Gram determinant is \(2h[16h^2+(2c-10)h+c]\): at \(c=0\), \(h=5/8\) corresponds to the recorded \(\Delta=-5/8\).}
\GSfield{Reconstruction}{The recorded polynomial is preserved under \(h=-\Delta\), \(\ell_n=-L_n\). The first-level assertion requires a minimum over coincident degeneracy conditions.}
\GSfield{Fate}{This repeats 1982-566.7 and GP-1982-03. \citet[Theorem~3.1 and Section~4]{feiginfuchs1982invariant} had already proved the criterion. The low-level check \oursComputed verifies the convention conversion. The polynomial is not itself a new error; the incompatible surrounding conventions must be repaired as a collection.}
\GSfield{Present status}{As of November 2026: the qualified statement is established; the notebook is an exposition, not a newly open question.}
\GSfield{Assessment}{literal statement: correct with missing hypotheses; publication match: exact; verification: read the result; development: annotated; confidence: high; problem \hyperref[prob:GP-1982-03]{GP-1982-03}}
\end{GSitem}
\begin{GSitem}{1982-573.7}{E}{Independent embeddings and worldsheet metrics}
\GSfield{Kind; source}{construction; \GSpage{P001058}}
\GSfield{The notes}{Equation (1) integrates independent \(\vec x(\xi)\) and \(g_{ab}(\xi)\) with the Dirichlet action and a cosmological area term; the margin emphasizes that the metric is not induced.}
\GSfield{Statement}{For a fixed smooth worldsheet \(\Sigma\), the Polyakov functional is \(S[X,g]=\frac12\int_\Sigma\sqrt g\,g^{ab}\partial_aX\cdot\partial_bX+\mu_0^2\int_\Sigma\sqrt g\). It treats the embedding and positive metric as independent variables. A quantum partition function additionally needs a topology, boundary conditions, regularization, target zero-mode normalization and a quotient by diffeomorphisms. The formal symbol is not by itself a finite measure.}
\GSfield{Reconstruction}{Fixed topology and the missing analytic data are explicit, without changing the recorded formal action.}
\GSfield{Fate}{The independent-metric formulation is in \citet[equations~(5)--(11)]{polyakov1981quantum}, published in July 1981. It recurs from records 567 and 569.}
\GSfield{Present status}{As of November 2026: the qualified statement is established; the notebook is an exposition, not a newly open question.}
\GSfield{Assessment}{literal statement: correct with missing hypotheses; publication match: exact; verification: read the result; development: annotated; confidence: high}
\end{GSitem}
\begin{GSitem}{1982-573.8}{E}{Diffeomorphism quotient and conformal representatives}
\GSfield{Kind; source}{method; \GSpage{P001058}}
\GSfield{The notes}{\(\delta g_{ab}=\nabla_a\varepsilon_b+\nabla_b\varepsilon_a\); the volume is factored by this group, and a representative \(g_{ab}=\rho\delta_{ab}\) is chosen.}
\GSfield{Statement}{Infinitesimal diffeomorphisms have the displayed Lie derivative. Locally a two-dimensional metric has isothermal coordinates. Globally on a fixed surface one instead writes \(g=e^\varphi\widehat g(m)\), allowing conformal moduli \(m\), and divides residual automorphisms. A fixed global Euclidean metric is unavailable on a general closed surface. For the ordinary geometric quotient, integration is over orbits as points of the quotient, not along their unreduced volumes.}
\GSfield{Reconstruction}{Moduli, reference metric and residual stabilizers are restored; the flat formula is interpreted locally.}
\GSfield{Fate}{\citet[equations~(9)--(14)]{polyakov1981quantum} is the prior formal gauge-fixing source. The local/global distinction also qualifies the repeated construction in 1982-569.9.}
\GSfield{Present status}{As of November 2026: the qualified statement is established; the notebook is an exposition, not a newly open question.}
\GSfield{Assessment}{literal statement: correct with missing hypotheses; publication match: partial; verification: read the result; development: annotated; confidence: high}
\end{GSitem}
\begin{GSitem}{1982-573.9}{E}{Scalar and ghost determinants after gauge fixing}
\GSfield{Kind; source}{method; \GSpage{P001058-P001059}}
\GSfield{The notes}{Equation (2) contains \([\det\Delta_g]^{-D/2}[\det\mathcal L]^{1/2}\), while (3) writes an extra upper \(a\) on the first covariant derivative. The later definition makes \(\Delta_g=\operatorname{div}\operatorname{grad}\).}
\GSfield{Statement}{For positive Euclidean metric the Gaussian matter determinant is \((\det{}^{\prime}(-\Delta_g))^{-D/2}\), with zero modes treated separately. Put \((P\varepsilon)_{ab}=\nabla_a\varepsilon_b+\nabla_b\varepsilon_a-g_{ab}\nabla^c\varepsilon_c\). With the standard \(L^2\) metrics, the notebook operator after removing the extra index is \(\mathcal L=-\frac12P^*P\). The positive ghost determinant is therefore expressed through \((\det{}^{\prime}P^*P)^{1/2}\), with convention-dependent constants, conformal Killing modes and moduli factors. Equation (2) also suppresses the cosmological weight.}
\GSfield{Reconstruction}{Positive operator, primed determinants, the covector operator and missing factors are specified; the doubtful upper index remains a palaeographic ambiguity.}
\GSfield{Fate}{\citet[equations~(12)--(16)]{polyakov1981quantum} gives the formal determinant mechanism. Integration by parts fixes the adjoint sign \ours; it is not a claim that the unregularized determinants exist.}
\GSfield{Present status}{As of November 2026: the qualified statement is established; the notebook is an exposition, not a newly open question.}
\GSfield{Assessment}{literal statement: false as written; publication match: partial; verification: read the result; development: annotated; confidence: high}
\end{GSitem}
\begin{GSitem}{1982-573.10}{E}{The induced Liouville action and its quantum measure}
\GSfield{Kind; source}{construction; \GSpage{P001059}}
\GSfield{The notes}{For \(D<26\), equation (4) gives a formal weight with coefficient \((26-D)/(48\pi)\), kinetic term \(\frac12(\partial\varphi)^2\), potential \(m_0^2e^\varphi\), and \(\rho=e^\varphi\).}
\GSfield{Statement}{This is the flat-reference formal Weyl-anomaly action. On a curved reference surface one needs a background-curvature term and a specified renormalized field measure. In the real probabilistic Liouville theory, \(0<\gamma<2\), \(\mu>0\), \(Q=2/\gamma+\gamma/2\), the exponential interaction is Gaussian-multiplicative-chaos renormalized; sphere insertions satisfy Seiberg bounds. The bare coefficient alone is not a rigorous measure definition and does not identify every \(D<26\) theory with this regime.}
\GSfield{Reconstruction}{Modern parameters, renormalization and curvature terms are qualifications, not formulas attributed to the 1982 notes.}
\GSfield{Fate}{The anomaly action was derived in \citet[equations~(15)--(18)]{polyakov1981quantum}. \citet[Theorems~3.2 and 3.11]{davidkupiainenrhodesvargas2016liouville}, first circulated in October 2014 and published in 2016, constructs a specified real-coupling sphere theory. These are different levels of assertion.}
\GSfield{Present status}{As of November 2026: the qualified statement is established; the notebook is an exposition, not a newly open question.}
\GSfield{Assessment}{literal statement: correct with missing hypotheses; publication match: exact; verification: read the result; development: annotated; confidence: high}
\end{GSitem}
\begin{GSitem}{1982-573.11}{F}{The claim that one integral would settle the string problem}
\GSfield{Kind; source}{remark; \GSpage{P001059}}
\GSfield{The notes}{\ru{Задача о струне будет решена, если вычислить этот интеграл}.}
\GSfield{Statement}{The sentence identifies the Liouville integral as a central remaining calculation, but it does not specify a worldsheet topology, observables, boundary conditions, a genus sum or the sense in which a string theory would be solved. It is not a definite assertion that evaluating one scalar partition function determines every string amplitude.}
\GSfield{Fate}{The modern sphere construction \citep{davidkupiainenrhodesvargas2016liouville} and bootstrap theorem \citep[Theorem~1.1]{guillarmoukupiainenrhodesvargas2024bootstrap} answer well-specified Liouville questions. They cannot resolve the unqualified phrase as a single mathematical problem.}
\GSfield{Present status}{As of November 2026: the broad string-theory assertion remains underdetermined; no new open problem is manufactured from it.}
\GSfield{Assessment}{literal statement: not determinable; publication match: thematic only; verification: read the result; development: annotated; confidence: medium}
\end{GSitem}
\begin{GSitem}{1982-573.12}{E}{Critical dimension and the remaining worldsheet integrations}
\GSfield{Kind; source}{claim; \GSpage{P001059}}
\GSfield{The notes}{For \(D=26\), the answer is said to equal the flat-metric answer, and dual models are called exact.}
\GSfield{Statement}{For \(D\) free bosons, matter and reparametrization ghosts have total central charge \(D-26\), so their local Weyl anomalies cancel at \(D=26\). This permits elimination of the Weyl factor in the critical theory, subject to the cosmological and regularization prescriptions. It does not remove moduli, ghost zero modes, vertex positions, boundary data or genus dependence; nor can every worldsheet be represented by one global flat metric.}
\GSfield{Reconstruction}{The exact local cancellation is separated from the unqualified global assertion.}
\GSfield{Fate}{This is the local anomaly-cancellation statement of \citet[equations~(15)--(18)]{polyakov1981quantum}. Reading the flat-metric shortcut as a global equality loses required integrations.}
\GSfield{Present status}{As of November 2026: the qualified statement is established; the notebook is an exposition, not a newly open question.}
\GSfield{Assessment}{literal statement: false as written; publication match: partial; verification: read the result; development: annotated; confidence: high}
\end{GSitem}
\begin{GSitem}{1982-573.13}{B}{Cylinder evolution and the spectral representation}
\GSfield{Kind; source}{method; \GSpage{P001059}}
\GSfield{The notes}{A drawing joins contours \(c_1,c_2\); equation (5) is \(\sum_nG(c_1,n)G(c_2,n)e^{-E_nt}\), annotated \(t\to\infty\).}
\GSfield{Statement}{For a specified self-adjoint Hamiltonian bounded below and suitable boundary vectors, \(Z(c_1,c_2;t)=\langle c_2,e^{-tH}c_1\rangle=\int e^{-tE}\,d\langle c_2,\Pi_Ec_1\rangle\). A sum is valid for a discrete orthonormal eigenbasis, with the appropriate dual or complex-conjugate amplitude. Large time isolates a lowest eigenstate only if one exists with nonzero overlap; at a continuous threshold one must retain a spectral integral and its threshold asymptotics.}
\GSfield{Reconstruction}{A Hilbert space, boundary vectors and spectral measure replace schematic amplitudes; no large-time gap is assumed.}
\GSfield{Fate}{The spectral-theorem interpretation is old; the Liouville realization in \citet[Sections~2.3--2.4 and 6]{guillarmoukupiainenrhodesvargas2024bootstrap} exhibits continuous momentum. The notebook cylinder is a propagation diagram, not evidence of a discrete Liouville spectrum.}
\GSfield{Present status}{As of November 2026: the identification and limits are as stated above.}
\GSfield{Assessment}{literal statement: correct with missing hypotheses; publication match: partial; verification: read the result; development: annotated; confidence: high}
\end{GSitem}
\begin{GSitem}{1982-573.14}{B}{Liouville stationary states as Virasoro representations}
\GSfield{Kind; source}{conjecture; \GSpage{P001059}}
\GSfield{The notes}{After the Liouville Lagrangian (6), the hypothesis says that its stationary spectrum is related to the \(L_n\) system on the preceding page.}
\GSfield{Statement}{A precise modern realization uses \(0<\gamma<2\), \(\mu>0\), \(Q=2/\gamma+\gamma/2\), and \(c_L=1+6Q^2>25\). Generalized primary states are labelled by \(\alpha=Q+iP\), \(P>0\), with left and right weights \((Q^2+P^2)/4\); their descendants have oscillator levels. In the Hamiltonian normalization of the cited preprint the energies are \((Q^2+P^2)/2+N+\bar N\), and the spectral resolution integrates over \(P\), rather than summing only over a discrete set. This realizes the recorded representation-theoretic relation in an explicitly constructed quantum theory.}
\GSfield{Reconstruction}{The note-supported relation is retained separately from this precise real-coupling implementation. No assertion is made about timelike coupling or all target dimensions.}
\GSfield{Fate}{\citet[Theorem~6.25, Sections~2.4--2.6 and 8]{guillarmoukupiainenrhodesvargas2024bootstrap} supplies the complete spectral resolution and its conformal-block application. The first preprint was 23 May 2020, the version covering \(0<\gamma<2\) is dated 11 November 2020, and publication is 2024. The bare 1982 action must first be renormalized and normalized as that theory.}
\GSfield{Present status}{As of November 2026: the stated real-coupling Liouville spectral relation is rigorously established; this is not a claim that the whole string integral is evaluated.}
\GSfield{Assessment}{literal statement: correct with missing hypotheses; publication match: partial; verification: read the result; development: annotated; confidence: high; problem \hyperref[prob:GP-1982-09]{GP-1982-09}}
\end{GSitem}
\begin{GSitem}{1982-573.15}{E}{The improved classical Liouville stress tensor}
\GSfield{Kind; source}{claim; \GSpage{P001059}}
\GSfield{The notes}{The purported conserved traceless tensor ends with \(\alpha\varphi_{ab}-\beta\delta_{ab}(\partial\varphi)^2\); it has no potential or Laplacian improvement term.}
\GSfield{Statement}{For \(S=\int[\frac12|\nabla\varphi|^2+V(\varphi)]\), \(V=m_0^2e^\varphi\), the equation is \(\Delta\varphi=V\). One consistent tensor is \[T_{ab}=\varphi_a\varphi_b-\delta_{ab}\bigl(\tfrac12|\nabla\varphi|^2+V\bigr)+2(\delta_{ab}\Delta\varphi-\varphi_{ab}).\] Its divergence is \((\Delta\varphi-V)\varphi_b\), and its trace is \(2(\Delta\varphi-V)\), so both vanish on shell. The printed ansatz cannot achieve both for arbitrary Liouville solutions by choosing constants \(\alpha,\beta\).}
\GSfield{Reconstruction}{The complete improved tensor is an editorial emendation; the notebook gives only a malformed ansatz.}
\GSfield{Fate}{This is an elementary classical check \ours of the notebook action, within the Liouville setting of \citet{polyakov1981quantum}; it is not a quantum stress-tensor construction. Differentiation and a separate trace calculation provide independent checks.}
\GSfield{Present status}{As of November 2026: the qualified statement is established; the notebook is an exposition, not a newly open question.}
\GSfield{Assessment}{literal statement: false as written; publication match: partial; verification: read the result; development: annotated; confidence: high}
\end{GSitem}
\begin{GSitem}{1982-573.16}{E}{Holomorphic stress components and the sign of the improvement}
\GSfield{Kind; source}{claim; \GSpage{P001059}}
\GSfield{The notes}{\(T=\varphi_z^2+2\varphi_{zz}\), \(\bar T=\varphi_{\bar z}^2+2\varphi_{\bar z\bar z}\), followed by \(\partial_{\bar z}T=0\) and a second \(\partial_zT=0\) without a bar.}
\GSfield{Statement}{With the recorded \(e^\varphi\) potential, \(\varphi_{z\bar z}=m_0^2e^\varphi/4\), so \[T=\varphi_z^2-2\varphi_{zz},\qquad \bar T=\varphi_{\bar z}^2-2\varphi_{\bar z\bar z}\] obey \(\partial_{\bar z}T=\partial_z\bar T=0\). The printed plus sign instead gives \(\partial_{\bar z}T=m_0^2e^\varphi\varphi_z\), generally nonzero. An overall normalization or simultaneous \(\varphi\mapsto-\varphi\) with a changed action is possible, but cannot preserve all printed formulas.}
\GSfield{Reconstruction}{The action fixes the relative sign; the missing bar is a separate emendation confirmed by the scan.}
\GSfield{Fate}{The chain-rule computation \ours directly checks conservation against equation (6). The general holomorphic stress-tensor mechanism is explained in \citet[Section~2]{belavinpolyakovzamolodchikov1984infinite}.}
\GSfield{Present status}{As of November 2026: the qualified statement is established; the notebook is an exposition, not a newly open question.}
\GSfield{Assessment}{literal statement: false as written; publication match: partial; verification: read the result; development: annotated; confidence: high}
\end{GSitem}
\begin{GSitem}{1982-573.17}{E}{Recovering Virasoro modes from the stress field}
\GSfield{Kind; source}{construction; \GSpage{P001059}}
\GSfield{The notes}{The final formula is \(T(z)=\sum_{n\in\mathbb Z}z^{n-2}L_n\), followed by the assertion that this is where the \(L_n\) come from.}
\GSfield{Statement}{For the modern generators, the stress-field convention is \[\ell_n=\frac{1}{2\pi i}\oint z^{n+1}T(z)\,dz,\qquad T(z)=\sum_{n\in\mathbb Z}\ell_nz^{-n-2}.\] Since the earlier notebook generators are \(L_n=-\ell_n\), its final coefficients cannot be silently identified with these modes. Keeping the written exponent \(n-2\) requires reindexing coefficients; keeping the earlier notebook generator identification gives \(T(z)=-\sum_nL_nz^{-n-2}\). Classically these modes have Poisson brackets. Quantum Virasoro commutators require a quantum stress field and its operator product and normal ordering; a Laurent expansion alone does not provide them.}
\GSfield{Reconstruction}{Both the sign identification \(\ell_n=-L_n\) and the required coefficient reindexing are stated explicitly. Classical extraction and quantum commutators are separated.}
\GSfield{Fate}{\citet[equations~(2.18)--(2.21)]{belavinpolyakovzamolodchikov1984infinite} makes the mode construction explicit; the Virasoro algebra itself is earlier. The distinction between extraction and quantization prevents the final sentence from being read as a proof of the spectral conjecture.}
\GSfield{Present status}{As of November 2026: the qualified statement is established; the notebook is an exposition, not a newly open question.}
\GSfield{Assessment}{literal statement: false as written; publication match: exact; verification: read the result; development: annotated; confidence: high}
\end{GSitem}

\GStalk{574}{1982-11-01}{A. A. Beilinson}{P001059-P001062}{Regular singular D-modules and gluing}
The actual entry begins on P001060, after the preceding Liouville discussion. Gelfand recalls the Lorentz-group reduction to linear algebra; Beilinson then describes regular singular modules on a punctured disk and gluing across a closed subvariety. The final relation \(vu=N\) is shorthand, not a complete definition. Beilinson’s later paper explicitly says it follows S. I. Gelfand’s 1982 notes \citep[introduction]{beilinson1987glue}.

\begin{GSitem}{1982-574.1}{E}{Lorentz-group representations reduced to linear algebra}
\GSfield{Kind; source}{remark; \GSpage{P001060}; I. M. Gelfand}
\GSfield{The notes}{Gelfand cites his joint work with Ponomarev, approximately 1968, on indecomposable representations of the Lorentz group.}
\GSfield{Statement}{For the real Lie algebra of \(SL(2,\mathbb C)\), the category considered consists of \(\mathfrak{su}(2)\)-locally finite modules with finite multiplicity of each compact type; its indecomposable Harish--Chandra modules admit an explicit classification by linear-algebra data. This does not classify arbitrary Hilbert-space representations without admissibility assumptions.}
\GSfield{Reconstruction}{The Harish--Chandra and finite-multiplicity conditions are supplied from the cited paper.}
\GSfield{Fate}{\citet[introduction and classification]{gelfandponomarev1968lorentz} is the exact 1968 publication recalled in the intervention.}
\GSfield{Present status}{As of November 2026: the stated result is established in the specified setting.}
\GSfield{Assessment}{literal statement: correct with missing hypotheses; publication match: exact; verification: read the result; development: annotated; confidence: high}
\end{GSitem}
\begin{GSitem}{1982-574.2}{E}{The flag-line model for SL(2,R)}
\GSfield{Kind; source}{method; \GSpage{P001060}}
\GSfield{The notes}{For \(SL(2,\mathbb R)\), classify \(D\)-modules on \(\mathbb P^1\) with singularities at two fixed points.}
\GSfield{Statement}{The complexified maximal compact torus of \(SL(2,\mathbb R)\) has two fixed points and one open orbit on the complex flag variety \(\mathbb P^1\). Regular holonomic twisted \(D\)-modules constructible for this orbit stratification give the geometric model, after imposing an infinitesimal character and equivariance appropriate to the representation category.}
\GSfield{Reconstruction}{The torus orbit stratification and the conditions for applying localization are editorial specifications, not a claim that all SL(2,R) representations equal all two-singularity modules.}
\GSfield{Fate}{The localization framework predates the talk \citep{beilinsonbernstein1981localisation}; the notebook does not specify a block or the precise equivariance.}
\GSfield{Present status}{As of November 2026: the stated result is established in the specified setting.}
\GSfield{Assessment}{literal statement: correct with missing hypotheses; publication match: partial; verification: publisher metadata checked; development: annotated; confidence: medium}
\end{GSitem}
\begin{GSitem}{1982-574.3}{A}{A regular singularity as two linear maps}
\GSfield{Kind; source}{theorem; \GSpage{P001060}}
\GSfield{The notes}{On a disk, the category is described by \(\Psi\mathrel{\mathop{\rightleftarrows}^{u}_{v}}\Phi\) with \(uv-\mathrm{Id}\) invertible; outside zero there is a bundle with connection.}
\GSfield{Statement}{For a small complex disk \(\Delta\), regular holonomic \(D_\Delta\)-modules whose singular support is contained in the zero section and \(T^*_0\Delta\) correspond, via Riemann--Hilbert, to finite-dimensional vector spaces \(\Psi,\Phi\) and maps \(u:\Psi\to\Phi\), \(v:\Phi\to\Psi\) for which \(1-vu\), equivalently \(1-uv\), is invertible. Morphisms commute with both maps.}
\GSfield{Reconstruction}{Finite-dimensionality, holonomicity and the choice of Riemann--Hilbert convention have been made explicit; the sign is consistent with Beilinson’s convention.}
\GSfield{Fate}{The disk equivalence is \citet[Corollary 3.2]{beilinson1987glue}, written from 1982 notes and published in 1987.}
\GSfield{Present status}{As of November 2026: the stated result is established in the specified setting.}
\GSfield{Assessment}{literal statement: correct with missing hypotheses; publication match: exact; verification: read the result; development: annotated; confidence: high}
\end{GSitem}
\begin{GSitem}{1982-574.4}{E}{Ordinary sheaf gluing across a closed subset}
\GSfield{Kind; source}{construction; \GSpage{P001060}}
\GSfield{The notes}{For \(i:Y\hookrightarrow X\), \(j:U=X\setminus Y\hookrightarrow X\), write \(\mathcal F_X=(\mathcal F_U,\mathcal F_Y,\mathcal F_Y\to i^*j_*\mathcal F_U)\).}
\GSfield{Statement}{For sheaves of vector spaces on a topological space, restriction to an open complement and its closed subset, together with the indicated specialization map, determines the sheaf by a fiber-product gluing construction. The analogous assertion for perverse sheaves cannot be obtained by using the same abelian restriction functors unchanged.}
\GSfield{Reconstruction}{The abelian sheaf category and the distinction from perverse gluing are specified editorially.}
\GSfield{Fate}{This standard recollement is the motivating elementary comparison for \citet[section 3]{beilinson1987glue}.}
\GSfield{Present status}{As of November 2026: the stated result is established in the specified setting.}
\GSfield{Assessment}{literal statement: correct with missing hypotheses; publication match: partial; verification: read the result; development: annotated; confidence: high}
\end{GSitem}
\begin{GSitem}{1982-574.5}{A}{Nearby and vanishing cycles provide gluing data}
\GSfield{Kind; source}{theorem; \GSpage{P001060}}
\GSfield{The notes}{The final tuple consists of \(M_U,M_Y\), maps to and from \(\Psi(M_U)\), and \(vu=N\), followed by \ru{что-то} (something).}
\GSfield{Statement}{For a regular function \(f:X\to\mathbb A^1\), \(Y=f^{-1}(0)\), \(U=X\setminus Y\), perverse sheaves are equivalent to the category of Beilinson gluing data: a perverse sheaf on \(U\), one on \(Y\), and two maps involving unipotent nearby cycles whose composite is its monodromy operator \(N\), with the Tate twist in the algebraic formulation.}
\GSfield{Reconstruction}{The function f, unipotent nearby-cycle convention and the meaning of N are supplied from the later text; the notebook alone does not define them.}
\GSfield{Fate}{\citet[section 3 and Proposition 3.1]{beilinson1987glue} supplies the precise equivalence; its 1982 circulation is explicit.}
\GSfield{Present status}{As of November 2026: the stated result is established in the specified setting.}
\GSfield{Assessment}{literal statement: correct with missing hypotheses; publication match: partial; verification: read the result; development: annotated; confidence: high}
\end{GSitem}

\GStalk{575}{1982-11-15}{I. M. Gelfand}{P001060-P001062}{Continuous-dimensional Grassmannians and a Schubert analogy}
The intervention occupies the top of P001061. Gelfand recalls finite Grassmannian cells, then asks about fixed-trace projections in a type \(\mathrm{II}_1\) factor, continuous flags, monotone parameters, topology and cohomology. The notebook expressly calls the problem imprecise, so these questions are preserved as a programme rather than assigned a fabricated solved/open theorem.

\begin{GSitem}{1982-575.1}{E}{Schubert cells and Grassmannian cohomology}
\GSfield{Kind; source}{theorem; \GSpage{P001061}}
\GSfield{The notes}{For \(G_{n,k}\), a cell is determined by \(d_j=\dim(V\cap F_j)\); its dimension is written \(\sum d_j\).}
\GSfield{Statement}{For a complete flag in \(\mathbb C^n\), let \(1\le a_1<\cdots<a_k\le n\) be the jump positions of \(d_j\). The corresponding Schubert cell is \(\mathbb C^{\sum_i(a_i-i)}\). The even real dimensional cells give a free integral cohomology basis by the dual Schubert classes. The written sum of the intersection dimensions is not this cell dimension.}
\GSfield{Reconstruction}{The field, jump convention and correct dimension are supplied; the notebook’s formula is retained in the correction ledger.}
\GSfield{Fate}{The classical cell formula is reviewed in \citet[section 47.2, Exercise 47.4 and Corollary 47.7]{etingof2024lie}; it also occurred as 1971-239.1--2 in the supplied index.}
\GSfield{Present status}{As of November 2026: the stated result is established in the specified setting.}
\GSfield{Assessment}{literal statement: false as written; publication match: exact; verification: read the result; development: annotated; confidence: high}
\end{GSitem}
\begin{GSitem}{1982-575.2}{E}{Fixed-trace projections as continuous Grassmannians}
\GSfield{Kind; source}{definition; \GSpage{P001061}}
\GSfield{The notes}{In a type \(\mathrm{II}_1\) factor, consider \(\alpha\)-dimensional projections.}
\GSfield{Statement}{For a type \(\mathrm{II}_1\) factor \((M,\tau)\) with normalized trace and \(0\le\alpha\le1\), define \(\mathrm{Gr}_\alpha(M)=\{p\in M:p=p^*=p^2,\ \tau(p)=\alpha\}\). Projections of the same trace are unitarily equivalent. The definition does not itself choose a topology on this set.}
\GSfield{Reconstruction}{A normalized finite trace and the projection equation are made explicit.}
\GSfield{Fate}{Continuous projection dimension is classical von Neumann geometry \citep{vonneumann1960continuous}; earlier factor background occurs at 1969-179.8--10.}
\GSfield{Present status}{As of November 2026: the stated result is established in the specified setting.}
\GSfield{Assessment}{literal statement: correct with missing hypotheses; publication match: partial; verification: publisher metadata checked; development: annotated; confidence: high}
\end{GSitem}
\begin{GSitem}{1982-575.3}{F}{Schubert-type strata for a continuous flag}
\GSfield{Kind; source}{question; \GSpage{P001061}}
\GSfield{The notes}{\ru{Есть ли аналог клетки Шуберта?} A continuous flag is indexed by \([0,1]\); the proposed description uses monotone functions.}
\GSfield{Statement}{Given a continuous increasing projection flag \((e_t)_{0\le t\le1}\) with \(\tau(e_t)=t\), one candidate invariant of \(p\in\mathrm{Gr}_\alpha(M)\) is \(t\mapsto\tau(p\wedge e_t)\). The notebook does not choose an acting group or say whether equal functions should define an orbit, a cell, or a closure class; hence no classification theorem is determined.}
\GSfield{Reconstruction}{The meet-trace function is a natural reconstruction of the monotonicity hint, not a uniquely recoverable formula.}
\GSfield{Fate}{The continuous-geometry background is identifiable \citep{vonneumann1960continuous}; the exact proposed Schubert equivalence relation is not.}
\GSfield{Present status}{As of November 2026: the notebook does not determine a unique proposition whose resolution can be assessed.}
\GSfield{Assessment}{literal statement: not determinable; publication match: thematic only; verification: unresolved lead; development: annotated; confidence: low}
\end{GSitem}
\begin{GSitem}{1982-575.4}{F}{Topology and cohomology of the proposed continuous Grassmannian}
\GSfield{Kind; source}{question; \GSpage{P001061}}
\GSfield{The notes}{\ru{какая топология? и какие когомологии?} follows \ru{Непр. когомологии!?}.}
\GSfield{Statement}{The question asks for a topology and a cohomology theory on fixed-trace projection spaces compatible with the continuous-flag analogy. Norm, strong-operator and trace-metric topologies lead to different questions; none is selected in the source.}
\GSfield{Fate}{No unique statement can be matched to a publication until the missing choices are made; continuous geometry alone does not supply the requested cohomology.}
\GSfield{Present status}{As of November 2026: this is an underdetermined historical programme, not evidence that every precise topology/cohomology variant is open.}
\GSfield{Assessment}{literal statement: not determinable; publication match: none found; verification: unresolved lead; development: annotated; confidence: low}
\end{GSitem}

\GStalk{576}{1982-11-15}{Yu. I. Manin}{P001060-P001065}{Twistor geometry and flag superspaces}
The entry occupies P001061--P001063. After an agenda extending to supersymmetric field equations, Manin recalls the Grassmannian model of compactified complex Minkowski space, formulates a classification problem for flag superspaces, and explains the functor of points and two examples. The surviving text stops at the opening of the Penrose correspondence; it contains no worked new Yang--Mills solution. The closely related 1983 chapter is identifiable, but its full classification has not been verified here \citep{manin1983flags}.

\begin{GSitem}{1982-576.1}{A}{The announced supertwistor programme}
\GSfield{Kind; source}{method; \GSpage{P001061}}
\GSfield{The notes}{The agenda lists flag supermanifolds, a super Radon--Penrose transform, non-self-dual Yang--Mills solutions, deformations and supergravity, then the index and characteristic classes. Only the first four are promised.}
\GSfield{Statement}{This is a programme for encoding gauge fields by holomorphic data on superflag correspondences. No specific non-self-dual solution, deformation theorem, or index formula is stated in the surviving notes.}
\GSfield{Fate}{The programme has a close contemporary publication in \citet[abstract]{manin1983flags}; the ordinary holomorphic framework had already appeared in \citet{manin1981holomorphic}.}
\GSfield{Present status}{As of November 2026: the programme and its publication are identified; the notebook does not justify adding any absent solution or index theorem.}
\GSfield{Assessment}{literal statement: not determinable; publication match: partial; verification: publisher metadata checked; development: annotated; confidence: medium}
\end{GSitem}
\begin{GSitem}{1982-576.2}{E}{Grassmannian compactification and null incidence}
\GSfield{Kind; source}{construction; \GSpage{P001061-P001062}}
\GSfield{The notes}{\(T=\mathbb C^4\), \(G(2,T)\) is compact complex Minkowski space; there are three-dimensional cones, two kinds of two-dimensional null planes, and null lines.}
\GSfield{Statement}{An affine chart of \(G(2,4)\) consists of graphs of \(2\times2\) matrices. The determinant gives the flat complex conformal quadratic form. At a point its rank-one tangent vectors form a three-dimensional null cone, ruled by the two families of totally null two-planes. A suitable real form and the affine chart recover real Minkowski space.}
\GSfield{Reconstruction}{The choice of real form, graph chart and determinant are supplied from the earlier source. The unexplained antiholomorphic involution on T plus T-dual is not promoted to a unique formula.}
\GSfield{Fate}{\citet[section~1, paragraph~2]{manin1981holomorphic} describes the graph construction before this talk.}
\GSfield{Present status}{As of November 2026: the stated result is established in the specified setting.}
\GSfield{Assessment}{literal statement: correct with missing hypotheses; publication match: exact; verification: read the result; development: annotated; confidence: high}
\end{GSitem}
\begin{GSitem}{1982-576.3}{E}{Tautological bundles and spinors}
\GSfield{Kind; source}{construction; \GSpage{P001062}}
\GSfield{The notes}{The diagram places two rank-two bundles above \(G(2,T)\).}
\GSfield{Statement}{On \(G(2,T)\), the tautological bundle \(S\) and quotient bundle \(Q=T/S\) have rank two and \(TG(2,T)\simeq S^*\otimes Q\). On a chart and after the relevant real and determinant choices, these are the two spinor factors of the complexified tangent bundle.}
\GSfield{Reconstruction}{The bundle identification is a standard reconstruction of the diagram; its handwritten bundle letters do not by themselves fix dual conventions.}
\GSfield{Fate}{This is the bundle geometry underlying \citet{manin1981holomorphic}.}
\GSfield{Present status}{As of November 2026: the stated result is established in the specified setting.}
\GSfield{Assessment}{literal statement: correct with missing hypotheses; publication match: partial; verification: read the result; development: annotated; confidence: high}
\end{GSitem}
\begin{GSitem}{1982-576.4}{A}{Flag superspaces with reduced space G(2,4)}
\GSfield{Kind; source}{problem; \GSpage{P001062-P001063}}
\GSfield{The notes}{\ru{Описать} homogeneous complex supermanifolds with reduced manifold \(G(2,T)\), requiring flags for simple Lie supergroups of classical type; only the beginning of a list follows.}
\GSfield{Statement}{Classify flag superspaces of classical simple complex Lie supergroups whose reduced complex manifold is isomorphic to \(G(2,4)\), specifying equivalence as homogeneous superspaces and the chosen class of flag stabilizers. The notebook supplies two families but does not claim its displayed list is complete.}
\GSfield{Reconstruction}{Equivalence as homogeneous superspaces is made explicit; an underlying supermanifold classification could identify different homogeneous presentations.}
\GSfield{Fate}{\citet{manin1983flags} is the close contemporary account; later general flag structure appears in \citet[sections~5--6]{voronovmanin1990supercells}. The complete restricted list was not read in the 1983 source and is not certified by this annotation.}
\GSfield{Present status}{As of November 2026: the classification programme has an identifiable 1983 publication; completeness of the precise restricted list remains unverified in this run, which is a verification limit, not an asserted open problem.}
\GSfield{Assessment}{literal statement: correct with missing hypotheses; publication match: partial; verification: publisher metadata checked; development: annotated; confidence: medium; problem \hyperref[prob:GP-1982-20]{GP-1982-20}}
\end{GSitem}
\begin{GSitem}{1982-576.5}{E}{Superflags require the functor of points}
\GSfield{Kind; source}{definition; \GSpage{P001062}}
\GSfield{The notes}{A naive flag of complex supervector spaces has too few coordinates; for every supercommutative algebra \(A\), use projective direct submodules and base change \(A\to B\).}
\GSfield{Statement}{The functor \(A\mapsto\{N^{2|0}\subset N^{2|N}\subset A^{4|N}\}\), with locally free quotients and locally direct summands, is represented by the flag superscheme. Scalar extension defines its covariant action on superalgebras. Restriction to ordinary complex points discards all odd directions.}
\GSfield{Reconstruction}{Locally free quotient conditions clarify the notebook’s projective direct-submodule shorthand.}
\GSfield{Fate}{The standard representable-functor definition is used in \citet[sections~1--2]{voronovmanin1990supercells}; it is background, rather than a new 1988 discovery.}
\GSfield{Present status}{As of November 2026: the stated result is established in the specified setting.}
\GSfield{Assessment}{literal statement: correct with missing hypotheses; publication match: exact; verification: read the result; development: annotated; confidence: high}
\end{GSitem}
\begin{GSitem}{1982-576.6}{A}{The standard N-extended superspace flag}
\GSfield{Kind; source}{construction; \GSpage{P001062}}
\GSfield{The notes}{\(F(2|0,2|N;\mathbb C^{4|N})\) has dimension \(4|4N\); its real big cell is the physicists\textquotesingle\ \(N\)-extended superspace.}
\GSfield{Statement}{The indicated flag superspace has reduced manifold \(G(2,4)\), even dimension \(4\) and odd dimension \(4N\). Its big cell admits four even coordinates and two sets of \(2N\) odd spinor coordinates. A real superspace requires an appropriate real structure; its compact flag ambient space supplies the stated conformal completion.}
\GSfield{Reconstruction}{The real structure is an extra choice, not specified by the complex flag alone.}
\GSfield{Fate}{This is the concrete superflag setting of \citet{manin1983flags}; the flag formalism is checked against \citet[section~1]{voronovmanin1990supercells}.}
\GSfield{Present status}{As of November 2026: the stated result is established in the specified setting.}
\GSfield{Assessment}{literal statement: correct with missing hypotheses; publication match: partial; verification: publisher metadata checked; development: annotated; confidence: medium}
\end{GSitem}
\begin{GSitem}{1982-576.7}{E}{Projective supercoordinates and their big cells}
\GSfield{Kind; source}{construction; \GSpage{P001062}}
\GSfield{The notes}{\(G(1|0,\mathbb C^{m|n})=\mathbb P^{m-1|n}\), with ratios by \(x_i\); the last odd ratio is written \(\xi_m/x_i\).}
\GSfield{Statement}{On the chart \(x_i\ne0\), the independent coordinates are \(x_j/x_i\) for \(j\ne i\) and \(\xi_a/x_i\) for \(1\le a\le n\). Thus each chart has dimension \((m-1)|n\), and these charts are the simplest superflag big cells.}
\GSfield{Reconstruction}{The omitted even coordinate equals 1; the last odd index is corrected from m to n.}
\GSfield{Fate}{The projective-bundle construction is part of the established superflag formalism \citep[section~1]{voronovmanin1990supercells}.}
\GSfield{Present status}{As of November 2026: the stated result is established in the specified setting.}
\GSfield{Assessment}{literal statement: false as written; publication match: exact; verification: read the result; development: annotated; confidence: high}
\end{GSitem}
\begin{GSitem}{1982-576.8}{F}{Classes of odd-codimension supercells}
\GSfield{Kind; source}{question; \GSpage{P001062}}
\GSfield{The notes}{An equation \(\xi_i=0\) gives codimension \(0|1\). The notes ask what cohomology contains such cells; ordinary differential and integral forms serve only extreme dimensions, while Schwarz’s intermediate objects have strange properties.}
\GSfield{Statement}{The local subsupermanifold cut out by one independent odd coordinate has codimension \(0|1\). The question seeks a cohomology or intersection theory representing such supercycles, but does not specify admissible cycles, equivalence, degree conventions, or required operations.}
\GSfield{Fate}{\citet[introduction]{voronovmanin1990supercells} still explains the need for a suitable theory in 1988 and names candidate constructions. That historical discussion does not establish either a universal solution or present-day openness of every formulation.}
\GSfield{Present status}{As of November 2026: the requested theory is insufficiently specified to attach a single resolved/open label.}
\GSfield{Assessment}{literal statement: not determinable; publication match: thematic only; verification: read the result; development: annotated; confidence: medium}
\end{GSitem}
\begin{GSitem}{1982-576.9}{A}{The odd-form isotropic example}
\GSfield{Kind; source}{construction; \GSpage{P001062-P001063}}
\GSfield{The notes}{The second entry is \(FI(2|0,2|2;\mathbb C^{4|4},\Pi\mathrm{Sp})\), dimension \(4|8\), with an odd symmetric form; further series and isolated examples are merely mentioned.}
\GSfield{Statement}{The notebook specifies an isotropic flag superspace for a nondegenerate odd symmetric form, with the displayed dimension. A complete interpretation must fix the periplectic form and isotropy conventions; the unnamed further examples cannot be reconstructed as an exhaustive list.}
\GSfield{Reconstruction}{The modern term periplectic is supplied for the odd symmetric form; no absent series is inserted.}
\GSfield{Fate}{The example belongs to the programme of \citet{manin1983flags}. The distinction between classical types and their isotropic flag conventions is developed in \citet[sections~2,5]{voronovmanin1990supercells}.}
\GSfield{Present status}{As of November 2026: the stated result is established in the specified setting.}
\GSfield{Assessment}{literal statement: correct with missing hypotheses; publication match: partial; verification: publisher metadata checked; development: annotated; confidence: medium}
\end{GSitem}
\begin{GSitem}{1982-576.10}{E}{The Penrose correspondence starts but is not recorded}
\GSfield{Kind; source}{method; \GSpage{P001063}}
\GSfield{The notes}{The last two symbols are \(M=G(2,T)\), \(P(T)=G(1,T)\).}
\GSfield{Statement}{These are the space-time and projective twistor spaces. Their usual incidence space is \(F(1,2;T)\), with projections to both Grassmannians; this incidence space is editorial reconstruction. No transform, cohomology degree, or field equation is stated on the surviving final page.}
\GSfield{Reconstruction}{The incidence space identifies the geometric context but is not drawn or defined in the final fragment.}
\GSfield{Fate}{The ordinary correspondence predates the talk \citep{manin1981holomorphic}; the fragment does not license importing the unrecorded supertransform.}
\GSfield{Present status}{As of November 2026: the stated result is established in the specified setting.}
\GSfield{Assessment}{literal statement: correct with missing hypotheses; publication match: partial; verification: read the result; development: annotated; confidence: high}
\end{GSitem}

\GStalk{577}{1982-11-22}{A. M. Polyakov}{P001063-P001071}{Liouville uniformization, accessory parameters and Virasoro blocks}
The talk occupies P001064--P001069. It moves from the Liouville action and cusp uniformization to the stress tensor, accessory parameters, and the proposed recovery of classical geometry from quantum Ward identities. The last pages give a Verma-module expansion and an oscillator realization. The classical generating-function conjecture was proved by Zograf and Takhtajan; the conformal-block programme appears in BPZ 1984, while rigorous sphere Liouville construction and bootstrap came much later. Several factors, powers and completeness relations in the notebook are inconsistent and are explicitly corrected below.

\begin{GSitem}{1982-577.1}{E}{Liouville action from the string conformal factor}
\GSfield{Kind; source}{construction; \GSpage{P001064}}
\GSfield{The notes}{In isothermal coordinates \(g_{ab}=e^\varphi\delta_{ab}\), \(S=\gamma^{-1}\int(\tfrac12|\nabla\varphi|^2+e^\varphi)\), with \(\gamma^{-1}=(26-D)/(48\pi)\).}
\GSfield{Statement}{The displayed expression is the formal conformal-gauge effective action of the bosonic string, with a choice of cosmological constant absorbed into \(\varphi\). It is an anomaly calculation with regularization and gauge fixing, not a definition of an infinite-dimensional Lebesgue measure. The coefficient is the notebook’s classical normalization.}
\GSfield{Reconstruction}{The words formal, effective, and the normalization qualification separate the physics derivation from a rigorous measure construction.}
\GSfield{Fate}{Polyakov’s 1981 construction is explicitly identified as the starting point by \citet[introduction]{davidkupiainenrhodesvargas2016liouville}. Their rigorous formulation also explains the quantum background-charge shift.}
\GSfield{Present status}{As of November 2026: the conformal anomaly construction is established physics; a rigorous subcritical sphere theory exists, without thereby solving the full bosonic string problem.}
\GSfield{Assessment}{literal statement: correct with missing hypotheses; publication match: partial; verification: read the result; development: annotated; confidence: high}
\end{GSitem}
\begin{GSitem}{1982-577.2}{E}{Conformal covariance and the purported unique local action}
\GSfield{Kind; source}{claim; \GSpage{P001064}}
\GSfield{The notes}{\(e^{\widetilde\varphi}=|w^{\prime}|^2e^{\varphi\circ w}\); the Liouville Lagrangian is called the only invariant local Lagrangian.}
\GSfield{Statement}{The stated transformation is the pullback law for a conformal metric, for locally univalent holomorphic \(w\). The kinetic-plus-exponential action has the appropriate covariance up to boundary and background terms. Without a derivative-order restriction the uniqueness claim is false: integrals of local scalar curvature invariants supply additional actions.}
\GSfield{Reconstruction}{No uniqueness theorem is asserted without a fixed order, boundary convention and admissible class of Lagrangians.}
\GSfield{Fate}{The correctly qualified covariance is part of the classical and quantum formulations \citep[section~1]{zograftakhtajan1988liouville}; \citep[Theorem~3.11]{davidkupiainenrhodesvargas2016liouville}.}
\GSfield{Present status}{As of November 2026: the stated result is established in the specified setting.}
\GSfield{Assessment}{literal statement: false as written; publication match: partial; verification: read the result; development: annotated; confidence: high}
\end{GSitem}
\begin{GSitem}{1982-577.3}{E}{Liouville equation and the developing map}
\GSfield{Kind; source}{theorem; \GSpage{P001064-P001065}}
\GSfield{The notes}{\(\Delta\varphi=e^\varphi\), with \(e^\varphi=|w^{\prime}|^2/(\operatorname{Im}w)^2\).}
\GSfield{Statement}{For the Euclidean \(\Delta=\partial_x^2+\partial_y^2\), a locally univalent holomorphic map into the upper half-plane gives \(e^\varphi=2|w^{\prime}|^2/(\operatorname{Im}w)^2\) satisfying \(\Delta\varphi=e^\varphi\). Without the factor 2 it satisfies \(\Delta\varphi=2e^\varphi\). The corresponding curvatures are respectively \(-1/2\) and \(-1\).}
\GSfield{Reconstruction}{The Laplacian convention is fixed to exhibit the notebook’s factor mismatch; w must have nonzero derivative and nonzero imaginary part.}
\GSfield{Fate}{This is classical uniformization geometry, reviewed explicitly as old material by \citet[section~1]{zograftakhtajan1988liouville}.}
\GSfield{Present status}{As of November 2026: the stated result is established in the specified setting.}
\GSfield{Assessment}{literal statement: false as written; publication match: exact; verification: read the result; development: annotated; confidence: high}
\end{GSitem}
\begin{GSitem}{1982-577.4}{E}{Complete cusp metrics on punctured spheres}
\GSfield{Kind; source}{construction; \GSpage{P001065}}
\GSfield{The notes}{A sphere with \(N\) horns has \(e^\varphi\sim2/(|z-z_i|^2\log^2|z-z_i|)\).}
\GSfield{Statement}{For \(N\ge3\) distinct punctures on \(\mathbb P^1\), there is a unique complete finite-area metric of curvature \(-1/2\) in its conformal class. In a puncture coordinate it has the displayed leading cusp asymptotic. Three points may be fixed at \(0,1,\infty\).}
\GSfield{Reconstruction}{The number of punctures, completeness and curvature normalization are supplied explicitly.}
\GSfield{Fate}{The classical complete hyperbolic setting and puncture asymptotics are reviewed in \citet[section~1]{zograftakhtajan1988liouville}.}
\GSfield{Present status}{As of November 2026: the stated result is established in the specified setting.}
\GSfield{Assessment}{literal statement: correct with missing hypotheses; publication match: exact; verification: read the result; development: annotated; confidence: high}
\end{GSitem}
\begin{GSitem}{1982-577.5}{B}{Quantum correlations and the classical saddle}
\GSfield{Kind; source}{method; \GSpage{P001065-P001067}}
\GSfield{The notes}{\(\Psi=\int\mathcal D\varphi e^{-S}\sim e^{-S_{cl}}\) as \(\gamma\to0\); later lines write this as an equality involving \(\Delta S_{cl}\).}
\GSfield{Statement}{The first formula is a formal saddle-point prescription with singular insertions and a regularized action. The later equality is only a schematic semiclassical statement: quantum correlations include fluctuation terms, and a scaling of insertion charges must be specified.}
\GSfield{Reconstruction}{The probabilistic parameter gamma is a modern convention and is not equated numerically with the notebook’s coupling.}
\GSfield{Fate}{\citet[Theorem~3.2]{davidkupiainenrhodesvargas2016liouville}, first circulated in 2014, constructs sphere correlations for coupling \(0<\gamma<2\), \(Q=2/\gamma+\gamma/2\), charges \(\alpha_i<Q\) and \(\sum\alpha_i>2Q\). It supplies a rigorous framework, not the unqualified exact saddle identity.}
\GSfield{Present status}{As of November 2026: rigorous subcritical sphere correlations exist; the notebook’s unscaled equality with a classical exponential is not an exact quantum formula.}
\GSfield{Assessment}{literal statement: correct with missing hypotheses; publication match: partial; verification: read the result; development: annotated; confidence: high}
\end{GSitem}
\begin{GSitem}{1982-577.6}{E}{Stress tensor and Schwarzian normalization}
\GSfield{Kind; source}{claim; \GSpage{P001065}}
\GSfield{The notes}{\(T=\tfrac12\varphi_z^2-\varphi_{zz}\) is holomorphic and is identified with a displayed Schwarzian having denominators \((w^{\prime})^2,(w^{\prime})^4\).}
\GSfield{Statement}{On a Liouville solution, \(\partial_{\bar z}T=0\). With the preceding developing map, \(T=-\{w,z\}\), where \(\{w,z\}=w^{\prime\prime\prime}/w^{\prime}-\tfrac32(w^{\prime\prime}/w^{\prime})^2\). Thus the written Schwarzian denominators and the sign identification are inconsistent.}
\GSfield{Reconstruction}{A new symbol q below separates the ODE coefficient from this stress tensor.}
\GSfield{Fate}{See \citet[section~1]{zograftakhtajan1988liouville} for the normalization of the projective connection.}
\GSfield{Present status}{As of November 2026: the stated result is established in the specified setting.}
\GSfield{Assessment}{literal statement: false as written; publication match: exact; verification: read the result; development: annotated; confidence: high}
\end{GSitem}
\begin{GSitem}{1982-577.7}{E}{Fuchsian equation and the accessory parameters}
\GSfield{Kind; source}{construction; \GSpage{P001065}}
\GSfield{The notes}{\(T=\tfrac14\sum(z-z_i)^{-2}+\sum c_i(z-z_i)^{-1}\), \(\psi^{\prime\prime}+T\psi=0\), \(w=\psi_1/\psi_2\); the \(c_i\) are said to be real.}
\GSfield{Statement}{Writing \(q=\tfrac12\{w,z\}\), the ratio of independent solutions of \(\psi^{\prime\prime}+q\psi=0\) is a developing map. Cusp singularities give double-pole coefficient \(1/4\); the simple-pole coefficients are accessory parameters, subject to the conditions at infinity. For arbitrary complex punctures they need not be real. Fuchsian uniformization selects monodromy conjugate into \(PSL(2,\mathbb R)\) with the appropriate developing-map geometry.}
\GSfield{Reconstruction}{q equals minus one-half of the stress tensor in the preceding item; the notebook uses T for both.}
\GSfield{Fate}{\citet[section~1.1]{zograftakhtajan1988liouville} supplies the classical formulation. Real monodromy alone, without the uniformizing branch and geometry, should not be read as a uniqueness assertion for arbitrary projective structures.}
\GSfield{Present status}{As of November 2026: the stated result is established in the specified setting.}
\GSfield{Assessment}{literal statement: false as written; publication match: exact; verification: read the result; development: annotated; confidence: high}
\end{GSitem}
\begin{GSitem}{1982-577.8}{E}{Varying punctures while preserving Fuchsian reality}
\GSfield{Kind; source}{remark; \GSpage{P001065}}
\GSfield{The notes}{When the punctures vary, monodromy changes but remains real, in \(SL(2,\mathbb R)\).}
\GSfield{Statement}{The family of uniformizing projective connections over punctured-sphere moduli has Fuchsian monodromy. This condition allows the conjugacy class of the representation to vary. It is therefore different from an isomonodromic deformation with a fixed monodromy representation.}
\GSfield{Reconstruction}{The projective monodromy and marking convention are supplied.}
\GSfield{Fate}{The deformation setting is made precise in \citet[sections~1--2]{zograftakhtajan1988liouville}.}
\GSfield{Present status}{As of November 2026: the stated result is established in the specified setting.}
\GSfield{Assessment}{literal statement: correct with missing hypotheses; publication match: exact; verification: read the result; development: annotated; confidence: high}
\end{GSitem}
\begin{GSitem}{1982-577.9}{B}{The Liouville action generates accessory parameters}
\GSfield{Kind; source}{conjecture; \GSpage{P001066-P001068}}
\GSfield{The notes}{The central claim is \(c_i=\partial S_{cl}/\partial z_i\), to be obtained through quantum theory and the classical limit.}
\GSfield{Statement}{On the configuration space of \(N\ge4\) distinct punctures with three fixed, the suitably regularized classical Liouville action generates the accessory parameters. In the Zograf--Takhtajan convention for the full Schwarzian projective connection, \(c_i=-(2\pi)^{-1}\partial S/\partial z_i\). The ODE coefficient q uses half these coefficients.}
\GSfield{Reconstruction}{The regularization, fixed punctures, sign and constant are indispensable; the notebook’s S and T normalizations do not determine them consistently.}
\GSfield{Fate}{Proved in \citet{zograftakhtajan1985action} and in detail in \citet[section~3, Theorem~1]{zograftakhtajan1988liouville}: announcement in 1985, Russian full paper in 1987, English in 1988. The latter introduction explicitly identifies Polyakov’s conjecture.}
\GSfield{Present status}{As of November 2026: the punctured-sphere generating-function conjecture is proved in the stated normalization.}
\GSfield{Assessment}{literal statement: correct with missing hypotheses; publication match: exact; verification: read the result; development: annotated; confidence: high; problem \hyperref[prob:GP-1982-21]{GP-1982-21}}
\end{GSitem}
\begin{GSitem}{1982-577.10}{E}{Canonical momentum in the chiral calculation}
\GSfield{Kind; source}{claim; \GSpage{P001066}}
\GSfield{The notes}{For the kinetic product \(\varphi_z\varphi_{\bar z}\), the derivative with respect to \(\varphi_{\bar z}\) is written \(p=\varphi_{\bar z}\), followed by a canonical bracket.}
\GSfield{Statement}{Treating the two derivatives as independent formal variables gives \(\partial\mathcal L/\partial\varphi_{\bar z}=\varphi_z\), with the action’s overall scale. A genuine equal-time Poisson bracket also requires a real slice, time variable and symplectic form; the two-dimensional complex-coordinate shorthand does not specify these.}
\GSfield{Reconstruction}{Only the formal derivative is corrected; a canonical structure absent from the notes is not invented.}
\GSfield{Fate}{The algebraic differentiation is elementary; the intended stress-tensor algebra has the established conformal-field-theory form \citep[section~2]{belavinpolyakovzamolodchikov1984infinite}.}
\GSfield{Present status}{As of November 2026: the stated result is established in the specified setting.}
\GSfield{Assessment}{literal statement: false as written; publication match: partial; verification: read the result; development: annotated; confidence: high}
\end{GSitem}
\begin{GSitem}{1982-577.11}{E}{The Virasoro cocycle and Laurent modes}
\GSfield{Kind; source}{theorem; \GSpage{P001066}}
\GSfield{The notes}{The stress bracket has a third derivative of delta. The mode formula omits \(\delta_{n+m,0}\), and the Laurent series omits conformal weight 2.}
\GSfield{Statement}{With \(T(z)=\sum_{n\in\mathbb Z}L_nz^{-n-2}\), the modes satisfy \([L_n,L_m]=(n-m)L_{n+m}+\tfrac{c}{12}(n^3-n)\delta_{n+m,0}\). The linear term in n depends on the chosen zero-mode shift. The third-derivative distribution is the local representative of the central cocycle.}
\GSfield{Fate}{The cohomology predates the lecture \citep{gelfandfuchs1968circle}; the exact mode convention is printed in \citet[section~2, equations~2.20--2.21]{belavinpolyakovzamolodchikov1984infinite}.}
\GSfield{Present status}{As of November 2026: the stated result is established in the specified setting.}
\GSfield{Assessment}{literal statement: false as written; publication match: exact; verification: read the result; development: annotated; confidence: high}
\end{GSitem}
\begin{GSitem}{1982-577.12}{B}{Anomaly coefficient versus quantum central charge}
\GSfield{Kind; source}{claim; \GSpage{P001066}}
\GSfield{The notes}{The notebook appends \(c=26-D\) to the Virasoro bracket.}
\GSfield{Statement}{For a consistent bosonic string, the matter, Liouville and ghost central charges must sum to zero, giving \(c_L=26-c_{matter}\). In quantum Liouville theory \(c_L=1+6Q^2\). This is not obtained by identifying an unrenormalized classical stress coefficient with the final quantum central charge without specifying the background-charge and field normalization.}
\GSfield{Reconstruction}{This separates the anomaly-balance claim from the quantum renormalization not recorded in the notebook.}
\GSfield{Fate}{\citet[Theorem~3.11]{davidkupiainenrhodesvargas2016liouville} proves the Weyl anomaly in the rigorous subcritical sphere construction. The interpretation of 26 originates in the pre-talk string calculation.}
\GSfield{Present status}{As of November 2026: the indicated anomaly formula is established for the specified Liouville construction; no unrestricted real-coupling construction for arbitrary D is inferred.}
\GSfield{Assessment}{literal statement: correct with missing hypotheses; publication match: partial; verification: read the result; development: annotated; confidence: high}
\end{GSitem}
\begin{GSitem}{1982-577.13}{A}{Primary fields and Ward operators}
\GSfield{Kind; source}{definition; \GSpage{P001066-P001069}}
\GSfield{The notes}{\(\delta_\varepsilon\psi=\varepsilon\psi^{\prime}+\Delta\varepsilon^{\prime}\psi\), but the mode relation repeatedly lacks \(z^n\) in the weight term.}
\GSfield{Statement}{A chiral primary of weight h obeys \([L_n,\psi(z)]=(z^{n+1}\partial_z+(n+1)hz^n)\psi(z)\). The local generator is the contour integral \((2\pi i)^{-1}\oint\varepsilon(z)T(z)\,dz\), not the unqualified area integral printed in the notes.}
\GSfield{Reconstruction}{The weight is denoted h to distinguish it from the Laplacian and the notebook’s other uses of Delta.}
\GSfield{Fate}{The primary-field formalism is part of \citet[sections~2--3]{belavinpolyakovzamolodchikov1984infinite}, received November 1983 and published in 1984.}
\GSfield{Present status}{As of November 2026: the stated result is established in the specified setting.}
\GSfield{Assessment}{literal statement: false as written; publication match: exact; verification: read the result; development: annotated; confidence: high}
\end{GSitem}
\begin{GSitem}{1982-577.14}{A}{Highest-weight modules and the four-point cross-ratio}
\GSfield{Kind; source}{construction; \GSpage{P001067-P001068}}
\GSfield{The notes}{Normalize four points to \(\infty,1,x,0\); \(|0\rangle=\psi(0)|v\rangle\) has \(L_n|0\rangle=0\) for \(n>0\), \(L_0|0\rangle=\Delta|0\rangle\).}
\GSfield{Statement}{Möbius covariance reduces a sphere four-point function to a function of the cross-ratio and its complex conjugate, times known coordinate factors. A chiral intermediate channel is modeled by a highest-weight module; its universal version is a Verma module. An actual irreducible theory may quotient null states and may require several or continuously many channels.}
\GSfield{Reconstruction}{The distinction between a chiral block and a full correlator is essential; the notebook calls both Psi.}
\GSfield{Fate}{\citet[sections~3--4]{belavinpolyakovzamolodchikov1984infinite} develops this algebraic framework.}
\GSfield{Present status}{As of November 2026: the stated result is established in the specified setting.}
\GSfield{Assessment}{literal statement: correct with missing hypotheses; publication match: exact; verification: read the result; development: annotated; confidence: high}
\end{GSitem}
\begin{GSitem}{1982-577.15}{A}{Descendant expansion requires the inverse Gram matrix}
\GSfield{Kind; source}{method; \GSpage{P001067-P001068}}
\GSfield{The notes}{\(I=\sum|n\rangle\langle n|\), later \(|n\rangle=L_{-n}|0\rangle\); the next line acknowledges nonorthogonality.}
\GSfield{Statement}{At level N, descendants are indexed by all partitions \(\lambda\vdash N\), not just one vector. If their contravariant Gram matrix \(G_N\) is invertible, the level projector is \(\sum_{\lambda,\mu}|\lambda\rangle(G_N^{-1})_{\lambda\mu}\langle\mu|\). At degenerate parameters pass to a suitable quotient or formulate a meromorphic block.}
\GSfield{Reconstruction}{This is the missing linear-algebra operation, not a claim that the author overlooked the problem; the notes explicitly register it.}
\GSfield{Fate}{The block machinery is explicit in \citet[section~4]{belavinpolyakovzamolodchikov1984infinite}; the contravariant-form and singular-vector structure was already treated in \citet[section~3]{feiginfuchs1982invariant}.}
\GSfield{Present status}{As of November 2026: the stated result is established in the specified setting.}
\GSfield{Assessment}{literal statement: false as written; publication match: exact; verification: read the result; development: annotated; confidence: high}
\end{GSitem}
\begin{GSitem}{1982-577.16}{A}{The first equal-weight block coefficient}
\GSfield{Kind; source}{claim; \GSpage{P001067}}
\GSfield{The notes}{The commutator gives \(\langle1|\psi(x)|0\rangle=\Delta x^{1-\Delta}\), followed by \(x^{-\Delta}(1+\Delta x/2+\cdots)\).}
\GSfield{Statement}{For a chiral block whose four external weights and intermediate weight all equal h, with normalized primary three-point factors, the level-one contribution is \(h^2/(2h)=h/2\) when h is nonzero. Thus the displayed coefficient is correct after inserting the missing inverse norm \(\langle L_{-1}h,L_{-1}h\rangle=2h\).}
\GSfield{Reconstruction}{The equal internal weight and normalization explain why this special coefficient is not a universal four-point coefficient.}
\GSfield{Fate}{The general block expansion is \citet[section~4]{belavinpolyakovzamolodchikov1984infinite}; the level-one verification here follows directly from the stated Virasoro and Ward relations.}
\GSfield{Present status}{As of November 2026: the stated result is established in the specified setting.}
\GSfield{Assessment}{literal statement: correct with missing hypotheses; publication match: exact; verification: read the result; development: annotated; confidence: high}
\end{GSitem}
\begin{GSitem}{1982-577.17}{F}{Coalescing punctures and accessory coefficients}
\GSfield{Kind; source}{question; \GSpage{P001067}}
\GSfield{The notes}{As \(z_i-z_j\to0\), does \(c_i\) go over to the preceding accessory coefficient?}
\GSfield{Statement}{A precise degeneration question must say which punctures remain, whether the colliding coefficient or a surviving one is examined, and what divergent principal part is subtracted. The source does not make these choices. Colliding cusp points cannot simply be substituted into the distinct-puncture formula.}
\GSfield{Fate}{\citet[section~3, Remark~7]{zograftakhtajan1988liouville} discusses degeneration as a further application, but is not itself the unqualified limit claimed here. No single modern open problem is inferred from the abbreviated question.}
\GSfield{Present status}{As of November 2026: the source’s unspecified limit cannot receive a unique solved/open status.}
\GSfield{Assessment}{literal statement: not determinable; publication match: partial; verification: read the result; development: annotated; confidence: medium}
\end{GSitem}
\begin{GSitem}{1982-577.18}{E}{Bosonic oscillators realize the Virasoro algebra}
\GSfield{Kind; source}{construction; \GSpage{P001068}}
\GSfield{The notes}{For positive n the notes write \(L_n=(\alpha+i\beta n)a_n+\sum_{k=1}^{n-1}a_ka_{n-k}+\sum_{k\ge1}a_k^+a_{n+k}\), with \([a_n,a_m^+]=n\delta_{nm}\).}
\GSfield{Statement}{The first quadratic sum requires a factor \(1/2\). With real \(\alpha,\beta\), \(L_{-n}=L_n^*\), \(L_0=\sum_{k\ge1}a_k^+a_k+h\), \(h=(\alpha^2+\beta^2)/2\), the corrected operators on finite-particle vectors have central charge \(c=1+12\beta^2\).}
\GSfield{Reconstruction}{The factor \(1/2\) and \(L_0\) are added to make the formula internally consistent. A vacuum level-two check is recorded in checks.py.}
\GSfield{Fate}{\citet[Theorem~2 and its preceding bosonic construction]{neretin1983highest} prints the associated group realization and attributes the Lie-algebra construction to Chodos--Thorn (1974). The pre-talk attribution is indirect through this primary 1983 paper.}
\GSfield{Present status}{As of November 2026: the stated result is established in the specified setting.}
\GSfield{Assessment}{literal statement: false as written; publication match: exact; verification: indirect; development: annotated; confidence: medium}
\end{GSitem}
\begin{GSitem}{1982-577.19}{E}{Oscillator and Verma bases agree only generically}
\GSfield{Kind; source}{claim; \GSpage{P001068}}
\GSfield{The notes}{The vectors \(a_{n_1}^+a_{n_2}^+\cdots|0\rangle\) and \(L_{-n_1}L_{-n_2}\cdots|0\rangle\) are said to be mutually expressible; the former are orthogonal.}
\GSfield{Statement}{Normalized occupation-number vectors form an orthogonal Fock basis. The change from Virasoro descendants is invertible at a fixed level only when the corresponding transition matrix has nonzero determinant. Singular vectors and special parameters can make the Verma-to-Fock map fail to be an isomorphism.}
\GSfield{Fate}{The need to separate generic and degenerate highest-weight parameters is established in \citet[section~3]{feiginfuchs1982invariant}; \citet[Theorem~2]{neretin1983highest} gives the Fock realization.}
\GSfield{Present status}{As of November 2026: the stated result is established in the specified setting.}
\GSfield{Assessment}{literal statement: correct with missing hypotheses; publication match: exact; verification: read the result; development: annotated; confidence: high}
\end{GSitem}
\begin{GSitem}{1982-577.20}{E}{Quadratic current and the Miura-type stress tensor}
\GSfield{Kind; source}{method; \GSpage{P001068}}
\GSfield{The notes}{\(T=p^{\prime}+p^2\), with a formal mode expansion of p.}
\GSfield{Statement}{The classical relation expresses a projective connection as a quadratic current plus a derivative. In the quantum oscillator realization the quadratic product must be normally ordered, and its derivative coefficient, zero mode and Laurent powers must match the chosen central charge. The written formula is schematic, not the complete normal-ordered operator identity.}
\GSfield{Reconstruction}{No numerical coefficient is forced from the notebook’s unnormalized p.}
\GSfield{Fate}{The corrected realization is checked through \citet[Theorem~2]{neretin1983highest}; it belongs to the pre-existing bosonic construction recalled there.}
\GSfield{Present status}{As of November 2026: the stated result is established in the specified setting.}
\GSfield{Assessment}{literal statement: correct with missing hypotheses; publication match: partial; verification: read the result; development: annotated; confidence: high}
\end{GSitem}
\begin{GSitem}{1982-577.21}{A}{Bargmann coordinates turn Ward identities into PDEs}
\GSfield{Kind; source}{method; \GSpage{P001069}}
\GSfield{The notes}{Put \(a_n=\partial_{z_n}\), \(a_n^+=nz_n\), then write a differential equation for \(F(x,z_1,z_2,\ldots)\).}
\GSfield{Statement}{On polynomials in oscillator coordinates these operators satisfy the stated Heisenberg relations. Substitution into the corrected \(L_n\) gives \([(\alpha+i\beta n)\partial_{z_n}+\tfrac12\sum_{k=1}^{n-1}\partial_{z_k}\partial_{z_{n-k}}+\sum_{k\ge1}kz_k\partial_{z_{n+k}}]F=(x^{n+1}\partial_x+(n+1)hx^n)F\), provided F is the corresponding highest-weight matrix element.}
\GSfield{Reconstruction}{The factors \(1/2\) and \(x^n\) are emendations; the function space and highest-weight input are editorial qualifications.}
\GSfield{Fate}{This is an explicit combination of the oscillator and Ward constructions subsequently treated in \citet{neretin1983highest} and \citet[sections~3--4]{belavinpolyakovzamolodchikov1984infinite}. It is a system of constraints; uniqueness needs initial and representation data.}
\GSfield{Present status}{As of November 2026: the stated result is established in the specified setting.}
\GSfield{Assessment}{literal statement: false as written; publication match: partial; verification: read the result; development: annotated; confidence: high}
\end{GSitem}
\begin{GSitem}{1982-577.22}{B}{Reconstructing four-point correlations from conformal blocks}
\GSfield{Kind; source}{claim; \GSpage{P001067-P001069}}
\GSfield{The notes}{The notes propose computing the coefficients algebraically and thereby determining \(\Psi\).}
\GSfield{Statement}{Ward identities determine descendants within a specified channel. A full four-point function additionally needs primary spectrum, three-point constants, a pairing of chiral and antichiral blocks, and an integration or summation prescription. The notebook does not supply those data.}
\GSfield{Reconstruction}{The modern theorem states a definite real-coupling range and channel conditions absent from the notes.}
\GSfield{Fate}{\citet[Theorem~1.1]{guillarmoukupiainenrhodesvargas2024bootstrap}, first circulated in 2020, proves the Liouville four-point bootstrap for \(0<\gamma<2\), \(\alpha_i<Q\), \(\alpha_1+\alpha_2>Q\), \(\alpha_3+\alpha_4>Q\), using the spectral integral and structure constants. This is a precise later realization, not a proof that symmetry alone specifies any CFT.}
\GSfield{Present status}{As of November 2026: the specified Liouville bootstrap theorem is proved; the statement is not enlarged to all parameter ranges or all conformal theories.}
\GSfield{Assessment}{literal statement: correct with missing hypotheses; publication match: partial; verification: read the result; development: annotated; confidence: high}
\end{GSitem}

\GStalk{578}{1982-11-29}{I. Yu. Kobzarev}{P001069-P001073}{Domain walls, vortex strings and the Nambu action}
The actual talk is P001070--P001073. A scalar double well supplies a planar wall; the Abelian Higgs model supplies a vortex string. Integrating a thin core yields the Nambu worldsheet action, with a collapsing circular loop as an illustration. The final page introduces the colour gauge Lagrangian. These are expository constructions predating the talk; the wall formula can be checked against Kobzarev’s own 1974 paper. The notes suppress couplings and contain a gauge-sign mismatch.

\begin{GSitem}{1982-578.1}{E}{Scalar field equations and polynomial interactions}
\GSfield{Kind; source}{definition; \GSpage{P001070}}
\GSfield{The notes}{\(\Box\psi+\mu^2\psi=0\), with interactions such as \(\lambda\varphi_1\varphi_2\varphi_3\).}
\GSfield{Statement}{A real scalar field on Minkowski space with quadratic kinetic and mass terms obeys the Klein--Gordon equation; polynomial potential terms change the Euler--Lagrange equation by the corresponding derivative. The date 1970 is a rhetorical contrast with soliton methods, not the historical invention of nonlinear field theory.}
\GSfield{Fate}{The interacting scalar setting is already explicit in \citet[section~2]{zeldovichkobzarevokun1975walls}.}
\GSfield{Present status}{As of November 2026: the stated result is established in the specified setting.}
\GSfield{Assessment}{literal statement: correct; publication match: partial; verification: read the result; development: annotated; confidence: high}
\end{GSitem}
\begin{GSitem}{1982-578.2}{E}{A double-well potential and two classical vacua}
\GSfield{Kind; source}{construction; \GSpage{P001070}}
\GSfield{The notes}{\(\mathcal L=\tfrac12(\partial\varphi)^2-\lambda^2(\varphi^2-\eta^2)^2\); the Hamiltonian is nonnegative, with vacua \(\varphi=\pm\eta\) exchanged by a discrete symmetry.}
\GSfield{Statement}{For positive \(\lambda,\eta\), with Minkowski kinetic convention, the static energy density is \(\tfrac12|\nabla\varphi|^2+\lambda^2(\varphi^2-\eta^2)^2\). Its homogeneous minima are exactly \(\pm\eta\); choosing one breaks the global \(\mathbb Z/2\) symmetry of the classical Lagrangian.}
\GSfield{Fate}{\citet[section~2]{zeldovichkobzarevokun1975walls} gives the same scalar potential before the lecture.}
\GSfield{Present status}{As of November 2026: the stated result is established in the specified setting.}
\GSfield{Assessment}{literal statement: correct with missing hypotheses; publication match: exact; verification: read the result; development: annotated; confidence: high}
\end{GSitem}
\begin{GSitem}{1982-578.3}{E}{The planar domain wall and its thickness}
\GSfield{Kind; source}{construction; \GSpage{P001070-P001071}}
\GSfield{The notes}{The field tends to \(\pm\eta\) as \(z\to\pm\infty\); the equation lacks a factor 4 and the solution is written \(\tanh(z\eta)\).}
\GSfield{Statement}{For the preceding normalization the field equation is \(\Box\varphi+4\lambda^2\varphi(\varphi^2-\eta^2)=0\). A static planar wall independent of t,x,y is \(\varphi(z)=\eta\tanh(\sqrt2\lambda\eta(z-z_0))\). Its characteristic width is of order \((\lambda\eta)^{-1}\); its energy is finite per unit area, not finite over the whole infinite plane.}
\GSfield{Reconstruction}{The translation parameter and the distinction between total and per-area energy are explicit. Substitution and a first-integral check are recorded independently in checks.py.}
\GSfield{Fate}{The correctly normalized profile and width are already in \citet[section~3]{zeldovichkobzarevokun1975walls}.}
\GSfield{Present status}{As of November 2026: the stated result is established in the specified setting.}
\GSfield{Assessment}{literal statement: false as written; publication match: exact; verification: read the result; development: annotated; confidence: high}
\end{GSitem}
\begin{GSitem}{1982-578.4}{E}{Local U(1) covariance of the Abelian Higgs action}
\GSfield{Kind; source}{definition; \GSpage{P001071}}
\GSfield{The notes}{The covariant derivative is \(\partial+iA\), but the printed transformation has \(\varphi\mapsto e^{i\alpha}\varphi\), \(A\mapsto A+d\alpha\).}
\GSfield{Statement}{With \(D=d+iA\), covariance requires \(A\mapsto A-d\alpha\). Then \(D\varphi\) transforms by the same phase as \(\varphi\), and the Maxwell term, scalar kinetic term and potential \((|\varphi|^2-\eta^2)^2\) are gauge invariant. A coupling can be restored consistently as \(D=d+ieA\).}
\GSfield{Fate}{The Abelian Higgs starting point is \citet[section~2, equation~2.1]{nielsenolesen1973vortex}.}
\GSfield{Present status}{As of November 2026: the stated result is established in the specified setting.}
\GSfield{Assessment}{literal statement: false as written; publication match: exact; verification: read the result; development: annotated; confidence: high}
\end{GSitem}
\begin{GSitem}{1982-578.5}{E}{Vortex winding and flux quantization}
\GSfield{Kind; source}{construction; \GSpage{P001071-P001072}}
\GSfield{The notes}{A transverse circuit rotates the vacuum vector once; \(\varphi_1=(x/r)f(r)\), \(\varphi_2=(y/r)f(r)\), approaching the vacuum at infinity.}
\GSfield{Statement}{A smooth static straight vortex of winding n has \(\varphi\sim\eta e^{in\theta}\) at large radius, vanishing core amplitude, and a gauge field with \(D\varphi\to0\). For \(D=d+ieA\), finite energy per unit length gives \(\int F=-2\pi n/e\) in this orientation. The profiles satisfy coupled nonlinear radial equations; the displayed ansatz is not an elementary closed-form solution.}
\GSfield{Reconstruction}{General winding, finite-energy boundary conditions and the gauge-field role complete the one-winding diagram.}
\GSfield{Fate}{\citet[section~2, equations~2.4--2.13]{nielsenolesen1973vortex} derives the flux and radial ansatz in 1973.}
\GSfield{Present status}{As of November 2026: the stated result is established in the specified setting.}
\GSfield{Assessment}{literal statement: correct with missing hypotheses; publication match: exact; verification: read the result; development: annotated; confidence: high}
\end{GSitem}
\begin{GSitem}{1982-578.6}{E}{The thin-core approximation produces a string action}
\GSfield{Kind; source}{method; \GSpage{P001072}}
\GSfield{The notes}{If the curvature radius greatly exceeds the thickness, \(S\simeq\nu\int dl\sqrt{1-v_\perp^2}\,dt\), where \(\nu\) is energy per length.}
\GSfield{Statement}{For a slowly varying vortex with core width much smaller than all relevant curvature radii, with internal excitations neglected, integrating the transverse profile gives the leading Nambu--Goto action \(S_{eff}=-\nu\int\sqrt{-\det h}\,d^2\sigma\), up to the overall action-sign convention. The approximation ceases to control regions with curvature on the core scale.}
\GSfield{Reconstruction}{The hypotheses are a physical effective-description regime, not an error estimate proved in the notebook.}
\GSfield{Fate}{The same reduction and transverse-velocity argument are \citet[section~3, equations~3.3--3.6]{nielsenolesen1973vortex}.}
\GSfield{Present status}{As of November 2026: the stated result is established in the specified setting.}
\GSfield{Assessment}{literal statement: correct with missing hypotheses; publication match: exact; verification: read the result; development: annotated; confidence: high}
\end{GSitem}
\begin{GSitem}{1982-578.7}{E}{Worldsheet area and extremality}
\GSfield{Kind; source}{definition; \GSpage{P001072-P001073}}
\GSfield{The notes}{The moving string sweeps a sheet; its area is called the Lagrangian, followed by minimum area.}
\GSfield{Statement}{The induced worldsheet metric is \(h_{ab}=\eta_{\mu\nu}\partial_aX^\mu\partial_bX^\nu\). The Nambu--Goto action is proportional to its Lorentzian area. Its Euler--Lagrange equations express stationarity, equivalently vanishing mean-curvature vector for a smooth timelike immersion, with suitable boundary conditions. They do not assert a Euclidean minimum.}
\GSfield{Reconstruction}{Action versus Lagrangian density and Lorentzian stationarity are distinguished.}
\GSfield{Fate}{The area and variation formulas are displayed in \citet[sections~1.2--1.2.2]{tong2009string}, as an exposition of the older construction.}
\GSfield{Present status}{As of November 2026: the stated result is established in the specified setting.}
\GSfield{Assessment}{literal statement: false as written; publication match: exact; verification: read the result; development: annotated; confidence: high}
\end{GSitem}
\begin{GSitem}{1982-578.8}{E}{A circular loop contracts in the effective model}
\GSfield{Kind; source}{claim; \GSpage{P001072}}
\GSfield{The notes}{A ring-shaped string contracts, quickly reaches the speed of light and becomes a small lump.}
\GSfield{Statement}{For an initially resting circular Nambu--Goto loop in flat space, \(R(t)=R_0\cos(t/R_0)\) until the first collapse, \(0\le t<\pi R_0/2\). Its inward speed tends to 1 at collapse. For a field-theory vortex, the thin-core approximation fails before zero radius, so this does not determine the final lump or radiation.}
\GSfield{Reconstruction}{The exact circular solution makes the qualitative claim precise only within the effective model; its conserved-energy and equation-of-motion checks are recorded in checks.py.}
\GSfield{Fate}{This follows from the classical action explained in \citet[section~1.2]{tong2009string}; the finite-core limitation is already part of \citet[section~3]{nielsenolesen1973vortex}.}
\GSfield{Present status}{As of November 2026: the stated result is established in the specified setting.}
\GSfield{Assessment}{literal statement: correct with missing hypotheses; publication match: exact; verification: read the result; development: annotated; confidence: high}
\end{GSitem}
\begin{GSitem}{1982-578.9}{F}{Monopoles as string endpoints}
\GSfield{Kind; source}{remark; \GSpage{P001072}}
\GSfield{The notes}{A variant places a monopole at each broken end of a string.}
\GSfield{Statement}{The diagram proposes a flux tube ending on magnetic sources. The smooth pure Abelian Higgs model on all of Minkowski space does not itself supply nonsingular monopole cores: one must add magnetic defects, boundary conditions, or an enlarged gauge theory. None of these is specified.}
\GSfield{Fate}{An identifiable earlier lead is \citet{kobayashi1974monopoles}, whose title and metadata match the idea. Its particular model cannot be assigned to this unspecified sketch.}
\GSfield{Present status}{As of November 2026: the source does not fix an endpoint model; this is not catalogued as an unsolved monopole-existence problem.}
\GSfield{Assessment}{literal statement: not determinable; publication match: thematic only; verification: publisher metadata checked; development: annotated; confidence: medium}
\end{GSitem}
\begin{GSitem}{1982-578.10}{E}{The colour gauge field and quark Lagrangian}
\GSfield{Kind; source}{definition; \GSpage{P001073}}
\GSfield{The notes}{\(SU(3)\) acts in the adjoint on the eight gauge components and in the triplet on quarks; two terms for u and d are displayed.}
\GSfield{Statement}{For a chosen convention, \(\mathcal L=-\tfrac14F^a_{\mu\nu}F^{a\mu\nu}+\sum_f\bar q_f(i\gamma^\mu D_\mu-m_f)q_f\), with \(D_\mu=\partial_\mu-igA_\mu^at^a\), \([t^a,t^b]=if^{abc}t^c\), and \(F^a_{\mu\nu}=\partial_\mu A^a_\nu-\partial_\nu A^a_\mu+gf^{abc}A^b_\mu A^c_\nu\). The adjoint dimension is 8 and the fundamental dimension 3. A two-flavour truncation is a model, not a claim that there are only two flavours.}
\GSfield{Reconstruction}{A consistent sign and coupling convention replaces the notebook’s mixed D and F formulas, and separate flavour masses are restored.}
\GSfield{Fate}{The colour gauge-theory framework predates the talk \citep{grosswilczek1973ultraviolet}. The present note records a definition; no confinement theorem is stated.}
\GSfield{Present status}{As of November 2026: the stated result is established in the specified setting.}
\GSfield{Assessment}{literal statement: correct with missing hypotheses; publication match: partial; verification: publisher metadata checked; development: annotated; confidence: high}
\end{GSitem}

\GStalk{579}{}{}{P001074-P001075}{Year 1983}
This MCCME record is a year heading. It has no separate mathematical content; the adjacent February material belongs to records 580 and 581.


\GStalk{580}{1983-02-07}{A. A. Kirillov}{P001074-P001075}{Near-unitary circle actions and Fock-space representations}
The record occupies the upper part of P001074. Following Gelfand’s methodological and supersymmetry remarks, Kirillov reports Neretin’s weighted pullbacks, their Hilbert--Schmidt failure of unitarity, and Gaussian/Fock unitarization. This construction had already appeared in 1982; extending the action on the original one-particle space and constructing a new unitary action on Fock space must be distinguished \citep{neretin1982complementary}.

\begin{GSitem}{1982-580.1}{F}{A call for infinite-dimensional symplectic geometry}
\GSfield{Kind; source}{remark; \GSpage{P001074}; I. M. Gelfand}
\GSfield{The notes}{Gelfand says that infinite-dimensional symplectic geometry is not yet familiar and calls Fock space a toy for mathematicians.}
\GSfield{Statement}{The intervention requests a geometric understanding beyond formal oscillator manipulations. It specifies neither a category of infinite-dimensional manifolds nor a theorem to prove.}
\GSfield{Fate}{The restricted Hilbert-space groups used immediately afterwards are identifiable in \citet[section 1]{neretin1982complementary}; they do not uniquely specify the broader programme.}
\GSfield{Present status}{As of November 2026: the notebook does not determine a unique proposition whose resolution can be assessed.}
\GSfield{Assessment}{literal statement: not determinable; publication match: thematic only; verification: unresolved lead; development: annotated; confidence: low}
\end{GSitem}
\begin{GSitem}{1982-580.2}{E}{Supersymmetric derivation of fixed point formulas}
\GSfield{Kind; source}{method; \GSpage{P001074}; I. M. Gelfand}
\GSfield{The notes}{Popular lectures by Witten are mentioned, with an arrow to the Atiyah--Bott fixed point theorem.}
\GSfield{Statement}{In supersymmetric quantum mechanics on a compact Riemannian manifold, a deformation of the de Rham/signature complex by a Killing field localizes an index calculation at its zeros. For isolated zeros this yields the corresponding fixed point contributions. The note does not identify a particular elliptic complex or a unique preprint.}
\GSfield{Reconstruction}{Compactness, the Killing-field setting and the chosen complexes are supplied from the source.}
\GSfield{Fate}{\citet[section 3, equations (44)--(45)]{witten1982supersymmetry} supplies this calculation in a 1982 paper, received 26 March 1982. The mathematical identification is stronger than the bibliographic identification of the unnamed lectures.}
\GSfield{Present status}{As of November 2026: the stated result is established in the specified setting.}
\GSfield{Assessment}{literal statement: correct with missing hypotheses; publication match: partial; verification: read the result; development: annotated; confidence: medium}
\end{GSitem}
\begin{GSitem}{1982-580.3}{E}{Weighted pullback extends the Möbius action}
\GSfield{Kind; source}{construction; \GSpage{P001074}}
\GSfield{The notes}{A fractional-linear action with a factor \(|\beta z+\delta|^\lambda\) is extended to diffeomorphisms by \(f\mapsto(f\circ\varphi^{-1})|D\varphi^{-1}|^{-\lambda/2}\).}
\GSfield{Statement}{For any real density weight a, \(T_a(\varphi)f=(f\circ\varphi^{-1})|(\varphi^{-1})^{\prime}|^a\) defines a representation of orientation-preserving smooth circle diffeomorphisms on smooth densities in a local trivialization. Its restriction to Möbius transformations has the familiar multiplier. Unitarity depends on the Hilbert form and parameter; it does not follow from the extension formula.}
\GSfield{Reconstruction}{A neutral weight a avoids identifying the notebook’s lambda with the publication’s lambda without the required reparametrization.}
\GSfield{Fate}{The weighted pullback is \citet[section 2, equation (2)]{neretin1982complementary}, with a different parameter convention.}
\GSfield{Present status}{As of November 2026: the stated result is established in the specified setting.}
\GSfield{Assessment}{literal statement: correct with missing hypotheses; publication match: exact; verification: read the result; development: annotated; confidence: high}
\end{GSitem}
\begin{GSitem}{1982-580.4}{E}{The complementary-series inner product is almost invariant}
\GSfield{Kind; source}{theorem; \GSpage{P001074}}
\GSfield{The notes}{The imaginary part of a double integral with kernel \(|x_1-x_2|^{-2+\lambda}/\Gamma(-1+\lambda)\) is Möbius invariant; under general diffeomorphisms the operator is marked \(U(1+\mathrm{HS})\).}
\GSfield{Statement}{For Neretin’s positive complementary-series Hilbert completions, with his parameter \(0<\rho<2\), \(\rho\ne1\), the weighted pullbacks are bounded invertible and have polar form \(O(1+K)\), with K Hilbert--Schmidt. In that range they are orthogonal exactly for Möbius transformations. Analytic continuation and a parameter-dependent normalization are required outside the convergent integral range.}
\GSfield{Reconstruction}{The positivity range is stated in the publication’s parameter, not asserted for every value of the handwritten lambda. The imaginary-part form is on the complexification of the positive real Hilbert space.}
\GSfield{Fate}{This is \citet[Theorem 1 and the following proposition]{neretin1982complementary}, already published in 1982.}
\GSfield{Present status}{As of November 2026: the stated result is established in the specified setting.}
\GSfield{Assessment}{literal statement: correct with missing hypotheses; publication match: exact; verification: read the result; development: annotated; confidence: high}
\end{GSitem}
\begin{GSitem}{1982-580.5}{E}{Gaussian unitarization on Fock space}
\GSfield{Kind; source}{construction; \GSpage{P001074}}
\GSfield{The notes}{The margin corrects the phrase unitarily: the weighted action is nonunitary, but unitary representations are constructed on Fock space using coherent states.}
\GSfield{Statement}{If a real Hilbert-space representation takes values in \(GL_0(H)=\{O(1+K):K\text{ Hilbert--Schmidt}\}\), the associated transformations of a Gaussian realization are quasi-invariant. Composition multiplied by the square root of the Radon--Nikodym derivative gives a unitary representation on Gaussian \(L^2\), equivalently the symmetric Fock space. This changes the representation space.}
\GSfield{Fate}{\citet[section 1 and the construction after Theorem 1]{neretin1982complementary} gives exactly this mechanism. Its restriction to the Möbius subgroup is a Fock representation, not the original single complementary-series representation.}
\GSfield{Present status}{As of November 2026: the stated result is established in the specified setting.}
\GSfield{Assessment}{literal statement: correct with missing hypotheses; publication match: exact; verification: read the result; development: annotated; confidence: high}
\end{GSitem}

\GStalk{581}{1983-02-07}{B. L. Feigin}{P001074-P001077}{String constraints and semi-infinite homology}
The actual notes occupy the lower P001074 and P001075. They pass from the loop-space symplectic form to oscillator quantization, explain the central-term obstruction to ordinary invariants, and replace invariants by the image of positive-mode invariants in negative-mode coinvariants. Feigin’s 1984 note identifies this physical-state construction with semi-infinite homology. Gelfand’s closing p-adic request has a separately documented answer in Neretin’s 1984 paper.

\begin{GSitem}{1982-581.1}{E}{The loop-space symplectic form modulo constants}
\GSfield{Kind; source}{construction; \GSpage{P001074}}
\GSfield{The notes}{For maps \(f,g:S^1\to\mathbb R^d\), with a Riemannian or Lorentzian target metric, use \(\omega(f,g)=\int_{S^1}(f,dg)\); quotient by constant maps.}
\GSfield{Statement}{For a nondegenerate symmetric bilinear form b on \(\mathbb R^d\), the form \(\omega(f,g)=\int b(f,g^{\prime})dt\) on smooth loops is skew by integration by parts. Its radical is exactly the constant loops, so it descends to a weak symplectic form on \(C^\infty(S^1,\mathbb R^d)/\mathbb R^d\). Orientation-preserving reparametrizations preserve it. Positivity of b is unnecessary for this algebraic assertion.}
\GSfield{Reconstruction}{Smoothness, orientation, nondegeneracy of the target metric and weak rather than strong symplectic topology are made explicit.}
\GSfield{Fate}{This is the classical oscillator phase space underlying the final construction in \citet{feigin1984semiinfinite}; the scalar endpoint realization is already in \citet[section 4]{neretin1982complementary}.}
\GSfield{Present status}{As of November 2026: the stated result is established in the specified setting.}
\GSfield{Assessment}{literal statement: correct with missing hypotheses; publication match: partial; verification: read the result; development: annotated; confidence: high}
\end{GSitem}
\begin{GSitem}{1982-581.2}{E}{The quotient carries a Poisson algebra, not automatically a smooth manifold}
\GSfield{Kind; source}{claim; \GSpage{P001074}}
\GSfield{The notes}{The space \(V^{\prime}/\operatorname{Diff}S^1\) is said to have a Hamiltonian structure.}
\GSfield{Statement}{Invariant observables in a symplectic Poisson algebra are closed under the bracket when the action preserves the symplectic structure. This defines a quotient Poisson algebra. A smooth Poisson manifold structure on the orbit space additionally requires a suitable regular quotient; the loop-space quotient in the notebook is generally singular and no such manifold theorem is established.}
\GSfield{Reconstruction}{The algebra of observables is the precise interpretation supported without additional regularity hypotheses.}
\GSfield{Fate}{The finite-dimensional reduction principle is older \citep{marsdenweinstein1974reduction}. Its infinite-dimensional application here is formal unless a class of admissible observables and a quotient category are chosen.}
\GSfield{Present status}{As of November 2026: the stated result is established in the specified setting.}
\GSfield{Assessment}{literal statement: correct with missing hypotheses; publication match: partial; verification: publisher metadata checked; development: annotated; confidence: high}
\end{GSitem}
\begin{GSitem}{1982-581.3}{E}{Symplectic reduction requires a moment-map level}
\GSfield{Kind; source}{theorem; \GSpage{P001074}; A. A. Kirillov}
\GSfield{The notes}{Kirillov describes a reduced symplectic manifold and says the invariant Poisson bracket is degenerate and one must take a quotient.}
\GSfield{Statement}{For a finite-dimensional Hamiltonian G-action with equivariant moment map \(\mu\), if \(\xi\) is a regular value and \(G_\xi\) acts freely and properly on \(\mu^{-1}(\xi)\), then \(\mu^{-1}(\xi)/G_\xi\) is symplectic, with reduced form characterized by pullback. The full orbit space \(M/G\) is not generally symplectic.}
\GSfield{Reconstruction}{The level set, Hamiltonian hypothesis and free proper regularity assumptions are restored editorially.}
\GSfield{Fate}{This is the regular reduction construction of \citet{marsdenweinstein1974reduction}, published in 1974.}
\GSfield{Present status}{As of November 2026: the stated result is established in the specified setting.}
\GSfield{Assessment}{literal statement: correct with missing hypotheses; publication match: partial; verification: publisher metadata checked; development: annotated; confidence: high}
\end{GSitem}
\begin{GSitem}{1982-581.4}{E}{Polarization and the symmetric oscillator space}
\GSfield{Kind; source}{construction; \GSpage{P001074-P001075}}
\GSfield{The notes}{Complexify \(V^{\prime}=H^+\oplus H^-\), using Fourier modes analytic inside and outside the circle; the notes display \(\Phi=\bigoplus_i S^iH_+\).}
\GSfield{Statement}{A positive/negative Fourier polarization gives two paired complex Lagrangian spaces. The algebraic bosonic Fock module is a symmetric algebra on the creation space; annihilation acts by contractions with the dual polarization. Which summand is labelled plus is conventional. A Hilbert completion and positive inner product need an additional choice; a Lorentzian target does not provide a positive Hilbert norm by itself.}
\GSfield{Reconstruction}{The duality implicit in the inconsistent plus/creation labels and the distinction between algebraic and Hilbert Fock spaces are stated explicitly.}
\GSfield{Fate}{The oscillator construction predates the lecture and is recalled in \citet[section 1]{neretin1982complementary}.}
\GSfield{Present status}{As of November 2026: the stated result is established in the specified setting.}
\GSfield{Assessment}{literal statement: correct with missing hypotheses; publication match: partial; verification: read the result; development: annotated; confidence: high}
\end{GSitem}
\begin{GSitem}{1982-581.5}{E}{The Schwinger term obstructs ordinary invariants}
\GSfield{Kind; source}{claim; \GSpage{P001074-P001075}}
\GSfield{The notes}{The map from vector fields through symplectic transformations to the Weyl representation is projective; a Schwinger term appears and invariants do not exist.}
\GSfield{Statement}{The quantized action factors through the Virasoro central extension. In a module on which its central element acts by a nonzero scalar, no nonzero vector is annihilated by every Virasoro generator including the center. For d free scalar oscillators with standard normal ordering, the matter central charge is d, also for a nondegenerate indefinite target form.}
\GSfield{Reconstruction}{The precise invariant obstruction and the normal-ordering convention replace the loose phrase invariants do not exist.}
\GSfield{Fate}{The central extension is older \citep{gelfandfuchs1968circle}; Feigin’s string construction explicitly treats the resulting central charge \citep[Example 1 and final paragraph]{feigin1984semiinfinite}.}
\GSfield{Present status}{As of November 2026: the stated result is established in the specified setting.}
\GSfield{Assessment}{literal statement: correct with missing hypotheses; publication match: partial; verification: read the result; development: annotated; confidence: high}
\end{GSitem}
\begin{GSitem}{1982-581.6}{A}{Physical states as an image of invariants in coinvariants}
\GSfield{Kind; source}{construction; \GSpage{P001075}}
\GSfield{The notes}{For \(L_+=\langle L_1,L_2,\ldots\rangle\), \(L_-=\langle L_{-1},L_{-2},\ldots\rangle\), take the image of \(\Phi^{L_+}\to\Phi/L_-\Phi\), then select an \(L_0\) condition.}
\GSfield{Statement}{For the string oscillator module, form \(Q^{\prime}=\operatorname{im}(H^0(L_+,\Phi)\to H_0(L_-,\Phi))\). Physical states require a further zero-mode condition, \(e_0v=-v\) in Feigin’s sign convention. Thus invariants are retained only modulo spurious negative-mode states; they are not the full Virasoro-invariant subspace.}
\GSfield{Reconstruction}{The zero-mode eigenvalue is supplied from the publication; the notebook’s subscript does not specify it unambiguously.}
\GSfield{Fate}{\citet[final physical-state paragraph]{feigin1984semiinfinite} states this exact construction and identifies it with degree-zero semi-infinite homology. Received 15 January 1983, published in 1984.}
\GSfield{Present status}{As of November 2026: the stated result is established in the specified setting.}
\GSfield{Assessment}{literal statement: correct with missing hypotheses; publication match: exact; verification: read the result; development: annotated; confidence: high}
\end{GSitem}
\begin{GSitem}{1982-581.7}{A}{The critical charge 26 in semi-infinite homology}
\GSfield{Kind; source}{theorem; \GSpage{P001075}}
\GSfield{The notes}{The resulting space is nontrivial at least for \(d=26\).}
\GSfield{Statement}{For Feigin’s graded category of Virasoro modules, the relative semi-infinite Tor functors of two modules with scalar central charges \(c_1,c_2\) vanish unless \(c_1+c_2=26\). In the string construction the fermionic complex contributes charge \(-26\), balancing the d-boson charge at \(d=26\). This necessary charge balance is not a blanket assertion that every module at charge 26 has nonzero physical states.}
\GSfield{Reconstruction}{The categorical assumptions and distinction between necessary charge cancellation and arbitrary-module nonvanishing are supplied from the paper.}
\GSfield{Fate}{\citet[Example 1, theorem (a), and final paragraph]{feigin1984semiinfinite} gives the precise homological statement.}
\GSfield{Present status}{As of November 2026: the stated result is established in the specified setting.}
\GSfield{Assessment}{literal statement: correct with missing hypotheses; publication match: partial; verification: read the result; development: annotated; confidence: high}
\end{GSitem}
\begin{GSitem}{1982-581.8}{B}{A p-adic analogue of diffeomorphism-group representations}
\GSfield{Kind; source}{problem; \GSpage{P001075}; I. M. Gelfand}
\GSfield{The notes}{Gelfand asks:\par\ru{Построить p-адич. аналог предст. группы диффеоморфизмов}.}
\GSfield{Statement}{Construct unitary representations for a p-adic analogue of the circle diffeomorphism group by replacing the real-circle almost-invariant Hilbert structure and Fock/Gaussian construction. The subsequently documented realization uses analytic diffeomorphisms of \(\mathbb P^1(\mathbb Q_p)\), embedded in the group of local automorphisms of the Bruhat--Tits tree boundary.}
\GSfield{Reconstruction}{The precise p-adic group and Hilbert construction come from the solution, not the one-line question. No p-adic Virasoro highest-weight theorem is claimed.}
\GSfield{Fate}{\citet[introduction and Theorem 1]{neretin1984padic} explicitly attributes this question to Gelfand and constructs the required near-orthogonal action and Gaussian unitary representations. Received 10 January 1984, published 1984.}
\GSfield{Present status}{As of November 2026: the analytic-projective-line version explicitly associated with Gelfand’s request is solved; other possible meanings of p-adic analogue are not identified with it.}
\GSfield{Assessment}{literal statement: correct with missing hypotheses; publication match: partial; verification: read the result; development: annotated; confidence: high; problem \hyperref[prob:GP-1982-22]{GP-1982-22}}
\end{GSitem}

\GStalk{582}{1983-02-14}{Yu. A. Neretin}{P001075-P001078}{Hilbert transforms, almost invariant kernels and circle representations}
The entry is P001076--P001077, including a name of Ginzburg crossed out before the final continuation question. The 1982 near-orthogonal construction is refined by endpoint complex structures and restricted symplectic actions, published in Neretin’s 1983 Doklady note. Gelfand’s pseudodifferential generalization remains unspecified. The written logarithmic convolution cannot yield the claimed derivative of a delta distribution; the Fock-space question needs a completed symmetric space rather than arbitrary analytic sequences.

\begin{GSitem}{1982-582.1}{A}{The almost invariant kernel needs a positivity qualification}
\GSfield{Kind; source}{theorem; \GSpage{P001076}}
\GSfield{The notes}{For every \(\rho\), a power-kernel inner product is almost preserved, with operator \(O(1+T)\), T Hilbert--Schmidt.}
\GSfield{Statement}{In the complementary interval the appropriately normalized and continued power kernel defines a positive Hilbert form and a near-orthogonal smooth diffeomorphism action. Outside it the unmodified form need not be positive. Neretin constructs positive forms for nonintegral parameters by adding a smooth kernel; their almost-invariant equivalence class is independent of that modification.}
\GSfield{Reconstruction}{The parameter restrictions, analytic continuation and smooth correction are indispensable additions to the phrase for every rho.}
\GSfield{Fate}{The interval case is already \citet[Theorem 1]{neretin1982complementary}; the qualified extension is \citet[Theorems 1--2]{neretin1983boson}, received 20 January 1983 and published later that year.}
\GSfield{Present status}{As of November 2026: the stated result is established in the specified setting.}
\GSfield{Assessment}{literal statement: false as written; publication match: partial; verification: read the result; development: annotated; confidence: high}
\end{GSitem}
\begin{GSitem}{1982-582.2}{E}{The periodic Hilbert transform is a complex structure modulo constants}
\GSfield{Kind; source}{construction; \GSpage{P001076}}
\GSfield{The notes}{The cotangent integral sends \(\cos nt\mapsto-\sin nt\), \(\sin nt\mapsto\cos nt\), and constants to zero.}
\GSfield{Statement}{Interpret the singular integral in the principal-value sense. With the displayed sign convention, H acts on nonzero Fourier modes by \(i\operatorname{sgn}(n)\), and \(H^2=-1\) on mean-zero functions or the quotient by constants. It therefore defines a real complex structure on the appropriate Sobolev completions.}
\GSfield{Reconstruction}{Principal value, removal of the zero mode and the Fourier multiplier make the formal integral meaningful.}
\GSfield{Fate}{The classical Hilbert-transform identity is used in \citet[Theorem 3]{neretin1983boson}; only its application to the restricted symplectic representation is new there.}
\GSfield{Present status}{As of November 2026: the stated result is established in the specified setting.}
\GSfield{Assessment}{literal statement: correct with missing hypotheses; publication match: partial; verification: read the result; development: annotated; confidence: high}
\end{GSitem}
\begin{GSitem}{1982-582.3}{E}{Two endpoint Hilbert forms}
\GSfield{Kind; source}{definition; \GSpage{P001076}}
\GSfield{The notes}{The endpoint cases are written as a logarithmic kernel and as a cotangent kernel with a derivative; one removes constants or imposes zero mean.}
\GSfield{Statement}{Up to positive normalization, the endpoint norms have Fourier expressions \(\sum_{n\ne0}|n|^{-1}|f_n|^2\) on mean-zero density coefficients and \(\sum_{n\ne0}|n||f_n|^2\) on functions modulo constants. The former is represented by the negative logarithmic sine kernel. The unmodified positive logarithm in the notes is negative on every nonzero mean-zero real Fourier mode.}
\GSfield{Reconstruction}{Fourier norms specify a positive convention without silently treating the handwritten logarithm as positive.}
\GSfield{Fate}{The endpoint constructions already appear in \citet[section 4]{neretin1982complementary}; \citet[section 2]{neretin1983boson} incorporates them into the common framework.}
\GSfield{Present status}{As of November 2026: the stated result is established in the specified setting.}
\GSfield{Assessment}{literal statement: false as written; publication match: exact; verification: read the result; development: annotated; confidence: high}
\end{GSitem}
\begin{GSitem}{1982-582.4}{A}{Endpoint actions preserve a symplectic form}
\GSfield{Kind; source}{theorem; \GSpage{P001076-P001077}}
\GSfield{The notes}{Extend the scalar product using H; its imaginary part is preserved by all diffeomorphisms, and \(\omega(f,g)=\langle f,Hg\rangle\).}
\GSfield{Statement}{For either endpoint Hilbert space with its specified density action, the Hilbert transform supplies a compatible complex structure; smooth orientation-preserving diffeomorphisms preserve the associated symplectic form and belong to the restricted symplectic group. On the function realization the form is, up to sign and normalization, \(\int f\,dg\). This stronger exact invariance is an endpoint statement, not a consequence of near-orthogonality at every parameter.}
\GSfield{Reconstruction}{The endpoint restriction and the respective function/density actions are restored.}
\GSfield{Fate}{This is \citet[Theorem 3]{neretin1983boson}; the paper also notes that the polar defect is trace class in its constructions.}
\GSfield{Present status}{As of November 2026: the stated result is established in the specified setting.}
\GSfield{Assessment}{literal statement: correct with missing hypotheses; publication match: exact; verification: read the result; development: annotated; confidence: high}
\end{GSitem}
\begin{GSitem}{1982-582.5}{E}{Normalized Hardy projections and Lagrangian subspaces}
\GSfield{Kind; source}{construction; \GSpage{P001077}}
\GSfield{The notes}{\(H\pm i\) are called projections, with Lagrangian images.}
\GSfield{Statement}{If \(H^2=-1\), the projections onto its \(\pm i\) eigenspaces in the complexification are \(P_\pm=(1\mp iH)/2\). Their images are the positive and negative Fourier polarizations and are Lagrangian for the complex bilinear extension of the symplectic form. The unscaled \(H\pm i\) have those images but are not idempotents.}
\GSfield{Reconstruction}{Projection normalization and complex bilinear, rather than Hermitian, isotropy are explicit.}
\GSfield{Fate}{This is the linear algebra of the Hilbert-transform structure used in \citet[Theorem 3 and section 3]{neretin1983boson}.}
\GSfield{Present status}{As of November 2026: the stated result is established in the specified setting.}
\GSfield{Assessment}{literal statement: false as written; publication match: partial; verification: read the result; development: annotated; confidence: high}
\end{GSitem}
\begin{GSitem}{1982-582.6}{F}{A general pseudodifferential invariance principle}
\GSfield{Kind; source}{question; \GSpage{P001076-P001077}; I. M. Gelfand}
\GSfield{The notes}{Gelfand asks for a general statement about pseudodifferential operators invariant under all diffeomorphisms, and for the full symbol series beyond \(c(1+\text{compact})\).}
\GSfield{Statement}{The proposed generalization does not specify the manifold dimension, density weights, operator order, topology of the remainder ideal, or whether equality is exact or modulo smoothing operators. These choices distinguish different classification and covariance problems.}
\GSfield{Fate}{Neretin’s concrete almost-invariant kernels are covered by \citet[section 2]{neretin1983boson}; that result does not resolve every possible pseudodifferential generalization.}
\GSfield{Present status}{As of November 2026: the notebook does not determine a unique proposition whose resolution can be assessed.}
\GSfield{Assessment}{literal statement: not determinable; publication match: thematic only; verification: unresolved lead; development: annotated; confidence: low}
\end{GSitem}
\begin{GSitem}{1982-582.7}{E}{The logarithmic Hilbert convolution is not a delta derivative}
\GSfield{Kind; source}{claim; \GSpage{P001077}}
\GSfield{The notes}{The integral \(\int\log|(t-s)/2|\,(s-\sigma)^{-1}ds\) is expected to give \(\delta^{\prime}\), with analytic continuation proposed.}
\GSfield{Statement}{After principal-value/tempered regularization, away from zero Fourier frequency the logarithmic kernel has multiplier proportional to \(|\xi|^{-1}\), while the Hilbert kernel has multiplier proportional to \(\operatorname{sgn}\xi\). Their product has order \(-1\), proportional to \(1/\xi\), whereas \(\delta^{\prime}\) has order \(+1\). Thus the displayed convolution cannot equal \(\delta^{\prime}\); at the symbolic level it is an antiderivative-type kernel. Zero-frequency ambiguities do not repair this mismatch.}
\GSfield{Reconstruction}{Regularization and exclusion of zero frequency are explicit; no globally defined unregularized convolution is asserted.}
\GSfield{Fate}{The endpoint formulas in \citet[section 2]{neretin1983boson} distinguish the inverse-derivative and derivative norms. The Fourier-order check here identifies a faithful error in the written intermediate calculation.}
\GSfield{Present status}{As of November 2026: the stated result is established in the specified setting.}
\GSfield{Assessment}{literal statement: false as written; publication match: partial; verification: read the result; development: annotated; confidence: high}
\end{GSitem}
\begin{GSitem}{1982-582.8}{E}{Analytic continuation and differentiation of homogeneous kernels}
\GSfield{Kind; source}{method; \GSpage{P001077}}
\GSfield{The notes}{Split \(\int|x|^\lambda f(x)dx\) into \(|x|<1\) and its complement, continue analytically, and differentiate in lambda to get a logarithm.}
\GSfield{Statement}{For Schwartz test functions the family \(|x|^\lambda\), initially \(\operatorname{Re}\lambda>-1\), extends meromorphically as a tempered distribution by subtracting Taylor coefficients at zero. At regular parameters, differentiation yields \(|x|^\lambda\log|x|\); at poles one must specify finite parts and residue terms. This procedure defines kernels but does not validate the preceding incorrect convolution identity.}
\GSfield{Reconstruction}{The test-function class and meromorphic, rather than everywhere holomorphic, continuation are supplied.}
\GSfield{Fate}{Analytic regularization is already explicitly used in \citet[section 2]{neretin1982complementary} and \citet[section 2]{neretin1983boson}, which cites the earlier Gelfand--Shilov distribution theory.}
\GSfield{Present status}{As of November 2026: the stated result is established in the specified setting.}
\GSfield{Assessment}{literal statement: correct with missing hypotheses; publication match: partial; verification: read the result; development: annotated; confidence: high}
\end{GSitem}
\begin{GSitem}{1982-582.9}{A}{Realizing circle diffeomorphisms on analytic symmetric sequences}
\GSfield{Kind; source}{question; \GSpage{P001077}}
\GSfield{The notes}{How does a diffeomorphism act on \(\{f_1(z),f_2(z_1,z_1),f_3(z_1,z_2,z_3),\ldots\}\)? For Möbius maps it agrees with the usual action.}
\GSfield{Statement}{Choose a positive-frequency one-particle Hilbert space H and its completed symmetric Fock space \(\widehat{\bigoplus}_{n\ge0}\operatorname{Sym}^nH\). Its vectors can be realized as symmetric holomorphic n-variable components satisfying the Fock norm condition. Restricted symplectic implementers act projectively on this completion; for Möbius transformations preserving the polarization this reduces to second quantization of the one-particle action. Arbitrary sequences of holomorphic functions are not automatically vectors of this representation.}
\GSfield{Reconstruction}{Symmetry, the vacuum component, Hilbert completion and the repeated second argument’s correction are editorial specifications.}
\GSfield{Fate}{\citet[section 3 and Theorems 4--5]{neretin1983boson} gives the implementation and endpoint highest-weight representations.}
\GSfield{Present status}{As of November 2026: the stated result is established in the specified setting.}
\GSfield{Assessment}{literal statement: correct with missing hypotheses; publication match: partial; verification: read the result; development: annotated; confidence: high}
\end{GSitem}
\begin{GSitem}{1982-582.10}{F}{Dependence on the smoothness of a diffeomorphism}
\GSfield{Kind; source}{question; \GSpage{P001077}}
\GSfield{The notes}{The final question asks what effect smoothness has.}
\GSfield{Statement}{The notebook does not choose a regularity scale or ask a specific threshold question: extension from smooth diffeomorphisms to Sobolev, Hölder, or weaker maps gives different problems. Neretin’s construction quoted here is continuous for the smooth topology, which does not settle an unspecified optimal regularity threshold.}
\GSfield{Fate}{\citet[section 1 and final remark in section 3]{neretin1983boson} states the smooth-topology construction. No unique stronger question follows from the notebook.}
\GSfield{Present status}{As of November 2026: the notebook does not determine a unique proposition whose resolution can be assessed.}
\GSfield{Assessment}{literal statement: not determinable; publication match: partial; verification: unresolved lead; development: annotated; confidence: low}
\end{GSitem}

\GStalk{583}{1983-02-14}{V. L. Ginzburg}{P001076-P001080}{Duistermaat--Heckman localization and orbital integrals}
The entry occupies P001078 and the upper part of P001079. It presents Hamiltonian torus actions, regular reduction, affine variation of reduced symplectic classes, polynomial densities and exact stationary phase, then relates these to convexity and characters. The 1982 theorem is expository here. The nonisolated fixed-component formula is abbreviated beyond validity: normal-bundle curvature terms are needed. The historical allusions to Semenov-Tian-Shansky and Vergne have identifiable pre-talk sources.

\begin{GSitem}{1982-583.1}{E}{Hamiltonian torus actions and the moment map}
\GSfield{Kind; source}{definition; \GSpage{P001078}}
\GSfield{The notes}{A torus acts symplectically, each infinitesimal field has a Hamiltonian, and the Hamiltonians commute; equivalently, give commuting functions with periodic flows. The map is \(\mu(m)(X)=H_X(m)\).}
\GSfield{Statement}{Let a torus T act on a compact connected symplectic manifold \((M,\omega)\). A Hamiltonian action has a T-invariant map \(\mu:M\to\mathfrak t^*\) satisfying \(d\langle\mu,X\rangle=\iota_{X_M}\omega\), with the sign convention fixed accordingly. Conversely, commuting complete Hamiltonian flows whose joint period lattice defines T give such an action. Individual periodicity alone does not specify an effective torus of the claimed rank.}
\GSfield{Reconstruction}{Connectedness, the period lattice and moment-map sign are supplied; the later paper is exposition of the named 1982 theorem.}
\GSfield{Fate}{The setting is explicit in \citet[section 6]{atiyahbott1984moment}, explaining the already published 1982 result \citep{duistermaatheckman1982variation}.}
\GSfield{Present status}{As of November 2026: the stated result is established in the specified setting.}
\GSfield{Assessment}{literal statement: correct with missing hypotheses; publication match: exact; verification: read the result; development: annotated; confidence: high}
\end{GSitem}
\begin{GSitem}{1982-583.2}{E}{The cocycle behind commuting Hamiltonian fields}
\GSfield{Kind; source}{claim; \GSpage{P001078}}
\GSfield{The notes}{Commuting fields need not come from Poisson-commuting Hamiltonians; the example uses y and minus x on the plane. The cocycle must be trivial.}
\GSfield{Statement}{If Hamiltonian vector fields commute on connected M, their Poisson brackets are constant. These constants form a Lie-algebra 2-cocycle; an equivariant moment map requires its vanishing after the allowed normalization. For an abelian Lie algebra a nonzero such cocycle cannot be removed merely by adding constants. The plane translations illustrate this obstruction, but are not a periodic torus action on compact M.}
\GSfield{Reconstruction}{The scope of the noncompact translation example and the distinction between a zero cocycle and an arbitrary cohomologically trivial representative are made explicit.}
\GSfield{Fate}{This is the standard moment-map equivariance obstruction underlying \citet[section 6]{atiyahbott1984moment}. In the compact torus setting, averaging Hamiltonians gives an invariant moment map when the individual infinitesimal contractions are exact.}
\GSfield{Present status}{As of November 2026: the stated result is established in the specified setting.}
\GSfield{Assessment}{literal statement: correct with missing hypotheses; publication match: partial; verification: read the result; development: annotated; confidence: high}
\end{GSitem}
\begin{GSitem}{1982-583.3}{E}{Orbit tangent spaces are isotropic}
\GSfield{Kind; source}{claim; \GSpage{P001078}}
\GSfield{The notes}{The local condition says that the fields lie in one Lagrangian subspace at each point.}
\GSfield{Statement}{For a Hamiltonian torus action, \(\omega(X_M,Y_M)=0\); hence each tangent space to a torus orbit is isotropic and can be extended to a Lagrangian linear subspace of \(T_mM\). This pointwise extension does not imply a globally defined Lagrangian foliation or polarization.}
\GSfield{Reconstruction}{The notebook’s expressly local qualification is preserved; no global integrability claim is added.}
\GSfield{Fate}{The identity follows from the commuting moment-map components in \citet[section 6]{atiyahbott1984moment}.}
\GSfield{Present status}{As of November 2026: the stated result is established in the specified setting.}
\GSfield{Assessment}{literal statement: correct with missing hypotheses; publication match: partial; verification: read the result; development: annotated; confidence: high}
\end{GSitem}
\begin{GSitem}{1982-583.4}{E}{Regular symplectic reduction by the torus}
\GSfield{Kind; source}{construction; \GSpage{P001078}}
\GSfield{The notes}{\(M_\xi=\mu^{-1}(\xi)\) has a degenerate form, which becomes nondegenerate on \(M_\xi/T\).}
\GSfield{Statement}{If \(\xi\) is a regular value and T acts freely on its level, \(M_\xi^{red}=\mu^{-1}(\xi)/T\) is a smooth symplectic manifold with reduced form \(\omega_\xi\). A locally free action gives the corresponding symplectic orbifold. The characteristic kernel on the level is the orbit tangent distribution. Singular values require a stratified reduction, absent from this elementary statement.}
\GSfield{Reconstruction}{Regularity, freeness and a separate notation for the reduced space are restored.}
\GSfield{Fate}{The regular construction predates the talk \citep{marsdenweinstein1974reduction}; it is recalled explicitly in \citet[introduction and section 2.1]{meinrenkensjamaar1999singular}.}
\GSfield{Present status}{As of November 2026: the stated result is established in the specified setting.}
\GSfield{Assessment}{literal statement: correct with missing hypotheses; publication match: partial; verification: read the result; development: annotated; confidence: high}
\end{GSitem}
\begin{GSitem}{1982-583.5}{E}{The reduced symplectic class varies affinely}
\GSfield{Kind; source}{theorem; \GSpage{P001078}}
\GSfield{The notes}{\(\omega_\xi\) is piecewise linear in \(\xi\), modulo exact differentials.}
\GSfield{Statement}{On a connected chamber of regular values, identify nearby reduced spaces by the regular family. Their de Rham classes vary affinely: \([\omega_\xi]=[\omega_{\xi_0}]+\langle\xi-\xi_0,c\rangle\), where c is the fixed curvature/Chern-class vector with a convention-dependent sign. This is a statement about cohomology under the family identifications, not pointwise linearity of differential forms on unrelated fibers.}
\GSfield{Reconstruction}{Affine variation, chamber restrictions and the identification of cohomology groups are explicit.}
\GSfield{Fate}{The 1982 theorem is explained in \citet[section 7, pp. 22--23]{atiyahbott1984moment}; original publication \citet{duistermaatheckman1982variation}.}
\GSfield{Present status}{As of November 2026: the stated result is established in the specified setting.}
\GSfield{Assessment}{literal statement: correct with missing hypotheses; publication match: exact; verification: read the result; development: annotated; confidence: high}
\end{GSitem}
\begin{GSitem}{1982-583.6}{E}{Piecewise polynomial reduced volume and density}
\GSfield{Kind; source}{theorem; \GSpage{P001078}}
\GSfield{The notes}{The corollary is written \(\operatorname{vol}(M_\xi)\) is piecewise polynomial, although \(M_\xi\) was defined as the unreduced level set.}
\GSfield{Statement}{For a \(2n\)-dimensional compact Hamiltonian manifold with effective torus dimension k, the reduced symplectic volume on a regular chamber is a polynomial of degree at most \(n-k\). With compatible Haar and Lebesgue normalizations, it is the density of the pushforward of Liouville measure by the moment map. The symplectic volume here belongs to the reduced \(2(n-k)\)-dimensional space, not the presymplectic level set.}
\GSfield{Reconstruction}{The notation is corrected to the quotient already described one line earlier.}
\GSfield{Fate}{This is the Duistermaat--Heckman polynomial-density conclusion \citep[section 7]{atiyahbott1984moment}; the original appeared in 1982 \citep{duistermaatheckman1982variation}.}
\GSfield{Present status}{As of November 2026: the stated result is established in the specified setting.}
\GSfield{Assessment}{literal statement: false as written; publication match: exact; verification: read the result; development: annotated; confidence: high}
\end{GSitem}
\begin{GSitem}{1982-583.7}{E}{Exact stationary phase for isolated fixed points}
\GSfield{Kind; source}{theorem; \GSpage{P001078}}
\GSfield{The notes}{The integral \(\int_Me^{iH_X}\,dm\) is exactly calculable by stationary phase.}
\GSfield{Statement}{For a compact symplectic manifold with Hamiltonian torus action, choose X so that its zero set consists of isolated nondegenerate fixed points. With Liouville amplitude, the stationary-phase sum is exact: it is a finite sum of \(e^{itH_X(p)}\) divided by the product of the linearized weights evaluated at \(tX\), including the phase and measure-normalization factors. This assertion is special to the Hamiltonian setting and its canonical amplitude.}
\GSfield{Reconstruction}{Isolation, nondegeneracy, canonical amplitude and scale factors are specified without pretending the notebook fixes the powers of 2 pi or i.}
\GSfield{Fate}{\citet[section 7, equation (7.5)]{atiyahbott1984moment} gives the isolated-point formula and identifies it with the 1982 theorem \citep{duistermaatheckman1982variation}.}
\GSfield{Present status}{As of November 2026: the stated result is established in the specified setting.}
\GSfield{Assessment}{literal statement: correct with missing hypotheses; publication match: exact; verification: read the result; development: annotated; confidence: high}
\end{GSitem}
\begin{GSitem}{1982-583.8}{E}{Positive-dimensional fixed components need the equivariant Euler class}
\GSfield{Kind; source}{theorem; \GSpage{P001078}}
\GSfield{The notes}{Equation (5) replaces each fixed component by its volume divided by the product of transverse weights.}
\GSfield{Statement}{For a fixed component F with normal bundle \(N_F\), the localization contribution is the top-degree part of \(e^{itH_X(F)}\int_F e^\omega/e_T(N_F)(tX)\), in a consistent equivariant convention. The Euler denominator includes curvature/Chern classes of the normal bundle as well as its weights. Consequently it cannot in general be replaced by \(\operatorname{vol}(F)/\prod\lambda_j(X)\); that simplification is valid for isolated points and in special cases when higher characteristic contributions vanish.}
\GSfield{Reconstruction}{This is a mathematical emendation of the handwritten formula, not a translation correction.}
\GSfield{Fate}{The full formula is \citet[Theorem (1.5)]{duistermaatheckman1983addendum} and \citet[equation (7.4)]{atiyahbott1984moment}. The Addendum circulated as a Utrecht 1982 preprint, explicitly listed in the latter’s reference 12, and appeared in February 1983.}
\GSfield{Present status}{As of November 2026: the stated result is established in the specified setting.}
\GSfield{Assessment}{literal statement: false as written; publication match: exact; verification: read the result; development: annotated; confidence: high}
\end{GSitem}
\begin{GSitem}{1982-583.9}{E}{Equivariant localization explains the improvements}
\GSfield{Kind; source}{remark; \GSpage{P001079}}
\GSfield{The notes}{Atiyah and Bott generalized and improved the result; Michèle Vergne did something else on the subject.}
\GSfield{Statement}{The identifiable common framework is localization in equivariant cohomology: restriction to the fixed locus becomes invertible after inverting nonzero weights, and equivariant integration is computed through the inverse normal Euler class. The one-line Vergne remark does not uniquely identify which announcement, preprint or formulation was discussed.}
\GSfield{Reconstruction}{The framework is secure, but no exact title is assigned to the words something else.}
\GSfield{Fate}{\citet[section 3, equation (3.8)]{atiyahbott1984moment} supplies the cohomological form. \citet[introduction and Theorem 1.6]{berlinevergne1983zeros} identifies its preceding 1982 C. R. announcement and gives the isolated-zero differential-form formula. Both full publications postdate the talk; the earlier announcement and the notebook’s report establish the prepublication context.}
\GSfield{Present status}{As of November 2026: the stated result is established in the specified setting.}
\GSfield{Assessment}{literal statement: correct with missing hypotheses; publication match: partial; verification: read the result; development: annotated; confidence: medium}
\end{GSitem}
\begin{GSitem}{1982-583.10}{E}{Convexity of the torus moment image}
\GSfield{Kind; source}{theorem; \GSpage{P001079}}
\GSfield{The notes}{The old Atiyah and Guillemin--Sternberg theorem says that the image of the moment map is convex.}
\GSfield{Statement}{For a compact connected symplectic manifold with Hamiltonian torus action, the moment image is the convex hull of the images of fixed components and is a convex polytope; moment-map fibers are connected. Compactness, connectedness and the Hamiltonian hypothesis are essential to this formulation.}
\GSfield{Reconstruction}{The connected-fiber assertion is supplied from the source to clarify the geometry, rather than attributed to the single notebook sentence.}
\GSfield{Fate}{\citet[Theorem 1]{atiyah1982convexity} is an exact primary source. The supplied earlier index already records the theorem at 1980-508.15 and 1981-556.12.}
\GSfield{Present status}{As of November 2026: the stated result is established in the specified setting.}
\GSfield{Assessment}{literal statement: correct with missing hypotheses; publication match: exact; verification: read the result; development: annotated; confidence: high}
\end{GSitem}
\begin{GSitem}{1982-583.11}{E}{Quadratic local phase and finite Fourier expansion force exactness}
\GSfield{Kind; source}{method; \GSpage{P001079}}
\GSfield{The notes}{The proof uses local quadratic form, rapid decay of the remainder, and piecewise polynomial density. The displayed formula inserts t in the exponent but leaves its denominator unchanged.}
\GSfield{Statement}{For isolated fixed points, equivariant Darboux coordinates make the Hamiltonian exactly quadratic and Liouville amplitude constant locally. The pushforward density is compactly supported and piecewise polynomial, so its Fourier transform has a finite exponential Laurent expansion. Comparison with stationary phase gives exactness. Each isolated contribution scales as \(t^{-n}\) on a \(2n\)-manifold; omitting this factor from the displayed t-dependent formula is incorrect.}
\GSfield{Reconstruction}{The argument here is restricted to isolated fixed points; the positive-dimensional case uses the characteristic-class expansion.}
\GSfield{Fate}{This is the proof mechanism explained in \citet[section 7]{atiyahbott1984moment}; \citet[Lemma 2.7, Proposition 2.8 and section 4]{semenovtianshansky1973kirillov} gives the earlier compact-orbit quadratic-phase case.}
\GSfield{Present status}{As of November 2026: the stated result is established in the specified setting.}
\GSfield{Assessment}{literal statement: false as written; publication match: exact; verification: read the result; development: annotated; confidence: high}
\end{GSitem}
\begin{GSitem}{1982-583.12}{E}{The compact orbital integral is a corrected character}
\GSfield{Kind; source}{claim; \GSpage{P001079}; A. A. Kirillov}
\GSfield{The notes}{The coadjoint-orbit integral is compared with the Harish-Chandra integral and Weyl’s formula. Kirillov says: this is the character.}
\GSfield{Statement}{For a compact connected semisimple group and an integral dominant highest weight \(\lambda\), the local Kirillov character formula relates \(\chi_\lambda(\exp X)\) to the Fourier transform of the canonical measure on the shifted orbit \(\mathcal O_{\lambda+\rho}\), multiplied by the appropriate inverse square root of the exponential-map Jacobian. It is not the uncorrected transform of an arbitrary orbit. With the canonical normalization, the shifted orbit volume is the representation dimension.}
\GSfield{Reconstruction}{Integrality, the rho shift, canonical measure and Jacobian factor supply the conditions hidden by the short intervention.}
\GSfield{Fate}{\citet[section 3, Theorem 2, equations (26)--(28)]{kirillov1968characters} gives these shifts and normalization before this talk.}
\GSfield{Present status}{As of November 2026: the stated result is established in the specified setting.}
\GSfield{Assessment}{literal statement: correct with missing hypotheses; publication match: exact; verification: read the result; development: annotated; confidence: high}
\end{GSitem}
\begin{GSitem}{1982-583.13}{E}{The ten-year-old Semenov-Tian-Shansky reference}
\GSfield{Kind; source}{remark; \GSpage{P001079}}
\GSfield{The notes}{A paper by Semenov-Tian-Shansky with similar content appeared about ten years earlier.}
\GSfield{Statement}{For compact semisimple coadjoint orbits, the leading stationary-phase orbital-integral formula is exact. The relevant 1973 paper proves the local symplectic quadratic normal form and compares the result with the compact character formula. It explicitly leaves noncompact regularization outside its treatment.}
\GSfield{Reconstruction}{The exact title and compact scope are recovered from the primary text; the ten-year estimate is accurate.}
\GSfield{Fate}{The allusion matches \citet[introduction, Lemma 2.7, Proposition 2.8 and section 4]{semenovtianshansky1973kirillov}, published in 1973.}
\GSfield{Present status}{As of November 2026: the stated result is established in the specified setting.}
\GSfield{Assessment}{literal statement: correct with missing hypotheses; publication match: exact; verification: read the result; development: annotated; confidence: high}
\end{GSitem}
\begin{GSitem}{1982-583.14}{E}{Vergne’s proof for discrete-series orbits}
\GSfield{Kind; source}{remark; \GSpage{P001079}}
\GSfield{The notes}{Michèle Vergne has a proof for noncompact orbits of the discrete series.}
\GSfield{Statement}{For a connected semisimple real group in the discrete-series setting with compact Cartan subgroup, the relevant regular elliptic orbit Fourier transform gives the corresponding discrete-series character in Rossmann’s formula, with the specified invariant measures and distributional interpretation. This does not extend compact stationary phase to arbitrary noncompact orbits without conditions.}
\GSfield{Reconstruction}{The orbit type, character family and distributional qualification come from the exact named primary paper.}
\GSfield{Fate}{\citet[sections 2--4]{vergne1979rossmann} proves Rossmann’s character formula using oscillators and intertwining operators. The manuscript was received in March 1979, revised in July 1979, and published in 1979.}
\GSfield{Present status}{As of November 2026: the stated result is established in the specified setting.}
\GSfield{Assessment}{literal statement: correct with missing hypotheses; publication match: partial; verification: read the result; development: annotated; confidence: high}
\end{GSitem}

\GStalk{584}{1983-02-14}{A. A. Kirillov}{P001078-P001082}{The SU(3) orbit hexagon and weight multiplicities}
The intervention begins in the lower part of P001079 and ends with a single request on P001080. It revisits the six-dimensional regular orbit, its torus-moment hexagon and changes in the fibers. The written map uses two conjugacy invariants and therefore cannot produce that hexagon; diagonal projection is the required interpretation. The final weight-multiplicity request recurs from 1981-556 and retains GP-1981-41.

\begin{GSitem}{1982-584.1}{E}{Regular SU(3) orbits of Hermitian matrices}
\GSfield{Kind; source}{construction; \GSpage{P001079}}
\GSfield{The notes}{Trace-zero Hermitian 3 by 3 matrices form an eight-dimensional space; an orbit has two invariants, \(\operatorname{tr}A^2\) and \(\operatorname{tr}A^3\), and dimension six.}
\GSfield{Statement}{The real vector space of trace-zero Hermitian 3 by 3 matrices has dimension eight. For three distinct eigenvalues with sum zero, the conjugation orbit is \(SU(3)/T^2\), hence has real dimension six. The two indicated trace invariants determine the characteristic polynomial and thus the orbit. Repeated eigenvalues give lower-dimensional orbits.}
\GSfield{Reconstruction}{Distinct eigenvalues are the missing hypothesis for dimension six.}
\GSfield{Fate}{This is the coadjoint-orbit realization underlying the Schur--Horn example in \citet[introduction]{atiyah1982convexity}; it already occurs as 1981-556.3 in the earlier index.}
\GSfield{Present status}{As of November 2026: the stated result is established in the specified setting.}
\GSfield{Assessment}{literal statement: correct with missing hypotheses; publication match: partial; verification: read the result; development: annotated; confidence: high}
\end{GSitem}
\begin{GSitem}{1982-584.2}{E}{The hexagon map is diagonal projection, not the invariant map}
\GSfield{Kind; source}{claim; \GSpage{P001079}}
\GSfield{The notes}{The diagram maps A to \((\operatorname{tr}A^2,\operatorname{tr}A^3)\) and says the image of its orbit is a hexagon.}
\GSfield{Statement}{Both trace polynomials are constant on a conjugation orbit, so the displayed map has a one-point image. The torus moment map is instead \(A\mapsto(A_{11},A_{22},A_{33})\) in the plane \(x_1+x_2+x_3=0\). Its image is the convex hull of the permutations of the eigenvalue triple.}
\GSfield{Reconstruction}{The diagonal map is a justified emendation needed to recover the drawn hexagon; the literal invariant map is preserved in the ledger.}
\GSfield{Fate}{\citet[introduction, Schur--Horn example]{atiyah1982convexity} identifies diagonal projection and its convex hull. Direct conjugation also independently checks that the trace invariants cannot vary.}
\GSfield{Present status}{As of November 2026: the stated result is established in the specified setting.}
\GSfield{Assessment}{literal statement: false as written; publication match: exact; verification: read the result; development: annotated; confidence: high}
\end{GSitem}
\begin{GSitem}{1982-584.3}{E}{Geometry of the SU(3) permutahedral hexagon}
\GSfield{Kind; source}{theorem; \GSpage{P001079}}
\GSfield{The notes}{The image is almost a regular hexagon: all angles are 120 degrees, and alternate sides have equal lengths.}
\GSfield{Statement}{For \(\lambda_1>\lambda_2>\lambda_3\) with sum zero, the convex hull of the six permutations of \(\lambda\) is an equiangular hexagon in the trace-zero plane. The two alternating side lengths are \(\sqrt2(\lambda_1-\lambda_2)\) and \(\sqrt2(\lambda_2-\lambda_3)\); it is regular exactly when these differences agree. Repeated eigenvalues collapse it to a triangle or a point.}
\GSfield{Reconstruction}{The eigenvalue inequalities and Euclidean metric on the trace-zero plane specify the notebook’s almost regular description.}
\GSfield{Fate}{The convex-hull description is the classical matrix example in \citet[introduction]{atiyah1982convexity}. The angle and length statements follow by subtracting adjacent permuted vertices. Earlier occurrence: 1981-556.3.}
\GSfield{Present status}{As of November 2026: the stated result is established in the specified setting.}
\GSfield{Assessment}{literal statement: correct with missing hypotheses; publication match: exact; verification: read the result; development: annotated; confidence: high}
\end{GSitem}
\begin{GSitem}{1982-584.4}{E}{Changing fibers and piecewise linear reduced volume}
\GSfield{Kind; source}{remark; \GSpage{P001079}}
\GSfield{The notes}{The fiber structure changes, the singularities are known, and the notes ask: polynomials?}
\GSfield{Statement}{On a regular six-dimensional orbit, a regular value of the rank-two moment map has a four-dimensional level set; its torus reduction has real dimension two. The reduced symplectic area, equivalently the normalized Duistermaat--Heckman density, is piecewise affine on regular chambers. The handwritten crossing lines do not uniquely label all singular strata, so no complete fiber classification is recovered from the sketch alone.}
\GSfield{Reconstruction}{Only the dimension and density consequences fixed by the regular-value hypotheses are asserted; labels absent from the drawing are not invented.}
\GSfield{Fate}{The degree bound follows from the 1982 polynomial-density theorem \citep[section 7]{atiyahbott1984moment}. The same level/reduced dimension distinction occurs at 1981-556.4 in the earlier index.}
\GSfield{Present status}{As of November 2026: the stated result is established in the specified setting.}
\GSfield{Assessment}{literal statement: correct with missing hypotheses; publication match: partial; verification: read the result; development: annotated; confidence: high}
\end{GSitem}
\begin{GSitem}{1982-584.5}{E}{Recovering weight multiplicities geometrically}
\GSfield{Kind; source}{problem; \GSpage{P001080}}
\GSfield{The notes}{\ru{По выписанным выше формулам вычислить кратности веса}: calculate weight multiplicities from the preceding formulas.}
\GSfield{Statement}{For an integral dominant \(SU(3)\) highest weight, recover each torus-weight multiplicity from the corresponding orbit/character data. Exact multiplicities require the integral lattice, character normalization and its shifts, or an equivariant quantization with the correctly shifted reduced line bundle; classical reduced area by itself is not the finite-dimensional multiplicity.}
\GSfield{Reconstruction}{An integral highest weight and the distinction between classical density and quantum index are required to make the one-line request precise.}
\GSfield{Fate}{Exact finite-dimensional multiplicities were already computable from the interlacing basis of \citet[section 3, formula (3)]{gelfandtsetlin1950unimodular} and the shifted orbital character formula \citep[section 3]{kirillov1968characters}. The geometric reduced-index formulation, including singular quotients, is supplied by \citet[Theorem 2.9 and Corollary 2.11]{meinrenkensjamaar1999singular}, first circulated July 1997, published 1999. This repeats 1981-556.1 and 1981-556.6.}
\GSfield{Present status}{As of November 2026: compact integral-orbit weight multiplicities are established; the later reduced-index theorem explains the geometric interpretation, including singular levels.}
\GSfield{Assessment}{literal statement: correct with missing hypotheses; publication match: exact; verification: read the result; development: annotated; confidence: high; problem \hyperref[prob:GP-1981-41]{GP-1981-41}}
\end{GSitem}

\GStalk{585}{1983-02-21}{V. I. Arnold}{P001080-P001088}{Mathieu gaps, Catalan diagrams and spectral perturbation}
The actual talk is P001081--P001086; P001080 belongs to Kirillov and P001087 begins Krichever. Arnold explains delayed splitting of Hill eigenvalues for trigonometric-polynomial perturbations by quiver paths, Catalan bracket fillings and a reversal involution. The paper received on 1 February 1983 and published later that year closely matches the talk. Its 30 March supplement acknowledges older results of Bloch, Levy--Keller and Hochstadt, so the estimates and coefficient identities are not all classified as new 1983 discoveries. Gelfand’s one-gap question remains underdetermined in this record.

\begin{GSitem}{1982-585.1}{E}{Instability tongues and periodic or antiperiodic spectra}
\GSfield{Kind; source}{construction; \GSpage{P001081}}
\GSfield{The notes}{\(\ddot x+(\omega^2+\varepsilon a(t))x=0\), \(a(t+2\pi)=a(t)\); the drawing shows tongues at \(\omega=1/2,1,3/2,\ldots\).}
\GSfield{Statement}{For real periodic a, the monodromy has determinant one. Instability occurs when one Floquet multiplier has modulus greater than one; the other then has modulus less than one. Tongue boundaries are periodic or antiperiodic eigenvalue loci. For the unperturbed equation their positive-frequency tips are \(\omega=N/2\), \(N\ge1\).}
\GSfield{Reconstruction}{Reality and the monodromy criterion are supplied; the diagram’s half-integer labeling is retained.}
\GSfield{Fate}{The classical spectral interpretation is recalled in \citet[section 3, proof of Theorem 3]{arnold1983mathieu}.}
\GSfield{Present status}{As of November 2026: the stated result is established in the specified setting.}
\GSfield{Assessment}{literal statement: correct with missing hypotheses; publication match: exact; verification: read the result; development: annotated; confidence: high}
\end{GSitem}
\begin{GSitem}{1982-585.2}{E}{Generic first-order opening of Hill gaps}
\GSfield{Kind; source}{claim; \GSpage{P001081}}
\GSfield{The notes}{For a coefficient a in general position the boundaries have first-order contact everywhere.}
\GSfield{Statement}{Fix \(N\ge1\). For a real bounded \(2\pi\)-periodic perturbation with nonzero Fourier coefficient \(a_N\), degenerate perturbation theory on \(\operatorname{span}(e^{iNt/2},e^{-iNt/2})\) gives a spectral gap \(2|\varepsilon a_N|+O(\varepsilon^2)\) in the eigenvalue variable. Thus a coefficient with every \(a_N\ne0\) has first-order opening at every fixed positive tip. This is not a uniform-in-N remainder statement.}
\GSfield{Reconstruction}{The nonvanishing Fourier coefficient and the chosen eigenvalue variable make the genericity statement precise.}
\GSfield{Fate}{The contrast between generic Hill equations and polynomial perturbations is explicit in \citet[introduction and section 3]{arnold1983mathieu}; the first coefficient is the restriction of A to the double eigenspace.}
\GSfield{Present status}{As of November 2026: the stated result is established in the specified setting.}
\GSfield{Assessment}{literal statement: correct with missing hypotheses; publication match: partial; verification: read the result; development: annotated; confidence: high}
\end{GSitem}
\begin{GSitem}{1982-585.3}{E}{Exponential gap decay requires analytic regularity}
\GSfield{Kind; source}{claim; \GSpage{P001081}}
\GSfield{The notes}{The generic Hill statement is followed by: the width decreases exponentially.}
\GSfield{Statement}{For a real smooth periodic potential, high-energy gap lengths decay faster than any inverse power of their index. Exponential decay in the index is the analytic-potential conclusion and is not guaranteed by smoothness or periodicity alone. This concerns the high-index limit at a fixed potential, distinct from the small-amplitude limit at a fixed gap.}
\GSfield{Reconstruction}{The missing regularity and limiting regime are restored; the notebook does not state analyticity.}
\GSfield{Fate}{\citet[section 1, historical equivalences and Theorems 1, 3]{poschel2004weighted} distinguishes the smooth and analytic cases and credits the pre-talk Hochstadt and Trubowitz results. Its 2004 preprint is used for verification, not as a claim of first discovery.}
\GSfield{Present status}{As of November 2026: the stated result is established in the specified setting.}
\GSfield{Assessment}{literal statement: false as written; publication match: partial; verification: read the result; development: annotated; confidence: high}
\end{GSitem}
\begin{GSitem}{1982-585.4}{E}{Mathieu’s cosine perturbation has increasingly high tangency}
\GSfield{Kind; source}{claim; \GSpage{P001081}}
\GSfield{The notes}{For \(a(t)=\cos t\), successive zone boundaries have increasing orders of contact; this is the degree-one example.}
\GSfield{Statement}{For the real Mathieu equation and each fixed \(N\ge1\), the Nth gap opens at order \(|\varepsilon|^N\), with a nonzero coefficient. The claim concerns small amplitude at fixed N and the standard positive-frequency branch; it should not be confused with a uniform estimate over all gaps.}
\GSfield{Reconstruction}{The fixed-index asymptotic and nonzero cosine coefficient are made explicit.}
\GSfield{Fate}{The classical Mathieu behavior is recalled in \citet[introduction]{arnold1983mathieu}; the paper’s supplement identifies the earlier Hill-gap literature \citep{levykeller1963instability,hochstadt1964instability}.}
\GSfield{Present status}{As of November 2026: the stated result is established in the specified setting.}
\GSfield{Assessment}{literal statement: correct with missing hypotheses; publication match: partial; verification: read the result; development: annotated; confidence: high}
\end{GSitem}
\begin{GSitem}{1982-585.5}{E}{The polynomial-degree gap bound uses a ceiling and an upper estimate}
\GSfield{Kind; source}{theorem; \GSpage{P001081-P001085}}
\GSfield{The notes}{Equation (3) writes the width as \(\asymp\varepsilon^{[N/p]}\) when \(\deg a=p\). The later proof only delays splitting.}
\GSfield{Statement}{For a fixed \(N\ge1\) and a trigonometric polynomial a of degree \(p\ge1\), the two periodic or antiperiodic eigenvalues issuing from the \(N/2\) double frequency satisfy \(|\lambda_N^+-\lambda_N^-|=O(|\varepsilon|^{\lceil N/p\rceil})\). For real a this bounds the Nth instability gap. A universal matching lower bound is false: coefficients can cancel and gaps can remain closed. The ceiling is \(-\lfloor-N/p\rfloor\), not \(\lfloor N/p\rfloor\).}
\GSfield{Reconstruction}{The bracket notation is unexplained in the notebook; the paper fixes the ceiling. Independently, the asymptotic equivalence is too strong even if the brackets meant ceiling.}
\GSfield{Fate}{The precise statements are \citet[Theorems 3 and 7]{arnold1983mathieu}. His Supplement dated 30 March 1983 credits \citet{levykeller1963instability} for the even case and \citet{hochstadt1964instability} for the general case. The original antecedent papers were metadata-checked; their theorem comparison is Arnold’s explicit report.}
\GSfield{Present status}{As of November 2026: the qualified gap bound is established and predates the talk; Arnold’s proof and its exact match to the notebook are verified.}
\GSfield{Assessment}{literal statement: false as written; publication match: exact; verification: read the result; development: annotated; confidence: high}
\end{GSitem}
\begin{GSitem}{1982-585.6}{F}{Perturbing a one-gap potential}
\GSfield{Kind; source}{question; \GSpage{P001081}; I. M. Gelfand}
\GSfield{The notes}{Gelfand observes that the constant potential is zero-gap and asks what happens near a one-gap potential.}
\GSfield{Statement}{The question suggests replacing the constant background by a real periodic one-gap potential and studying perturbed spectral gaps. It does not specify the perturbation class, which spectral gaps, a norm, or whether the target is a first variation, a uniform asymptotic, or persistence of finite-gap structure. These lead to distinct problems.}
\GSfield{Fate}{The next meeting, items 1982-586.3--6, supplies a thematic continuation: \citet[Theorems 1--3, Lemma 3 and equations (6), (13)--(19)]{krichever1983hessians} gives a precise first-order finite-gap perturbation framework, including fixed-period constraints, Hessian diagonalization and gap-opening modes. The paper was presented on 19 October 1982, received on 23 November 1982 and published in 1983; those administrative dates do not establish public circulation. The spectral framework and finite-gap density in \citet[Theorems 1.1--1.2]{grebertkappelerposchel2004gaps} are another related result. Neither publication resolves all possible readings of the notebook’s sentence. Current primary-source searches did not identify the missing formulation.}
\GSfield{Present status}{As of November 2026: the historical question needs a specified perturbation and asymptotic regime before it can be designated solved or open; the link to the following meeting is thematic and does not establish an identical problem.}
\GSfield{Assessment}{literal statement: not determinable; publication match: thematic only; verification: unresolved lead; development: annotated; confidence: low}
\end{GSitem}
\begin{GSitem}{1982-585.7}{E}{Rayleigh--Schrödinger expansion at a simple eigenvalue}
\GSfield{Kind; source}{method; \GSpage{P001081}}
\GSfield{The notes}{Equation (4) expands the eigenvector and eigenvalue in powers of epsilon for \(\Lambda+\varepsilon A\).}
\GSfield{Statement}{For a finite-dimensional diagonalizable operator with a simple eigenvalue \(\lambda_0\), analytic perturbation gives a normalized local eigenvector and eigenvalue power series. Equating coefficients recursively computes them after fixing the component of the eigenvector in the original eigenspace. A formal expansion can also be considered without claiming convergence.}
\GSfield{Reconstruction}{Simplicity, normalization and finite-dimensional analytic validity are supplied.}
\GSfield{Fate}{\citet[section 1]{arnold1983mathieu} starts from this classical expansion and later acknowledges \citet{bloch1958perturbations} for general diagrammatic coefficient formulas.}
\GSfield{Present status}{As of November 2026: the stated result is established in the specified setting.}
\GSfield{Assessment}{literal statement: correct with missing hypotheses; publication match: exact; verification: read the result; development: annotated; confidence: high}
\end{GSitem}
\begin{GSitem}{1982-585.8}{E}{Catalan closed form and the erroneous seventh entry}
\GSfield{Kind; source}{definition; \GSpage{P001082}}
\GSfield{The notes}{The sequence is \(1,1,2,5,14,42,138\), and \(a_n=(2n)!/(n!(n+1)!)\).}
\GSfield{Statement}{With indexing \(C_0=1\), the Catalan numbers are \(C_n=(2n)!/(n!(n+1)!)\); thus \(C_6=132\), not 138. The closed form in the notes agrees with the corrected sequence.}
\GSfield{Reconstruction}{The starting index is specified to align the two initial ones.}
\GSfield{Fate}{\citet[section 4, Catalan remark]{arnold1983mathieu} prints 132 and the same closed formula. The value was independently checked from both the factorial formula and convolution recurrence.}
\GSfield{Present status}{As of November 2026: the stated result is established in the specified setting.}
\GSfield{Assessment}{literal statement: false as written; publication match: exact; verification: read the result; development: annotated; confidence: high}
\end{GSitem}
\begin{GSitem}{1982-585.9}{E}{The Catalan convolution has incorrect bounds}
\GSfield{Kind; source}{claim; \GSpage{P001082}}
\GSfield{The notes}{Equation (5) is \(a_{n+1}=\sum_{j=1}^{n-1}a_ja_{n-j}\).}
\GSfield{Statement}{For the stated indexing, \(C_{n+1}=\sum_{j=0}^{n}C_jC_{n-j}\). The omitted endpoint terms are necessary. The generating series \(C(z)=\sum_{n\ge0}C_nz^n\) satisfies \(C=1+zC^2\), with the branch determined by \(C(0)=1\).}
\GSfield{Reconstruction}{The generating equation is the explicit version of the notebook’s generating-function instruction.}
\GSfield{Fate}{The Catalan objects in \citet[section 4]{arnold1983mathieu} obey this standard recurrence. The notebook’s printed bounds already fail at n equal to one.}
\GSfield{Present status}{As of November 2026: the stated result is established in the specified setting.}
\GSfield{Assessment}{literal statement: false as written; publication match: partial; verification: read the result; development: annotated; confidence: high}
\end{GSitem}
\begin{GSitem}{1982-585.10}{E}{The polygon interpretation has an index shift}
\GSfield{Kind; source}{claim; \GSpage{P001082}}
\GSfield{The notes}{The number \(a_n=(2n)!/(n!(n+1)!)\) is called the number of simplicial subdivisions of an n-gon.}
\GSfield{Statement}{The number \(C_n\) counts triangulations of a convex \((n+2)\)-gon by noncrossing diagonals, with no added vertices. A convex n-gon has \(C_{n-2}\) such triangulations. Allowing arbitrary additional vertices would not give this finite count.}
\GSfield{Reconstruction}{Convexity, triangulation and the prohibition on new vertices are restored.}
\GSfield{Fate}{\citet[section 4, Catalan remark]{arnold1983mathieu} explicitly prints the \((n+2)\)-gon.}
\GSfield{Present status}{As of November 2026: the stated result is established in the specified setting.}
\GSfield{Assessment}{literal statement: false as written; publication match: exact; verification: read the result; development: annotated; confidence: high}
\end{GSitem}
\begin{GSitem}{1982-585.11}{E}{Dimension of curve moduli}
\GSfield{Kind; source}{remark; \GSpage{P001082}}
\GSfield{The notes}{Riemann surfaces of genus g have \(3g-3\) moduli; branch points on the sphere are invoked.}
\GSfield{Statement}{For smooth compact complex curves of genus \(g\ge2\), the moduli space has complex dimension \(3g-3\). The formula requires separate treatment in genera zero and one. In the subsequent genus-\(2k\) covering example, \(6k\) simple branch points modulo target projective transformations give dimension \(6k-3=3g-3\).}
\GSfield{Reconstruction}{Genus restriction and dimension over the complex numbers are editorial specifications; the branch count belongs to the following special example.}
\GSfield{Fate}{The specific covering example is identified by \citet[section 4, footnote]{arnold1983mathieu}; the dimension comparison follows from Riemann--Hurwitz and the target automorphism quotient.}
\GSfield{Present status}{As of November 2026: the stated result is established in the specified setting.}
\GSfield{Assessment}{literal statement: correct with missing hypotheses; publication match: partial; verification: read the result; development: annotated; confidence: high}
\end{GSitem}
\begin{GSitem}{1982-585.12}{E}{Catalan many minimal pencils on a general even-genus curve}
\GSfield{Kind; source}{theorem; \GSpage{P001082}}
\GSfield{The notes}{Griffiths and Harris count ways of obtaining one and the same surface as a covering; the answer is the Catalan sequence.}
\GSfield{Statement}{A general smooth complex curve of genus \(2k\), \(k\ge1\), has \(C_k\) pencils of degree \(k+1\), equivalently maps to the projective line modulo target automorphisms, with \(6k\) simple branch points in the generic case. Special curves need not have this finite count.}
\GSfield{Reconstruction}{Even genus, general curve, degree and equivalence under target automorphisms are recovered from the footnote rather than the abbreviated notebook.}
\GSfield{Fate}{\citet[section 4, footnote]{arnold1983mathieu} states exactly this example and identifies Griffiths--Harris, Principles of Algebraic Geometry (1978), as the antecedent. The book’s result was not independently opened, so this historical attribution is indirect through Arnold’s primary paper.}
\GSfield{Present status}{As of November 2026: the stated result is established in the specified setting.}
\GSfield{Assessment}{literal statement: correct with missing hypotheses; publication match: partial; verification: indirect; development: annotated; confidence: high}
\end{GSitem}
\begin{GSitem}{1982-585.13}{A}{Bracket symbols as Catalan objects}
\GSfield{Kind; source}{definition; \GSpage{P001082}}
\GSfield{The notes}{A properly formed string of letters and parentheses has length equal to letters plus opening parentheses; the number of shapes is Catalan.}
\GSfield{Statement}{Consider balanced parenthesis words with nonempty contents in each pair, and letters serving as indistinguishable ordered slots. Give each letter and each opening parenthesis weight one. The shapes of length n are counted by \(C_n\); length two gives \(ab,(a)\), and length three gives five shapes. Letters become independently chosen quiver vertices only in a later filling.}
\GSfield{Reconstruction}{Distinguishing abstract shapes from fillings prevents counting infinitely many alphabet choices.}
\GSfield{Fate}{This exact notation and its counting appear in \citet[section 4]{arnold1983mathieu}, received 1 February and published later in 1983. The Supplement identifies the equivalent Catalan staircase diagrams in Bloch’s 1958 approach; A marks this particular presentation, not priority for Catalan perturbation theory.}
\GSfield{Present status}{As of November 2026: the stated result is established in the specified setting.}
\GSfield{Assessment}{literal statement: correct with missing hypotheses; publication match: exact; verification: read the result; development: annotated; confidence: high}
\end{GSitem}
\begin{GSitem}{1982-585.14}{A}{The spectral quiver and its chronological products}
\GSfield{Kind; source}{construction; \GSpage{P001082-P001083}}
\GSfield{The notes}{The eigenspaces \(V_i\) are vertices and \(A_{ij}:V_i\to V_j\) are arrows. Products are taken following the arrows.}
\GSfield{Statement}{Let \(\Lambda\) act as the distinct scalar \(\lambda_i\) on each block \(V_i\). Draw an arrow \(i\to j\) when the perturbing block \(A_{ij}\) is nonzero. For a path \(i_0\to i_1\to\cdots\to i_k\), its chronological product is \(A_{i_{k-1}i_k}\cdots A_{i_0i_1}:V_{i_0}\to V_{i_k}\), with composition read right to left.}
\GSfield{Reconstruction}{The direction convention is fixed from the handwritten map, so no transposition of the block indices is silently introduced.}
\GSfield{Fate}{\citet[section 1, definitions]{arnold1983mathieu} gives this exact quiver presentation. Diagrammatic perturbation itself is older, as the paper’s Supplement explicitly acknowledges.}
\GSfield{Present status}{As of November 2026: the stated result is established in the specified setting.}
\GSfield{Assessment}{literal statement: correct with missing hypotheses; publication match: exact; verification: read the result; development: annotated; confidence: high}
\end{GSitem}
\begin{GSitem}{1982-585.15}{E}{Local finiteness controls formal coefficients, not convergence by itself}
\GSfield{Kind; source}{claim; \GSpage{P001083}}
\GSfield{The notes}{Everything is assumed finite or locally finite; for Mathieu-type operators the quiver is locally finite, and the perturbed subspace is called analytic.}
\GSfield{Statement}{A locally finite quiver has finitely many paths of any fixed length issuing from a fixed vertex, so the prescribed perturbation coefficient is a finite sum. This is a formal statement. In finite dimensions a spectrally isolated cluster has an analytic spectral projection; for Hill operators bounded polynomial perturbations and the elliptic unperturbed operator give analytic cluster perturbation. Local finiteness alone does not prove convergence for an arbitrary infinite operator.}
\GSfield{Reconstruction}{The analytic claim is confined to a specified operator setting, not inferred from the graph alone.}
\GSfield{Fate}{\citet[section 1, Remark 2]{arnold1983mathieu} explicitly makes this distinction: it treats formal series or finite matrices, and invokes ellipticity for Mathieu convergence.}
\GSfield{Present status}{As of November 2026: the stated result is established in the specified setting.}
\GSfield{Assessment}{literal statement: correct with missing hypotheses; publication match: partial; verification: read the result; development: annotated; confidence: high}
\end{GSitem}
\begin{GSitem}{1982-585.16}{E}{The perturbed invariant subspace as a graph}
\GSfield{Kind; source}{construction; \GSpage{P001083}}
\GSfield{The notes}{\(V_0(\varepsilon)\) is the graph of \(B(\varepsilon):V_0\to\bigoplus_{i\ne0}V_i\), and equation (7) imposes invariance.}
\GSfield{Statement}{For a finite-dimensional diagonalizable \(\Lambda\) and an isolated eigenvalue \(\lambda_0\), the small perturbed spectral subspace is the graph of an analytic \(B(\varepsilon)\) with \(B(0)=0\), relative to the fixed spectral complement. The invariance equation is \((\Lambda+\varepsilon A)(\iota+B)=(\iota+B)S(\varepsilon)\), where \(\iota:V_0\hookrightarrow V\). This notation avoids treating a projection and an inclusion as the same typed map.}
\GSfield{Reconstruction}{Isolation, fixed complement and the inclusion convention are supplied.}
\GSfield{Fate}{\citet[section 2]{arnold1983mathieu} writes the graph equation. The general perturbation method predates the talk; the Supplement credits \citet{bloch1958perturbations}.}
\GSfield{Present status}{As of November 2026: the stated result is established in the specified setting.}
\GSfield{Assessment}{literal statement: correct with missing hypotheses; publication match: exact; verification: read the result; development: annotated; confidence: high}
\end{GSitem}
\begin{GSitem}{1982-585.17}{E}{The secular operator preserves the cluster spectrum}
\GSfield{Kind; source}{theorem; \GSpage{P001083}}
\GSfield{The notes}{The operator \(\lambda_0+\varepsilon\alpha^1+\cdots\) is called secular; restriction to \(V_0(\varepsilon)\) is similar to it.}
\GSfield{Statement}{The effective operator \(S(\varepsilon)=(\iota+B)^{-1}(\Lambda+\varepsilon A)|_{\operatorname{graph}B}(\iota+B)\) acts on the fixed space \(V_0\) and has exactly the perturbed cluster eigenvalues with algebraic multiplicity. Analyticity of the matrix S does not by itself imply analyticity of each individual eigenvalue for a general non-self-adjoint family.}
\GSfield{Reconstruction}{The inverse is restricted to the graph and multiplicity is specified; no unsupported eigenvalue-analyticity claim is added.}
\GSfield{Fate}{The similarity and spectrum assertion is explicit in \citet[section 2]{arnold1983mathieu}; the exact notebook formulation matches it.}
\GSfield{Present status}{As of November 2026: the stated result is established in the specified setting.}
\GSfield{Assessment}{literal statement: correct with missing hypotheses; publication match: exact; verification: read the result; development: annotated; confidence: high}
\end{GSitem}
\begin{GSitem}{1982-585.18}{E}{Closed paths give universal coefficient formulas}
\GSfield{Kind; source}{theorem; \GSpage{P001083}}
\GSfield{The notes}{\(\alpha_k=\sum_{|\gamma|=k}\Pi_\gamma(A)Q_\gamma(\Lambda)\), where the loops start and end at zero and Q is rational in eigenvalues.}
\GSfield{Statement}{In the graph-normalized effective operator, the coefficient of \(\varepsilon^k\) is a sum over length-k closed paths at the chosen eigenspace. The factor \(\Pi_\gamma(A)\) is the ordered block product, and \(Q_\gamma\) is a universal rational function of the eigenvalue differences. If several parenthesis fillings yield the same loop, their denominator contributions must be added in \(Q_\gamma\).}
\GSfield{Reconstruction}{Graph normalization and the aggregation of different fillings specify the otherwise ambiguous rational coefficient.}
\GSfield{Fate}{This is \citet[Theorems 1--2 and section 4, Remark 1]{arnold1983mathieu}. His dated Supplement explicitly credits Bloch’s 1958 antecedent for these coefficient results \citep{bloch1958perturbations}; Arnold’s proof was read directly.}
\GSfield{Present status}{As of November 2026: the stated result is established in the specified setting.}
\GSfield{Assessment}{literal statement: correct with missing hypotheses; publication match: exact; verification: read the result; development: annotated; confidence: high}
\end{GSitem}
\begin{GSitem}{1982-585.19}{A}{Filled bracket diagrams remember the denominator data}
\GSfield{Kind; source}{construction; \GSpage{P001083-P001084}}
\GSfield{The notes}{Insert vertices into a bracket symbol; opening parentheses become zeros, closing parentheses are omitted from the solid path, and dashed arrows are then added. Different fillings can have one loop.}
\GSfield{Statement}{Fill each letter slot by a nonzero quiver vertex. Put zero at both ends and at every opening parenthesis to form a solid path; require its arrows to be allowed. Add a dashed arrow from the vertex immediately preceding each closing parenthesis to the zero of its matching opening parenthesis. A symbol of length \(k-1\) yields a loop of k solid arrows. Dashed arrows encode information lost by the solid loop alone.}
\GSfield{Reconstruction}{Admissibility and the exclusion of zero from letter slots are restored from the published definitions.}
\GSfield{Fate}{The definitions and examples are an exact match to \citet[section 4]{arnold1983mathieu}. This is Arnold’s particular diagrammatic construction; the underlying coefficient expansion is older.}
\GSfield{Present status}{As of November 2026: the stated result is established in the specified setting.}
\GSfield{Assessment}{literal statement: correct with missing hypotheses; publication match: exact; verification: read the result; development: annotated; confidence: high}
\end{GSitem}
\begin{GSitem}{1982-585.20}{E}{The denominator rule for each filled diagram}
\GSfield{Kind; source}{theorem; \GSpage{P001084}}
\GSfield{The notes}{Equation (11) sums \(\Pi_{\gamma(I)}(A)/q_I(\Lambda)\); equation (12) multiplies \(\lambda_0-\lambda_i\) for solid arrows into nonzero i and \(\lambda_j-\lambda_0\) for dashed arrows from nonzero j.}
\GSfield{Statement}{For each admissible filling I of length \(k-1\), take exactly the indicated signed product \(q_I\), counting arrows with multiplicity. Its contribution to \(\alpha_k\) is the chronological solid-loop product divided by \(q_I\). Summing over all fillings gives the coefficient. In particular the filling \((i)j\) has denominator \((\lambda_0-\lambda_i)(\lambda_0-\lambda_j)(\lambda_i-\lambda_0)\), agreeing with the notebook example.}
\GSfield{Reconstruction}{The exact sign and multiplicities are retained; no denominator factor is changed.}
\GSfield{Fate}{This is \citet[Theorem 4 and the denominator definition]{arnold1983mathieu}. The Supplement credits the equivalent coefficient theorem to Bloch 1958, while recording Arnold’s distinct parenthesis notation \citep{bloch1958perturbations}.}
\GSfield{Present status}{As of November 2026: the stated result is established in the specified setting.}
\GSfield{Assessment}{literal statement: correct with missing hypotheses; publication match: exact; verification: read the result; development: annotated; confidence: high}
\end{GSitem}
\begin{GSitem}{1982-585.21}{E}{Fourier blocks of a trigonometric perturbation}
\GSfield{Kind; source}{construction; \GSpage{P001085}}
\GSfield{The notes}{\(V_k=\langle e_k,e_{-k}\rangle\), \(\Lambda=d^2/dt^2\), and equation (16) gives a two by two matrix marked with a question mark.}
\GSfield{Statement}{For \(i,j>0\) and \(A_{ij}:V_i\to V_j\), in the ordered bases \((e_i,e_{-i})\), \((e_j,e_{-j})\), multiplication by \(a(t)=\sum a_se^{ist}\) has block \(\begin{pmatrix}a_{j-i}&a_{i+j}\\a_{-i-j}&a_{i-j}\end{pmatrix}\). The printed signs are correct for that convention. The zero-frequency block has dimension one; odd-numbered tongues instead use the half-integer Fourier lattice.}
\GSfield{Reconstruction}{The zero block and antiperiodic sector are specified; the handwritten question mark does not require a false correction.}
\GSfield{Fate}{\citet[section 3, Lemma 1 and preceding setup]{arnold1983mathieu} gives the same matrix and treats both integer and half-integer frequencies.}
\GSfield{Present status}{As of November 2026: the stated result is established in the specified setting.}
\GSfield{Assessment}{literal statement: correct with missing hypotheses; publication match: exact; verification: read the result; development: annotated; confidence: high}
\end{GSitem}
\begin{GSitem}{1982-585.22}{E}{Remote blocks are scalar for an even polynomial}
\GSfield{Kind; source}{claim; \GSpage{P001085}}
\GSfield{The notes}{If \(i+j>p\), off-diagonal entries vanish; for \(a_{-s}=a_s\), the block is scalar. Products of scalar blocks remain scalar.}
\GSfield{Statement}{If a has Fourier support \(|s|\le p\) and \(i+j>p\), the off-diagonal block entries vanish. If a is even, the two remaining diagonal entries coincide, so \(A_{ij}=a_{j-i}I\). For a general non-even polynomial the remote block is only diagonal; equality of the summed effective diagonal entries needs the later involution.}
\GSfield{Reconstruction}{The distinction between evenness and reality is explicit: an even polynomial may have complex coefficients.}
\GSfield{Fate}{This is \citet[section 3, Lemma 2 and Remark 2]{arnold1983mathieu}; it is the elementary ingredient of the older even-gap estimate.}
\GSfield{Present status}{As of November 2026: the stated result is established in the specified setting.}
\GSfield{Assessment}{literal statement: correct with missing hypotheses; publication match: exact; verification: read the result; development: annotated; confidence: high}
\end{GSitem}
\begin{GSitem}{1982-585.23}{A}{Reflection of a short loop delays splitting}
\GSfield{Kind; source}{method; \GSpage{P001085}}
\GSfield{The notes}{The combinatorial lemma says \(\tau p<2k\) implies \(\alpha_\tau\) is scalar. A loop is reflected across zero at a nonremote arrow.}
\GSfield{Statement}{A loop based at k on the nonnegative frequency lattice, with \(\tau\) arrows each of length at most p, cannot contain an arrow \(i\to j\) with \(i+j\le p\) when \(\tau p<2k\): reflecting the subsequent path would give a path from k to minus k of total length at most \(\tau p\). Hence all factors are remote. For even a, this makes \(\alpha_\tau\) scalar and prevents splitting to that order.}
\GSfield{Reconstruction}{The reflected path is a length argument; the notebook’s parenthetical homotopy language is not needed for a stronger topological theorem.}
\GSfield{Fate}{The proof is exactly \citet[section 3, Lemma 3 and proof of Theorem 3]{arnold1983mathieu}. A labels this short combinatorial proof; the resulting spectral bound itself is E because of Levy--Keller’s 1963 result.}
\GSfield{Present status}{As of November 2026: the stated result is established in the specified setting.}
\GSfield{Assessment}{literal statement: correct with missing hypotheses; publication match: exact; verification: read the result; development: annotated; confidence: high}
\end{GSitem}
\begin{GSitem}{1982-585.24}{A}{Dual bracket fillings equalize the diagonal corrections}
\GSfield{Kind; source}{method; \GSpage{P001086}}
\GSfield{The notes}{Opening parentheses split a word into fragments; reverse letters within each fragment with their attached closing parentheses, then restore the opening parentheses. The contributions exchange the diagonal entries.}
\GSfield{Statement}{In the remote-loop regime, this operation is an involution on admissible fillings. Reversed arrows interchange the diagonal entries of the commuting Fourier blocks, while the signed denominator remains unchanged because the letter multiset and its closing-parenthesis multiplicities are preserved. Pairing the contributions therefore makes every effective coefficient below the first allowed splitting order scalar, even for a non-even complex trigonometric polynomial.}
\GSfield{Reconstruction}{The remote-loop restriction and denominator invariance are supplied from the proof. Complex coefficients are allowed, so no Hermitian or positivity assertion is inferred.}
\GSfield{Fate}{\citet[section 5, Lemmas 2--4 and proof of Theorem 7]{arnold1983mathieu} matches the final page. The Supplement attributes the spectral estimate to Hochstadt 1964; A marks the verified involution proof, not a new first proof of the estimate.}
\GSfield{Present status}{As of November 2026: the stated result is established in the specified setting.}
\GSfield{Assessment}{literal statement: correct with missing hypotheses; publication match: exact; verification: read the result; development: annotated; confidence: high}
\end{GSitem}

\GStalk{586}{1983-02-28}{I. M. Krichever}{P001086-P001090}{Hessians and perturbations of finite gap KdV solutions}
The actual notes start on P001087 and end above Wodzicki\textquotesingle s heading on P001089. They describe the spectral diagonalization of a stationary KdV Hessian and a basis for first order perturbation theory; Gelfand asks for higher variations. The catalogue range includes the end of Arnold and the start of Wodzicki.

\begin{GSitem}{1982-586.1}{E}{KdV Hamiltonians, stationary finite gap solutions, and Lenard recursion}
\GSfield{Kind; source}{construction; \GSpage{P001087}}
\GSfield{The notes}{The notes write the perturbed KdV equation, commuting integrals \(I_k\), \(H=I_{n+2}-\sum c_kI_k\), and \(\delta H/\delta u=0\); the recursion is explicitly qualified as being up to coefficients.}
\GSfield{Statement}{For smooth real periodic potentials, the stationary equations of the KdV hierarchy are Euler--Lagrange equations for linear combinations of its local integrals. With suitable normalization they characterize the nonsingular finite gap potentials of genus at most \(n\). The two compatible Hamiltonian operators produce the Lenard recursion. In the displayed convention \(u_t=\partial_x(\tfrac14u_{xx}-\tfrac34u^2)\), a Hamiltonian is \(\int(-\tfrac18u_x^2-\tfrac14u^3)\,dx\).}
\GSfield{Reconstruction}{\oursComputed{} Smooth periodicity, reality, and the Hamiltonian normalization have been supplied; the notebook conflates a density with its integral. The elementary identity is checked in the accompanying \nolinkurl{independent_checks.py}; this is a verification, not a novelty claim.}
\GSfield{Fate}{This is background already used in \citet[pp.~1312--1313]{krichever1983hessians}; the paper distinguishes it from its new Hessian theorem.}
\GSfield{Present status}{As of November 2026, this is an established result with the hypotheses stated above.}
\GSfield{Assessment}{literal statement: false as written; publication match: exact; verification: read the result; development: annotated; confidence: high}
\end{GSitem}
\begin{GSitem}{1982-586.2}{E}{Band edges and normalized quasimomentum}
\GSfield{Kind; source}{construction; \GSpage{P001088}}
\GSfield{The notes}{For \(L=-\partial_x^2+u_0\), the notes draw \(n\) gaps and double points and give \(i\,dp=Q(E)dE/\sqrt{R(E)}\), with vanishing gap periods.}
\GSfield{Statement}{For a real smooth periodic finite gap Hill operator with \(n\) open gaps, its desingularized spectral curve is \(y^2=\prod_{j=1}^{2n+1}(E-E_j)\). Its quasimomentum differential is the unique second kind differential with the prescribed principal part at infinity and zero \(a\)-periods. Its numerator has one zero in each open gap. The remaining closed gap points are double eigenvalues of the periodic or antiperiodic problems.}
\GSfield{Reconstruction}{The overall factor, including the printed \(i\), depends on whether the Bloch multiplier is written as \(e^{ipT}\) or \(e^{pT}\); no normalization is silently identified across these conventions.}
\GSfield{Fate}{The precise differential and periodicity constraints are recalled in \citet[pp.~1313--1314, eq.~(5)]{krichever1983hessians}; this spectral geometry predates the talk.}
\GSfield{Present status}{As of November 2026, this is an established result with the hypotheses stated above.}
\GSfield{Assessment}{literal statement: correct with missing hypotheses; publication match: exact; verification: read the result; development: none; confidence: high}
\end{GSitem}
\begin{GSitem}{1982-586.3}{A}{Spectral form of the stationary Hessian}
\GSfield{Kind; source}{theorem; \GSpage{P001088}}
\GSfield{The notes}{Equation (8) gives \(\delta^2H=\tfrac12\sum Q(E_i)(\delta E_i)^2+\tfrac14\sum_sQ(e_s)(\delta e_s)^2\), where \(\delta e_s\) is an opened gap width.}
\GSfield{Statement}{At a stationary real finite gap potential, use Krichever\textquotesingle s normalization of \(H\) and his convention for its second variation. Let \(E_i\) be old band edges, \(e_s\) closed-gap double points, and \(\delta e_s\) the leading full widths of the newly opened gaps. Then \[\delta^2H=\frac12\sum_{i=1}^{2n+1}Q(E_i)(\delta E_i)^2+\frac14\sum_sQ(e_s)(\delta e_s)^2.\] This is a quadratic variation, not a finite-perturbation identity. The old-edge shifts satisfy the differentiated fixed-period constraints; the signs of \(Q\) therefore cannot be used to count negative directions as if all edge coordinates were independent.}
\GSfield{Reconstruction}{The statement concerns the quadratic term at the stationary potential, not an exact identity for finite perturbations. The overall normalization of the hierarchy integrals must be inherited from the publication; the notebook\textquotesingle s preceding \(I_3\) normalization is inconsistent.}
\GSfield{Fate}{\citet[Theorems~1--3]{krichever1983hessians} is the exact short publication: presented 19 October 1982, received 23 November 1982, published in Doklady 270(6) in 1983. Submission establishes priority before this talk, but is not evidence of public distribution.}
\GSfield{Present status}{As of November 2026, this is an established result with the hypotheses stated above.}
\GSfield{Assessment}{literal statement: correct with missing hypotheses; publication match: exact; verification: read the result; development: none; confidence: high}
\end{GSitem}
\begin{GSitem}{1982-586.4}{A}{First order opening of a closed spectral gap}
\GSfield{Kind; source}{claim; \GSpage{P001088}}
\GSfield{The notes}{The squared gap width is written as \(|\langle\delta u\,\Psi^2\rangle|^2/\langle|\Psi|^2\rangle\); the constant potential case is identified with Fourier coefficients.}
\GSfield{Statement}{Let a double periodic or antiperiodic eigenvalue have an orthonormal basis \(\phi_1,\phi_2\) of its two-dimensional eigenspace. For a real perturbation \(\varepsilon v\), its two first order eigenvalue shifts are the eigenvalues of \(M_{ij}=\langle\phi_i,v\phi_j\rangle\). Thus the leading squared gap width is \(\varepsilon^2[(M_{11}-M_{22})^2+4|M_{12}|^2]\). For the free operator this becomes \(4\varepsilon^2|\widehat v(2k)|^2\) with normalized plane waves.}
\GSfield{Reconstruction}{\ours{} The displayed matrix statement fixes the Hilbert-space normalization explicitly; its eigenvalue separation is checked in \nolinkurl{independent_checks.py}. Rescaling \(\Psi\) changes the notebook\textquotesingle s right side but cannot change the gap, so the formula is not correct without an unstated fixed normalization.}
\GSfield{Fate}{\citet[Lemma~1 and Theorem~2]{krichever1983hessians} gives the spectral perturbation form. \ours{} The invariant two-dimensional compression is our reconstruction, which makes the missing normalization visible.}
\GSfield{Present status}{As of November 2026, this is an established result with the hypotheses stated above.}
\GSfield{Assessment}{literal statement: false as written; publication match: partial; verification: read the result; development: none; confidence: high}
\end{GSitem}
\begin{GSitem}{1982-586.5}{F}{Gelfand on third and higher variations}
\GSfield{Kind; source}{question; \GSpage{P001088}; I. M. Gelfand}
\GSfield{The notes}{Gelfand asks: \ru{посчитать 3-ю вариацию и т.п.}}
\GSfield{Statement}{At a fixed stationary finite gap potential, calculate third and higher variations of the chosen functional \(H\). A third Fr\'echet derivative in potential coordinates is well-defined once \(H\) and the function space are fixed. An analogue of the preceding spectral diagonalization requires additional choices of spectral coordinates and how a double point is unfolded.}
\GSfield{Reconstruction}{The context suggests a spectral-coordinate formula, but the notes do not say that. We do not turn a request to calculate higher derivatives into an invariant open conjecture.}
\GSfield{Fate}{The located exact paper concerns the Hessian and first order perturbation theory, not a general higher-variation theorem. Searches for the author, Hessians, higher/third variations, and finite gap perturbations did not identify a paper with the unspecified spectral-coordinate conclusion.}
\GSfield{Present status}{As of November 2026, ordinary higher variations are calculable from the differential-polynomial Hamiltonians; the stronger requested form cannot be identified from these notes.}
\GSfield{Assessment}{literal statement: not determinable; publication match: none found; verification: unresolved lead; development: none; confidence: medium}
\end{GSitem}
\begin{GSitem}{1982-586.6}{A}{Gap-opening modes and their adjoint pairing}
\GSfield{Kind; source}{method; \GSpage{P001089 (upper portion)}}
\GSfield{The notes}{The notes expand the first order perturbation in \(w_s^\pm\), obtained by opening one additional gap, and append band-edge and divisor derivatives. They give a logarithmic theta expression and an orthogonality formula with mismatched indices.}
\GSfield{Statement}{Differentiation of an exact family of KdV solutions gives a solution of the linearized equation. At a finite gap background, gap-opening modes together with tangent vectors to the finite gap family give Krichever\textquotesingle s first order perturbation expansion; squares of the corresponding Bloch solutions solve the adjoint equation and give the projections. The parameter variations must satisfy the fixed-period constraints.}
\GSfield{Reconstruction}{The complete expansion includes slow variation of the background spectral parameters. The notebook\textquotesingle s single-theta formula would vanish in genus zero, although the free Fourier modes do not.}
\GSfield{Fate}{\citet[Lemmas~2--3, eqs.~(13)--(19), Theorem~3]{krichever1983hessians} gives the exact basis and forced first order equations. Its mode formula uses a theta quotient with shifts \(\pm2A(e_s)\), differentiated twice, rather than the notebook\textquotesingle s logarithm of a single shifted theta function.}
\GSfield{Present status}{As of November 2026, this is an established result with the hypotheses stated above.}
\GSfield{Assessment}{literal statement: false as written; publication match: partial; verification: read the result; development: none; confidence: high}
\end{GSitem}

\GStalk{587}{1983-02-28}{Mariusz Wodzicki}{P001088-P001094}{Finite gap Hill zeta functions and regularized determinants}
From the lower part of P001089 through P001092, Wodzicki presents zeta values for finite gap Hill operators, universal polynomials, and a Lam\'e example. Interventions concern elementary zeta values, arithmetic, determinant multiplicativity, and a Clausen-series formula. The core became a three-page 1984 paper; the peripheral questions have different evidentiary status.

\begin{GSitem}{1982-587.1}{A}{Shifted spectral zeta function of a Hill operator}
\GSfield{Kind; source}{definition; \GSpage{P001089 (lower portion)}}
\GSfield{The notes}{\(\zeta_{L-w}(s)=\sum(\lambda-w)^{-s}\), initially \(\Re s>1/2\); the question concerns positive integer values for a finite gap potential and spectral \(w\).}
\GSfield{Statement}{For smooth real \(T\)-periodic \(u\), let \(L=-\partial_x^2+u\) have periodic or antiperiodic boundary conditions on \([0,T]\). Remove zero eigenvalues of \(L-w\) and choose a cut for negative eigenvalues when needed. The spectral series converges for \(\Re s>1/2\) and continues meromorphically. At spectral shifts the deleted zero eigenspace and boundary condition are part of the data.}
\GSfield{Reconstruction}{Boundary conditions, a spectral cut, and the prime on the sum are essential additions. The short publication treats the simple and double spectra of both boundary problems.}
\GSfield{Fate}{\citet[Theorems~1--2]{wodzicki1984values}, received 25 July 1983 and published in 1984, explicitly records an earlier presentation at Manin\textquotesingle s seminar in November 1982.}
\GSfield{Present status}{As of November 2026, this is an established result with the hypotheses stated above.}
\GSfield{Assessment}{literal statement: correct with missing hypotheses; publication match: exact; verification: read the result; development: none; confidence: high}
\end{GSitem}
\begin{GSitem}{1982-587.2}{E}{Hurwitz and Epstein comparisons and a binomial expansion}
\GSfield{Kind; source}{remark; \GSpage{P001089}; I. M. Gelfand}
\GSfield{The notes}{Gelfand compares \(\sum(n+a)^{-s}\) and \(\sum(n^2+a)^{-s}\), mentions functional equations, and writes a malformed expansion.}
\GSfield{Statement}{For \(a>0\), the Hurwitz zeta function is \(\zeta(s,a)=\sum_{n\ge0}(n+a)^{-s}\). For \(|a|<n^2\), \((n^2+a)^{-s}=n^{-2s}\sum_{j\ge0}(-1)^j(s)_j a^j/(j!n^{2j})\). Classical Epstein functional equations concern positive definite quadratic forms; an arbitrary homogeneous polynomial is not covered merely by homogeneity.}
\GSfield{Fate}{These are older special-function comparisons, not the new finite gap theorem. The authoritative formulas for Hurwitz zeta and its continuation are collected in \citet[\S25.11]{nist2026dlmf}; the expansion also follows directly from the binomial theorem. The handwritten date 1899 is not used as a priority claim.}
\GSfield{Present status}{As of November 2026, this is an established result with the hypotheses stated above.}
\GSfield{Assessment}{literal statement: false as written; publication match: partial; verification: read the result; development: none; confidence: high}
\end{GSitem}
\begin{GSitem}{1982-587.3}{F}{Arithmetic of shifted spectral values}
\GSfield{Kind; source}{question; \GSpage{P001089}}
\GSfield{The notes}{The discussion asks what happens for rational \(a\) and refers to arithmetic features of the resulting zeta series.}
\GSfield{Statement}{No particular property---rationality, algebraicity, irrationality, special value, or functional equation---is fixed, and the shift and boundary conditions are not specified consistently enough to extract a single arithmetic problem.}
\GSfield{Fate}{Neither the located finite gap paper nor its Lam\'e example states the blanket arithmetic question suggested by the discussion. We retain the question as unidentifiable, without attaching a famous irrationality conjecture.}
\GSfield{Present status}{As of November 2026, no definite proposition can be assigned to this wording.}
\GSfield{Assessment}{literal statement: not determinable; publication match: none found; verification: unresolved lead; development: none; confidence: low}
\end{GSitem}
\begin{GSitem}{1982-587.4}{A}{Universal polynomials and positive integer zeta values}
\GSfield{Kind; source}{construction; \GSpage{P001090-P001091}}
\GSfield{The notes}{The notes define \(F(u)=1+\sum c_i u^i/i!\), polynomial coefficients from \((u\partial_u+p)F^p\), renormalized \(\widehat Q_{p,k}\), and local derivatives of the quasimomentum; equation (4) expresses zeta values using them and Bernoulli numbers.}
\GSfield{Statement}{In the finite gap setting the Taylor coefficients of the normalized quasimomentum at a branch point determine all positive integer shifted zeta values by universal polynomials. A normalization-independent computational formulation is to differentiate the regularized spectral determinant: off the spectrum, \(\partial_w^k\zeta_{L-w}\prime(0)=(k-1)!\zeta_{L-w}(k)\). At an eigenvalue, remove the appropriate zero factor before taking the finite part. The four simple/double and periodic/antiperiodic cases have different coefficients.}
\GSfield{Reconstruction}{The derivative identity gives an exact modern statement without guessing ambiguous factorials in the written polynomial definition. The printed \(p=0\) convention and the power-series convention must be read together.}
\GSfield{Fate}{\citet[\S2, Theorem~2 and its proof]{wodzicki1984values} gives the universal-polynomial formulas and derives them from Theorem~1. It does not support using the notebook\textquotesingle s displayed coefficients independently of its unrecorded normalization.}
\GSfield{Present status}{As of November 2026, this is an established result with the hypotheses stated above.}
\GSfield{Assessment}{literal statement: not determinable; publication match: partial; verification: read the result; development: none; confidence: high}
\end{GSitem}
\begin{GSitem}{1982-587.5}{E}{The free circle and Bernoulli normalization}
\GSfield{Kind; source}{claim; \GSpage{P001091}}
\GSfield{The notes}{The notes say that \(u=0,T=2\pi\) gives \(2\zeta(2s)\), then state a Bernoulli formula without the corresponding period factors.}
\GSfield{Statement}{For \(-\partial_x^2\) on a circle of length \(T\), with the zero mode removed, \(\zeta_L(s)=2(T/2\pi)^{2s}\zeta_{\mathbb R}(2s)\). Consequently \(\zeta_L(m)=T^{2m}|B_{2m}|/(2m)!\) for positive integers \(m\). In particular, at \(T=2\pi\) the factor is \((2\pi)^{2m}\), not \(2^{2m}\).}
\GSfield{Reconstruction}{\oursComputed{}  The elementary identity is checked in the accompanying \nolinkurl{independent_checks.py}; this is a verification, not a novelty claim.}
\GSfield{Fate}{This is the elementary free check of the formulas in \citet[Theorem~2]{wodzicki1984values}; Euler\textquotesingle s value formula is given in \citet[\S25.6]{nist2026dlmf}.}
\GSfield{Present status}{As of November 2026, this is an established result with the hypotheses stated above.}
\GSfield{Assessment}{literal statement: false as written; publication match: exact; verification: read the result; development: annotated; confidence: high}
\end{GSitem}
\begin{GSitem}{1982-587.6}{A}{One-gap Lamé quasimomentum}
\GSfield{Kind; source}{claim; \GSpage{P001091}}
\GSfield{The notes}{For \(u=2\wp(x)\), the notes write a cubic spectral curve, a Weierstrass-zeta Bloch exponent, and a quasimomentum differential.}
\GSfield{Statement}{Choose a real nonsingular periodic translate \(u(x)=2\wp(x+x_0)\). If the elliptic curve for \(\wp\) is \(y^2=4\prod(z-e_i)\), the Hill spectral curve uses \(E=-z\), hence its branch points are \(-e_i\). The Bloch exponent is obtained from \(T\zeta(a)-\eta a\), with \(\wp(a)=-E\), and the differential and zeta values follow after this coordinate change.}
\GSfield{Reconstruction}{The real-axis translate avoids poles. A \(+g_3\) in the spectral cubic and a \(-g_3\) in the Weierstrass cubic are compatible under \(E=-z\); the sign is not automatically a transcription error.}
\GSfield{Fate}{\citet[\S4]{wodzicki1984values} explicitly distinguishes the Weierstrass curve from the reflected spectral curve and gives the simple and double spectral values.}
\GSfield{Present status}{As of November 2026, this is an established result with the hypotheses stated above.}
\GSfield{Assessment}{literal statement: correct with missing hypotheses; publication match: exact; verification: read the result; development: none; confidence: high}
\end{GSitem}
\begin{GSitem}{1982-587.7}{B}{Multiplicativity of the zeta determinant}
\GSfield{Kind; source}{question; \GSpage{P001092}}
\GSfield{The notes}{Gelfand asks whether regularized determinants of differential operators are multiplicative.}
\GSfield{Statement}{For commuting, positive, invertible, classical elliptic pseudodifferential operators \(A,B\) of positive orders \(a,b\), on a closed smooth manifold and on the same Hermitian vector bundle, define \(\det_\zeta A=\exp[-\zeta_A\prime(0)]\). In general \(\det_\zeta(AB)\ne\det_\zeta A\det_\zeta B\). In this commuting positive setting the logarithmic anomaly is \[\frac{\operatorname{res}\bigl(\log(A^bB^{-a})\bigr)^2}{2ab(a+b)}.\] Noncommuting or nonpositive operators require the appropriate more general formula and compatible admissible cuts.}
\GSfield{Fate}{\citet[\S1.1]{kontsevichvishik1994geometry} explains the anomaly and the spectral-cut hypotheses; its June 1994 preprint is an abridgement of their April 1994 work. The commuting formula is explicitly stated with its hypotheses in \citet[eq.~(1)]{castillogarate2015anomaly}.}
\GSfield{Present status}{As of November 2026, general multiplicativity is false. Local residue anomaly formulas and special vanishing classes are established; there is no remaining blanket conjecture here.}
\GSfield{Assessment}{literal statement: correct with missing hypotheses; publication match: exact; verification: read the result; development: none; confidence: high; problem \hyperref[prob:GP-1982-40]{GP-1982-40}}
\end{GSitem}
\begin{GSitem}{1982-587.8}{E}{Euler, Apéry, and regularized sums}
\GSfield{Kind; source}{remark; \GSpage{P001092}}
\GSfield{The notes}{The discussion mentions Euler\textquotesingle s even values, Ap\'ery\textquotesingle s \(\zeta(3)\), and assigns values to \(\sum n^2\), \(\sum n^3\).}
\GSfield{Statement}{Euler\textquotesingle s formula is \(\zeta(2m)=(-1)^{m+1}B_{2m}(2\pi)^{2m}/[2(2m)!]\). Ap\'ery proved that \(\zeta(3)\) is irrational. Analytic continuation gives \(\zeta(-2)=0\) and \(\zeta(-3)=1/120=-B_4/4\); these are zeta-regularized values, not sums of convergent positive-term series.}
\GSfield{Fate}{\citet{apery1979irrationalite} published the irrationality proof in 1979, before the talk. The continuation identities are recorded in \citet[\S25.6]{nist2026dlmf}.}
\GSfield{Present status}{As of November 2026, this is an established result with the hypotheses stated above.}
\GSfield{Assessment}{literal statement: false as written; publication match: exact; verification: read the result; development: none; confidence: high}
\end{GSitem}
\begin{GSitem}{1982-587.9}{F}{The D-module sketch for a complex power}
\GSfield{Kind; source}{remark; \GSpage{P001092}}
\GSfield{The notes}{A short diagram separates integer and noninteger exponents of \(x^{-s}\), with arrows and subquotients.}
\GSfield{Statement}{On the punctured complex line, \(x^{-s}\) generates the rank-one connection \((x\partial_x+s)f=0\) with monodromy \(e^{-2\pi i s}\). At integral exponents, extensions across zero can have punctual subquotients. A specific extension functor or module presentation is needed to identify the diagram.}
\GSfield{Reconstruction}{The rank-one connection is supplied only to locate the subject. It is not an assertion that the diagram equals a specified minimal or maximal extension.}
\GSfield{Fate}{The local connection is an elementary reconstruction; the notebook does not specify the extension or its arrows sufficiently to match an exact module theorem. No fate is inferred from a general reference on D-modules.}
\GSfield{Present status}{As of November 2026, the local connection is standard; the displayed extension diagram remains unidentifiable.}
\GSfield{Assessment}{literal statement: not determinable; publication match: none found; verification: unresolved lead; development: none; confidence: low}
\end{GSitem}
\begin{GSitem}{1982-587.10}{F}{Clausen series and a negative-even zeta derivative}
\GSfield{Kind; source}{claim; \GSpage{P001092}}
\GSfield{The notes}{The notes express \(\zeta_{L-w}\prime(-2)\) using Clausen functions and branch-count corrections; they write \(\operatorname{Cl}_5(0)=\zeta(-5)\).}
\GSfield{Statement}{With \(\operatorname{Cl}_{2k+1}(\theta)=\sum_{n\ge1}\cos(n\theta)/n^{2k+1}\), one has \(\operatorname{Cl}_5(0)=\zeta(5)\). For the free circle, \(\zeta_L\prime(-2)=48\zeta(5)/T^4\), obtained by differentiating \(2(T/2\pi)^{2s}\zeta(2s)\). The general shifted formula needs the spectral branch and the counting function for negative eigenvalues.}
\GSfield{Reconstruction}{\oursComputed{} Only the explicitly checked free specialization and the definition of the Clausen function are asserted here.}
\GSfield{Fate}{The located Wodzicki paper proves positive integer values, not the entire Clausen expression on this page. The free specialization follows from \citet[\S25.4]{nist2026dlmf}; the complete displayed identity has not been verified in an exact source.}
\GSfield{Present status}{As of November 2026, the free specialization is established; the full recorded branch-sensitive formula remains unverified.}
\GSfield{Assessment}{literal statement: false as written; publication match: partial; verification: read the result; development: none; confidence: medium}
\end{GSitem}

\GStalk{588}{1983-03-21}{S. P. Novikov}{P001092-P001095}{Many-valued actions and transgression}
The talk occupies P001093 and the upper part of P001094. It constructs a globally defined variation from locally defined action functionals, illustrated by the Dirac monopole and mapping spaces. The 1982 survey is already a pre-talk source.

\begin{GSitem}{1982-588.1}{E}{A closed one-form as a many-valued action}
\GSfield{Kind; source}{construction; \GSpage{P001093}}
\GSfield{The notes}{The action is \(S=S_0+S_1\), where local \(q\)-form primitives of a closed \((q+1)\)-form define \(S_1\). The variation is asserted to be single-valued.}
\GSfield{Statement}{Let \(N\) be a closed oriented smooth \(q\)-manifold and \(\Omega\in\Omega^{q+1}(M)\) closed. On the smooth mapping space the transgression \(\alpha_\phi(v)=\int_N\phi^*(\iota_v\Omega)\) is a closed one-form. Locally it is the derivative of an action; its periods are integrals of \(\Omega\) over the cycles swept out by loops of maps. Adding a single-valued functional \(S_0\) does not change this period obstruction.}
\GSfield{Reconstruction}{Closedness and orientation of the source, and the transgression formula, are explicit modern hypotheses. The notebook does not develop a complete differential-cohomology model.}
\GSfield{Fate}{\citet[\S5]{novikov1982hamiltonian} gives the many-valued functional construction. It was received 22 April 1982 and published in 1982, before this March 1983 exposition.}
\GSfield{Present status}{As of November 2026, this is an established result with the hypotheses stated above.}
\GSfield{Assessment}{literal statement: correct with missing hypotheses; publication match: exact; verification: read the result; development: none; confidence: high}
\end{GSitem}
\begin{GSitem}{1982-588.2}{E}{Dirac monopole and the phase of the action}
\GSfield{Kind; source}{claim; \GSpage{P001093-P001094}}
\GSfield{The notes}{The example takes \(N=S^1\), \(M=S^2\), and a closed two-form that is not globally exact.}
\GSfield{Statement}{For the magnetic two-form of a monopole on \(S^2\), local magnetic potentials give a multi-valued action on loops. The curvature can have a global primitive only if its integral over \(S^2\) vanishes. If the transgression periods lie in \(2\pi\mathbb Z\), exponentiating the action gives a single-valued phase on the chosen component, after a reference normalization.}
\GSfield{Fate}{\citet[\S\S4--5]{novikov1982hamiltonian} explicitly treats the Dirac monopole and the resulting multivalued action. The integral-period phase formulation is a reconstruction of the obstruction, not a separate theorem claimed from the seminar.}
\GSfield{Present status}{As of November 2026, this is an established result with the hypotheses stated above.}
\GSfield{Assessment}{literal statement: correct with missing hypotheses; publication match: exact; verification: read the result; development: none; confidence: high}
\end{GSitem}
\begin{GSitem}{1982-588.3}{E}{Local actions and constant transition differences}
\GSfield{Kind; source}{claim; \GSpage{P001094 (upper portion)}}
\GSfield{The notes}{The notes choose primitive patches \(U_\alpha\) in \(M\), use maps whose entire images lie in one patch, and say \(S_\alpha-S_\beta\) is constant. A tentative Witten/chiral-field reference follows.}
\GSfield{Statement}{On an overlap of mapping-space domains where two local actions have the same derivative, their difference is constant on each connected component. Domains consisting of maps into one primitive patch need not cover the entire mapping space; a general construction uses local homotopies to reference maps or subdivisions. The chiral-field attribution does not state a further definite claim.}
\GSfield{Reconstruction}{Locally constant replaces globally constant on a disconnected overlap. The tentative attribution is retained as a bibliographic lead, not an exact later solution.}
\GSfield{Fate}{\citet[\S5, the local-action construction]{novikov1982hamiltonian} explains why the variations agree. The scope repair follows from the closed transgression form, without asserting that every map fits into one target chart.}
\GSfield{Present status}{As of November 2026, this is an established result with the hypotheses stated above.}
\GSfield{Assessment}{literal statement: correct with missing hypotheses; publication match: exact; verification: read the result; development: none; confidence: high}
\end{GSitem}

\GStalk{589}{1983-03-21}{M. A. Semenov-Tian-Shansky}{P001093-P001097}{Classical r-matrices and factorization}
From the second heading on P001094 through P001095 the notes introduce an operator \(R\), two Lie--Poisson brackets, invariant Hamiltonians, modified Yang--Baxter identities and factorization. Gelfand asks for the \(\mathfrak{sl}_2\) classification. The final quadratic bracket anticipates the exact 1983 article.

\begin{GSitem}{1982-589.1}{A}{Double Lie algebras and invariant Hamiltonians}
\GSfield{Kind; source}{definition; \GSpage{P001094 (lower portion)}}
\GSfield{The notes}{\([X,Y]_R=[RX,Y]+[X,RY]\) is required to be a Lie bracket; invariant functions commute for both associated brackets, and their Hamiltonian equation has a Lax form. The heading cites Drinfeld, Doklady volume 268, no.5, under the words Lie bialgebra.}
\GSfield{Statement}{If a linear map \(R:\mathfrak g\to\mathfrak g\) makes \([\ ,\ ]_R\) a Lie bracket, the two brackets define Lie--Poisson structures on \(\mathfrak g^*\). Adjoint-invariant smooth functions are in involution for both. For such \(h\), the second Hamiltonian vector field is \(\dot L=\operatorname{ad}^*_{R\,dh(L)}L\), with the coadjoint sign convention fixed consistently. A double Lie algebra alone is not the same datum as a Lie bialgebra.}
\GSfield{Reconstruction}{Finite dimensionality is sufficient for the stated differential-geometric formulation; infinite-dimensional examples require a specified pairing and domains.}
\GSfield{Fate}{\citet{semenov1983classical} was received 15 April 1983, revised 16 May, and published in volume 17(4) later in 1983; this is a publication after the March talk, not evidence that the theorem was new on that evening. The exact result is \citet[\S1, Theorem~1; \S2]{semenov1983classical}. The distinct bibliographic pointer is \citet[\S4, Theorem~2]{drinfeld1983hamiltonian}, which defines a Lie bialgebra by a cobracket satisfying the one-cocycle compatibility, equivalently a compatible invariant double on \(\mathfrak g\oplus\mathfrak g^*\). Its actual issue is 268(2), pages 285--287, received 4 June 1982 and published 1983. This background is not the same assertion as the operator double Lie algebra in the talk.}
\GSfield{Present status}{As of November 2026, this is an established result with the hypotheses stated above.}
\GSfield{Assessment}{literal statement: correct with missing hypotheses; publication match: exact; verification: read the result; development: none; confidence: high}
\end{GSitem}
\begin{GSitem}{1982-589.2}{A}{Modified Yang–Baxter equation, the double, and factorization}
\GSfield{Kind; source}{method; \GSpage{P001094-P001095}}
\GSfield{The notes}{The defect \([RX,RY]-R([RX,Y]+[X,RY])\) is set equal to \(-[X,Y]\). The maps \(R\pm1\), a decomposition in the double, and the factorization of \(\exp(t\,dh)\) are used to integrate the flow.}
\GSfield{Statement}{The modified identity implies that \(R\pm I:\mathfrak g_R\to\mathfrak g\) are Lie homomorphisms. Their joint image in \(\mathfrak g\oplus\mathfrak g\) is complementary to the diagonal. In the corresponding local groups, the factorization problem, with its common-quotient constraint, is locally soluble. For the unhalved bracket used here, the exponential is \(\exp(2t\,dh(L_0))\); the resulting coadjoint action integrates the invariant Hamiltonian flow.}
\GSfield{Reconstruction}{Global factorization is not automatic. The notebook omits the quotient compatibility when the two images intersect and loses the factor of two associated with its bracket convention.}
\GSfield{Fate}{\citet[\S\S3--4, Proposition~9 and Theorem~11, eq.~(18)]{semenov1983classical} gives the exact local result, including the factor \(2\). \citet{semenov1983classical} was received 15 April 1983, revised 16 May, and published in volume 17(4) later in 1983; this is a publication after the March talk, not evidence that the theorem was new on that evening.}
\GSfield{Present status}{As of November 2026, this is an established result with the hypotheses stated above.}
\GSfield{Assessment}{literal statement: false as written; publication match: exact; verification: read the result; development: none; confidence: high}
\end{GSitem}
\begin{GSitem}{1982-589.3}{B}{Gelfand’s classification problem for sl(2)}
\GSfield{Kind; source}{question; \GSpage{P001095}}
\GSfield{The notes}{Gelfand asks to find all \(R\) satisfying the displayed modified Yang--Baxter equation on \(\mathfrak{sl}_2\).}
\GSfield{Statement}{Over \(\mathbb C\), classify all linear maps \(R\) satisfying \([Rx,Ry]-R([Rx,y]+[x,Ry])=-[x,y]\), modulo Lie-algebra automorphisms. No skew-adjoint restriction is imposed by the question. Writing \(B=(R-I)/2\) identifies this exactly with Rota--Baxter operators of weight one: \([Bx,By]=B([Bx,y]+[x,By]+[x,y])\).}
\GSfield{Reconstruction}{\oursComputed{} The ground field and equivalence relation are supplied. Real forms require a separate classification, which the notebook does not select. The elementary identity is checked in the accompanying \nolinkurl{independent_checks.py}; this is a verification, not a novelty claim.}
\GSfield{Fate}{An explicit classification and proof are in \citet[Proposition~2.5 and its proof]{gubarevkozlov2023rota}; its source cases are closed under conjugacy and the involution \(B\mapsto-B-I\), with \(0,-I\) included separately. Received September 2022, published 2023. This answers the complex operator problem without replacing it by the narrower skew tensor problem. With \([h,e]=2e\), \([h,f]=-2f\), \([e,f]=h\), representative weight-one operators are \(B(e)=-e,B(f)=B(h)=0\); \(B(e)=-(e+h),B(f)=B(h)=0\); and \(B(e)=-e,B(f)=0,B(h)=th\), \(t\ne0\), together with \(0,-I\), conjugacy and \(B\mapsto-I-B\).}
\GSfield{Present status}{As of November 2026, the complex finite-dimensional operator classification is established. The primary proof was read; this item does not claim a new classification.}
\GSfield{Assessment}{literal statement: correct with missing hypotheses; publication match: exact; verification: read the result; development: annotated; confidence: high; problem \hyperref[prob:GP-1982-41]{GP-1982-41}}
\end{GSitem}
\begin{GSitem}{1982-589.4}{A}{Quadratic Poisson bracket and the Hamiltonian hierarchy}
\GSfield{Kind; source}{construction; \GSpage{P001095}}
\GSfield{The notes}{A quadratic bracket with terms \(R(XL)\) and \(R(LX)\) is followed by a hierarchy of invariant Hamiltonians.}
\GSfield{Statement}{Let \(\mathfrak g\) be an associative algebra with nondegenerate symmetric invariant trace pairing, and let \(R\) be skew with respect to it and satisfy the modified Yang--Baxter equation. Then the quadratic bracket of Semenov-Tian-Shansky\textquotesingle s equation (22) is Poisson and compatible with the linear \(R\)-bracket. Invariant functions commute for it, and trace powers produce the associated Lenard hierarchy. In formal pseudodifferential examples the trace is the integrated formal residue and the differential/negative-order splitting supplies \(R\).}
\GSfield{Reconstruction}{Skew-adjointness, the trace pairing and associative multiplication are needed for the particular theorem verified here; the bare assumption that the linear \(R\)-bracket satisfies Jacobi is insufficient.}
\GSfield{Fate}{\citet[\S5, Theorem~13 and eqs.~(22)--(23)]{semenov1983classical} gives precisely these hypotheses. This is related to earlier Gelfand--Dikii discussions (1977-426.5 and 1977-426.10), but not an exact recurrence of the earlier question about the first bracket.}
\GSfield{Present status}{As of November 2026, this is an established result with the hypotheses stated above.}
\GSfield{Assessment}{literal statement: correct with missing hypotheses; publication match: exact; verification: read the result; development: none; confidence: high}
\end{GSitem}

\GStalk{590}{1983-03-28}{Mariusz Wodzicki}{P001095-P001099}{APS boundary index theory and local spectral invariants}
The actual record is P001096--P001098. It reviews all three Atiyah--Patodi--Singer papers, includes Fedosov\textquotesingle s symbol-projector construction, and connects heat coefficients, eta regularity and secondary invariants. Gelfand asks about formal cohomology of metrics and symbols. Several questions recur verbatim in the 1981--82 index.

\begin{GSitem}{1982-590.1}{E}{Spectral boundary conditions and the APS index formula}
\GSfield{Kind; source}{theorem; \GSpage{P001096}}
\GSfield{The notes}{For \(D=\partial_u+A\) near the boundary, the projector onto the nonnegative spectrum of self-adjoint \(A\) is imposed; the index is \(\int_X\alpha_0-(h+\eta(0))/2\).}
\GSfield{Statement}{Let a Dirac-type elliptic operator on a compact manifold with boundary have product structure in a collar, with self-adjoint tangential operator \(A\). Impose \(P_{\ge0}(A)(s|_{\partial X})=0\). The realization is Fredholm and \(\operatorname{ind}D_{APS}=\int_X\alpha_0-(\eta_A(0)+\dim\ker A)/2\). The normal bundle isomorphism in the usual Dirac decomposition is suppressed in the notes.}
\GSfield{Fate}{\citet[Theorem~3.10 and \S4]{atiyahpatodisinger1975spectral1} is the exact pre-talk source (received March 1974, published 1975). This repeats 1981-558.11--12 and 1981-559.2.}
\GSfield{Present status}{As of November 2026, this is an established result with the hypotheses stated above.}
\GSfield{Assessment}{literal statement: correct with missing hypotheses; publication match: exact; verification: read the result; development: none; confidence: high}
\end{GSitem}
\begin{GSitem}{1982-590.2}{E}{Replacing self-adjointness by a spectral-cut hypothesis}
\GSfield{Kind; source}{question; \GSpage{P001096}}
\GSfield{The notes}{The notes say self-adjointness may be replaced by absence of spectrum near the imaginary axis.}
\GSfield{Statement}{Complex powers of a general elliptic classical operator require an Agmon ray: the ray avoids both the nonzero principal-symbol eigenvalues and the operator spectrum. For boundary problems the corresponding sectorial projection must also define an elliptic boundary condition. Avoiding the imaginary axis only at the level of the discrete global spectrum does not supply these microlocal hypotheses.}
\GSfield{Reconstruction}{The statement separates the established spectral-cut calculus from a blanket boundary Fredholm theorem, which the one-line remark does not formulate.}
\GSfield{Fate}{The non-self-adjoint closed-operator calculus was already available in \citet[pp.~96--98]{seeley1968powers}, a 1968 Russian translation explicitly identified as a preprint. It requires the principal-symbol ray condition, constructs complex powers and meromorphic diagonal kernels, and gives the shifted-zeta index identity. \citet[\S1]{kontsevichvishik1994geometry} restates admissibility precisely. Wodzicki\textquotesingle s \citet{wodzicki1982spectral} is also pre-talk, but only publisher metadata was checked. This is the earlier problem GP-1981-44; the shorthand about boundary projections does not by itself state a complete boundary theorem.}
\GSfield{Present status}{As of November 2026, this is an established result with the hypotheses stated above.}
\GSfield{Assessment}{literal statement: correct with missing hypotheses; publication match: partial; verification: read the result; development: none; confidence: high; problem \hyperref[prob:GP-1981-44]{GP-1981-44}}
\end{GSitem}
\begin{GSitem}{1982-590.3}{E}{Fedosov’s idempotents for an elliptic symbol}
\GSfield{Kind; source}{construction; \GSpage{P001096-P001097}; B. V. Fedosov}
\GSfield{The notes}{Fedosov constructs a matrix projector from cutoffs \(f_1,f_2\) and \(\sigma,\sigma^*\), then gives a pair \(P_1,P_2\) for general bundles.}
\GSfield{Statement}{An elliptic symbol \(\sigma:\pi^*E_1\to\pi^*E_2\), invertible away from the zero section, defines a compactly supported \(K^0(T^*X)\) class. In a unitary trivialization the block matrix \(\left(\begin{smallmatrix}f_1^2&f_1f_2\sigma\\f_1f_2\sigma^*&f_2^2\end{smallmatrix}\right)\) is an idempotent when \(\sigma\sigma^*=I\) and \(f_1^2+f_2^2=1\). For the later nonorthogonal idempotent pair, the auxiliary symbol and its parametrix must be mutual inverses outside a compact set for their difference to represent a compactly supported class.}
\GSfield{Reconstruction}{\oursComputed{} The algebraic projector identities and the support condition are separate issues. We do not assert that the later pair fails idempotence; the potential error is its cutoff behavior at infinity. The elementary identity is checked in the accompanying \nolinkurl{independent_checks.py}; this is a verification, not a novelty claim.}
\GSfield{Fate}{The symbol-class and Chern-character framework is the classical index construction described in \citet[Chapter~3]{gilkey1984invariance}. The written projector identities were checked directly; the exact displayed cutoff formula was not attributed to an independently located Fedosov paper.}
\GSfield{Present status}{As of November 2026, this is an established result with the hypotheses stated above.}
\GSfield{Assessment}{literal statement: false as written; publication match: partial; verification: read the result; development: annotated; confidence: medium}
\end{GSitem}
\begin{GSitem}{1982-590.4}{E}{Chern character, Todd form, and the symbol index}
\GSfield{Kind; source}{method; \GSpage{P001097}; B. V. Fedosov}
\GSfield{The notes}{The displayed integral is a difference of traces of \(P_i\exp(\omega_i/2\pi i)\) multiplied by the Todd class; \(\omega_i=P_i dP_i dP_i\).}
\GSfield{Statement}{For smooth idempotents representing a compactly supported symbol class, the Grassmann connection has curvature \(P(dP)^2\). Its Chern character difference is paired with \(\pi^*\operatorname{Td}(T_\mathbb C X)\), using the cotangent orientation and the conventional index sign, to obtain the analytic index. Curvature normalization and complexification matter in comparing \(R/(e^R-1)\) with \(R/(1-e^{-R})\).}
\GSfield{Reconstruction}{The support, orientation and curvature conventions have been supplied; the notebook\textquotesingle s two Todd conventions are not treated as an unambiguous sign error.}
\GSfield{Fate}{This is classical topological index theory; \citet[Chapter~3]{gilkey1984invariance} supplies the Chern--Weil formulation. No new universal Todd identity is inferred from the crossed-out curvature notation.}
\GSfield{Present status}{As of November 2026, this is an established result with the hypotheses stated above.}
\GSfield{Assessment}{literal statement: correct with missing hypotheses; publication match: exact; verification: read the result; development: none; confidence: high}
\end{GSitem}
\begin{GSitem}{1982-590.5}{E}{Index from shifted zeta functions and local parametrices}
\GSfield{Kind; source}{method; \GSpage{P001097}}
\GSfield{The notes}{Wodzicki sets \(DD^*=\Delta_2\), \(D^*D=\Delta_1\), writes an index difference of shifted zeta values, and a local symbol integral for \(\alpha_0\).}
\GSfield{Statement}{On a closed manifold, \(\operatorname{ind}D=\zeta_{D^*D+1}(0)-\zeta_{DD^*+1}(0)\), because the positive spectra cancel and each zero mode becomes the eigenvalue one. For a positive elliptic differential operator \(A\), \(\zeta_A(0)\) with zero modes removed equals its integrated constant heat coefficient minus \(\dim\ker A\). Resolvent-symbol contour integrals compute the heat coefficients, with an admissible cut and the chosen Fourier normalization.}
\GSfield{Reconstruction}{The exact shifted identity concerns a closed elliptic problem, or paired boundary realizations with the required spectral cancellation. A generic pseudodifferential heat expansion may include logarithmic terms, so the differential-operator hypothesis is retained.}
\GSfield{Fate}{\citet[\S1.10 and Chapter~3]{gilkey1984invariance} gives the heat/zeta formulation. This is a recurrence of 1981-558.3--4, 1981-559.10--11; the first argument on the scan is an ambiguous crossed zero, not a new parameter. The exact shifted spectral cancellation is already written in \citet[p.~98]{seeley1968powers}, attributed there to Atiyah and Bott.}
\GSfield{Present status}{As of November 2026, this is an established result with the hypotheses stated above.}
\GSfield{Assessment}{literal statement: correct with missing hypotheses; publication match: exact; verification: read the result; development: none; confidence: high}
\end{GSitem}
\begin{GSitem}{1982-590.6}{E}{Local index densities for signature and Dirac operators}
\GSfield{Kind; source}{theorem; \GSpage{P001098}}
\GSfield{The notes}{The notebook says the local density agrees with the topological formula for signature and Dirac operators and mentions equivariant versions.}
\GSfield{Statement}{For geometric signature and Dirac-type complexes on a closed Riemannian manifold, the supertrace heat-kernel constant term is the appropriate characteristic-form density, with the twisting Chern character when present. For an orientation-preserving compact-group symmetry preserving the operator data, the equivariant index is described by fixed-point local contributions.}
\GSfield{Fate}{The pre-talk local index theorem is proved in \citet{atiyahbottpatodi1973heat}; \citet[Chapter~3]{gilkey1984invariance} provides the corresponding invariant-theoretic treatment. This repeats 1981-558.1 and 1981-558.7.}
\GSfield{Present status}{As of November 2026, this is an established result with the hypotheses stated above.}
\GSfield{Assessment}{literal statement: correct with missing hypotheses; publication match: exact; verification: read the result; development: none; confidence: high}
\end{GSitem}
\begin{GSitem}{1982-590.7}{F}{Formal cohomology of metrics and symbols}
\GSfield{Kind; source}{question; \GSpage{P001098}; I. M. Gelfand}
\GSfield{The notes}{Gelfand proposes calculating formal cohomology for the metric and the symbol and asks whether classes other than the index occur; Gelfand--Fuks cohomology and maximal topological invariants are mentioned.}
\GSfield{Statement}{The note does not specify the Lie algebra of formal coordinate changes, coefficient module, equivalence relation on symbols, or degree. Those choices distinguish Chern--Weil invariant polynomials, local scalar invariants, and Gelfand--Fuks cohomology, so no single formal-cohomology calculation is determined.}
\GSfield{Fate}{\citet[\S\S2.4--2.6]{gilkey1984invariance} classifies specified local invariant problems; it is thematic evidence, not a solution of the unspecified combined metric-and-symbol problem.}
\GSfield{Present status}{As of November 2026, no definite open question can be extracted without choosing the formal complex.}
\GSfield{Assessment}{literal statement: not determinable; publication match: thematic only; verification: read the result; development: none; confidence: low}
\end{GSitem}
\begin{GSitem}{1982-590.8}{E}{Signature defect and equivariant eta series}
\GSfield{Kind; source}{theorem; \GSpage{P001098}}
\GSfield{The notes}{The signature formula is \(\operatorname{sign}X=\int_XL-\eta(0)\); an equivariant eta series inserts \(\operatorname{Tr}(g|E_\lambda)\).}
\GSfield{Statement}{For a compact oriented \(4k\)-manifold with product metric near its boundary, the APS signature theorem expresses the signature as the integral of the Hirzebruch \(L\)-form minus the eta invariant of the boundary odd signature operator in the APS convention. If \(g\) is a symmetry commuting with the self-adjoint tangential operator, \(\eta_A(s,g)=\sum_{\lambda\ne0}\operatorname{sign}(\lambda)\operatorname{Tr}(g|E_\lambda)|\lambda|^{-s}\) initially converges in a right half-plane.}
\GSfield{Reconstruction}{The signature normalization differs from the general index expression with \((\eta+h)/2\); its kernel convention cannot simply be interchanged.}
\GSfield{Fate}{\citet[\S4]{atiyahpatodisinger1975spectral1} gives the signature theorem; the boundary operator uses the Hodge star, not merely the formal adjoint of \(d\). This continues 1981-558.8.}
\GSfield{Present status}{As of November 2026, this is an established result with the hypotheses stated above.}
\GSfield{Assessment}{literal statement: correct with missing hypotheses; publication match: exact; verification: read the result; development: none; confidence: high}
\end{GSitem}
\begin{GSitem}{1982-590.9}{E}{Eta regularity without a bounding manifold}
\GSfield{Kind; source}{theorem; \GSpage{P001098}}
\GSfield{The notes}{Part III is summarized as regularity of \(\eta(0)\) in odd dimension without requiring the operator to be a boundary operator.}
\GSfield{Statement}{For a self-adjoint classical elliptic operator of positive order on a closed smooth manifold, the global eta function is regular at zero. APS III proves the odd-dimensional result by its topological argument; the general-dimensional theorem does not imply that the local eta residue density vanishes pointwise.}
\GSfield{Fate}{\citet[Theorem~4.5]{atiyahpatodisinger1976spectral3} proves the odd-dimensional statement (received January 1975, published 1976). \citet[Theorem~0.1]{gilkey1981residue} proves the general statement before this 1983 talk. This settles the exact earlier problem GP-1981-43.}
\GSfield{Present status}{As of November 2026, the global eta-regularity question is settled; local residue assertions require separate hypotheses.}
\GSfield{Assessment}{literal statement: correct with missing hypotheses; publication match: exact; verification: read the result; development: none; confidence: high; problem \hyperref[prob:GP-1981-43]{GP-1981-43}}
\end{GSitem}
\begin{GSitem}{1982-590.10}{E}{Flat bundles, secondary K-theory, and reduced eta}
\GSfield{Kind; source}{construction; \GSpage{P001098 (bottom)}}
\GSfield{The notes}{A unitary representation of \(\pi_1X\), a class with \(\mathbb R/\mathbb Z\) coefficients, and a map to \(\mathbb R/\mathbb Z\) are drawn next to the symbol class.}
\GSfield{Statement}{A rank-\(n\) flat unitary bundle defines a secondary class relative to the trivial rank-\(n\) bundle. Pairing an appropriate \(K\)-theory class with it gives a secondary index in \(\mathbb R/\mathbb Z\), represented analytically by the difference of reduced eta invariants \(\bar\eta(A_E)-n\bar\eta(A)\) modulo integers. The base class lies on \(X\); its product with a compactly supported cotangent symbol is formed only after pullback.}
\GSfield{Reconstruction}{Reduced eta means \((\eta(0)+\dim\ker)/2\). We supply the rank subtraction and distinguish a class on the base from its pullback to the cotangent bundle.}
\GSfield{Fate}{\citet[\S\S4.3--4.4]{gilkey1984invariance} gives the secondary-index and eta interpretation of APS theory. The notebook\textquotesingle s diagram is schematic; the exact coefficient grading is not recoverable from its arrows alone.}
\GSfield{Present status}{As of November 2026, this is an established result with the hypotheses stated above.}
\GSfield{Assessment}{literal statement: correct with missing hypotheses; publication match: partial; verification: read the result; development: none; confidence: medium}
\end{GSitem}

\GStalk{591}{1983-03-28}{George Wilson}{P001097-P001102}{Loop Grassmannians, integrable flows, and the tau function}
The actual record is P001099--P001100. Wilson announces joint work with Segal and relates loop groups, Grassmannians, KdV, modified KdV, determinant lines and affine central extensions. Krichever adds the zero-curvature KP interpretation. This is the prospectus for the detailed April lecture.

\begin{GSitem}{1982-591.1}{A}{Loops, Hardy polarization, and the restricted Grassmannian}
\GSfield{Kind; source}{construction; \GSpage{P001099}}
\GSfield{The notes}{\(H^{(n)}=L^2(S^1,\mathbb C^n)=H_+\oplus H_-\), \(W_\gamma=\gamma H_+\); the positive projection is Fredholm and the negative one compact. The Grassmannian is compared with \(\Omega U(n)\).}
\GSfield{Statement}{For a continuous invertible matrix loop \(\gamma\), multiplication on \(L^2\) is bounded invertible, \(W_\gamma\) is closed, and its Hardy projections have the stated properties. The full restricted Grassmannian is larger than the loop orbit. The subspace with \(zW\subset W\) is the appropriate loop Grassmannian; continuous based unitary loops inject into a dense subspace of it, with a completion admitting measurable loops.}
\GSfield{Fate}{\citet[\S2, Proposition~2.7 and the following discussion]{segalwilson1985loop} states the continuous-loop inclusion and the distinction between the completed Grassmannian and \(\Omega U(n)\). \citet{segalwilson1985loop} develops the joint work reported here; publication was in 1985. The present seminar notes themselves document circulation in March--April 1983.}
\GSfield{Present status}{As of November 2026, this is an established result with the hypotheses stated above.}
\GSfield{Assessment}{literal statement: correct with missing hypotheses; publication match: exact; verification: read the result; development: none; confidence: high}
\end{GSitem}
\begin{GSitem}{1982-591.2}{E}{Which measurable loops give the Fredholm and compact projections?}
\GSfield{Kind; source}{question; \GSpage{P001099}}
\GSfield{The notes}{The notes call it an unsolved problem to characterize all loops for which conditions (i) and (ii) hold.}
\GSfield{Statement}{For continuous invertible loops the conditions already hold. A precise wider version takes an essentially bounded unitary matrix loop \(\gamma\). Then \(W_\gamma\to H_+\) is Fredholm and \(W_\gamma\to H_-\) compact exactly when the matrix entries lie in \(QC=(H^\infty+C)\cap\overline{H^\infty+C}\). Indeed Hartman\textquotesingle s theorem identifies compact Hankel blocks; the Toeplitz block is essentially unitary exactly when both off-diagonal blocks are compact.}
\GSfield{Reconstruction}{We explicitly restrict the reconstruction to essentially bounded unitary loops. \ours{} Write the unitary multiplier as \(U=\left(\begin{smallmatrix}a&b\\c&d\end{smallmatrix}\right)\). If \(c\) is compact, then \(a^*a=1\) in the Calkin algebra; Fredholmness makes \(a\) invertible there, whence \(aa^*=1\) and \(bb^*=0\), so \(b\) is compact. Conversely compact \(b,c\) make \(a\) unitary modulo compacts and hence Fredholm. Hartman\textquotesingle s criterion then gives \(QC\) entrywise.}
\GSfield{Fate}{\citet[pp.~221--223]{page1975compact} proves the compact-Hankel characterization (Hartman\textquotesingle s theorem dates to 1958). \ours{} The unitary-loop equivalence is the elementary block-operator consequence checked here. The notes do not specify this ambient measurable class, so no unresolved general \(GL_n\) loop problem is asserted.}
\GSfield{Present status}{As of November 2026, the unitary measurable reconstruction has the stated classical criterion. The literal phrase all loops lacks a function class.}
\GSfield{Assessment}{literal statement: not determinable; publication match: partial; verification: read the result; development: none; confidence: medium}
\end{GSitem}
\begin{GSitem}{1982-591.3}{A}{Grassmannian flows and Krichever’s algebro-geometric locus}
\GSfield{Kind; source}{construction; \GSpage{P001099-P001100}}
\GSfield{The notes}{The notes map a Grassmannian to solutions of the \(n\)th KdV hierarchy and place the Krichever construction inside it, followed by a dressing-method comparison.}
\GSfield{Statement}{For scalar polarized \(L^2(S^1)\), the reduction \(z^nW\subset W\) yields an order-\(n\) scalar differential Lax operator and its Gelfand--Dikii flows. General Grassmannian points give KP. The algebro-geometric construction supplies a distinguished locus, and polynomial/rational finite-dimensional approximations are dense in the relevant Grassmannian topology; this is not a classification of all solutions of arbitrary growth.}
\GSfield{Reconstruction}{The reduction and analytic class delimit the word solutions; no universal surjectivity onto every smooth KdV initial value is inferred.}
\GSfield{Fate}{\citet[\S\S2,5--7]{segalwilson1985loop} verifies the reduction and its algebraic loci. \citet{segalwilson1985loop} develops the joint work reported here; publication was in 1985. The present seminar notes themselves document circulation in March--April 1983.}
\GSfield{Present status}{As of November 2026, this is an established result with the hypotheses stated above.}
\GSfield{Assessment}{literal statement: correct with missing hypotheses; publication match: exact; verification: read the result; development: none; confidence: high}
\end{GSitem}
\begin{GSitem}{1982-591.4}{E}{Modified hierarchies and the Miura flag}
\GSfield{Kind; source}{construction; \GSpage{P001100}}
\GSfield{The notes}{A commutative diagram relates a flag manifold to modified KdV and the Grassmannian to KdV; the Miura fibre is written \(G/B\).}
\GSfield{Statement}{For a local monic scalar differential operator of order \(n\), factorizations into first order operators correspond to suitable complete flags in its local solution space. The Miura transformation carries modified hierarchy solutions to the ordinary hierarchy. Analytic factorizations near an ordinary point form the open set of flags with nonvanishing required Wronskians; allowing meromorphic factors recovers the full flag variety.}
\GSfield{Reconstruction}{Local solution fields and meromorphic versus analytic factorizations are distinguished; global periodic fibres need not equal a single flag variety.}
\GSfield{Fate}{\citet[\S3.7 and \S6]{drinfeldsokolov1985lie} develops the Miura transformation and generalized hierarchies. Its introduction identifies \citet{drinfeldsokolov1981equations} as the earlier Doklady announcement; the latter’s publication metadata was checked. Thus this is pre-talk mathematics despite the full survey’s 1984 Russian and 1985 English publication dates.}
\GSfield{Present status}{As of November 2026, this is an established result with the hypotheses stated above.}
\GSfield{Assessment}{literal statement: correct with missing hypotheses; publication match: exact; verification: read the result; development: none; confidence: high}
\end{GSitem}
\begin{GSitem}{1982-591.5}{A}{Determinant lines, affine central extensions, and tau}
\GSfield{Kind; source}{construction; \GSpage{P001100}}
\GSfield{The notes}{The tau function is introduced as a determinant, with an affine Lie-algebra extension and a theta-function special case.}
\GSfield{Statement}{The determinant line over the index-zero restricted Grassmannian produces a tau function after an admissible frame and a lift of the positive loop action are chosen. The restricted linear group acts through a central extension; on smooth matrix loops its infinitesimal cocycle is a scalar multiple of \((2\pi i)^{-1}\oint\operatorname{tr}(X\,dY)\). The scalar tau function is not canonically independent of all trivializations.}
\GSfield{Fate}{\citet[\S\S3--4 and 10]{segalwilson1985loop} gives the determinant construction and central extension. The affine algebra itself was older; the specific analytic tau model is the joint work announced here.}
\GSfield{Present status}{As of November 2026, this is an established result with the hypotheses stated above.}
\GSfield{Assessment}{literal statement: correct with missing hypotheses; publication match: exact; verification: read the result; development: none; confidence: high}
\end{GSitem}
\begin{GSitem}{1982-591.6}{E}{Krichever’s KP zero-curvature and tau reconstruction}
\GSfield{Kind; source}{method; \GSpage{P001100 (lower portion)}; I. M. Krichever}
\GSfield{The notes}{Krichever replaces commuting ordinary operators by compatible \(\partial_y-L\), \(\partial_t-L_1\), and writes \(u=-2\partial_x^2\log\tau\) for \(L=-\partial_x^2+u\).}
\GSfield{Statement}{For a KP auxiliary system, the compatibility condition is \([\partial_y-L,\partial_t-L_1]=0\). In the convention \(L=-\partial_x^2+u\), the potential reconstructed from tau is \(u=-2\partial_x^2\log\tau\), up to any explicitly chosen constant background. The sign reverses for the standard monic \(\partial_x^2+u\) convention.}
\GSfield{Reconstruction}{The relation is a compatibility equation in two evolution variables, not ordinary commutativity of time-dependent \(L,L_1\).}
\GSfield{Fate}{\citet[\S5, especially eq.~(5.19)]{segalwilson1985loop} gives the monic convention and its tau reconstruction. The notebook\textquotesingle s sign is consistent with its negative leading term.}
\GSfield{Present status}{As of November 2026, this is an established result with the hypotheses stated above.}
\GSfield{Assessment}{literal statement: correct with missing hypotheses; publication match: exact; verification: read the result; development: none; confidence: high}
\end{GSitem}

\GStalk{592}{1983-04-04}{George Wilson}{P001100-P001109}{The determinant tau function and the KP Grassmannian}
P001101--P001107 gives the detailed Segal--Wilson construction: frames and determinants, Schur expansions, algebraic curves, Baker functions and Miwa shifts, formal Lax equations, determinant-line extensions, and local Miura factorizations. Gelfand\textquotesingle s broad request about theta functions of all abelian varieties is kept separate from the precise Jacobian results.

\begin{GSitem}{1982-592.1}{A}{Index-zero Grassmannian and admissible frames}
\GSfield{Kind; source}{construction; \GSpage{P001101-P001102}}
\GSfield{The notes}{The restricted Grassmannian is defined by Fredholm/compact projections. Frames \(w:H_+\to W\) have \(w_+=I+\text{trace class}\), \(w_-\) compact; the text says all points have such frames.}
\GSfield{Statement}{The Fredholm index labels the connected components of the restricted Grassmannian. The stated frame condition parametrizes the index-zero component. For \(g\in\Gamma_+\), holomorphic and nonvanishing on a neighbourhood of the closed disc, choose its block decomposition \(g^{-1}=\left(\begin{smallmatrix}a&b\\0&d\end{smallmatrix}\right)\). The trace-class block needed for the determinant is available in this analytic loop class.}
\GSfield{Reconstruction}{The scan first notes the full Grassmannian, with all Fredholm indices, and also displays the index-zero qualification. The later frame construction is interpreted on that index-zero component; the unrestricted assertion about every point would be false.}
\GSfield{Fate}{\citet[\S\S2--3]{segalwilson1985loop} gives the component and frame construction. \citet{segalwilson1985loop} develops the joint work reported here; publication was in 1985. The present seminar notes themselves document circulation in March--April 1983.}
\GSfield{Present status}{As of November 2026, this is an established result with the hypotheses stated above.}
\GSfield{Assessment}{literal statement: correct with missing hypotheses; publication match: exact; verification: read the result; development: none; confidence: high}
\end{GSitem}
\begin{GSitem}{1982-592.2}{A}{The Fredholm determinant tau function and its zero set}
\GSfield{Kind; source}{definition; \GSpage{P001102}}
\GSfield{The notes}{\(\tau_W(g)=\det(w_++a^{-1}bw_-)\); transversality is stated as both invertibility of the positive projection and \(W\cap H_-=0\).}
\GSfield{Statement}{On the index-zero component, the operator inside the determinant is identity plus trace class, so its Fredholm determinant is defined. Its vanishing is equivalent to failure of \(g^{-1}W\to H_+\) to be invertible. There, injectivity is equivalent to invertibility because the index is zero. On the big cell the graph frame gives \(\tau_W(1)=1\); elsewhere this normalization is unavailable.}
\GSfield{Reconstruction}{An arbitrary Fredholm operator has no ordinary determinant. The frame normalization and trace-class subtraction are essential.}
\GSfield{Fate}{\citet[\S3, determinant-line discussion]{segalwilson1985loop} gives this construction.}
\GSfield{Present status}{As of November 2026, this is an established result with the hypotheses stated above.}
\GSfield{Assessment}{literal statement: correct with missing hypotheses; publication match: exact; verification: read the result; development: none; confidence: high}
\end{GSitem}
\begin{GSitem}{1982-592.3}{A}{Nonvanishing, Schur functions, and Plücker coefficients}
\GSfield{Kind; source}{theorem; \GSpage{P001102-P001103}}
\GSfield{The notes}{The tau function is never identically zero, even on \(g=e^{xz}\). Coordinate subspaces \(H_S\) give Schur functions, and general tau has a Pl\"ucker-coordinate expansion.}
\GSfield{Statement}{For an index-zero Grassmannian point, \(\tau_W(t)=\sum_S w^S F_S(t)\), where \(F_S\) is the Schur function associated with the charge-zero indexing set \(S\), with the sign convention determined by \(g^{-1}\). For \(t=(x,0,\ldots)\), the leading power has nonzero coefficient and degree equal to the codimension of the stratum containing \(W\). Schur functions also encode polynomial \(GL_n\) characters after a finite-variable specialization.}
\GSfield{Reconstruction}{The coefficient conventions use the charge-zero partition, not an arbitrary increasing set without its eventual indexing condition.}
\GSfield{Fate}{\citet[Propositions~8.2--8.3 and 8.6]{segalwilson1985loop} proves the expansion and the stronger one-variable nonvanishing. The representation-theoretic identification answers the brief \(GL(n)\) aside.}
\GSfield{Present status}{As of November 2026, this is an established result with the hypotheses stated above.}
\GSfield{Assessment}{literal statement: correct with missing hypotheses; publication match: exact; verification: read the result; development: none; confidence: high}
\end{GSitem}
\begin{GSitem}{1982-592.4}{A}{The Krichever map and its Fredholm index}
\GSfield{Kind; source}{construction; \GSpage{P001103}}
\GSfield{The notes}{A pointed curve, a disc at infinity, a line bundle (rank-one torsion-free sheaf for a singular curve), and a trivialization produce \(W\); the index is \(\chi(\mathcal L)-1\).}
\GSfield{Statement}{Let \(X\) be an integral complete algebraic curve, smooth at a marked point \(p\), with a local coordinate and a trivialized rank-one torsion-free sheaf \(\mathcal L\) near \(p\). The closure of boundary values of sections on the complement defines a Grassmannian point. With \(H_+=\operatorname{span}\{1,z,z^2,\ldots\}\) and \(z\) a pole coordinate at \(p\), the projection index is \(\chi(\mathcal L(-p))=\chi(\mathcal L)-1\).}
\GSfield{Reconstruction}{The coordinate, local freeness at the marked point, and boundary-value completion are supplied. Genus and degree need not be guessed from the index.}
\GSfield{Fate}{\citet[\S6, Proposition~6.1]{segalwilson1985loop} is the exact analytic realization of the older Krichever construction.}
\GSfield{Present status}{As of November 2026, this is an established result with the hypotheses stated above.}
\GSfield{Assessment}{literal statement: correct with missing hypotheses; publication match: exact; verification: read the result; development: none; confidence: high}
\end{GSitem}
\begin{GSitem}{1982-592.5}{A}{Theta factorization of an algebro-geometric tau function}
\GSfield{Kind; source}{theorem; \GSpage{P001103-P001104}}
\GSfield{The notes}{The notes relate \(\Gamma_+\) to \(H^1(X,\mathcal O)\) and write \(\tau=e^{\alpha_{W,0}+Q}\theta(\bar f-\beta_W)\), with \(Q\) depending only on the curve and disc.}
\GSfield{Statement}{For a smooth pointed curve with fixed local coordinate and line-bundle data, the Krichever tau function is a theta function on its Jacobian composed with the linear map from flow times, multiplied by an exponential with linear and quadratic terms. The quadratic form is determined by the pointed curve and normalization; the line bundle changes the translate and the linear factor.}
\GSfield{Fate}{\citet[\S9, especially Theorem~9.11]{segalwilson1985loop} gives the exact theta comparison. This continues the curve-based theta constructions indexed in 1977-431.3.}
\GSfield{Present status}{As of November 2026, this is an established result with the hypotheses stated above.}
\GSfield{Assessment}{literal statement: correct with missing hypotheses; publication match: exact; verification: read the result; development: none; confidence: high}
\end{GSitem}
\begin{GSitem}{1982-592.6}{F}{Gelfand’s request for theta functions of all abelian varieties}
\GSfield{Kind; source}{question; \GSpage{P001104}; I. M. Gelfand}
\GSfield{The notes}{\ru{Построить теорию интегрирования, связанную с \(\theta\)-функцией \(\forall\) абелевых многообразий}; the parenthesis says only Jacobians of curves occur here.}
\GSfield{Statement}{The notes ask for an integration theory attached to arbitrary abelian varieties but do not specify equations, a notion of integrability, or a universal construction. The exact Jacobian theta solution of KP is a different, narrower statement.}
\GSfield{Fate}{The exact curve construction is in \citet[\S9]{segalwilson1985loop}. \citet{shiota1986characterization} is a later Jacobian-characterization reference, not a solution of the all-abelian-varieties programme; only its publisher metadata and bibliography were checked here. That bibliography records 1983 Kyoto seminar notes separately from the 1986 publication.}
\GSfield{Present status}{As of November 2026, this broad programme does not specify one checkable open problem; no full solution is attributed to the Schottky/KP literature.}
\GSfield{Assessment}{literal statement: not determinable; publication match: thematic only; verification: publisher metadata checked; development: none; confidence: medium}
\end{GSitem}
\begin{GSitem}{1982-592.7}{A}{Baker function, tau ratio, and Miwa shifts}
\GSfield{Kind; source}{construction; \GSpage{P001104-P001105}}
\GSfield{The notes}{\(\psi_W(g,z)=g(z)(1+\sum a_i(g)z^{-i})\in W\); multiplying by \(q_\zeta=1-z/\zeta\) gives its tau quotient and a vertex operator.}
\GSfield{Statement}{Whenever \(g^{-1}W\to H_+\) is invertible, there is a unique Baker function with the stated asymptotic form and membership in \(W\). Its reduced form is \(\widetilde\psi_W(g,\zeta)=\tau_W(gq_\zeta)/\tau_W(g)\), \(|\zeta|>1\). If \(g=\exp(\sum t_i z^i)\), this is the shift \(t_i\mapsto t_i-1/(i\zeta^i)\), and hence a differential-exponential vertex-operator formula. The coefficients extend meromorphically through the tau zero set.}
\GSfield{Reconstruction}{Uniqueness is first asserted on the big cell; the meromorphic extension does not give a finite regular Baker value at every tau zero.}
\GSfield{Fate}{\citet[Proposition~5.1 and formulas~(5.14)--(5.16)]{segalwilson1985loop} provides the exact result.}
\GSfield{Present status}{As of November 2026, this is an established result with the hypotheses stated above.}
\GSfield{Assessment}{literal statement: correct with missing hypotheses; publication match: exact; verification: read the result; development: none; confidence: high}
\end{GSitem}
\begin{GSitem}{1982-592.8}{A}{Formal KP Lax operator and the Gelfand–Dikii reduction}
\GSfield{Kind; source}{construction; \GSpage{P001105}}
\GSfield{The notes}{\(Q=D+\sum q_iD^{-i}\), \(\partial_{t_r}Q=[(Q^r)_+,Q]\); when \(Q^n=L\) is differential this becomes the \(n\)th KdV hierarchy. The condition is \(z^nW\subset W\).}
\GSfield{Statement}{For the Baker dressing operator \(K=1+\sum a_iD^{-i}\), define \(Q=KDK^{-1}\). Then \(\partial_{t_r}Q=[(Q^r)_+,Q]\). If \((Q^n)_-=0\), the differential operator \(L=Q^n\) obeys \(\partial_{t_r}L=[(L^{r/n})_+,L]\). In the Grassmannian construction this reduction is precisely \(z^nW\subset W\). For \(L=D^2+u\), \(u=2\partial_x^2\log\tau\).}
\GSfield{Reconstruction}{The positive part is taken after forming \(Q^r\). The factor 2 in the tau potential is required in the monic convention.}
\GSfield{Fate}{\citet[Propositions~5.5 and 5.9, Corollary~5.12, eq.~(5.19)]{segalwilson1985loop} is the exact source. The hierarchy formalism itself is credited there to Sato and earlier Lax theory.}
\GSfield{Present status}{As of November 2026, this is an established result with the hypotheses stated above.}
\GSfield{Assessment}{literal statement: false as written; publication match: exact; verification: read the result; development: none; confidence: high}
\end{GSitem}
\begin{GSitem}{1982-592.9}{A}{Pole-order bounds and regular singular scalar operators}
\GSfield{Kind; source}{theorem; \GSpage{P001105}}
\GSfield{The notes}{The notes say the resulting differential operator has only regular singular points in \(x\).}
\GSfield{Statement}{For a Segal--Wilson solution and fixed remaining times, the Baker coefficient \(a_i(x)\) has poles of order at most \(i\). In the reduced order-\(n\) operator \(L=D^n+\sum_{i=0}^{n-2}u_iD^i\), each coefficient has pole order at most \(n-i\) at finite poles. Thus these finite singular points are regular singular in the scalar Fuchs criterion.}
\GSfield{Reconstruction}{No assertion about regular singularity at infinity is added.}
\GSfield{Fate}{\citet[Proposition~5.17 and Corollary~5.18 and their proofs in \S8]{segalwilson1985loop} proves precisely the finite-plane bounds.}
\GSfield{Present status}{As of November 2026, this is an established result with the hypotheses stated above.}
\GSfield{Assessment}{literal statement: correct with missing hypotheses; publication match: exact; verification: read the result; development: none; confidence: high}
\end{GSitem}
\begin{GSitem}{1982-592.10}{A}{The determinant line and its nontrivial central extension}
\GSfield{Kind; source}{construction; \GSpage{P001106}}
\GSfield{The notes}{\(\mathcal P\) is the frame bundle; \(\mathcal T\) is the trace-class change-of-frame group; pairs \((g,q)\) form \(\mathcal E\). Quotienting \(\mathcal E\) by \(\ker\det\) gives a central \(\mathbb C^*\)-extension acting on \(\mathrm{Det}\).}
\GSfield{Statement}{The determinant line is associated to the admissible-frame bundle by the Fredholm determinant character of \(I+\mathcal L^1\). On the identity component of the restricted linear group, pairs \((g,q)\) with \(a q^{-1}-I\) trace class form an extension; quotienting by determinant-one frame transformations gives a central \(\mathbb C^*\)-extension. It acts on the determinant line, and generally does not admit a continuous global splitting.}
\GSfield{Reconstruction}{The restricted operator ideals and the index-zero component are part of the definition; a lift is needed for the action on the line.}
\GSfield{Fate}{\citet[\S3, especially the discussion preceding formula~(3.5)]{segalwilson1985loop} gives the nontrivial extension and its local cocycle.}
\GSfield{Present status}{As of November 2026, this is an established result with the hypotheses stated above.}
\GSfield{Assessment}{literal statement: correct with missing hypotheses; publication match: exact; verification: read the result; development: none; confidence: high}
\end{GSitem}
\begin{GSitem}{1982-592.11}{E}{Local Miura factorization and the flag of solutions}
\GSfield{Kind; source}{construction; \GSpage{P001107}}
\GSfield{The notes}{\(D^2+u=(D-v)(D+v)\), \(u=v_x-v^2\); the general factorization is related to analytic initial values and meromorphic flags \(SL(n)/B\).}
\GSfield{Statement}{Let \(L\) be monic of order \(n\) with analytic coefficients near an ordinary point and no \(D^{n-1}\) term. Complete flags in the \(n\)-dimensional local solution space determine factorizations into monic first order operators with meromorphic coefficients. Flags with all necessary initial Wronskians nonzero give analytic factors; they form an open chart with \(n(n-1)/2\) initial parameters. The second-order factorization obeys the stated Riccati relation.}
\GSfield{Reconstruction}{\oursComputed{} This is local around an ordinary point; global meromorphic or periodic coefficients impose further monodromy constraints. The phrase every solution is not extended to global factorizations. The elementary identity is checked in the accompanying \nolinkurl{independent_checks.py}; this is a verification, not a novelty claim.}
\GSfield{Fate}{\citet[\S3.7, Miura transformation]{drinfeldsokolov1985lie} supplies the hierarchy connection. The local flag statement follows from successive Wronskian factorization and explains the distinction drawn in the notebook.}
\GSfield{Present status}{As of November 2026, this is an established result with the hypotheses stated above.}
\GSfield{Assessment}{literal statement: correct with missing hypotheses; publication match: exact; verification: read the result; development: annotated; confidence: high}
\end{GSitem}
\begin{GSitem}{1982-592.12}{F}{Differential fields and the proposed Galois analogy}
\GSfield{Kind; source}{question; \GSpage{P001107 (bottom)}}
\GSfield{The notes}{The notes compare fields of coefficients, roots, and solutions, with \(B\) and \(SL(n)\) drawn beneath them, while observing that Miura fibres are not a group.}
\GSfield{Statement}{The Picard--Vessiot differential Galois group is a subgroup of \(SL_n\) when the \(D^{n-1}\) coefficient vanishes; a complete solution flag has a stabilizer in a Borel subgroup. Which differential field is generated by a chosen factorization depends on the operator and flag. The diagram does not specify a universal equality of Galois groups or a question to prove.}
\GSfield{Fate}{The Lie-theoretic Miura formalism of \citet[\S3]{drinfeldsokolov1985lie} is relevant background, but it is not an exact solution of an unstated differential-Galois problem.}
\GSfield{Present status}{As of November 2026, the analogy is meaningful but the notes do not fix one differential-field assertion.}
\GSfield{Assessment}{literal statement: not determinable; publication match: thematic only; verification: read the result; development: none; confidence: medium}
\end{GSitem}

\GStalk{594}{1983-04-11}{George Wilson}{P001107-P001114}{KP and its scalar reductions: concluding note}
Only the lower strip of P001108 belongs to this brief conclusion. It repeats the formal KP hierarchy and its scalar differential reductions, and identifies \(2q_1\) as a KP solution. Ginzburg begins on the next page.

\begin{GSitem}{1982-594.1}{E}{The first KP coefficient and the two-time equation}
\GSfield{Kind; source}{claim; \GSpage{P001108}}
\GSfield{The notes}{The note repeats \(Q=D+\sum q_iD^{-i}\), its Lax equations and the reduction, then says \(2q_1(x,t_2,t_3)\) solves KP.}
\GSfield{Statement}{With \(x=t_1\), \(y=t_2\), \(t=t_3\) and \(u=2q_1\), the formal KP Lax equations imply \((4u_t-6uu_x-u_{xxx})_x=3u_{yy}\). The condition \(Q^n\) differential gives the Gelfand--Dikii reduction. Other signs result from explicit changes of time or leading-operator conventions.}
\GSfield{Reconstruction}{The explicit PDE is supplied to make the coefficient assertion checkable; it is not written out on the source page.}
\GSfield{Fate}{This is the same formalism as 1982-592.8, proved in \citet[\S5]{segalwilson1985loop}; the short conclusion does not introduce another new theorem.}
\GSfield{Present status}{As of November 2026, this is an established result with the hypotheses stated above.}
\GSfield{Assessment}{literal statement: correct with missing hypotheses; publication match: exact; verification: read the result; development: none; confidence: high}
\end{GSitem}

\GStalk{593}{1983-04-11}{V. L. Ginzburg}{P001107-P001109}{Equivariant localization and the Duistermaat–Heckman theorem}
The actual report occupies P001109--P001112, after Wilson\textquotesingle s conclusion on P001108. It compares the Duistermaat--Heckman theorem, the Atiyah--Bott equivariant interpretation, Witten\textquotesingle s parity complex, and the explicit primitive attributed to Vergne. The last page records a successful correction of an initially impossible closed angular form.

\begin{GSitem}{1982-593.1}{E}{Exact stationary phase and the fixed-component denominator}
\GSfield{Kind; source}{theorem; \GSpage{P001109}}
\GSfield{The notes}{For a compact symplectic manifold and Hamiltonian with compact or periodic flow, the notes write \(\int_M e^{iH}=\sum_j\operatorname{vol}(M_j)e^{iH(M_j)}/\lambda_j\), with \(\lambda_j\) a normal Pfaffian. The Fourier transform is said to have an exact stationary-phase answer.}
\GSfield{Statement}{Let \((M^{2n},\omega)\) be compact and let a Hamiltonian circle action have generator \(\xi\), with \(\iota_\xi\omega=dH\). For the Cartan differential \(d+u\iota_\xi\), the equivariantly closed form is \(\omega-uH\). Localization gives \[\int_M e^{\omega-uH}=\sum_{F\subset M^{S^1}}\int_F\frac{e^{\omega_F-uH_F}}{e_{S^1}(\nu F)(u)}.\] Its left side means \(\int_M e^{-uH}\omega^n/n!\), with compatible Euler-class conventions. For isolated fixed points the denominators are products of weights and powers of \(u\). For positive-dimensional fixed components the full equivariant Euler class, including curvature terms, is needed.}
\GSfield{Reconstruction}{The Liouville measure, deformation parameter, sign convention, and full fixed-component denominator are supplied. The notebook\textquotesingle s volume divided by a constant Pfaffian is not the unrestricted fixed-submanifold formula.}
\GSfield{Fate}{The underlying theorem was published by \citet{duistermaatheckman1982variation} in 1982. \citet[\S\S3,6--7, especially formula~(3.8)]{atiyahbott1984moment}, received 20 December 1982 and published in 1984, gives this cohomological form and distinguishes isolated points from fixed submanifolds. Its bibliography records the 1982 Berline--Vergne and Duistermaat--Heckman addendum preprints.}
\GSfield{Present status}{As of November 2026, this is an established result with the hypotheses stated above.}
\GSfield{Assessment}{literal statement: false as written; publication match: partial; verification: read the result; development: none; confidence: high}
\end{GSitem}
\begin{GSitem}{1982-593.2}{E}{Polynomial density on regular-value chambers}
\GSfield{Kind; source}{theorem; \GSpage{P001109}}
\GSfield{The notes}{The measure pushed forward by \(H\) is first called piecewise polynomial, then a polynomial on \(\mathbb R\). The reduced symplectic form is called linear in the level.}
\GSfield{Statement}{For a compact Hamiltonian torus manifold, the pushforward of Liouville measure has a polynomial density on each chamber of regular values, relative to Lebesgue measure on the moment-map image. Under the natural local identification of reduced spaces, the cohomology class of the reduced symplectic form varies affinely with the moment value. This cohomological variation gives the polynomial volume; walls can change the polynomial.}
\GSfield{Reconstruction}{Affineness concerns the reduced cohomology class after identification, not a global choice of local coordinate expressions for every reduced two-form.}
\GSfield{Fate}{\citet{duistermaatheckman1982variation} is the 1982 source; \citet[\S7]{atiyahbott1984moment} gives the localization and Fourier-transform interpretation. This is not one polynomial valid across all critical values.}
\GSfield{Present status}{As of November 2026, this is an established result with the hypotheses stated above.}
\GSfield{Assessment}{literal statement: false as written; publication match: exact; verification: read the result; development: none; confidence: high}
\end{GSitem}
\begin{GSitem}{1982-593.3}{F}{Gelfand’s question about stationary eigenvalues}
\GSfield{Kind; source}{question; \GSpage{P001109}; I. M. Gelfand}
\GSfield{The notes}{Gelfand remarks that the eigenvalues of the stationary action have not yet been used.}
\GSfield{Statement}{The normal isotropy weights are precisely the linearized eigenvalue data entering the equivariant Euler denominator. No further proposed identity or missing conclusion is specified by the question.}
\GSfield{Fate}{\citet[\S\S3 and 7]{atiyahbott1984moment} explains where those weights occur. The question is an expository challenge to the abbreviated derivation, not a separately identifiable unresolved theorem.}
\GSfield{Present status}{As of November 2026, the role of the normal weights is explicit in localization; no distinct open proposition is stated here.}
\GSfield{Assessment}{literal statement: not determinable; publication match: partial; verification: read the result; development: none; confidence: medium}
\end{GSitem}
\begin{GSitem}{1982-593.4}{E}{Borel equivariant cohomology and localization}
\GSfield{Kind; source}{construction; \GSpage{P001110}}
\GSfield{The notes}{\(H_G^*(M)=H^*(EG\times_GM)\), as a module over \(H^*(BG)\); for the circle \(H^*(BS^1;\mathbb Q)=\mathbb Q[u]\), \(\deg u=2\). Restriction to the fixed locus becomes an isomorphism over the fraction field. The notes recall Atiyah--Segal in K-theory and Quillen in cohomology.}
\GSfield{Statement}{For a compact smooth circle manifold, restriction induces \[H_{S^1}^*(M;\mathbb Q)\otimes_{\mathbb Q[u]}\mathbb Q(u)\simeq H^*(M^{S^1};\mathbb Q)\otimes\mathbb Q(u).\] More generally a torus requires localization of its polynomial coefficient ring. The analogous statement for an arbitrary nonabelian compact group cannot be obtained just by replacing the circle fixed locus by the full group fixed locus.}
\GSfield{Reconstruction}{Rational coefficients and the circle or torus restriction are explicit. A polynomial description of \(H^*(BG)\) needs the corresponding coefficient and compact-Lie-group hypotheses.}
\GSfield{Fate}{\citet[\S\S2--3]{atiyahbott1984moment} reviews the earlier Borel/Atiyah--Segal localization and proves the integration formula. The lecture\textquotesingle s later explicit restriction to \(S^1\) supplies the relevant hypothesis.}
\GSfield{Present status}{As of November 2026, this is an established result with the hypotheses stated above.}
\GSfield{Assessment}{literal statement: correct with missing hypotheses; publication match: exact; verification: read the result; development: none; confidence: high}
\end{GSitem}
\begin{GSitem}{1982-593.5}{E}{The invariant Cartan complex and the square of its differential}
\GSfield{Kind; source}{construction; \GSpage{P001111}}
\GSfield{The notes}{\(d_1=d+u\iota_\xi\) is first written on all forms; the notes then insist on invariant forms from the start. They print \(d_1^2=u^2L_\xi\). Goncharov compares the construction with Chern--Weil theory.}
\GSfield{Statement}{The Cartan model is \((\Omega(M)^{S^1}\otimes\mathbb C[u],d+u\iota_\xi)\), \(\deg u=2\). Cartan\textquotesingle s homotopy identity gives \((d+u\iota_\xi)^2=uL_\xi\), which vanishes on invariant forms. Its cohomology identifies with Borel equivariant cohomology over \(\mathbb C\); the Chern--Weil comparison uses a connection on the universal bundle.}
\GSfield{Reconstruction}{The later correction to invariant forms is retained as part of the source argument, rather than counted as an unresolved omission.}
\GSfield{Fate}{\citet[\S4]{atiyahbott1984moment} gives the differential-form model and its relation with the topological construction. The invariant-form qualification is already repaired in the notebook; only the extra power of \(u\) remains an actual formula error.}
\GSfield{Present status}{As of November 2026, this is an established result with the hypotheses stated above.}
\GSfield{Assessment}{literal statement: false as written; publication match: exact; verification: read the result; development: none; confidence: high}
\end{GSitem}
\begin{GSitem}{1982-593.6}{E}{Parity localization and Vergne’s explicit contracting form}
\GSfield{Kind; source}{method; \GSpage{P001111-P001112}}
\GSfield{The notes}{After putting \(u=1\), the notes identify the parity cohomology with that of the fixed set. They first seek an invariant angular one-form \(\alpha\) with \(\iota_\xi\alpha=1\) and \(d\alpha=0\), then replace it by \(\beta=\alpha(1+d\alpha)^{-1}\) and write \(\omega=d_1(\beta\wedge\omega)+\beta\wedge d_1\omega\).}
\GSfield{Statement}{On the complement of the zero set of a Killing field generating a compact action, choose an invariant metric and \(\alpha=\xi^\flat/\|\xi\|^2\). Then \(\iota_\xi\alpha=1\), \(L_\xi\alpha=0\). The finite exterior-algebra series \(\beta=\alpha\sum_{j\ge0}(-d\alpha)^j\) satisfies \((d+\iota_\xi)\beta=1\). Wedge multiplication by \(\beta\) contracts the invariant parity complex away from the fixed set. For compact \(M\), restriction therefore gives its \(\mathbb Z/2\)-graded cohomology as \(H^*(M^{S^1};\mathbb C)\) with parity grading; boundary integration near the fixed set yields the localization residues.}
\GSfield{Reconstruction}{\ours{} A direct check is \(D\alpha=1+d\alpha\), \(D(d\alpha)=L_\xi\alpha-d(\iota_\xi\alpha)=0\); hence \(D\beta=1\), and the graded Leibniz rule gives the written identity. This verifies the recorded proof, not a new localization theorem. The earlier aspiration \(d\alpha=0\) is explicitly abandoned.}
\GSfield{Fate}{\citet[\S\S4--5 and 8]{atiyahbott1984moment} treats this model and the relation to Witten and Berline--Vergne. Its reference [5] records the Rennes Berline--Vergne preprint of 1982, before this talk. The notebook itself contains the valid primitive and the contracting identity.}
\GSfield{Present status}{As of November 2026, this is an established result with the hypotheses stated above.}
\GSfield{Assessment}{literal statement: correct with missing hypotheses; publication match: exact; verification: read the result; development: annotated; confidence: high}
\end{GSitem}
\begin{GSitem}{1982-593.7}{F}{A group-free formulation through the double fibration}
\GSfield{Kind; source}{question; \GSpage{P001111}; I. M. Gelfand}
\GSfield{The notes}{Gelfand asks to formulate the same theorem without \(G\), in the language of the two projections of \(EG\times_GM\).}
\GSfield{Statement}{The diagram alone does not specify what structures should replace the group action, the polynomial coefficient ring, the fixed set, or the localization multiplicative set. Without those data it is not a single theorem about arbitrary double fibrations.}
\GSfield{Fate}{The Borel and Cartan comparison in \citet[\S\S2--4]{atiyahbott1984moment} explains the diagram, but does not solve the additional unspecified group-free programme. No later groupoid or foliated theorem is treated as an exact answer.}
\GSfield{Present status}{As of November 2026, the wording remains insufficient to identify one falsifiable open problem.}
\GSfield{Assessment}{literal statement: not determinable; publication match: thematic only; verification: read the result; development: none; confidence: low}
\end{GSitem}

\GStalk{596}{1983-04-18}{V. L. Ginzburg}{P001112-P001115}{Moment-map convexity and density chambers}
The actual report is the upper and middle portion of P001113, ending at Beilinson\textquotesingle s heading. A preceding plan mentions reading Levinson on the zeta function but supplies no theorem. P001115 starts 5 September 1983 and belongs to the next academic year, despite the catalogue overlap.

\begin{GSitem}{1982-596.1}{E}{Hamiltonian torus actions and the moment map}
\GSfield{Kind; source}{definition; \GSpage{P001113}}
\GSfield{The notes}{The notes start with a compact symplectic torus manifold, then insert the assignment \(X\mapsto H_X\) compatible with Poisson brackets and define \(\mu(m)(X)=H_X(m)\).}
\GSfield{Statement}{For a Hamiltonian action of a torus \(T\) on \((M,\omega)\), choose a linear assignment \(X\mapsto H_X\) with \(dH_X=\iota_{X_M}\omega\), compatible with the Lie and Poisson brackets. The map \(\mu:M\to\mathfrak t^*\) is defined by evaluation. A symplectic torus action need not be Hamiltonian; the inserted Hamiltonian assignment is an essential assumption.}
\GSfield{Reconstruction}{The inserted Poisson representation is treated as part of the notes, not silently supplied by symplecticity alone.}
\GSfield{Fate}{\citet[\S1]{atiyah1982convexity} states the Hamiltonian setup before the convexity theorem. This continues the moment-map discussion in 1981-556.12 and 1980-508.15; it is not the exact orbital-integral question of 1982-568.}
\GSfield{Present status}{As of November 2026, this is an established result with the hypotheses stated above.}
\GSfield{Assessment}{literal statement: correct with missing hypotheses; publication match: exact; verification: read the result; development: none; confidence: high}
\end{GSitem}
\begin{GSitem}{1982-596.2}{E}{Convexity of the moment image}
\GSfield{Kind; source}{theorem; \GSpage{P001113}}
\GSfield{The notes}{\(\mu(M)\) is asserted to be a convex polytope; the margin credits Atiyah--Bott and Guillemin--Sternberg independently.}
\GSfield{Statement}{If a torus acts in a Hamiltonian manner on a compact connected symplectic manifold, its moment-map image is the convex hull of the finitely many images of connected fixed-point components, and its fibres are connected. Compactness and connectedness are part of the theorem.}
\GSfield{Reconstruction}{The connected-fibre conclusion is included to identify Atiyah\textquotesingle s precise theorem, although it is not separately stated in this short note.}
\GSfield{Fate}{\citet[Theorem~1]{atiyah1982convexity}, received 27 March 1981 and published in 1982, proves this theorem. The independent convexity work is by Guillemin--Sternberg, as also described in \citet[\S6]{atiyahbott1984moment}; the source margin conflates Atiyah\textquotesingle s convexity paper with the joint equivariant-localization work.}
\GSfield{Present status}{As of November 2026, this is an established result with the hypotheses stated above.}
\GSfield{Assessment}{literal statement: correct with missing hypotheses; publication match: exact; verification: read the result; development: none; confidence: high}
\end{GSitem}
\begin{GSitem}{1982-596.3}{E}{Critical walls and the Duistermaat–Heckman density}
\GSfield{Kind; source}{theorem; \GSpage{P001113}}
\GSfield{The notes}{The density is piecewise polynomial on \(\mu(M)\); transitions are called faces where \(d\mu\) degenerates and are said to lie in planes.}
\GSfield{Statement}{The moment-map critical values lie in finitely many affine subspaces coming from isotropy strata. The regular-value chambers carry polynomial Duistermaat--Heckman densities. Critical walls may occur in the interior of the moment polytope; they are not restricted to its boundary faces.}
\GSfield{Reconstruction}{The word faces is read as a wall decomposition, not as an assertion that every critical value belongs to the boundary of the image.}
\GSfield{Fate}{This is the torus version of 1982-593.2. \citet{duistermaatheckman1982variation} and \citet[\S\S6--7]{atiyahbott1984moment} supply the precise chamber picture.}
\GSfield{Present status}{As of November 2026, this is an established result with the hypotheses stated above.}
\GSfield{Assessment}{literal statement: correct with missing hypotheses; publication match: exact; verification: read the result; development: none; confidence: high}
\end{GSitem}

\GStalk{595}{1983-04-18}{A. A. Beilinson}{P001112-P001114}{Arithmetic zeta functions and the global continuation problem}
The talk starts in the lower portion of P001113 and continues on P001114, with a short Wodzicki interruption before Beilinson resumes. The heading promises special values up to a rational factor and exact formulations, but the recorded mathematics concerns local factors, global continuation, a Langlands-type correspondence, and the elliptic functional equation. No rank or leading-coefficient formula for Birch--Swinnerton-Dyer is actually recorded.

\begin{GSitem}{1982-595.1}{F}{Special-value programme and the limits of the surviving statement}
\GSfield{Kind; source}{remark; \GSpage{P001113 (lower portion)}}
\GSfield{The notes}{The heading distinguishes formulations modulo \(\mathbb Q^*\) from exact formulations, followed by the name Birch--Swinnerton-Dyer.}
\GSfield{Statement}{The notes contain no regulator map, motivic cohomology group, rank equality, leading coefficient, or specified integer special value. They announce a programme rather than state a precise special-value conjecture.}
\GSfield{Fate}{The modern distinctions between Hasse--Weil continuation and stronger arithmetic conjectures are reviewed in \citet[\S1]{gee2026modularity}. The earlier regulator occurrence 1979-476.11 is related context, not an exact recurrence inferable from this heading.}
\GSfield{Present status}{As of November 2026, no particular regulator or Birch--Swinnerton-Dyer formula can be identified from this entry alone.}
\GSfield{Assessment}{literal statement: not determinable; publication match: thematic only; verification: read the result; development: none; confidence: high}
\end{GSitem}
\begin{GSitem}{1982-595.2}{E}{Arithmetic models and the local point-count zeta function}
\GSfield{Kind; source}{definition; \GSpage{P001113 (bottom)-P001114 (upper portion)}}
\GSfield{The notes}{The notes choose equations supposedly giving smooth reduction at every prime, write \(\zeta(X,s)=\prod_p\zeta(X_{\mathbb F_p},s)\), and define each local factor by \(\sum_n\#X(\mathbb F_{p^n})p^{-ns}\). The point example then gives the usual Riemann Euler product.}
\GSfield{Statement}{For a smooth projective geometrically connected curve \(X/\mathbb Q\), there is a smooth proper model over \(\mathbb Z[1/S]\) for a finite set of bad primes \(S\); changing equations cannot generally remove bad reduction. At a good prime, \[Z(X_p,t)=\exp\!\left(\sum_{n\ge1}\#X_p(\mathbb F_{p^n})\frac{t^n}{n}\right),\qquad t=p^{-s}.\] Bad local factors require the chosen arithmetic zeta or cohomological L-function convention. A point has \(Z(t)=(1-t)^{-1}\); the unexponentiated series is the logarithmic derivative data, not the local zeta function.}
\GSfield{Reconstruction}{The finite bad set, exponential, and divisor \(n\) are indispensable. The word rational in the margin is not interpreted as a genus-zero hypothesis because the continuation discussion explicitly addresses general curves.}
\GSfield{Fate}{\citet[\S1]{deligne1974weil} gives the finite-field zeta and cohomological factors. \citet[\S1.1]{boxercalegarigeepilloni2021surfaces} explicitly separates the global Euler product and finitely many excluded factors.}
\GSfield{Present status}{As of November 2026, this is an established result with the hypotheses stated above.}
\GSfield{Assessment}{literal statement: false as written; publication match: exact; verification: read the result; development: none; confidence: high}
\end{GSitem}
\begin{GSitem}{1982-595.3}{E}{Weil rationality and the correct location of finite-field zeros}
\GSfield{Kind; source}{theorem; \GSpage{P001114 (upper portion)}}
\GSfield{The notes}{The local function is rational in \(p^{-s}\), with numerator and denominator factors; a damaged formula places zeros at \(\operatorname{Re}z=p^{?/2}\). The notes remark that the curve case is proved.}
\GSfield{Statement}{For a smooth projective geometrically connected curve of genus \(g\) over \(\mathbb F_p\), \[Z(X,t)=\frac{\prod_{j=1}^{2g}(1-\alpha_jt)}{(1-t)(1-pt)},\qquad |\alpha_j|=p^{1/2}.\] Thus the numerator zeros in the \(t\)-plane have modulus \(p^{-1/2}\); in the variable \(t=p^{-s}\), they lie on \(\Re s=1/2\). These statements concern a local finite-field factor, not the global Riemann hypothesis.}
\GSfield{Reconstruction}{The damaged exponent is not guessed. The three distinct variables---Frobenius eigenvalue, local zero, and complex s---are separated.}
\GSfield{Fate}{\citet[Theorem~1.6]{deligne1974weil} proves purity in every cohomological degree for smooth projective varieties, not merely curves. The paper records July 1973 lectures and was published in 1974; both precede this 1983 exposition.}
\GSfield{Present status}{As of November 2026, this is an established result with the hypotheses stated above.}
\GSfield{Assessment}{literal statement: false as written; publication match: exact; verification: read the result; development: none; confidence: high}
\end{GSitem}
\begin{GSitem}{1982-595.4}{F}{Which coefficients define reasonable Euler factors?}
\GSfield{Kind; source}{question; \GSpage{P001114}; I. M. Gelfand}
\GSfield{The notes}{Gelfand asks which \(a_n(p)\) in \(\prod_p\sum_n a_n(p)p^{-ns}\) should be called reasonable, observing that \(a_n(p)=1\) for the Riemann zeta function.}
\GSfield{Statement}{The question does not impose local rationality, integrality, purity, conductor bounds, compatibility across primes, or an automorphic or motivic origin. These choices give inequivalent classification problems.}
\GSfield{Fate}{\citet[\S1]{gee2026modularity} specifies the geometric and Galois data in a modern Hasse--Weil question. It does not supply a necessary-and-sufficient characterization of the undefined word reasonable.}
\GSfield{Present status}{As of November 2026, the wording does not identify one precise open classification problem.}
\GSfield{Assessment}{literal statement: not determinable; publication match: thematic only; verification: read the result; development: none; confidence: medium}
\end{GSitem}
\begin{GSitem}{1982-595.5}{C}{Global continuation and functional equations for curves}
\GSfield{Kind; source}{problem; \GSpage{P001114}}
\GSfield{The notes}{The notes say analytic continuation is not proved even for curves over \(\mathbb Q\), contrast the finite-field case, and locate the obstacle in compatibility among different primes.}
\GSfield{Statement}{For every smooth projective geometrically connected curve over \(\mathbb Q\), prove meromorphic continuation of its completed arithmetic zeta function and the expected functional equation, with the local bad factors and archimedean factors specified. Equivalently, one may formulate continuation for the Hasse--Weil L-function of its \(H^1\), after separating the Tate factors. This is a global assertion; finite-field rationality alone does not imply it.}
\GSfield{Reconstruction}{The precise completion and the distinction between the arithmetic zeta function and its \(H^1\) L-factor are supplied. The historical sentence is read as a general conjecture, not as denying all previously known special curves.}
\GSfield{Fate}{For genus one over \(\mathbb Q\), \citet[Theorem~A]{breuilconraddiamondtaylor2001modularity} proves modularity and hence the analytic assertion (2001). For every genus-two curve over a totally real field, \citet[Theorem~1.1.2]{boxercalegarigeepilloni2021surfaces} proves meromorphic continuation and the functional equation, by potential automorphy in Theorem~1.1.3. Their work circulated at the end of 2018, was received 25 December 2018, and appeared in 2021. \citet[\S1 and Remark~2.7]{gee2026modularity} gives the November 2026 research perspective.}
\GSfield{Present status}{As of November 2026, the unrestricted curve statement remains open. Genus one over \(\mathbb Q\) and genus two over totally real fields have the stated continuation results. Potential automorphy gives meromorphic continuation over the original field; it does not establish holomorphy of every genus-two \(H^1\) L-function.}
\GSfield{Assessment}{literal statement: correct with missing hypotheses; publication match: exact; verification: read the result; development: none; confidence: high; problem \hyperref[prob:GP-1982-42]{GP-1982-42}}
\GSfield{Worth developing}{Worth pursuing in explicit higher-genus families because general global continuation remains unresolved and local point-count data alone do not control the Euler product.}
\end{GSitem}
\begin{GSitem}{1982-595.6}{C}{Matching geometric H1 factors with automorphic L-functions}
\GSfield{Kind; source}{conjecture; \GSpage{P001114}}
\GSfield{The notes}{A strong conjecture says the two classes of geometric zeta functions and Langlands zeta functions coincide.}
\GSfield{Statement}{A precise direction relevant to a curve is that the compatible Galois representations on its \(H^1\) should correspond to suitable algebraic automorphic data, allowing isobaric decompositions when the motive is reducible, with matching good Euler factors and compatible completions. Literal equality with every automorphic L-function is too broad: algebraicity, weight, and motivic constraints matter.}
\GSfield{Reconstruction}{The formulation is one modern geometric-to-automorphic direction extracted from the broad notebook sentence, not a claim that the notes formulated a complete Langlands reciprocity conjecture.}
\GSfield{Fate}{\citet[Theorem~1.1.3]{boxercalegarigeepilloni2021surfaces} gives potential automorphy for all genus-two curves over totally real fields. Stronger actual modularity is proved under explicit extra hypotheses by \citet[Theorem~B]{boxercalegarigeepilloni2025modularity}: for a genus-two curve \(y^2=f(x)\) over \(\mathbb Q\), \(\deg f=5\), surjective mod 3 representation on its Jacobian and good ordinary reduction at 2 and 3 imply modularity. The preprint was posted 28 February 2025. \citet[Remark~2.7]{gee2026modularity} places this advance in the continuing programme.}
\GSfield{Present status}{As of November 2026, general actual automorphy for arbitrary curve \(H^1\) remains open. The 2025 theorem has the stated genus-two, rational-Weierstrass-point, residual-image and ordinary-reduction hypotheses; it is not modularity of every abelian surface.}
\GSfield{Assessment}{literal statement: correct with missing hypotheses; publication match: exact; verification: read the result; development: none; confidence: high; problem \hyperref[prob:GP-1982-43]{GP-1982-43}}
\GSfield{Worth developing}{Worth developing for carefully specified motives or explicit curves because actual automorphy is stronger than potential automorphy and can yield holomorphy as well as continuation.}
\end{GSitem}
\begin{GSitem}{1982-595.7}{B}{The elliptic completion and what the BSD heading does not say}
\GSfield{Kind; source}{claim; \GSpage{P001114 (bottom, after Wodzicki)}}
\GSfield{The notes}{Under the resumed Birch--Swinnerton-Dyer heading, only continuation and symmetry of \((2\pi)^{-s}\Gamma(s)\zeta(X,s)\) under \(s\mapsto2-s\) are written; the strip \(0\le\Re s\le2\) is drawn.}
\GSfield{Statement}{For an elliptic curve \(E/\mathbb Q\) with conductor \(N\), the completed Hasse--Weil L-function is \(\Lambda(E,s)=N^{s/2}(2\pi)^{-s}\Gamma(s)L(E,s)\), and \(\Lambda(E,s)=w_E\Lambda(E,2-s)\), where \(w_E\in\{\pm1\}\). The whole arithmetic zeta function of the curve contains additional Tate factors and has a different completion. Neither the equality of analytic and algebraic ranks nor the BSD leading-term formula appears in these notes.}
\GSfield{Reconstruction}{The conductor, root number, and identification of the displayed series with \(L(H^1(E),s)\) are necessary emendations; genus one is inferred from the single gamma factor and the named conjecture, with that inference made explicit.}
\GSfield{Fate}{\citet[Theorem~A]{breuilconraddiamondtaylor2001modularity} establishes modularity of every elliptic curve over \(\mathbb Q\), hence the continuation and functional equation written here after these corrections. The broader BSD conjecture is not thereby solved and is not manufactured as a separate occurrence.}
\GSfield{Present status}{As of November 2026, this is an established result with the hypotheses stated above.}
\GSfield{Assessment}{literal statement: false as written; publication match: exact; verification: read the result; development: none; confidence: high}
\end{GSitem}

\GStalk{598}{1983-04-18}{Mariusz Wodzicki}{P001113-P001115}{Shintani-type zeta series and a proposed operator realization}
This is the short interruption in the middle of P001114, between Beilinson\textquotesingle s strong correspondence conjecture and his resumed functional-equation paragraph. The catalogue\textquotesingle s extension to P001115 must not absorb the September 1983 Deligne talk.

\begin{GSitem}{1982-598.1}{E}{Dedekind zeta and Shintani cone series}
\GSfield{Kind; source}{method; \GSpage{P001114 (middle interruption)}}
\GSfield{The notes}{Wodzicki asks which zeta functions can be computed and cites Shintani, writing \(\zeta_K(s)=\sum_{\mathfrak a}N\mathfrak a^{-s}=\sum_n a_n n^{-s}\).}
\GSfield{Statement}{For a number field \(K\), the sum is over nonzero integral ideals, \(N\mathfrak a=|\mathcal O_K/\mathfrak a|\), and converges for \(\Re s>1\). For totally real \(K\), Shintani\textquotesingle s cone decomposition expresses partial Dedekind zeta functions as finite sums of shifted positive-cone Dirichlet series associated with products of real linear forms; these yield values at nonpositive integers. The totally real assumption belongs to this particular construction.}
\GSfield{Reconstruction}{The nonzero-integral-ideal condition and totally real scope are supplied; the notes do not specify which special values or which field signature.}
\GSfield{Fate}{\citet[opening pages and \S1]{cassounogues1982series} explicitly states Shintani\textquotesingle s 1976 result and develops the associated polynomial Dirichlet series. It was a 1981 arithmetic-days report, published in 1982, before this talk. The original 1976 paper\textquotesingle s bibliographic identity is confirmed by that primary research account; its full text was not obtained here.}
\GSfield{Present status}{As of November 2026, this is an established result with the hypotheses stated above.}
\GSfield{Assessment}{literal statement: correct with missing hypotheses; publication match: partial; verification: read the result; development: none; confidence: high}
\end{GSitem}
\begin{GSitem}{1982-598.2}{F}{Constant-coefficient operators on a torus: unresolved identification}
\GSfield{Kind; source}{claim; \GSpage{P001114 (middle interruption)}}
\GSfield{The notes}{The Shintani construction is described as reducing the Dedekind zeta function to zeta functions of pseudodifferential operators on a torus with constant coefficients.}
\GSfield{Statement}{A cone-restricted polynomial lattice sum can be viewed formally as a spectral sum after choosing a Fourier multiplier and a projection onto a cone of modes. That does not by itself make it the zeta function of an ordinary positive classical elliptic operator on the whole torus. The symbol, ellipticity, multiplicities, cone projection and zero-mode convention must be specified.}
\GSfield{Reconstruction}{The Fourier-multiplier interpretation is supplied only to explain what additional data an exact statement would require, and is not claimed as a recovered historical theorem.}
\GSfield{Fate}{The primary account \citet{cassounogues1982series} verifies the cone Dirichlet series but does not identify the unrestricted torus pseudodifferential assertion written here. Searches combining Shintani, Dedekind zeta, torus, constant coefficients and pseudodifferential operators did not locate an exact operator theorem. This is an unresolved identification, not evidence that the arithmetic method was undeveloped.}
\GSfield{Present status}{As of November 2026, the Shintani arithmetic construction is established; the exact operator realization intended by the note has not been identified.}
\GSfield{Assessment}{literal statement: not determinable; publication match: thematic only; verification: unresolved lead; development: none; confidence: low}
\end{GSitem}

\section{Developments}\label{dev:1982}
\subsection{The four-unitary free energy: a controlled local calculation}
\label{dev:1982-migdal}

\paragraph{The question actually recorded.}
Item \hyperref[item:1982-565.6]{1982-565.6}, problem
\hyperref[prob:GP-1982-02]{GP-1982-02}, is introduced on
\GSpage{P001040}, equation~(19), by \ru{Надо вычислить}.
Both the Russian transcription and an independent inspection of the scan show
a negative sign in the exponential, an ordered sum over \(\mu,\nu\), and an
inserted \(N^{-2}\log\). The notes ask for the limit; they do not supply a
proof of its existence, a domain for the coupling, or a normalization of Haar
measure. We make the last choice explicit by taking probability Haar measure.
The strongest formulation directly supported by the page is: evaluate the
displayed large-matrix limit on its domain of existence. No continuum-limit
claim follows simply from writing this integral.

For fixed \(d\geq2\), independent Haar matrices \(U_1,\ldots,U_d\in U(N)\),
and real \(\beta\), put
\begin{align}
 S_{d,N}(U)&=\sum_{\mu,\nu=1}^{d}
       \operatorname{Tr}(U_\mu U_\nu U_\mu^*U_\nu^*),\label{eq:1982-S}\\
 F_{d,N}(\beta)&=\frac1{N^2}\log
       \mathbb E\exp\{-\beta N S_{d,N}\}.\label{eq:1982-F}
\end{align}
The notebook case is \(d=4\). The general \(d\) notation is our extension
for checking coefficients. Pairing \((\mu,\nu)\) with \((\nu,\mu)\) makes
\(S_{d,N}\) real. The diagonal terms give \(dN\), exactly. Thus removing
them adds \(d\beta\) to \(F_{d,N}\); reversing the action sign replaces
\(\beta\) by \(-\beta\). These changes cannot be suppressed when comparing
formulas. In particular, the usual positive plaquette coupling corresponds
to negative \(\beta\) in \eqref{eq:1982-F}.

\paragraph{What later work establishes.}
The reduced lattice model itself predates this seminar:
\citet{eguchikawai1982reduction} was published on 19~April~1982.
That identification does not evaluate the four-matrix integral.
For our exact mathematical normalization,
\citet[Theorem~0.1]{collinsguionnetmaurel2009asymptotics}
is decisive: for a fixed self-adjoint noncommutative polynomial potential in
finitely many independent Haar unitaries, the normalized logarithmic integral
has a large-\(N\) limit for sufficiently small real coupling; that limit is
analytic near zero and its Taylor coefficients are the limits of the formal
finite-\(N\) coefficients. There are no external matrices here, so the
boundedness and joint-moment assumptions on those matrices are vacuous.
Take the potential to be the negative of the ordered commutator sum in
\eqref{eq:1982-S}. The preprint first circulated on 8~August~2006; publication
followed in 2009.

The distinction between an unmodified reduced model and a valid reduction
of an extended gauge theory remains material. Center symmetry and the precise
quenching prescription are discussed in
\citet[Secs.~II--III]{neuberger2020quenched}; neither quenching nor twisting
is present in the notebook's integral. The mathematical local convergence
theorem therefore applies directly to the displayed integral, while a
physical equivalence requires separate hypotheses. The search behind this
chapter located no verified all-coupling evaluation of this exact four-unitary
integral. This is a bounded literature finding as of November~2026, not a
proof that no such evaluation exists. The sources and failed leads are
recorded in the research audit.
The checked sources do not establish that the exact integral was abandoned.
They do explain why its evaluation and its use as a reduction of QCD are
different research tasks: the latter depends on symmetry and on the chosen
reduction prescription, while the former still involves four interacting
unitary matrices.

\paragraph{The stronger target and the scope of this development.}
A plausible stronger formulation is existence and determination of
\(F_4(\beta)=\lim_{N\to\infty}F_{4,N}(\beta)\) for every real \(\beta\),
including any points where it fails to be differentiable. This is our
formulation of the global target; it is not a theorem asserted by the notes
or established in this chapter. We supply exact finite-\(N\) coefficients,
elementary global bounds, and the rigorously justified local expansion.
These are reproducible bounded deductions, not claims of a new solution
of the model or of priority for the coefficient formulas.

\begin{proposition}[Normalization, convexity, and compactness \ours]
For every \(d,N\) and every real \(\beta\), \(F_{d,N}\) is well-defined,
real analytic, convex, and \(d^2\)-Lipschitz. It satisfies
\[
 -d\beta-d(d-1)|\beta|\ \leq\ F_{d,N}(\beta)
 \ \leq\ -d\beta+d(d-1)|\beta|.
\]
Every sequence of matrix sizes tending to infinity has a subsequence along
which \(F_{d,N}\) converges uniformly on every compact real interval to a
convex \(d^2\)-Lipschitz function vanishing at zero.
\end{proposition}
\begin{proof}
Each commutator is unitary, hence its trace has absolute value at most \(N\).
The off-diagonal sum is real and has absolute value at most \(d(d-1)N\).
Applying these pointwise bounds inside the positive integral proves the
inequalities. At any real coupling the integrand is bounded on a compact
group, so differentiation under the integral is legitimate to every order.
The partition function is entire in complex coupling and strictly positive
on the real axis, hence its logarithm is analytic in a neighborhood of each
real point. It need not have an entire logarithm. With expectation under
the tilted probability measure denoted by \(\mathbb E_\beta\),
\[
 F_{d,N}'(\beta)=-\frac1N\mathbb E_\beta S_{d,N},\qquad
 F_{d,N}''(\beta)=\operatorname{Var}_\beta(S_{d,N})\geq0.
\]
The bound \(|S_{d,N}|\leq d^2N\) gives the Lipschitz constant.
Arzel\`a--Ascoli on each interval \([-k,k]\), followed by a diagonal
subsequence, gives compact convergence. Convexity and the Lipschitz bound
pass to that limit. This argument gives subsequences only, not uniqueness
of the limiting free energy.
\end{proof}

\begin{lemma}[One plaquette and two adjacent plaquettes \ours]
Let \(U,V,W\) be independent probability-Haar matrices in \(U(N)\),
\(N\geq2\). Write \(C(U,V)=\operatorname{Tr}(UVU^*V^*)\) and
\(X(U,V)=C(U,V)+\overline{C(U,V)}\). Then
\begin{align*}
 \mathbb E C(U,V)&=N^{-1},&
 \mathbb E C(U,V)^2&=\frac4{N^2-1},&
 \mathbb E|C(U,V)|^2&=\frac{N^2}{N^2-1},\\
 \operatorname{Var}X(U,V)&=
       \frac{2(N^2+4)}{N^2-1}-\frac4{N^2},&
 \operatorname{Cov}(X(U,V),X(U,W))&=\frac4{N^2}.
\end{align*}
Plaquettes involving disjoint pairs of matrices are independent.
\end{lemma}
\begin{proof}
The first Haar identity is
\(\mathbb E_U UAU^*=N^{-1}\operatorname{Tr}(A)I\).
For fixed matrices \(A,B,C,D\), the second identity gives
\begin{align*}
 &\mathbb E_U\bigl[\operatorname{Tr}(UAU^*B)
                    \operatorname{Tr}(UCU^*D)\bigr]\\
 &\quad=\frac{\operatorname{Tr}A\operatorname{Tr}B
                    \operatorname{Tr}C\operatorname{Tr}D
                    +\operatorname{Tr}(AC)\operatorname{Tr}(BD)}{N^2-1}\\
 &\qquad\ -\frac{\operatorname{Tr}(AC)\operatorname{Tr}B\operatorname{Tr}D
                    +\operatorname{Tr}A\operatorname{Tr}C\operatorname{Tr}(BD)}
                    {N(N^2-1)}.
\end{align*}
This is the degree-two specialization of
\citet[Corollary~2.4]{collinssniady2006integration}; it also follows by
projecting \(A\otimes C\) onto the span of the identity and the tensor flip,
then solving the two equations obtained by taking its trace and flip-trace.
The determinant of those two equations is \(N^2(N^2-1)\), explaining
the restriction \(N\geq2\).

Put \(t=\operatorname{Tr}V\) and \(s=\operatorname{Tr}(V^2)\).
Schur orthogonality, using
\(t^2=\chi_{(2)}+\chi_{(1,1)}\) and
\(s=\chi_{(2)}-\chi_{(1,1)}\), yields
\[
 \mathbb E|t|^2=1,\quad \mathbb E|t|^4=2,\quad
 \mathbb E|s|^2=2,\quad \mathbb E(s\overline t^{\,2})=0.
\]
Both degree-two characters exist and are distinct for \(N\geq2\).
First condition on \(V\). Substituting \(A=C=V\), \(B=D=V^*\)
into the second identity gives
\[
 \mathbb E_U C(U,V)^2
 =\frac{|t|^4+|s|^2}{N^2-1}
   -\frac{s\overline t^{\,2}+\overline s t^2}{N(N^2-1)}.
\]
For the absolute square, use
\(\overline{C(U,V)}=\operatorname{Tr}(UV^*U^*V)\); the result is
\[
 \mathbb E_U|C(U,V)|^2
     =\frac{|t|^4+N^2-2|t|^2}{N^2-1}.
\]
Averaging the displayed formulas proves the three moments and hence the
variance. Finally, conditioning on the shared matrix \(U\), the two
plaquettes are independent and each has conditional mean
\(2|\operatorname{Tr}U|^2/N\). Their covariance is consequently
\(4\operatorname{Var}(|\operatorname{Tr}U|^2)/N^2=4/N^2\).
Independence of disjoint pairs is immediate from product Haar measure.
\end{proof}

\begin{proposition}[Exact first and second derivatives \ours]
For \(d\geq2\) and \(N\geq2\),
\begin{align}
 F_{d,N}'(0)&=-d-\frac{d(d-1)}{N^2},\label{eq:1982-first}\\
 F_{d,N}''(0)&=d(d-1)\left\{
      \frac{N^2+4}{N^2-1}+\frac{4d-10}{N^2}\right\}.
       \label{eq:1982-second}
\end{align}
For \(N=1\), the exact answer instead is
\(F_{d,1}(\beta)=-d^2\beta\).
\end{proposition}
\begin{proof}
The mean of \(S_{d,N}\) is
\(dN+d(d-1)/N\), which gives \eqref{eq:1982-first}.
There are \(m=d(d-1)/2\) unordered plaquettes and
\(a=d\binom{d-1}{2}\) unordered pairs of distinct plaquettes sharing a matrix.
The constant diagonal terms contribute no variance. The preceding lemma and
the factor of two for covariance terms in a variance give
\begin{align*}
 \operatorname{Var}S_{d,N}
 &=m\left\{\frac{2(N^2+4)}{N^2-1}-\frac4{N^2}\right\}
       +2a\frac4{N^2}\\
 &=d(d-1)\left\{\frac{N^2+4}{N^2-1}+\frac{4d-10}{N^2}\right\}.
\end{align*}
Since \(F_{d,N}''(0)=\operatorname{Var}S_{d,N}\), the formula follows.
At \(N=1\) every scalar commutator is one, so the integral is
\(\exp(-\beta d^2)\). No substitution into the singular \(N\geq2\)
formula is permissible.
\end{proof}

\begin{corollary}[The justified local answer for the notebook convention \ours]
There is an \(\varepsilon>0\) such that the large-\(N\) free energy in
the notebook exists for real \(|\beta|<\varepsilon\), has an analytic
extension near zero, and
\[
                 F_4(\beta)=-4\beta+6\beta^2+O(\beta^3).
\]
More generally \(F_d(\beta)=-d\beta+\tfrac12d(d-1)\beta^2+O(\beta^3)\)
locally. The error here is a statement about the analytic limiting function,
not a uniform finite-\(N\) error inferred from a finite computation.
\end{corollary}
\begin{proof}
Apply the local unitary-integral theorem stated above, then take the limits
of \eqref{eq:1982-first} and \eqref{eq:1982-second}.
The theorem, rather than the numerical or symbolic checks, justifies the
identification of these limits with derivatives of \(F_d\).
\end{proof}

\begin{remark}[Verification \oursComputed]
The independent calculation in \nolinkurl{code/haar_check.py} expands products of traces
into matrix indices and applies the degree-one and degree-two Weingarten
formulas directly. It counts free index loops with a union--find algorithm.
This does not use the conditional-covariance argument of the proof. Six
basic mixed moments agree as rational functions of the symbolic integer
\(N\geq2\); summing all 584 ordered plaquette pairs for
\(d=2,3,4,5\) independently reproduces \eqref{eq:1982-second}.
Additional exact rational evaluations use \(N=2,3,4,7,11\).
For the four-matrix model the second derivatives at \(N=2,3,4\) are
\(50,55/2,41/2\), respectively, and tend to \(12\).
The separate scalar check gives zero second derivative at \(N=1\).
The check confirms the finite-rank restriction and both coefficients;
it makes no inference from finite samples to a global analytic limit.
An independent review also rederived the single-plaquette moments from
\(\mathbb E\chi_\rho([U,V])=1/\dim\rho\): Schur averaging first in
\(U\) gives \(|\chi_\rho(V)|^2/\dim\rho\), and averaging in \(V\)
gives the identity. Applying it to the symmetric and alternating squares,
and to the decomposition of the fundamental representation tensored with
its dual, reproduces the same second moments for every \(N\geq2\).
That review separately checked the adjacent-pair count and the exact scope
of the cited local-limit theorem; it found no mathematical correction.
The source, symbolic output, and machine-readable rational results are
included in the archive.
\end{remark}

\paragraph{What remains and what another run could do.}
Neither the two coefficients nor subsequential compactness select a global
large-\(N\) branch. The remaining obstacle is control of the interacting
unitary ensemble outside the small-potential region, including competing
states and any transition points. A complete physical reduction is an
additional question. One further bounded run could calculate higher connected
Haar cumulants, provide explicit certified bounds on a small-coupling interval,
and compare conventions against the original lattice strong-coupling series.
The precise next calculation is the third and fourth Taylor coefficients of
the same ordered-sum integral with their finite-\(N\) corrections, followed by
a justified limiting argument. It would enlarge the controlled local answer;
there is no basis here for promising an all-coupling solution.

\section{Problem index}\label{problems:1982}
Earlier-year occurrence identifiers are retained as text; current-year identifiers link to their entries.

\subsection{GP-1982-01: Quantization of cyclic invariant charges of the closed string}\label{prob:GP-1982-01}
Construct and describe an associative formal quantum deformation of the Poisson algebra of cyclic iterated integrals of closed-string chiral currents, with prescribed classical bracket.
\par\noindent\textit{Occurrences:} \hyperref[item:1982-562.9]{1982-562.9}.
\par\noindent\textit{Status (2026-11):} As of November 2026, abstract existence follows from quasi-Lie-bialgebra quantization, with the string application stated by Bahns--Meinecke. Naive canonical normal ordering fails. A preferred explicit presentation or positive representation is a stronger target than the recorded existence question.

\subsection{GP-1982-02: Global free energy of the recorded four-unitary matrix integral}\label{prob:GP-1982-02}
For four independent probability-Haar matrices in \(U(N)\), set \[Z_N(\beta)=\mathbb E\exp\!\left[-\beta N\sum_{\mu,\nu=1}^4\operatorname{tr}(U_\mu U_\nu U_\mu^\dagger U_\nu^\dagger)\right].\] Determine the existence and value of \(F(\beta)=\lim_{N\to\infty}N^{-2}\log Z_N(\beta)\) for all real \(\beta\), including analytic branches and transition values. The full ordered sum and negative exponent are retained.
\par\noindent\textit{Occurrences:} \hyperref[item:1982-565.6]{1982-565.6}.
\par\noindent\textit{Status (2026-11):} As of November 2026, rigorous small-potential theory supplies a local analytic limit; the chapter computes \(F(\beta)=-4\beta+6\beta^2+O(\beta^3)\). The all-coupling formula and transition structure for this precise convention were not located. This is a bounded search conclusion, not a claim that no relevant unpublished solution exists.

\subsection{GP-1982-03: Reducibility and first singular vectors of Virasoro Verma modules}\label{prob:GP-1982-03}
Determine all pairs \((h,c)\) for which the complex Virasoro Verma module is reducible, and the least positive level of a singular vector.
\par\noindent\textit{Occurrences:} \hyperref[item:1982-566.7]{1982-566.7}, \hyperref[item:1982-573.6]{1982-573.6}.
\par\noindent\textit{Status (2026-11):} As of November 2026, solved before the September 1982 talk by the Kac determinant criterion, with a proof in Feigin--Fuchs (1982). Later appearances in this year reuse this identifier.

\subsection{GP-1982-04: Recover a crossing value from hyperbola orbital integrals}\label{prob:GP-1982-04}
For compactly supported smooth \(f\) on \(\mathbb R^2\), recover \(f(0,0)\) from the family \(\int f(x,y)\delta(xy-c)\,dx\,dy\), \(c\ne0\).
\par\noindent\textit{Occurrences:} \hyperref[item:1982-568.5]{1982-568.5}.
\par\noindent\textit{Status (2026-11):} As of November 2026, an elementary solved singular-asymptotic problem: the leading coefficient is \(-2f(0,0)\log|c|\). The notebook already refers to the older generalized-functions treatment.

\subsection{GP-1982-05: Nonminimal finite-energy SU(2) Yang–Mills–Higgs critical points}\label{prob:GP-1982-05}
In the Prasad--Sommerfield limit on \(\mathbb R^3\), with \(|\Phi|\to v>0\), do finite-energy stationary points exist outside the Bogomolny minimum equations? A fixed magnetic-charge sector must be specified separately.
\par\noindent\textit{Occurrences:} \hyperref[item:1982-571.2]{1982-571.2}.
\par\noindent\textit{Status (2026-11):} As of November 2026: existence was published in September 1982. Later sectorwise multiplicity is covered by Taubes with his 1989 correction. This Higgs-field problem on \(\mathbb R^3\) differs from the earlier pure four-dimensional Yang--Mills problem GP-1977-04.

\subsection{GP-1982-06: Existence of a unit instanton when the second cohomology vanishes}\label{prob:GP-1982-06}
For every smooth metric on a closed oriented simply connected four-manifold with \(H^2(M;\mathbb R)=0\), does the \(SU(2)\) bundle of charge one admit an anti-self-dual connection?
\par\noindent\textit{Occurrences:} \hyperref[item:1982-571.3]{1982-571.3}.
\par\noindent\textit{Status (2026-11):} As of November 2026: yes, by the Taubes result reported in 1982. Verification is indirect through Donaldson, Section 3, citing the exact Taubes theorem; the incorrect notebook journal/year is separately logged.

\subsection{GP-1982-07: Asymptotic completeness for three quantum particles with Coulomb forces}\label{prob:GP-1982-07}
For three nonrelativistic quantum particles in \(\mathbb R^3\) with positive finite masses and pair Coulomb interactions, remove the center of mass. Do the Dollard-modified channel wave operators exhaust the continuous spectral subspace of the resulting Schrödinger Hamiltonian?
\par\noindent\textit{Occurrences:} \hyperref[item:1982-571.20]{1982-571.20}.
\par\noindent\textit{Status (2026-11):} As of November 2026: established in Enss 1989, equations (1.12)--(1.13). The 1983 chapter is short-range and does not verify the Coulomb claim. Earlier 1971-252.3 is a thematic wave-operator discussion without an identical problem ID.

\subsection{GP-1982-08: Primary correlators, conformal blocks and null-state equations}\label{prob:GP-1982-08}
Determine the constraints and calculational recursions for primary-field matrix elements from Virasoro covariance and specified null states, distinguishing the additional channel and three-point data needed to obtain a full correlation function.
\par\noindent\textit{Occurrences:} \hyperref[item:1982-573.4]{1982-573.4}.
\par\noindent\textit{Status (2026-11):} As of November 2026, the BPZ program, received in November 1983 and published in 1984, gives conformal blocks and degenerate-field differential equations. The notebook data alone do not determine structure constants or a unique full conformal theory; this is an underdetermination, not a newly posed open classification problem.

\subsection{GP-1982-09: The Virasoro representation content of the Liouville spectrum}\label{prob:GP-1982-09}
Relate the stationary spectrum of a quantized and renormalized Liouville field theory to highest-weight Virasoro representations. A precise modern version fixes \(0<\gamma<2\), \(\mu>0\), and the reflection-positive probabilistic Liouville Hamiltonian.
\par\noindent\textit{Occurrences:} \hyperref[item:1982-573.14]{1982-573.14}.
\par\noindent\textit{Status (2026-11):} As of November 2026, the real-coupling implementation has a complete continuous spectral resolution with Virasoro descendant states, proved in the 2020 preprint of Guillarmou--Kupiainen--Rhodes--Vargas and published in 2024. This does not identify the unrenormalized notebook action at every \(D<26\) with that theorem.

\subsection{GP-1982-20: Flag superspaces with reduced Grassmannian G(2,4)}\label{prob:GP-1982-20}
Classify classical homogeneous flag superspaces whose reduced manifold is \(G(2,4)\), with equivalence and flag-stabilizer convention specified.
\par\noindent\textit{Occurrences:} \hyperref[item:1982-576.4]{1982-576.4}.
\par\noindent\textit{Status (2026-11):} As of November 2026: Manin’s 1983 chapter is the close publication match. Its complete classification has not been independently read; this entry does not assert that the question is still open.

\subsection{GP-1982-21: Polyakov accessory-parameter conjecture}\label{prob:GP-1982-21}
For the complete hyperbolic metric on a sphere with distinct punctures, does a suitably regularized classical Liouville action generate the uniformizing accessory parameters under motion of the punctures?
\par\noindent\textit{Occurrences:} \hyperref[item:1982-577.9]{1982-577.9}.
\par\noindent\textit{Status (2026-11):} As of November 2026: proved by Zograf and Takhtajan; 1985 announcement, 1987 Russian full paper, 1988 English translation. With their projective-connection normalization, \(c_i=-(2\pi)^{-1}\partial_{z_i}S\).

\subsection{GP-1982-22: A p-adic analogue of the circle-diffeomorphism representations}\label{prob:GP-1982-22}
Construct a p-adic version of the almost-invariant-kernel and Fock/Gaussian construction of unitary diffeomorphism-group representations. The documented realization takes analytic diffeomorphisms of \(\mathbb P^1(\mathbb Q_p)\).
\par\noindent\textit{Occurrences:} \hyperref[item:1982-581.8]{1982-581.8}.
\par\noindent\textit{Status (2026-11):} As of November 2026: Neretin’s 1984 construction solves the explicitly credited analytic-projective-line version \citep[introduction and Theorem 1]{neretin1984padic}. The group specification is supplied by that solution.

\subsection{GP-1981-41: Geometric realization of compact-group weight multiplicities}\label{prob:GP-1981-41}
For an integral compact coadjoint orbit quantizing an irreducible representation, express its torus-weight multiplicities geometrically through the correctly normalized orbital character or the index of the corresponding reduced space, with the shifted quotient line bundle.
\par\noindent\textit{Occurrences:} 1981-556.1, 1981-556.6, \hyperref[item:1982-584.5]{1982-584.5}.
\par\noindent\textit{Status (2026-11):} As of November 2026: finite-dimensional multiplicities are classical; the reduced-space index interpretation extends to singular reductions by \citet[Theorem 2.9 and Corollary 2.11]{meinrenkensjamaar1999singular} (July 1997 preprint; 1999 publication). A level set is not itself the reduced space, and reduced volume alone is not the finite-dimensional multiplicity.

\subsection{GP-1982-40: Multiplicativity of zeta-regularized determinants}\label{prob:GP-1982-40}
For admissible elliptic operators, determine whether \(\det_\zeta(AB)=\det_\zeta A\det_\zeta B\), and if not identify its defect.
\par\noindent\textit{Occurrences:} \hyperref[item:1982-587.7]{1982-587.7}.
\par\noindent\textit{Status (2026-11):} November 2026: false in general; local multiplicative-anomaly formulas solve the well-posed operator versions. See \citep{kontsevichvishik1994geometry,castillogarate2015anomaly}.

\subsection{GP-1982-41: Classical modified Yang–Baxter operators on sl(2)}\label{prob:GP-1982-41}
Classify complex-linear \(R:\mathfrak{sl}_2(\mathbb C)\to\mathfrak{sl}_2(\mathbb C)\) with modified Yang--Baxter defect \(-[x,y]\), up to Lie-algebra automorphism.
\par\noindent\textit{Occurrences:} \hyperref[item:1982-589.3]{1982-589.3}.
\par\noindent\textit{Status (2026-11):} November 2026: solved via the equivalent weight-one Rota--Baxter classification; no skew restriction may be silently imposed. See \citep{gubarevkozlov2023rota,semenov1983classical}.

\subsection{GP-1981-44: Spectral zeta constructions without self-adjointness}\label{prob:GP-1981-44}
Extend complex-power, zeta and spectral-asymmetry constructions beyond self-adjoint elliptic operators under suitable spectral-cut hypotheses.
\par\noindent\textit{Occurrences:} 1981-559.8, 1981-559.13, 1981-559.14, \hyperref[item:1982-590.2]{1982-590.2}.
\par\noindent\textit{Status (2026-11):} November 2026: complex powers and zeta functions are established for admissible elliptic operators; the cut must avoid principal-symbol spectrum as well as global spectrum. Boundary realizations need their own ellipticity conditions. See \citep{kontsevichvishik1994geometry,wodzicki1982spectral,seeley1968powers}.

\subsection{GP-1981-43: Eta regularity without a bounding manifold}\label{prob:GP-1981-43}
Is the global eta function regular at zero for a self-adjoint elliptic operator without assuming it bounds an elliptic complex?
\par\noindent\textit{Occurrences:} 1981-559.9, \hyperref[item:1982-590.9]{1982-590.9}.
\par\noindent\textit{Status (2026-11):} November 2026: yes for closed classical elliptic operators of positive order; APS III treats odd dimension and Gilkey 1981 supplies the general-dimensional result. Local residue densities need not vanish. See \citep{atiyahpatodisinger1976spectral3,gilkey1981residue}.

\subsection{GP-1982-42: Global continuation of curve zeta functions}\label{prob:GP-1982-42}
For every smooth projective geometrically connected curve over \(\mathbb Q\), prove the expected meromorphic continuation and functional equation of its completed arithmetic zeta function, including the correct bad and archimedean factors.
\par\noindent\textit{Occurrences:} \hyperref[item:1982-595.5]{1982-595.5}.
\par\noindent\textit{Status (2026-11):} November 2026: proved for genus one over \(\mathbb Q\) by modularity and for genus two over totally real fields by potential automorphy; open in unrestricted higher genus. See \citep{breuilconraddiamondtaylor2001modularity,boxercalegarigeepilloni2021surfaces,gee2026modularity}.

\subsection{GP-1982-43: Automorphy of the H1 of arbitrary curves}\label{prob:GP-1982-43}
Associate suitable algebraic automorphic data to the compatible \(H^1\) Galois representations of arbitrary smooth projective curves over \(\mathbb Q\), matching local Euler factors and completions.
\par\noindent\textit{Occurrences:} \hyperref[item:1982-595.6]{1982-595.6}.
\par\noindent\textit{Status (2026-11):} November 2026: general actual automorphy remains open. Potential automorphy is known for genus-two curves over totally real fields; the 2025 actual modularity theorem proves a specified residual-surjective ordinary genus-two class over \(\mathbb Q\). See \citep{boxercalegarigeepilloni2021surfaces,boxercalegarigeepilloni2025modularity,gee2026modularity}.

\section{Corrections to the transcription and translation}\label{corr:1982}
The evidence distinguishes a transcription correction from a mathematical emendation of wording faithfully preserved in the notebook.
\subsection*{P001027 (both; faithful error): Generalized Cartan conditions and relation (2)}
\noindent\textit{Transcription reads.} \(a_{ij}<0\), and \(e_i^{1-a_{ij}}e_j=f_i^{1-a_{ij}}f_j=0\).\par
\noindent\textit{Critical reading.} Allow \(a_{ij}=0\) for disconnected vertices; write adjoint powers for \(i\ne j\).\par
\noindent\textit{Evidence.} Scan inspected: strict inequality and omission of ad are present in the handwriting. These are mathematical emendations, not mistranscriptions.\par
\subsection*{P001027 (both; faithful error): Finite/affine positivity criterion}
\noindent\textit{Transcription reads.} Positivity of \((a_{ij})\) without a symmetrizer.\par
\noindent\textit{Critical reading.} For indecomposable symmetrizable A test the symmetric matrix DA with D positive diagonal.\par
\noindent\textit{Evidence.} The displayed B2 and G2 matrices are nonsymmetric. Scan and Russian agree; the qualification is necessary.\par
\subsection*{P001028 (both; faithful error): First affine matrix}
\noindent\textit{Transcription reads.} \(\left(\begin{smallmatrix}2&-2\\-2&-2\end{smallmatrix}\right)\).\par
\noindent\textit{Critical reading.} \(\left(\begin{smallmatrix}2&-2\\-2&2\end{smallmatrix}\right)\).\par
\noindent\textit{Evidence.} Scan inspected: lower-right -2 is clear. Kirillov repeats the correct +2 matrix lower on the same page.\par
\subsection*{P001028 (both; faithful error): Loop cocycle versus h2}
\noindent\textit{Transcription reads.} \(+\operatorname{Res}\operatorname{tr}(A\,dB)K\) and \(h_2=-h_1+K\).\par
\noindent\textit{Critical reading.} Change the cocycle to \(\operatorname{Res}\operatorname{tr}((dA)B)K\), or change h2 to -h1-K.\par
\noindent\textit{Evidence.} Scan inspected; differentiation of \(t^{-1}\) gives central term -K with the printed convention.\par
\subsection*{P001029 (both; faithful error): Principal Heisenberg bracket}
\noindent\textit{Transcription reads.} \(\delta(k+l)(2l+1)K\).\par
\noindent\textit{Critical reading.} \(\delta_{k+l+1,0}(2l+1)K\) for the printed cocycle.\par
\noindent\textit{Evidence.} Scan inspected. \(\operatorname{tr}(B_{2k+1}dB_{2l+1})=(2l+1)t^{k+l}dt\), so the residue condition is \(k+l=-1\).\par
\subsection*{P001029 (both; faithful error): Oscillator action indices and vacuum}
\noindent\textit{Transcription reads.} Positive B modes multiply; negative modes differentiate; indices alternate between \(x_k\) and \(x_{2k+1}\).\par
\noindent\textit{Critical reading.} Fix one indexing convention and match raising/annihilation operators to the sign of the chosen B modes and highest-weight vector.\par
\noindent\textit{Evidence.} Scan inspected. For the highest-weight convention, B1=e1+e2 annihilates v0, so it cannot act as multiplication by a nonzero variable on the vacuum 1.\par
\subsection*{P001030 (both; faithful error): Lorentzian area integrand}
\noindent\textit{Transcription reads.} \(\sqrt{\dot X^2(X^{\prime})^2-(\dot X\cdot X^{\prime})^2}\).\par
\noindent\textit{Critical reading.} Use \(\sqrt{(\dot X\cdot X^{\prime})^2-\dot X^2(X^{\prime})^2}\) for a timelike worldsheet.\par
\noindent\textit{Evidence.} Russian and scan 1030 checked. The displayed induced determinant is negative for Lorentzian signature.\par
\subsection*{P001030 (both; faithful error): Legendre-transform constraint}
\noindent\textit{Transcription reads.} \(p^2-(X^{\prime})^2=0\).\par
\noindent\textit{Critical reading.} \(p^2+T^2(X^{\prime})^2=0\), T=1 in the notes.\par
\noindent\textit{Evidence.} Scans 1030–31 checked; expanding both \((p+X^{\prime})^2=(p-X^{\prime})^2=0\) forces the plus sign.\par
\subsection*{P001031 (both; faithful error): Fourier canonical delta}
\noindent\textit{Transcription reads.} Opposite Fourier signs with \(\delta_{n+m,0}\).\par
\noindent\textit{Critical reading.} Use \(\delta_{nm}\), or same-sign expansions.\par
\noindent\textit{Evidence.} Russian and scan 1031 checked; direct Fourier orthogonality.\par
\subsection*{P001031 (both; faithful error): Reality and oscillator indices}
\noindent\textit{Transcription reads.} Individual modes Hermitian; equation~18 repeats n on the left, \(\delta_{nm}\) on the right.\par
\noindent\textit{Critical reading.} Use \(x_n^\dagger=x_{-n}\), \(p_n^\dagger=p_{-n}\), distinct n,m indices, target metric and consistent mode factors.\par
\noindent\textit{Evidence.} Scan 1031 inspected; a real periodic field has conjugate opposite modes.\par
\subsection*{P001032 (both; faithful error): Vacuum condition}
\noindent\textit{Transcription reads.} \(a_n\Psi_0=0\) for \(n\ge0\).\par
\noindent\textit{Critical reading.} Positive oscillators annihilate the vacuum; specify zero-mode momentum separately.\par
\noindent\textit{Evidence.} Scan 1032 checked against the center-of-mass momentum on the same page.\par
\subsection*{P001032 (both; faithful error): Quantum charge invariance}
\noindent\textit{Transcription reads.} Normal ordering gives \([Q,L_n]=0\) unconditionally.\par
\noindent\textit{Critical reading.} Retain as a historical assertion and annotate the anomalous canonical relations.\par
\noindent\textit{Evidence.} Scan 1032 inspected; Bahns 2004 sections~4–5 directly analyze the prescription.\par
\subsection*{P001033 (both; faithful error): Gauge transformation (5)}
\noindent\textit{Transcription reads.} \(D=d+A,\ \phi^g=g\phi,\ A^g=gAg^{-1}+dg\,g^{-1}\).\par
\noindent\textit{Critical reading.} For these conventions use \(A^g=gAg^{-1}-dg\,g^{-1}\).\par
\noindent\textit{Evidence.} Scan 1033 and Russian checked: the plus is present. Expansion of Dg(g phi) fixes the sign.\par
\subsection*{P001035 (both; faithful error): Alice-string criterion (14)--(16)}
\noindent\textit{Transcription reads.} A reversing element merely in \(\widetilde G\) implies a topological string.\par
\noindent\textit{Critical reading.} Require reversing holonomy in a disconnected component of the stabilizer H; specify simply connected gauge group for pi1(G/H)=pi0(H).\par
\noindent\textit{Evidence.} Scan 1035 and Russian checked; the omission is in the notebook. Schwarz 1982 abstract explicitly supplies disconnected unbroken subgroup.\par
\subsection*{P001037 (both; ambiguity): Agenda item 7, large-rank equality}
\noindent\textit{Transcription reads.} \(u(\infty)=sp(\infty)=O(\infty)\).\par
\noindent\textit{Critical reading.} Interpret only as an informal leading-planar comparison after fixing theories and coupling normalizations; do not identify the direct-limit classical Lie algebras.\par
\noindent\textit{Evidence.} Scan P001037 and Russian checked. The equality occurs in the physics agenda. The classical matrix series have different Lie-algebra structures and nonorientable corrections; no exact algebra isomorphism follows.\par
\subsection*{P001038 (both; faithful error): Feigenbaum constant (2)}
\noindent\textit{Transcription reads.} \(\delta\approx4.6992\).\par
\noindent\textit{Critical reading.} \(\delta\approx4.6692\).\par
\noindent\textit{Evidence.} Scan 1038 inspected:4.6992 is clear, agreeing with Russian. A faithful numerical error.\par
\subsection*{P001038 (both; faithful error): Flow equation (3)}
\noindent\textit{Transcription reads.} \(\dot x_i=f(x_1,\ldots,x_r,\mu)\) for every i.\par
\noindent\textit{Critical reading.} Use component functions \(f_i\), or write \(\dot x=f(x,\mu)\) with vector-valued f.\par
\noindent\textit{Evidence.} Scan 1038 checked; the component index on f is omitted in the notebook.\par
\subsection*{P001038 (both; faithful error): Polynomial solution of (7)}
\noindent\textit{Transcription reads.} A polynomial of degree ten solves the functional equation.\par
\noindent\textit{Critical reading.} A degree-ten polynomial approximates the fixed point; validated analytic and numerical error bounds establish an exact nearby solution.\par
\noindent\textit{Evidence.} Scan 1038 and Lanford 1981 preprint section~3 checked. This is a mathematical qualification, not a mistranscription.\par
\subsection*{P001038 (both; faithful error): Euclidean matrix action and real form}
\noindent\textit{Transcription reads.} Anti-Hermitian \(\mathfrak u(N)\) field with positive \(\operatorname{tr}(\partial\varphi)^2\).\par
\noindent\textit{Critical reading.} Choose Hermitian matrix fields, or use the positive form \(-\operatorname{tr}(AB)\) on anti-Hermitian matrices.\par
\noindent\textit{Evidence.} Russian and scan 1038 inspected. For nonzero anti-Hermitian A, tr(A squared) is negative. This is a sign convention repair.\par
\subsection*{P001039 (both; faithful error): White-noise covariance and Gibbs weight}
\noindent\textit{Transcription reads.} Unit covariance with drift \(-\nabla S\) and stationary weight \(e^{-S}\).\par
\noindent\textit{Critical reading.} Use covariance \(2\delta\) for this drift, or halve the drift; otherwise the stationary weight is \(e^{-2S}\).\par
\noindent\textit{Evidence.} Russian and scan 1039 inspected; finite-dimensional stationary Fokker--Planck equation checked.\par
\subsection*{P001039-P001040 (both; ambiguity): Conjugated noise and spacetime covariance}
\noindent\textit{Transcription reads.} Finite matrix noise transported as \(J(x,\tau)=e^{ipx}\bar J(\tau)e^{-ipx}\) replaces the spacetime noise.\par
\noindent\textit{Critical reading.} Keep this as a regulated quenched planar prescription; it is not an exact realization of independent spacetime white noise at finite \(N\).\par
\noindent\textit{Evidence.} Each fixed pair of matrix entries has a single spatial phase \(e^{i(p_i-p_j)x}\), so its covariance cannot equal a spatial delta distribution before the specified large-rank and momentum-averaging limits.\par
\subsection*{P001040 (both; faithful error): Gaussian momentum variance}
\noindent\textit{Transcription reads.} Density \(e^{-p^2/\Lambda^2}\), variance \(1/\Lambda^2\).\par
\noindent\textit{Critical reading.} Variance \(\Lambda^2/2\) per real component for the written density.\par
\noindent\textit{Evidence.} Russian and scan 1040 agree; direct Gaussian integral gives the variance.\par
\subsection*{P001040 (both; faithful error): Reduced drift in equation (16)}
\noindent\textit{Transcription reads.} \(U(\varphi)\) in the Langevin drift.\par
\noindent\textit{Critical reading.} \(U^{\prime}(\varphi)\) for the action potential named earlier.\par
\noindent\textit{Evidence.} Scan 1040 inspected: no derivative is visible. Variation of tr U gives tr U prime delta phi.\par
\subsection*{P001040 (both; faithful error): Formal reduced expectation in equations (17)--(18)}
\noindent\textit{Transcription reads.} Unnormalized integral over \(U(N)\) with the displayed positive exponent.\par
\noindent\textit{Critical reading.} Specify the matrix integration domain and stable action convention, and divide by its partition function for an expectation.\par
\noindent\textit{Evidence.} Scan 1040 and Russian inspected. Earlier phi is a Lie-algebra field, whereas U(N) denotes the compact group; the notebook does not resolve the change.\par
\subsection*{P001040 (both; provenance): Name after external-field equation (20)}
\noindent\textit{Transcription reads.} \ru{Пиризен} [uncertain first vowel], Gross; then Fateev.\par
\noindent\textit{Critical reading.} Probable identification: Brézin--Gross (1980); retain uncertainty in the written surname and in the Fateev attribution.\par
\noindent\textit{Evidence.} Brézin author academy bibliography confirms the exact external-field paper with Gross. Handwriting does not justify silently replacing the ambiguous surname.\par
\subsection*{P001042 (both; faithful error): Zero-mode commutators at the top of the page}
\noindent\textit{Transcription reads.} \(e_0=[e_{-1},e_1]\), and comparison with \([e_{-2},e_2]\).\par
\noindent\textit{Critical reading.} \([e_{-1},e_1]=2e_0\), \([e_{-2},e_2]=4e_0+z/2\).\par
\noindent\textit{Evidence.} Scan 1042 inspected; substitute indices in the exact bracket on 1041.\par
\subsection*{P001042 (both; faithful error): Highest weight following the Feigin--Fuchs reference}
\noindent\textit{Transcription reads.} \(h=\mu(\mu-2\lambda+1)/2\).\par
\noindent\textit{Critical reading.} \(h=\mu(\mu-2\lambda-1)/2\) for the stated charge-zero tail.\par
\noindent\textit{Evidence.} Scan 1042 visibly has +1 and a question mark. Feigin--Fuchs 1982 Proposition~3.1 Russian p59 has -1 in this exact convention.\par
\subsection*{P001042 (both; faithful error): Equation (11), determinant factors}
\noindent\textit{Transcription reads.} Minus \((c-13)(k_i^2-1)/24\) in both factors.\par
\noindent\textit{Critical reading.} Plus \((c-13)(k_i^2-1)/24\) in both factors.\par
\noindent\textit{Evidence.} Scan 1042 checked; Feigin--Fuchs 1982 Theorem~3.1 Russian p58 prints plus. At c0, r1,s2, the corrected factor has root h0, agreeing with the level 2 singular vector.\par
\subsection*{P001042 (both; faithful error): Equations (14)--(16), purported oscillator realization}
\noindent\textit{Transcription reads.} \(F_k=\sum_i x_i\partial_{x_{i+k}}\) used as the quadratic construction modes.\par
\noindent\textit{Critical reading.} Use actual multiplication/differentiation Heisenberg modes and specify zero-mode, normal ordering and background-charge coefficient consistently.\par
\noindent\textit{Evidence.} Scan 1042 checked. Every written \(F_k\) kills \(1\), so their commutator on \(1\) is \(0\) and cannot equal a nonzero scalar; Neretin (1983), Theorem 2 gives an independently printed oscillator construction.\par
\subsection*{P001044 (both; faithful error): Unqualified lattice-walk convergence assertion}
\noindent\textit{Transcription reads.} The sum over all walks with weight \(e^{-mN}\) converges.\par
\noindent\textit{Critical reading.} Specify step normalization and mass range; \(2de^{-m}<1\) gives the uniform geometric-series regime for unnormalized nearest-neighbor walks.\par
\noindent\textit{Evidence.} Russian page 1044 read. There are \((2d)^N\) unrestricted N-step walks; fixed-endpoint critical behavior has a separate dimensional qualification.\par
\subsection*{P001044 (both; faithful error): Endpoint identity and equations (15)--(18)}
\noindent\textit{Transcription reads.} \(\sum n_i=(x_1-x_2)/\Delta\); Fourier delta integral has no Jacobian.\par
\noindent\textit{Critical reading.} \(\sum n_i=(x_N-x_0)/\Delta\); use density \(\Delta^{-d}G_N\), or retain factor \(\Delta^d\) in its Fourier transform.\par
\noindent\textit{Evidence.} Russian page 1044 read; endpoint notation conflicts with the immediately defined R. Change of variables in the delta function gives the Jacobian.\par
\subsection*{P001044-P001045 (both; faithful error): Small-momentum expansion and equation (20)}
\noindent\textit{Transcription reads.} Linear isotropic term \(b\vec p\) in the exponent.\par
\noindent\textit{Critical reading.} First nonzero term \(-\Delta^2|p|^2/(2d)\) for normalized isotropic steps; tune step fugacity on the same scale.\par
\noindent\textit{Evidence.} Scan 1045 and Russian page 1044 checked. Odd sphere moments vanish and covariance is I/d.\par
\subsection*{P001045 (both; faithful error): Gaussian spring paths, equations (22)--(23)}
\noindent\textit{Transcription reads.} Gaussian lengths, unnormalized increments and immediate identification with the preceding proper-time convention.\par
\noindent\textit{Critical reading.} Gaussian displacement vectors with normalized kernels; lengths have a chi-type distribution. Kinetic coefficient \(1/2\) gives positive kinetic symbol \(|p|^2/2\), requiring time/mass rescaling to match (21).\par
\noindent\textit{Evidence.} Scan 1045 inspected. Direct Gaussian convolution gives the factors; Brownian kinetic action is symbolic.\par
\subsection*{P001045 (both; ambiguity): Function name after equation (24)}
\noindent\textit{Transcription reads.} Inverse Euclidean propagator called a Hankel function.\par
\noindent\textit{Critical reading.} Use modified Bessel K at positive real mass/distance, or state imaginary-argument analytic continuation.\par
\noindent\textit{Evidence.} Scan 1045 inspected. The heat-kernel integral decays exponentially, unlike a real-argument oscillatory Hankel kernel.\par
\subsection*{P001046 (both; faithful error): Classical minimization after equation (28)}
\noindent\textit{Transcription reads.} Arbitrary nonzero cosmological term reproduces the original minimal-surface equations.\par
\noindent\textit{Critical reading.} Classical equivalence uses zero cosmological term; the quantum term is a renormalization counterterm.\par
\noindent\textit{Evidence.} Russian 1046 and Polyakov 1981 pp. 208--209 read. The two-dimensional metric equation trace isolates lambda.\par
\subsection*{P001046 (both; ambiguity): Resolution definition and real illustration}
\noindent\textit{Transcription reads.} Generic bijectivity and equal dimension alone define desingularization; the cylinder picture is said to be exactly the Springer example.\par
\noindent\textit{Critical reading.} An algebraic resolution is proper birational with smooth source; retain the cylinder as a real topological rank-one model and distinguish complex exceptional fibers.\par
\noindent\textit{Evidence.} Scans 1046--47 inspected. Incidence projection is proper because the Borel variety is projective. The real exceptional fiber is RP1, the complex fiber CP1.\par
\subsection*{P001047 (both; faithful error): Power expansion and limit of hyperbola integrals}
\noindent\textit{Transcription reads.} Expansion of \(\delta(xy-c)\) in powers of c; direct limiting axis integral.\par
\noindent\textit{Critical reading.} Include logarithmic terms and finite-part axis distributions; the leading full-plane term is \(-2\log|c|\,\delta_{(0,0)}\).\par
\noindent\textit{Evidence.} Scan 1047 checked. Direct integration f(x,c/x)/|x| with f constant near 0 gives the logarithmic divergence, so a plain Taylor series is impossible.\par
\subsection*{P001047 (both; ambiguity): Equation (4), relation between masses}
\noindent\textit{Transcription reads.} \(m^2=m_0-\varkappa\).\par
\noindent\textit{Critical reading.} Treat as a schematic entropy subtraction until the lattice spacing, proper-time rescaling and diffusion coefficient have been fixed.\par
\noindent\textit{Evidence.} Scan 1047 inspected. Length and Brownian time scale differently; the text itself states regularization dependence.\par
\subsection*{P001048 (both; ambiguity): Equation (6), metric on metric variations}
\noindent\textit{Transcription reads.} Unrestricted coefficient c in a squared norm.\par
\noindent\textit{Critical reading.} Require \(c>-1/2\) in two dimensions for positivity on all symmetric variations.\par
\noindent\textit{Evidence.} Scan 1048 inspected; trace/traceless decomposition gives |h0|²+(c+1/2)(tr h)².\par
\subsection*{P001048 (both; faithful error): Equation (7), covariant derivative}
\noindent\textit{Transcription reads.} \(\nabla_a\varepsilon_b=\partial_a\varepsilon_b+\Gamma^c_{ab}\varepsilon_c\).\par
\noindent\textit{Critical reading.} \(\nabla_a\varepsilon_b=\partial_a\varepsilon_b-\Gamma^c_{ab}\varepsilon_c\).\par
\noindent\textit{Evidence.} The scan and diplomatic Russian explicitly have plus. A covector carries the dual connection; direct metric-compatibility check requires minus.\par
\subsection*{P001048 (both; faithful error): Sentence after equations (8)--(9)}
\noindent\textit{Transcription reads.} The integral is now well defined.\par
\noindent\textit{Critical reading.} The quotient is a formal gauge-invariant prescription, pending regularization, zero-mode treatment and measure construction.\par
\noindent\textit{Evidence.} The scan agrees with the transcription; an infinite-dimensional field-space metric does not supply a locally finite flat volume. Later rigorous sphere construction requires renormalization and Seiberg bounds.\par
\subsection*{P001048 (both; ambiguity): Equation (10), residual conformal group}
\noindent\textit{Transcription reads.} Conformally flat representative on a plane; remaining group finite-dimensional.\par
\noindent\textit{Critical reading.} Conformal flatness is local; retain conformal moduli and distinguish compact-surface global automorphisms from local conformal transformations.\par
\noindent\textit{Evidence.} Scan 1048 inspected; the global and local statements have different hypotheses.\par
\subsection*{P001048 (both; faithful error): Sentence and equation (11), equal determinants}
\noindent\textit{Transcription reads.} Ghost and scalar determinants are identical; determinant exponent (26-d)/2.\par
\noindent\textit{Critical reading.} The scalar and ghost operators differ. Their conformal anomalies combine with coefficient 26-d, while zero modes, determinant normalizations and moduli factors remain.\par
\noindent\textit{Evidence.} Scan 1048 inspected. Polyakov 1981 equations~13--24 distinguish matter and ghost operators before anomaly calculation.\par
\subsection*{P001050 (both; provenance): Speaker heading and catalogue attribution}
\noindent\textit{Transcription reads.} Catalogue V. L. Ginzburg; notebook V. Ginzburg.\par
\noindent\textit{Critical reading.} Retain V. L. Ginzburg as the catalogue attribution and V. Ginzburg as the notebook reading. Do not expand the manuscript initial to V. A. without direct identity evidence.\par
\noindent\textit{Evidence.} The scan reads only \ru{В. Гинзбург} and labels the presentation as following Guillemin--Sternberg. The cited 1983 paper is by V. A. Ginzburg, but its authorship and mathematical relevance do not identify the speaker. No independent dated identification was found.\par
\subsection*{P001050 (both; ambiguity): Lagrangian relation in a product cotangent bundle}
\noindent\textit{Transcription reads.} \(\Lambda\subset T^*X\times T^*X\), with no symplectic sign convention.\par
\noindent\textit{Critical reading.} Use the product form \(\omega\oplus(-\omega)\), or equivalently reverse the second covector when treating a kernel conormal.\par
\noindent\textit{Evidence.} Scan 1050 inspected. Even the diagonal is Lagrangian only for the reversed second symplectic form.\par
\subsection*{P001051 (both; faithful error): Product example, normal directions}
\noindent\textit{Transcription reads.} Covectors normal to \(X_1\), along \(X_2\), followed by a diagonal in \(T^*X_1\).\par
\noindent\textit{Critical reading.} The formula and operator conclusion require covectors annihilating \(TX_2\), hence \(\Sigma=T^*X_1\times0_{X_2}\).\par
\noindent\textit{Evidence.} Scan 1051 confirms the verbal factor order and the incompatible displayed formula. A covector normal to the first factor instead has first cotangent component zero and exchanges the two operator roles.\par
\subsection*{P001051 (both; faithful error): Minimax and sector transition}
\noindent\textit{Transcription reads.} \(\max_t\min_{\rm curve}\); a monopole--antimonopole pair after a charge-one definition.\par
\noindent\textit{Critical reading.} \(\inf_{\gamma\in\Gamma}\sup_{t\in S^1}E(\gamma(t))\). The pair has total charge zero.\par
\noindent\textit{Evidence.} Scan P001051 confirms the order and pair. The initial subscript is doubtful. The variational order is an emendation; degree addition fixes total charge.\par
\subsection*{P001052 (both; provenance): Taubes unit-instanton reference}
\noindent\textit{Transcription reads.} Taubes, Comm. Math. Phys., 1981.\par
\noindent\textit{Critical reading.} Taubes, Self-dual Yang--Mills connections on non-self-dual four-manifolds, JDG 17 (1982), 139--170.\par
\noindent\textit{Evidence.} Scan P001052 confirms journal/year. Donaldson reference 6 and Taubes Part I reference 24 identify the JDG work; full original proof not read.\par
\subsection*{P001052 (both; ambiguity): Moduli dimension and boundary}
\noindent\textit{Transcription reads.} The moduli space is a five-manifold with boundary \(M\).\par
\noindent\textit{Critical reading.} The dimension is virtual; regularity needs unobstructedness or perturbation. The boundary comes from bubbling compactification.\par
\noindent\textit{Evidence.} Donaldson Lemmas 1--2 state generic perturbation and a product end with compact complement. Scan P001052 inspected.\par
\subsection*{P001052 (both; ambiguity): Pontryagin class and bounding}
\noindent\textit{Transcription reads.} Boundary iff the Pontryagin class vanishes.\par
\noindent\textit{Critical reading.} For closed oriented \(M\), oriented null-cobordism iff \(\langle p_1(TM),[M]\rangle=3\sigma(M)=0\).\par
\noindent\textit{Evidence.} Thom IV.13--IV.14 checked. Orientation and characteristic-number interpretation matter; connected oriented top-degree classes are detected by the number.\par
\subsection*{P001052 (both; faithful error): Donaldson manifold equality}
\noindent\textit{Transcription reads.} The manifold equals a generalized \(\mathbb{CP}^2\).\par
\noindent\textit{Critical reading.} Its integral form is \(\operatorname{diag}(1,\ldots,1)\); no diffeomorphism classification follows.\par
\noindent\textit{Evidence.} Scan P001052 confirms the wording with nearby cancellation. Donaldson page 81 states diagonalization.\par
\subsection*{P001052-P001053 (both; faithful error): Conjugation and equation (2)}
\noindent\textit{Transcription reads.} \(d_t=e^{-tf}de^{tf}\), then \(d_t=d-t\,df\wedge\), \(d_t^*=d^*-t\,\iota_{\nabla f}\).\par
\noindent\textit{Critical reading.} Both later signs are plus for the given conjugation.\par
\noindent\textit{Evidence.} Both scanned minuses are clear. Leibniz gives \(e^{-tf}d(e^{tf}\omega)=d\omega+t\,df\wedge\omega\).\par
\subsection*{P001053 (both; faithful error): Creation and annihilation prose}
\noindent\textit{Transcription reads.} The unstarred \(a_i\) is exterior multiplication.\par
\noindent\textit{Critical reading.} \(a_i=\iota_{e_i}\) contracts; \(a_i^*=e_i^\flat\wedge\) creates.\par
\noindent\textit{Evidence.} Kirillov explicitly corrects the earlier phrase in scan P001053; the vacuum convention agrees.\par
\subsection*{P001054 (both; faithful error): Equation (9) and marginal coefficient}
\noindent\textit{Transcription reads.} \(t\sum_{i<j}[a_i^*,a_j]c_{ij}\), with marginal coefficient \(4t\).\par
\noindent\textit{Critical reading.} For \(f=\sum_i s_i\lambda_i y_i^2\), use \(2t\sum_i s_i\lambda_i[a_i^*,a_i]\).\par
\noindent\textit{Evidence.} Scan P001054 confirms \(i<j\) and four. The Hessian entries are \(2s_i\lambda_i\), independently agreeing with oscillator cancellation.\par
\subsection*{P001054 (both; faithful error): Equation (10), annihilated form}
\noindent\textit{Transcription reads.} \(dy_{n-p+1}\wedge\cdots\wedge dy_n\) alone.\par
\noindent\textit{Critical reading.} \(e^{-t\sum_i\lambda_i y_i^2}dy_{n-p+1}\wedge\cdots\wedge dy_n\).\par
\noindent\textit{Evidence.} The scan has only the wedge. Substitution and square-integrability require the Gaussian.\par
\subsection*{P001054 (both; ambiguity): Poincaré group parenthesis}
\noindent\textit{Transcription reads.} The group is \ru{простая}, translated simple.\par
\noindent\textit{Critical reading.} If uncomplicated is intended, say so; the group is not algebraically simple.\par
\noindent\textit{Evidence.} Translations form a proper normal subgroup. The colloquial Russian meaning remains possible; scan inspected.\par
\subsection*{P001055 (both; ambiguity): Equation (1), height-domain notation}
\noindent\textit{Transcription reads.} \(p\subset\mathbb Z^2\).\par
\noindent\textit{Critical reading.} \(p\in\mathbb Z^2\), or let \(p\) denote a plaquette indexed by \(\mathbb Z^2\).\par
\noindent\textit{Evidence.} Scan P001055 has a subset-shaped sign; the surrounding height map and plaquette explanation show the intended indexing.\par
\subsection*{P001056 (both; ambiguity): Equation (5), tail equivalence}
\noindent\textit{Transcription reads.} \(\Pr(|h_p|=n)\sim e^{-{\rm const}\,n}\).\par
\noindent\textit{Critical reading.} An exponential upper bound is justified by the stated Peierls argument; a sharp asymptotic requires an additional theorem.\par
\noindent\textit{Evidence.} Scan P001056 confirms the tilde. Fröhlich--Spencer Section 7 records bounded low-temperature fluctuations, not the exact tail constant asserted here.\par
\subsection*{P001056 (both; faithful error): Continuous symmetry of integer heights}
\noindent\textit{Transcription reads.} \(h_p\mapsto h_p+\mathrm{const}\) is continuous for the SOS measure.\par
\noindent\textit{Critical reading.} Exact microscopic shifts require an integer constant; the real shift belongs to a continuous-field representation or limit.\par
\noindent\textit{Evidence.} Scan P001055 fixes \(h_p\in\mathbb Z\); P001056 calls the shift continuous. Noninteger shifts leave the configuration space.\par
\subsection*{P001056 (both; ambiguity): Equation (9), combinatorial justification}
\noindent\textit{Transcription reads.} \(a_\gamma(n)\sim C_\gamma^n n^{\delta(\gamma)}\) follows from general combinatorial theorems.\par
\noindent\textit{Critical reading.} The identified 1983 theorem gives \(\log a_\gamma(n)=\beta_0 n+o(n)\), with no general power-law refinement.\par
\noindent\textit{Evidence.} Scan P001056 confirms the stronger formula. The checked Durhuus--Fröhlich--Jónsson abstract states only exponential growth; \(C^ne^{\sqrt n}\) illustrates the logical gap.\par
\subsection*{P001057 (both; ambiguity): Equations (12)--(13), overwritten endpoints}
\noindent\textit{Transcription reads.} Possibly \(\beta\ge\beta_0\) and \(\beta\le\beta_0\).\par
\noindent\textit{Critical reading.} Use strict \(\beta>\beta_0\) for convergence and \(\beta<\beta_0\) for divergence; leave equality undecided.\par
\noindent\textit{Evidence.} Scan P001057 is overwritten. The root test alone says nothing at the critical value.\par
\subsection*{P001057 (both; faithful error): Density conclusion after equation (13)}
\noindent\textit{Transcription reads.} \(Z_{\beta,\gamma}=\infty\) implies positive occupied-plaquette density.\par
\noindent\textit{Critical reading.} A finite-volume normalized measure and a specified limit are needed before asserting a density.\par
\noindent\textit{Evidence.} The scan makes this implication explicitly. An infinite normalizing constant does not define the probability law used in the conclusion.\par
\subsection*{P001057 (both; faithful error): Roughening label at large inverse temperature}
\noindent\textit{Transcription reads.} \(\beta>\beta_R\): roughening; almost flat, nearly constant shape.\par
\noindent\textit{Critical reading.} Call the described large-\(\beta\) phase smooth or rigid. Roughening occurs on the smaller-\(\beta\) side.\par
\noindent\textit{Evidence.} Scan P001057 confirms the English word and contradictory explanatory prose; the preceding SOS statements have the correct temperature order.\par
\subsection*{P001057 (both; ambiguity): Equation (14), handle factor and scale}
\noindent\textit{Transcription reads.} The handle factor is in the denominator, \(e^{-cg}\), with \(\prod_e d\ell_e/\ell_e\).\par
\noindent\textit{Critical reading.} The displayed sign gives \(e^{cg}\). Retain it unless a suppression convention is explicitly intended; specify a scale constraint or regulator before normalization.\par
\noindent\textit{Evidence.} The fraction bar is clear in scan P001057. Global scaling gives \(dr/r\), unchanged deficits and area \(r^2A\), proving small-scale divergence of the bare closed-triangulation integral.\par
\subsection*{P001057 (both; provenance): Regge date and curvature normalization}
\noindent\textit{Transcription reads.} Regge 64; unqualified equality of hinge sum and scalar-curvature integral.\par
\noindent\textit{Critical reading.} Regge 1961. For the convention \(\mathrm{Scal}=2K\), the hinge sum has an overall factor two; state any alternative normalization.\par
\noindent\textit{Evidence.} Scan P001057 shows 64. The opened publisher record gives February 1961. In dimension two the deficit sum is total Gaussian curvature, half total usual scalar curvature.\par
\subsection*{P001057 (both; ambiguity): Curvature theorem versus functional-integral limit}
\noindent\textit{Transcription reads.} Curvature convergence implies that equation (14) tends to Polyakov theory.\par
\noindent\textit{Critical reading.} Keep the proposed measure limit separate from deterministic weak curvature convergence.\par
\noindent\textit{Evidence.} Cheeger--Müller--Schrader Theorem 5.1 assumes fat metric approximations and proves curvature-measure convergence, not convergence of the displayed edge-length probability measure.\par
\subsection*{P001058 (both; faithful error): Opening Virasoro commutator}
\noindent\textit{Transcription reads.} \((m-n)L_{n+m}+(c/12)\delta_{n+m,0}\mathbf1\).\par
\noindent\textit{Critical reading.} Retain the reverse Witt term for notebook \(L_n\), but restore \(\frac{c}{12}(n^3-n)\delta_{n+m,0}\mathbf1\). Modern generators are \(\ell_n=-L_n\), with \(h=-\Delta\).\par
\noindent\textit{Evidence.} Scan checked: no cubic factor is present. The written central term violates antisymmetry even at \(n=m=0\). The reverse Witt sign itself is a convention, not an independent error.\par
\subsection*{P001058 (both; faithful error): Sentence after the primary commutator}
\noindent\textit{Transcription reads.} The operator is uniquely determined by this formula.\par
\noindent\textit{Critical reading.} The covariance law needs module channels and a normalization; it is unchanged by scalar multiplication.\par
\noindent\textit{Evidence.} Russian and scan explicitly say uniquely. Multiplying any solution by a scalar is an immediate counterexample to unrestricted uniqueness.\par
\subsection*{P001058 (both; ambiguity): Neretin reducibility polynomial convention}
\noindent\textit{Transcription reads.} Both factors begin with \(-\Delta\).\par
\noindent\textit{Critical reading.} Retain this polynomial with \(h=-\Delta\) under \(\ell_n=-L_n\) for the reverse Witt bracket; translate coherently if the opening bracket and weight are renormalized.\par
\noindent\textit{Evidence.} Scan confirms both minus signs. At \((\alpha,\beta)=(1,2)\) it gives \(16\Delta^2-2(c-5)\Delta+c=0\), exactly the standard level-two condition with \(h=-\Delta\). The polynomial alone is not a faithful mathematical error.\par
\subsection*{P001058 (both; ambiguity): First degeneration level}
\noindent\textit{Transcription reads.} Degeneration occurs at level \(\alpha\beta\).\par
\noindent\textit{Critical reading.} That condition supplies a singular-vector level; the first is the minimum of \(rs\) over all matching positive integer pairs.\par
\noindent\textit{Evidence.} Multiple reducibility conditions may hold simultaneously; Feigin--Fuchs determinant gives factors at every level and their multiplicities.\par
\subsection*{P001058 (both; faithful error): Equation (2), matter and ghost determinants}
\noindent\textit{Transcription reads.} \((\det\Delta_g)^{-D/2}(\det\mathcal L)^{1/2}\), with no primes or cosmological weight.\par
\noindent\textit{Critical reading.} Use positive operators \(-\Delta_g\) and \(P^*P\), primed determinants and explicit zero-mode/moduli factors; retain the cosmological weight.\par
\noindent\textit{Evidence.} The Laplacian definition on P001059 is nonpositive. Integration by parts gives \(\mathcal L=-P^*P/2\). The area exponential in equation (1) does not disappear under the Gaussian matter integration.\par
\subsection*{P001058 (both; faithful error): Primary covariance sign with the reverse Witt bracket}
\noindent\textit{Transcription reads.} \([L_n,\varphi]=[z^{n+1}\partial_z+(n+1)\Delta z^n]\varphi\).\par
\noindent\textit{Critical reading.} For notebook \(L_n=-\ell_n\), use the negative of the standard primary differential operator, with the field weight specified separately.\par
\noindent\textit{Evidence.} Jacobi applied to two mode commutators gives the standard \((n-m)\) bracket for the printed positive primary law, contrary to the notebook \((m-n)\) term. The sign conversion repairs this independently of the scalar-normalization ambiguity.\par
\subsection*{P001059 (both; ambiguity): Equation (3), first derivative}
\noindent\textit{Transcription reads.} \(\nabla_a^a\varepsilon_b\).\par
\noindent\textit{Critical reading.} \(\nabla_a\varepsilon_b\), with the upper mark recorded as doubtful or cancelled.\par
\noindent\textit{Evidence.} The scan has an overwritten upper mark. The free indices of the symmetric trace-free derivative permit no extra upper \(a\).\par
\subsection*{P001059 (both; faithful error): Critical-dimension flat-metric statement}
\noindent\textit{Transcription reads.} At \(D=26\) the answer equals that for \(g_{ab}=\delta_{ab}\).\par
\noindent\textit{Critical reading.} The local matter-plus-ghost Weyl anomaly cancels; residual moduli, automorphisms and insertion integrations remain.\par
\noindent\textit{Evidence.} The scan states an unqualified equality. A general closed surface has no single global Euclidean metric; the local cancellation is the precise prior result.\par
\subsection*{P001059 (both; faithful error): Stress tensor after equation (6)}
\noindent\textit{Transcription reads.} \(\varphi_a\varphi_b-\frac12\delta_{ab}|\nabla\varphi|^2+\alpha\varphi_{ab}-\beta\delta_{ab}|\nabla\varphi|^2\).\par
\noindent\textit{Critical reading.} For \(V=m_0^2e^\varphi\), a consistent tensor is \(\varphi_a\varphi_b-\delta_{ab}(\frac12|\nabla\varphi|^2+V)+2(\delta_{ab}\Delta\varphi-\varphi_{ab})\).\par
\noindent\textit{Evidence.} Scan confirms a repeated gradient-square final term. The corrected divergence and trace vanish by \(\Delta\varphi=V\); the printed ansatz cannot have both properties for general solutions.\par
\subsection*{P001059 (both; faithful error): Holomorphic and antiholomorphic stress formulas}
\noindent\textit{Transcription reads.} \(\varphi_z^2+2\varphi_{zz}\) and its barred counterpart.\par
\noindent\textit{Critical reading.} \(\varphi_z^2-2\varphi_{zz}\) and its barred counterpart, for the recorded \(e^\varphi\) action.\par
\noindent\textit{Evidence.} The plus signs are clear in the scan. Using \(\varphi_{z\bar z}=m_0^2e^\varphi/4\), the printed \(\partial_{\bar z}T\) equals \(m_0^2e^\varphi\varphi_z\), generally nonzero.\par
\subsection*{P001059 (both; faithful error): Second holomorphic conservation equation}
\noindent\textit{Transcription reads.} \(\partial_zT=0\).\par
\noindent\textit{Critical reading.} \(\partial_z\bar T=0\).\par
\noindent\textit{Evidence.} The scan has no bar on this \(T\); it is an omission already in the notebook, not an English translation error.\par
\subsection*{P001059 (both; ambiguity): Final Laurent series}
\noindent\textit{Transcription reads.} \(T(z)=\sum z^{n-2}L_n\).\par
\noindent\textit{Critical reading.} For modern \(\ell_n\), use \(T(z)=\sum_n\ell_nz^{-n-2}\). With earlier notebook \(L_n=-\ell_n\), this is \(-\sum_nL_nz^{-n-2}\); retaining the final exponent requires reindexing the coefficients.\par
\noindent\textit{Evidence.} Scan confirms \(n-2\). Standard residue extraction is \(\ell_n=(2\pi i)^{-1}\oint z^{n+1}T(z)\,dz\). One cannot use the same labels unaltered in all three notebook formulas.\par
\subsection*{P001060 (both; ambiguity): Final gluing theorem}
\noindent\textit{Transcription reads.} \(vu=N\), \ru{что-то}.\par
\noindent\textit{Critical reading.} Retain the uncertainty; do not replace \(N\) silently by a fully specified monodromy logarithm.\par
\noindent\textit{Evidence.} Scan inspected. The word \ru{что-то} is present. The precise operator depends on the nearby-cycle convention; see \citet[section 3]{beilinson1987glue}.\par
\subsection*{P001061 (both; faithful error): Schubert-cell dimension}
\noindent\textit{Transcription reads.} \(\dim\text{cell}=\sum_jd_j\), where \(d_j=\dim(V\cap F_j)\).\par
\noindent\textit{Critical reading.} For jumps \(a_i\), the cell dimension is \(\sum_i(a_i-i)\).\par
\noindent\textit{Evidence.} The scan and both transcriptions agree on the written sum. For the zero-dimensional cell \(V=F_k\), \(d_j=\min(j,k)\) has positive sum. This is a faithful mathematical error, not a translation correction.\par
\subsection*{P001062 (both; faithful error): Projective superspace chart}
\noindent\textit{Transcription reads.} Last odd coordinate \(\xi_m/x_i\).\par
\noindent\textit{Critical reading.} \(\xi_n/x_i\); omit \(x_i/x_i=1\) from independent even coordinates.\par
\noindent\textit{Evidence.} Scan and Russian agree on m, while the preceding homogeneous list has n odd variables. The dimension and coordinate list force the correction.\par
\subsection*{P001064 (both; faithful error): Uniqueness of the local invariant Lagrangian}
\noindent\textit{Transcription reads.} The only local invariant Lagrangian.\par
\noindent\textit{Critical reading.} Qualify by admissible derivative order, total divergences and background terms; unrestricted uniqueness is false.\par
\noindent\textit{Evidence.} The scan says local without an order bound. For example a curvature-squared metric action provides a higher-derivative local invariant.\par
\subsection*{P001065 (both; faithful error): Equations 5--6}
\noindent\textit{Transcription reads.} \(\Delta\varphi=e^\varphi\), \(e^\varphi=|w^\prime|^2/(\operatorname{Im}w)^2\).\par
\noindent\textit{Critical reading.} For Euclidean Delta insert 2 in the numerator, or replace the equation by \(\Delta\varphi=2e^\varphi\).\par
\noindent\textit{Evidence.} Both scans agree; \(w(z)=z\) gives \(\Delta(-2\log y)=2/y^2\). The cusp formula on the same page already contains 2.\par
\subsection*{P001065 (both; faithful error): Equations 9--12}
\noindent\textit{Transcription reads.} Schwarzian denominators \((w^\prime)^2,(w^\prime)^4\) and a common symbol T.\par
\noindent\textit{Critical reading.} \(\{w,z\}=w^{\prime\prime\prime}/w^\prime-\tfrac32(w^{\prime\prime}/w^\prime)^2\); the stress tensor is minus this, while the ODE coefficient is half this.\par
\noindent\textit{Evidence.} Scan confirms the powers. Direct differentiation and the ratio-of-solutions Schwarzian identity independently fix the normalizations.\par
\subsection*{P001065 (both; faithful error): Reality of accessory parameters}
\noindent\textit{Transcription reads.} The accessory coefficients must be real.\par
\noindent\textit{Critical reading.} The uniformizing monodromy is conjugate into PSL(2,R); coefficients need not be real for arbitrary complex punctures.\par
\noindent\textit{Evidence.} Russian and English agree; the surrounding text does not restrict punctures to the real axis. The classical projective-connection formula makes the distinction.\par
\subsection*{P001065 (both; faithful error): Solutions following equation~13}
\noindent\textit{Transcription reads.} Two independent solutions of 13.\par
\noindent\textit{Critical reading.} Two independent solutions of 12.\par
\noindent\textit{Evidence.} Equation 12 is the second-order ODE; 13 is their quotient. Scan confirms the self-reference.\par
\subsection*{P001066 (both; faithful error): Equation 16}
\noindent\textit{Transcription reads.} \(\partial\mathcal L/\partial\varphi_{\bar z}=\varphi_{\bar z}\).\par
\noindent\textit{Critical reading.} \(\partial\mathcal L/\partial\varphi_{\bar z}=\varphi_z\), with the overall normalization.\par
\noindent\textit{Evidence.} The separate bar on the final z is visible in the scan; elementary differentiation of the preceding kinetic product reverses it.\par
\subsection*{P001066 (both; faithful error): Equations 18--19}
\noindent\textit{Transcription reads.} Laurent modes without weight 2 and central term without Kronecker delta.\par
\noindent\textit{Critical reading.} Use \(T(z)=\sum L_nz^{-n-2}\) and central term \(c(n^3-n)\delta_{n+m,0}/12\).\par
\noindent\textit{Evidence.} The scan contains no delta. Antisymmetry already fails without it, for example \(n=m=2\); BPZ section 2 confirms the conventional modes.\par
\subsection*{P001066 (both; faithful error): Equation 20}
\noindent\textit{Transcription reads.} \(T_\varepsilon=\int\varepsilon T\,d^2z\).\par
\noindent\textit{Critical reading.} For the holomorphic charge use \((2\pi i)^{-1}\oint\varepsilon T\,dz\).\par
\noindent\textit{Evidence.} Scan says d²z; a holomorphic mode generator is a contour charge. Other real-time conventions require an explicitly specified slice.\par
\subsection*{P001066 (both; faithful error): Equation 22}
\noindent\textit{Transcription reads.} \((n+1)\Delta\psi\) or \((n+1)\Delta F\).\par
\noindent\textit{Critical reading.} Insert \(z^n\), or \(x^n\) in the x-coordinate version.\par
\noindent\textit{Evidence.} All three scan lines omit the factor; substituting \(\epsilon=z^{n+1}\) into the immediately preceding primary variation requires it.\par
\subsection*{P001067-P001068 (both; faithful error): Resolution of the identity}
\noindent\textit{Transcription reads.} \(I=\sum|n\rangle\langle n|\), \(|n\rangle=L_{-n}|0\rangle\).\par
\noindent\textit{Critical reading.} Sum over all partition descendants with the inverse Gram matrix; restrict or quotient when degenerate.\par
\noindent\textit{Evidence.} The source itself notes nonorthogonality. At level 2 there are two descendants and at level 1 the norm is \(2\Delta\), so the displayed resolution is not literal.\par
\subsection*{P001068 (both; faithful error): Primary-field axiom}
\noindent\textit{Transcription reads.} \((n+1)\Delta\psi\) or \((n+1)\Delta F\).\par
\noindent\textit{Critical reading.} Insert \(z^n\), or \(x^n\) in the x-coordinate version.\par
\noindent\textit{Evidence.} All three scan lines omit the factor; substituting \(\epsilon=z^{n+1}\) into the immediately preceding primary variation requires it.\par
\subsection*{P001068-P001069 (both; faithful error): Bosonic quadratic sum}
\noindent\textit{Transcription reads.} \(\sum_{k=1}^{n-1}a_ka_{n-k}\) and the corresponding second-derivative sum.\par
\noindent\textit{Critical reading.} Insert a factor \(1/2\).\par
\noindent\textit{Evidence.} Scan confirms its absence. At \(\alpha=\beta=0\), \(n=2\), the vacuum commutator gives 2 instead of the required \(c/2=1/2\); inserting \(1/2\) gives the correct value.\par
\subsection*{P001069 (both; faithful error): Right-hand Ward operator}
\noindent\textit{Transcription reads.} \((n+1)\Delta\psi\) or \((n+1)\Delta F\).\par
\noindent\textit{Critical reading.} Insert \(z^n\), or \(x^n\) in the x-coordinate version.\par
\noindent\textit{Evidence.} All three scan lines omit the factor; substituting \(\epsilon=z^{n+1}\) into the immediately preceding primary variation requires it.\par
\subsection*{P001070 (both; faithful error): Equation 6}
\noindent\textit{Transcription reads.} \(\Box\varphi+\lambda^2\varphi(\varphi^2-\eta^2)=0\).\par
\noindent\textit{Critical reading.} Insert 4 before \(\lambda^2\).\par
\noindent\textit{Evidence.} Scan confirms no 4. Differentiate the preceding potential; independent confirmation is the 1974 Kobzarev coauthored wall paper, section~3.\par
\subsection*{P001070 (both; faithful error): After equation~7}
\noindent\textit{Transcription reads.} Independent of x,y,z.\par
\noindent\textit{Critical reading.} Independent of t,x,y for the displayed z-dependent static wall.\par
\noindent\textit{Evidence.} The source list includes z, contradicting the asymptotic condition immediately before it; the intended planar wall fixes the emendation.\par
\subsection*{P001071 (both; faithful error): Wall profile and thickness}
\noindent\textit{Transcription reads.} \(\tanh(z\eta)\), width of order \(1/\eta\).\par
\noindent\textit{Critical reading.} For equation~2 use \(\eta\tanh(\sqrt2\lambda\eta z)\), width of order \((\lambda\eta)^{-1}\).\par
\noindent\textit{Evidence.} The amplitude and coupling are absent in the scan. The printed potential and boundary values determine them; a rescaling would have to be stated.\par
\subsection*{P001071 (both; faithful error): Equations 9, 12, 13}
\noindent\textit{Transcription reads.} D=d+iA and A maps to A+dalpha for phase exp(ialpha).\par
\noindent\textit{Critical reading.} With that phase and D, use A maps to A-dalpha.\par
\noindent\textit{Evidence.} All three plus signs are visible in the scan. Direct differentiation produces an extra \(2i\,d\alpha\) unless one sign is changed.\par
\subsection*{P001073 (both; faithful error): Top line}
\noindent\textit{Transcription reads.} Minimum area.\par
\noindent\textit{Critical reading.} Stationary Lorentzian worldsheet area.\par
\noindent\textit{Evidence.} The scan says min. The timelike Nambu action has an indefinite second variation; its equation is an extremality condition.\par
\subsection*{P001073 (both; ambiguity): Equations 23, 24, 26}
\noindent\textit{Transcription reads.} D has no coupling factor while F has g; sign conventions are not specified.\par
\noindent\textit{Critical reading.} Choose one consistent generator, D and F convention.\par
\noindent\textit{Evidence.} Scan confirms the unequal appearances of g. Either absorb it throughout or retain it in both the covariant derivative and curvature.\par
\subsection*{P001074 (both; ambiguity): Kirillov’s phrase extends unitarily}
\noindent\textit{Transcription reads.} The \(SL(2,\mathbb R)\) representation extends unitarily to \(\operatorname{Diff}(S^1)\).\par
\noindent\textit{Critical reading.} The weighted pullback extends algebraically and is generally not unitary; the new Fock/Gaussian representation is unitary.\par
\noindent\textit{Evidence.} The scan itself inserts \ru{это неунитарное, но} beside the pullback formula. The two representation spaces must not be conflated.\par
\subsection*{P001076 (both; faithful error): Power kernel for every rho}
\noindent\textit{Transcription reads.} For every rho the displayed expression is an inner product.\par
\noindent\textit{Critical reading.} Restrict to a positive parameter range or add the appropriate smooth correction and regularization.\par
\noindent\textit{Evidence.} The scan says every rho. Neretin (1983), Theorems 1--2, explicitly distinguishes positive complementary kernels from indefinite kernels and their corrected positive forms.\par
\subsection*{P001076 (both; faithful error): Logarithmic scalar product}
\noindent\textit{Transcription reads.} \(\iint\log|\sin((t-s)/2)|f(t)f(s)dt\,ds\).\par
\noindent\textit{Critical reading.} Insert an overall minus sign for a positive norm on real mean-zero functions.\par
\noindent\textit{Evidence.} The scan has a positive logarithm. Its nonzero Fourier coefficients are negative; testing a cosine gives a negative quadratic form. The sign is a mathematical emendation.\par
\subsection*{P001077 (both; faithful error): Projections}
\noindent\textit{Transcription reads.} \(H\pm i\).\par
\noindent\textit{Critical reading.} \(P_\pm=(1\mp iH)/2\), with the same eigenspace images.\par
\noindent\textit{Evidence.} The scan explicitly has H plus/minus i. Squaring proves these operators are not idempotent; H squared equals minus the identity on the quotient.\par
\subsection*{P001077 (both; faithful error): Expected logarithmic convolution}
\noindent\textit{Transcription reads.} The displayed log/PV(1/x) convolution should give \(\delta^{\prime}\).\par
\noindent\textit{Critical reading.} The displayed kernels have product Fourier order minus one, not plus one; additional derivatives or inverse kernels would be required.\par
\noindent\textit{Evidence.} The prime on delta is clearly visible in the scan. Independent Fourier-symbol and circle Fourier-mode checks agree on the order mismatch.\par
\subsection*{P001077 (both; ambiguity): Analytic sequence second component}
\noindent\textit{Transcription reads.} \(f_2(z_1,z_1)\).\par
\noindent\textit{Critical reading.} In the symmetric Fock interpretation use \(f_2(z_1,z_2)\).\par
\noindent\textit{Evidence.} The second subscript resembles 1 in the scan, already flagged by the transcription. The n-variable sequence supplies the intended two independent variables; preserve the diplomatic reading and record this as an interpretive emendation.\par
\subsection*{P001078 (both; faithful error): Polynomial volume corollary}
\noindent\textit{Transcription reads.} \(M_\xi=\mu^{-1}(\xi)\), then \(\operatorname{vol}(M_\xi)\) is piecewise polynomial.\par
\noindent\textit{Critical reading.} The symplectic volume is that of \(\mu^{-1}(\xi)/T\), at regular values, with the reduced form.\par
\noindent\textit{Evidence.} The level-set definition and later volume notation are both visible in the scan. The intervening line already descends the form to the quotient, which is the intended object.\par
\subsection*{P001078 (both; faithful error): Equation (5), nonisolated fixed components}
\noindent\textit{Transcription reads.} \(\sum_j\operatorname{vol}(M_j)e^{iH_X(M_j)}/\prod_k\lambda_{jk}(X)\).\par
\noindent\textit{Critical reading.} Use the equivariant integral over each fixed component with the inverse Euler class of its normal bundle; normal Chern/curvature terms generally contribute.\par
\noindent\textit{Evidence.} The scan explicitly includes positive-dimensional fixed components and only their volume. Duistermaat--Heckman’s Addendum, Theorem (1.5), and Atiyah--Bott equation (7.4) show the omitted terms.\par
\subsection*{P001079 (both; faithful error): Stationary-phase display with parameter t}
\noindent\textit{Transcription reads.} The exponent contains \(tH_X\), but the denominator remains \(\prod_k\lambda_{jk}(X)\).\par
\noindent\textit{Critical reading.} For isolated fixed points on a \(2n\)-dimensional manifold the denominator must scale with \(\prod_k\lambda_{jk}(tX)=t^n\prod_k\lambda_{jk}(X)\), with the chosen phase factors.\par
\noindent\textit{Evidence.} The missing t is confirmed in the scan. Scaling in the quadratic Gaussian integral and the exact localization formula both require the power; positive-dimensional components require the additional characteristic terms.\par
\subsection*{P001079 (both; faithful error): SU(3) orbit projection diagram}
\noindent\textit{Transcription reads.} \(A\mapsto(\operatorname{tr}A^2,\operatorname{tr}A^3)\) on an orbit has hexagonal image.\par
\noindent\textit{Critical reading.} These invariants are constant on the orbit. The hexagonal moment image belongs to the diagonal projection \(A\mapsto(A_{11},A_{22},A_{33})\), with sum zero.\par
\noindent\textit{Evidence.} The two trace expressions are visibly written in the scanned diagram. Conjugacy invariance proves its image is a point; the preceding orbit discussion and Schur--Horn convex hull identify the intended diagonal moment map.\par
\subsection*{P001081 (both; faithful error): Equation (3), order of polynomial gap opening}
\noindent\textit{Transcription reads.} Width \(\asymp\varepsilon^{[N/p]}\).\par
\noindent\textit{Critical reading.} The generally valid conclusion is \(O(|\varepsilon|^{\lceil N/p\rceil})\); a matching lower bound requires additional nonvanishing conditions.\par
\noindent\textit{Evidence.} The scan confirms the two-sided asymptotic symbol and the unexplained brackets. Arnold’s Theorems 3 and 7 explicitly use an upper bound and minus the floor of minus N/p. This separates an ambiguous bracket convention from the false universal asymptotic equivalence.\par
\subsection*{P001081 (both; faithful error): Generic Hill high-index width}
\noindent\textit{Transcription reads.} The width decreases exponentially for the unspecified generic periodic coefficient.\par
\noindent\textit{Critical reading.} Exponential high-index decay requires analytic regularity; smooth real periodic potentials generally give rapid polynomial decay instead.\par
\noindent\textit{Evidence.} The scan and Russian text say exp decreases without an analytic hypothesis. Pöschel’s section 1 distinguishes the smooth and analytic spectral equivalences; the small-amplitude and high-index limits also need separation.\par
\subsection*{P001082 (both; faithful error): Catalan sequence final number}
\noindent\textit{Transcription reads.} 138\par
\noindent\textit{Critical reading.} 132\par
\noindent\textit{Evidence.} The scan clearly writes 138. The adjacent factorial formula and the convolution recurrence independently give 132, as printed in Arnold’s section 4.\par
\subsection*{P001082 (both; faithful error): Equation (5), Catalan recurrence}
\noindent\textit{Transcription reads.} \(a_{n+1}=\sum_{j=1}^{n-1}a_ja_{n-j}\).\par
\noindent\textit{Critical reading.} With \(a_n=C_n\), use \(C_{n+1}=\sum_{j=0}^{n}C_jC_{n-j}\).\par
\noindent\textit{Evidence.} The scan clearly shows the upper bound n minus one and lower bound one. The listed initial values force the endpoint terms; at n equal to one the written empty sum contradicts \(C_2=2\).\par
\subsection*{P001082 (both; faithful error): Triangulation interpretation}
\noindent\textit{Transcription reads.} The nth Catalan number counts simplicial subdivisions of an n-gon.\par
\noindent\textit{Critical reading.} It counts triangulations of a convex (n+2)-gon by noncrossing diagonals, with no new vertices.\par
\noindent\textit{Evidence.} The scan writes n-gon. Arnold’s section 4 explicitly writes (n+2)-gon, and the triangle/quadrilateral counts independently fix the shift.\par
\subsection*{P001087 (both; faithful error): equation (3), variational annotation}
\noindent\textit{Transcription reads.} \(I_3=\tfrac14(u_x^2+u^3),\quad \delta I_3/\delta u=\tfrac14u_{xx}-\tfrac34u^2\)\par
\noindent\textit{Critical reading.} For the written density the derivative is \(-\tfrac12u_{xx}+\tfrac34u^2\); the displayed flow instead uses \(\int(-u_x^2/8-u^3/4)\,dx\).\par
\noindent\textit{Evidence.} Scan and both texts agree; integration by parts gives the discrepancy.\par
\subsection*{P001088 (both; faithful error): equation (10)}
\noindent\textit{Transcription reads.} \((\delta e_s)^2=|\langle\delta u\Psi^2\rangle|^2/\langle|\Psi|^2\rangle\)\par
\noindent\textit{Critical reading.} Specify normalization and use the two-dimensional eigenvalue perturbation matrix; a scale-invariant single-mode formula needs the squared norm in the denominator.\par
\noindent\textit{Evidence.} Scan confirms no outer square. Rescaling the Bloch eigenfunction changes the written expression; free Fourier perturbation also fixes the full-gap factor 4.\par
\subsection*{P001089 (both; faithful error): equation (12)}
\noindent\textit{Transcription reads.} Second derivative of log of a single shifted theta function times an exponential\par
\noindent\textit{Critical reading.} Use the theta quotient and the Abel shifts in Krichever 1983 Lemma 3, eq.(13).\par
\noindent\textit{Evidence.} The notebook and transcriptions agree. The printed lemma and genus-zero limit disprove the literal formula.\par
\subsection*{P001089 (both; ambiguity): equation (13)}
\noindent\textit{Transcription reads.} Both factors carry \(s\), while the condition is \(s\ne s^\prime\).\par
\noindent\textit{Critical reading.} Distinguish the paired mode indices, as in Krichever 1983 Lemma 3.\par
\noindent\textit{Evidence.} Scan crop orthogonality.png confirms the condition and the unresolved repeated index; no confident claim of an OCR error.\par
\subsection*{P001089 (both; faithful error): binomial expansion}
\noindent\textit{Transcription reads.} \(n^{-2s}(-as/n^{2s+2}+\cdots)\)\par
\noindent\textit{Critical reading.} \(n^{-2s}(1-sa/n^2+s(s+1)a^2/(2n^4)-\cdots)\)\par
\noindent\textit{Evidence.} Scan and transcriptions omit the constant term and duplicate a power of n; binomial expansion checks the correction.\par
\subsection*{P001091 (both; faithful error): free spectral zeta normalization}
\noindent\textit{Transcription reads.} \(T=2\pi,\quad 2^{2m}|B_{2m}|/(2m)!\)\par
\noindent\textit{Critical reading.} For period \(T\): \(T^{2m}|B_{2m}|/(2m)!\).\par
\noindent\textit{Evidence.} The eigenvalues are \((2\pi n/T)^2\); the printed normalization conflicts with the preceding period.\par
\subsection*{P001092 (both; faithful error): Euler and formal sums}
\noindent\textit{Transcription reads.} Even zeta formula without \(\pi\); \(\sum n^3=-B_2\).\par
\noindent\textit{Critical reading.} \(\zeta(2m)=(-1)^{m+1}B_{2m}(2\pi)^{2m}/[2(2m)!]\); \(\zeta(-3)=-B_4/4\).\par
\noindent\textit{Evidence.} Both errors are present in the scan; these are analytic-continuation values, not ordinary sums.\par
\subsection*{P001092 (both; faithful error): Clausen value}
\noindent\textit{Transcription reads.} \(\mathrm{Cl}_5(0)=\zeta(-5)\)\par
\noindent\textit{Critical reading.} \(\mathrm{Cl}_5(0)=\zeta(5)\)\par
\noindent\textit{Evidence.} The displayed cosine-series definition on the same page gives the positive argument.\par
\subsection*{P001094 (both; faithful error): local action covering}
\noindent\textit{Transcription reads.} Every map is represented by a map into one primitive patch\par
\noindent\textit{Critical reading.} Use local domains in the mapping space; one-target-patch domains need not cover it. Differences of local primitives are locally constant.\par
\noindent\textit{Evidence.} The scan shows a target-space covering. The general construction requires local homotopies or subdivision.\par
\subsection*{P001094 (both; faithful error): factorization exponential}
\noindent\textit{Transcription reads.} \(\exp(t\,dh(L_0))\)\par
\noindent\textit{Critical reading.} \(\exp(2t\,dh(L_0))\) for the unhalved \(R\)-bracket.\par
\noindent\textit{Evidence.} Semenov 1983 Theorem 11, equation (18), explicitly has \(2t\); the derivative at \(t=0\) gives the same check.\par
\subsection*{P001094 (both; faithful error): Drinfeld journal pointer}
\noindent\textit{Transcription reads.} Doklady 268, no. 5.\par
\noindent\textit{Critical reading.} Doklady 268, no. 2 (1983), 285--287.\par
\noindent\textit{Evidence.} Drinfeld 1983 publisher metadata and original three-page paper both opened; issue 5 is the notebook reading.\par
\subsection*{P001096 (both; faithful error): Fedosov block projector}
\noindent\textit{Transcription reads.} \(f_1+f_2=1\)\par
\noindent\textit{Critical reading.} \(f_1^2+f_2^2=1\), with a unitary symbol for the displayed adjoint blocks.\par
\noindent\textit{Evidence.} Multiplying the displayed block matrix gives \(P^2=(f_1^2+f_2^2)P\) under unitary normalization. The scan confirms the unsquared written partition.\par
\subsection*{P001097 (both; ambiguity): auxiliary symbol a and cutoff}
\noindent\textit{Transcription reads.} \(a=f_1\sigma,\quad\tau=f_2\sigma^{-1}\)\par
\noindent\textit{Critical reading.} If the reading \(f_1\) is retained, replace its infinity behavior: \(a\) and \(\tau\) must be inverse outside a compact set. A likely intended choice is \(a=f_2\sigma\).\par
\noindent\textit{Evidence.} The \(f_1\) subscript is ambiguous in \nolinkurl{cutoff 1097b.png}. Under the preceding cutoff definitions, the literal \(f_1\) choice destroys compact support of \(P_1-P_2\); idempotence itself still holds.\par
\subsection*{P001097 (both; ambiguity): first shifted zeta argument}
\noindent\textit{Transcription reads.} \(\zeta_{\Delta_1+1}(\theta)\)\par
\noindent\textit{Critical reading.} \(\zeta_{\Delta_1+1}(0)\)\par
\noindent\textit{Evidence.} The mark is a crossed or overwritten zero. Only zero gives the index formula stated in context.\par
\subsection*{P001098 (both; ambiguity): boundary signature operator}
\noindent\textit{Transcription reads.} \(\pm(d^*-d)\)\par
\noindent\textit{Critical reading.} Use the odd signature operator built from \(d\) and the Hodge star, with degree-dependent signs.\par
\noindent\textit{Evidence.} The scan has a star near \(d\), but the shorthand does not specify Hodge star versus formal adjoint. The APS signature formula fixes its intended geometric meaning.\par
\subsection*{P001101-P001102 (both; ambiguity): full Grassmannian versus admissible frames}
\noindent\textit{Transcription reads.} Every point of \(\mathrm{Gr}\) has a frame with \(w_+=I+\text{trace class}\).\par
\noindent\textit{Critical reading.} The stated frames cover the index-zero component; the preceding full Grassmannian has all integer indices.\par
\noindent\textit{Evidence.} The scan contains both the all-index discussion and an index-zero annotation; Segal–Wilson section 3 explicitly restricts components.\par
\subsection*{P001103 (both; ambiguity): all Grassmannian points and KdV}
\noindent\textit{Transcription reads.} Any \(W\) gives the stated KdV construction.\par
\noindent\textit{Critical reading.} Any \(W\) gives KP; the \(n\)th KdV reduction requires \(z^nW\subseteq W\).\par
\noindent\textit{Evidence.} The restriction is made explicit later on P001105 and in Segal–Wilson Corollary 5.12.\par
\subsection*{P001105 (both; faithful error): \(n=2\) tau reconstruction}
\noindent\textit{Transcription reads.} \(u=D^2\log\tau\)\par
\noindent\textit{Critical reading.} \(u=2D^2\log\tau\) for \(L=D^2+u\).\par
\noindent\textit{Evidence.} The scan confirms the missing 2; Segal--Wilson equation (5.19) and the coefficient \(2q_1\) in \(Q^2\) give the correction.\par
\subsection*{P001109 (both; faithful error): stationary-phase formula for fixed submanifolds}
\noindent\textit{Transcription reads.} \(\operatorname{vol}(M_j)\exp(iH(M_j))/\lambda_j\)\par
\noindent\textit{Critical reading.} Use the full inverse equivariant Euler class on a positive-dimensional fixed component.\par
\noindent\textit{Evidence.} The normal curvature terms contribute after integration; Atiyah–Bott formula (3.8), sections 6–7 distinguish this from isolated fixed points.\par
\subsection*{P001109 (both; faithful error): pushforward and variation of symplectic form}
\noindent\textit{Transcription reads.} Polynomial on \(\mathbb R\); the symplectic form is linear.\par
\noindent\textit{Critical reading.} Polynomial on regular-value chambers; reduced cohomology class varies affinely.\par
\noindent\textit{Evidence.} The same page first says piecewise polynomial; Duistermaat–Heckman gives the chamberwise and cohomological statement.\par
\subsection*{P001111 (both; faithful error): square of the Cartan differential}
\noindent\textit{Transcription reads.} \(d_1^2=u^2L_\xi\)\par
\noindent\textit{Critical reading.} \(d_1^2=uL_\xi\)\par
\noindent\textit{Evidence.} Both versions and the scan agree on \(u^2\); expand \((d+u\iota_\xi)^2\) using Cartan homotopy.\par
\subsection*{P001112 (both; ambiguity): angular form and the subsequent repair}
\noindent\textit{Transcription reads.} Initially \(d_1\alpha=1\) is sought.\par
\noindent\textit{Critical reading.} The successful form is \(\beta=\alpha(1+d\alpha)^{-1}\), with \(d_1\beta=1\).\par
\noindent\textit{Evidence.} The scan explicitly gives this correction and the contracting identity; preserve the sequence, not an unresolved-error label.\par
\subsection*{P001113 (both; faithful error): convexity attribution}
\noindent\textit{Transcription reads.} Atiyah–Bott and Guillemin–Sternberg independently\par
\noindent\textit{Critical reading.} Atiyah and Guillemin–Sternberg independently for convexity; Atiyah–Bott for the equivariant localization treatment.\par
\noindent\textit{Evidence.} Atiyah 1982 Theorem 1 and Atiyah–Bott 1984 section 6.\par
\subsection*{P001113 (both; faithful error): good integral equations at every prime}
\noindent\textit{Transcription reads.} Choose sufficiently good equations so every reduction is smooth\par
\noindent\textit{Critical reading.} A smooth proper model exists away from finitely many bad primes; bad reduction is not generally removed by coordinate choice.\par
\noindent\textit{Evidence.} Arithmetic-model setup in BCGP 2021 section 1.1; scan contains the universal prime claim.\par
\subsection*{P001114 (both; faithful error): local zeta definition (2)}
\noindent\textit{Transcription reads.} \(\sum_n \#X(\mathbb F_{p^n})p^{-ns}\)\par
\noindent\textit{Critical reading.} \(\exp(\sum_{n\ge1}\#X(\mathbb F_{p^n})p^{-ns}/n)\)\par
\noindent\textit{Evidence.} The point example on the same page requires the exponential and \(1/n\); Deligne 1974, section 1, gives the standard definition.\par
\subsection*{P001114 (both; ambiguity): finite-field zero location}
\noindent\textit{Transcription reads.} \(\operatorname{Re}z=p^{?/2}\)\par
\noindent\textit{Critical reading.} \(|\alpha_i|=\sqrt p\); local \(t\)-zeros have \(|t|=p^{-1/2}\); with \(t=p^{-s}\), \(\operatorname{Re}s=1/2\).\par
\noindent\textit{Evidence.} Scan exponent is ambiguous but real part versus modulus and the variable are genuinely confused; no conjectural digit restored.\par
\subsection*{P001114 (both; faithful error): completed elliptic functional equation}
\noindent\textit{Transcription reads.} \((2\pi)^{-s}\Gamma(s)\zeta(X,s)\) symmetric under \(s\mapsto2-s\).\par
\noindent\textit{Critical reading.} \(\Lambda(E,s)=N^{s/2}(2\pi)^{-s}\Gamma(s)L(E,s)\), with \(\Lambda(E,s)=w_E\Lambda(E,2-s)\).\par
\noindent\textit{Evidence.} The conductor factor and root number are absent; the arithmetic curve zeta function and the \(H^1\) \(L\)-function are distinct.\par
\subsection*{P001115 (both; provenance): catalogue range for records 596 and 598}
\noindent\textit{Transcription reads.} Page P001115 included by overlapping catalogue ranges\par
\noindent\textit{Critical reading.} The page begins 5 September 1983, Deligne, and lies outside the 1982–83 academic-year records.\par
\noindent\textit{Evidence.} Date and speaker are clearly visible in the scan and both transcriptions.\par
